% !TEX program = pdflatex
\documentclass[11pt]{article}

\usepackage[T1]{fontenc}
\usepackage{lmodern}
\usepackage{amsmath,amssymb,amsthm}
\usepackage{geometry}
\usepackage{xcolor}
\usepackage{booktabs,longtable,array}
\usepackage{enumitem}
\usepackage{fancyhdr}
\usepackage{hyperref}

\geometry{a4paper,margin=25mm,headheight=15pt}
\setlength{\emergencystretch}{2em}
\hypersetup{
  colorlinks=true,
  linkcolor=blue!55!black,
  citecolor=green!40!black,
  urlcolor=blue!60!black,
  pdftitle={Simon's multidimensional L2 conjecture and the failure of fixed-potential selected-channel entropy},
  pdfsubject={Pluecker--Schur reductions and a reciprocal Riesz--Talbot construction},
  pdfauthor={Siam Siraji}
}

\pagestyle{fancy}
\fancyhf{}
\lhead{Selected-channel entropy}
\rhead{S.~Siraji}
\cfoot{\thepage}

\definecolor{proved}{RGB}{20,110,65}
\definecolor{open}{RGB}{170,80,15}
\definecolor{failed}{RGB}{165,35,35}
\definecolor{navy}{RGB}{20,45,75}
\definecolor{pale}{RGB}{242,247,251}

\newcommand{\statusproved}{\textcolor{proved}{\textsc{proved here}}}
\newcommand{\statusknown}{\textcolor{proved}{\textsc{known}}}
\newcommand{\statusopen}{\textcolor{open}{\textsc{open}}}
\newcommand{\statusconditional}{\textcolor{open}{\textsc{conditional}}}
\newcommand{\statusfailed}{\textcolor{failed}{\textsc{false}}}
\newcommand{\statusdiagnostic}{\textcolor{navy}{\textsc{diagnostic}}}

\newcommand{\R}{\mathbb R}
\newcommand{\C}{\mathbb C}
\newcommand{\N}{\mathbb N}
\newcommand{\Sd}{\mathbb S^{d-1}}
\newcommand{\K}{\mathcal K}
\newcommand{\cH}{\mathcal H}
\newcommand{\cE}{\mathcal E}
\newcommand{\cF}{\mathcal F}
\newcommand{\cP}{\mathcal P}
\newcommand{\cQ}{\mathcal Q}
\newcommand{\HS}{\mathfrak S_2}
\newcommand{\TC}{\mathfrak S_1}
\newcommand{\ac}{\mathrm{ac}}
\newcommand{\Id}{\mathrm I}
\newcommand{\dd}{\,\mathrm d}
\newcommand{\eps}{\varepsilon}
\newcommand{\wedgeN}{\bigwedge\nolimits^N}
\DeclareMathOperator{\tr}{tr}
\DeclareMathOperator{\rank}{rank}
\DeclareMathOperator{\spec}{spec}
\DeclareMathOperator{\esssupp}{ess\,supp}
\DeclareMathOperator{\dist}{dist}
\DeclareMathOperator{\Ran}{Ran}
\DeclareMathOperator{\diag}{diag}

\newtheorem{theorem}{Theorem}[section]
\newtheorem{proposition}[theorem]{Proposition}
\newtheorem{lemma}[theorem]{Lemma}
\newtheorem{corollary}[theorem]{Corollary}
\theoremstyle{definition}
\newtheorem{definition}[theorem]{Definition}
\newtheorem{problem}[theorem]{Proof obligation}
\newtheorem{programme}[theorem]{Attack module}
\theoremstyle{remark}
\newtheorem{remark}[theorem]{Remark}

\newcommand{\statusbox}[1]{%
  \begin{center}
  \fcolorbox{open}{pale}{\begin{minipage}{0.91\textwidth}\small #1\end{minipage}}
  \end{center}}

\title{\textbf{Simon's Multidimensional $L^2$ Conjecture and the Failure of Fixed-Potential Selected-Channel Entropy}\\[4pt]
\large Pl\"ucker--Schur reductions and a reciprocal Riesz--Talbot construction}
\author{Siam Siraji\\[3pt]
\small The University of Adelaide, Adelaide, South Australia, Australia\\
\small \href{mailto:siam.siraji@adelaide.edu.au}{siam.siraji@adelaide.edu.au}}
\date{24 July 2026}

\begin{document}
\maketitle

\begin{abstract}
Let $H=-\Delta+V$ on $L^2(\mathbb R^d)$, $d\geq2$, with $V$ bounded,
real, and square integrable against the radial weight $|x|^{1-d}$.
We examine a finite-angular-channel entropy strategy for Simon's
multidimensional $L^2$ conjecture.  For a finite angular model we first
derive an exact Pl\"ucker--Jacobi identity which expresses the selected
spectral density through the reciprocal volume of an analytic inverse Jost
frame.  We also record the corresponding full-column Hilbert--Schmidt cost,
strong-resolvent completion, and double-buffered high-channel estimates.

The main result is negative for the proposed proof mechanism.  In dimension
three we construct a single bounded, real, mean-zero zonal potential
$V\in X_3$ and a compact energy interval $K\Subset(0,\infty)$ such that, for
every cofinal positively smoothed, radially truncated, double-buffered
Galerkin diagonal,
\[
 \int_K\log^-\|J_n(k)^{-1}e_0\|^2\,\dd k\longrightarrow+\infty .
\]
The construction combines a translated-barrier Hardy identity, zonal
quarter-Talbot Riesz products, fractional-moment synchronization of a
reciprocal M\"obius phase, and uniform lacunary placement of radial cells at
summable charged cost.  It disproves the universal fixed-potential
selected-channel entropy estimate and the stronger endpoint estimates that
would imply it.  Several natural alternatives are also decided: unrestricted
matrix low-column estimates, coefficientwise quadratic expansion, raw
reflection bounds, and one-point Poisson--Jensen replacements fail, whereas
the exact Pl\"ucker, Riccati, interpolation, and buffered-completion
identities remain valid.

This result does not determine the absolutely continuous spectral
multiplicity of the limiting Schr\"odinger operator.  In particular, it
neither proves nor disproves Simon's conjecture: spectral density may persist
in angular directions which move with the approximation.  Any continuation
of this entropy programme must therefore be basis-free, moving-probe, or
full-operator, or must use a different spectral mechanism.
\end{abstract}

\noindent\textbf{2020 Mathematics Subject Classification.}
Primary 35P25, 47A55; Secondary 81Q10, 42A55.

\noindent\textbf{Keywords.}
Schr\"odinger operator, absolutely continuous spectrum, Jost matrix,
Pl\"ucker coordinates, entropy, Riesz product, Talbot effect.

\statusbox{\textbf{Scope of the claims.}
The unconditional main conclusion is the failure of the fixed selected-channel
entropy route.  Statements used only conditionally are labelled as such at the
point of use, and failed intermediate estimates are retained as explicit
counterexamples.  Simon's multidimensional $L^2$ conjecture remains open.}

\tableofcontents

\section{The conjecture and the exact standard of proof}

Let $V\in L^\infty(\R^d;\R)$ and
\begin{equation}
  H=-\Delta+V,
  \qquad \operatorname{dom}(H)=H^2(\R^d).
  \label{eq:H}
\end{equation}
Set
\begin{equation}
 \|V\|_{X_d}^2
 :=\int_{\R^d}|x|^{1-d}|V(x)|^2\dd x.
 \label{eq:Xd}
\end{equation}
In polar coordinates $x=r\omega$,
\begin{equation}
 \|V\|_{X_d}^2
 =\int_0^\infty\|V(r,\cdot)\|_{L^2(\Sd)}^2\dd r.
 \label{eq:polar-cost}
\end{equation}
The bounded-real formulation of Simon's Conjecture~20.2 is
\begin{equation}
  m_{\ac}(E;H)=\infty
  \quad\text{for Lebesgue-a.e. }E>0.
  \label{eq:simon}
\end{equation}
It does not assert that the spectrum on $(0,\infty)$ is purely absolutely continuous.  Embedded point or singular-continuous spectrum may coexist with the asserted absolutely continuous component.

The unitary polar flattening
\begin{equation}
  (\mathcal U f)(r,\omega)=r^{(d-1)/2}f(r\omega)
  \label{eq:polar-unitary}
\end{equation}
gives, on the exterior half-line,
\begin{equation}
 \widetilde H
 =-\partial_r^2+\frac{A_d}{r^2}+M_{V(r)},
 \qquad
 A_d=-\Delta_{\Sd}+\frac{(d-1)(d-3)}4,
 \label{eq:polar-H}
\end{equation}
acting on $L^2((1,\infty);\K)$ with $\K=L^2(\Sd)$.  Formula~\eqref{eq:polar-H} exposes the central tension: $M_{V(r)}$ is controlled in angular $L^2$, whereas its operator norm is $\|V(r,\cdot)\|_\infty$.

The primary literature still describes the conjecture as unresolved; see
Simon~\cite{Simon2019}, Denisov's survey~\cite{DenisovSurvey2012}, and his
later operator-valued discussion~\cite{Denisov2022}.  The radial case follows
from the one-dimensional $L^2$ theorem of Deift and
Killip~\cite{DeiftKillip1999}.  Known multidimensional theorems require
additional decay, sign, angular regularity, operator-norm control, or a
generic coupling parameter~\cite{LNS2005,Perelman2005,LaptevSafronov2024,Denisov2025}.

\section{Audit baseline and non-negotiable obstructions}

The earlier finite-channel audit establishes the following ledger.

\begin{longtable}{>{\raggedright\arraybackslash}p{0.30\textwidth}>{\raggedright\arraybackslash}p{0.18\textwidth}>{\raggedright\arraybackslash}p{0.43\textwidth}}
\toprule
Statement & Status & Consequence \\
\midrule
\endhead
Strong-resolvent convergence of controlled radial/angular truncations & \statusproved & Herglotz matrices and matrix spectral measures converge vaguely. \\
Finite-model Jost/Gram formula & \statusproved & The canonical boundary density is positive definite; compact-probe density is positive definite a.e. and has finite local entropy. \\
Lower semicontinuity of $\int\log^-\det W$ & \statusproved & A uniform entropy estimate survives the truncation limit. \\
Rank/countability finish & \statusconditional & Uniform $N$-channel entropy for every $N$ implies~\eqref{eq:simon}. \\
Leakage-only abstract FCEN & \statusfailed & Schr\"odinger locality and propagation are indispensable. \\
Truncated leakage $(P_L-P_N)M_VP_N$ & \statusdiagnostic & Angular frequencies above $L$ can alias into retained high channels. \\
Uniform full-leakage $\mathrm{PSR}_N$ endpoint at $\alpha=1$ & \statusfailed & Already for $N=1$, a cutoff-edge scalar rotator and pure-degree barrier have bounded full-column cost but divergent quotient entropy. \\
Coefficientwise quotient estimate $\mathrm{QBC}_1$ & \statusfailed & A separated radial two-slab family has bounded $L^2$ cost but a divergent positive quadratic coefficient. \\
Exact scalar physical-coupling entropy & \statusknown & The one-dimensional transmission trace inequality controls the radial $m=N=1$ Jost entropy uniformly in support radius. \\
Dual inverse frame and differential quotient flux & \statusproved & $\Gamma_P=\|\wedge^N(J^{-1}P)\|^{-1}$, and the high--high cocycle generator cancels exactly. \\
Unrestricted cocycle low-column bound & \statusfailed & A high block preloads amplification $e^H$, while a later quarter-turn costs only $\pi^2/(4L)$. \\
Physical Riccati descent & \statusproved & The flux-normalized outgoing frame gives an exact Riccati formula for $\Gamma_P$ and its differential quotient flux. \\
Cutoff-compatible scalar Riccati/Feshbach replacement & \statusfailed & The fixed Riesz--Talbot potential defeats every cofinal positive double-buffered diagonal, so a fixed selected probe is still insufficient. \\
Positive fixed-potential buffered diagonal and completion & \statusproved & Positive angular smoothing preserves $L^\infty$ and $X_3$, the buffered models converge strongly resolvent, and one liminf entropy estimate suffices. \\
Fixed-potential buffered selected $N=1$ entropy & \statusfailed & One bounded real mean-zero $X_3$ potential has divergent selected entropy along every cofinal positive smoothing, radial truncation, and double-buffered Galerkin diagonal. \\
Outgoing high-channel edge removal & \statusproved & If $J/R\to\infty$, the outgoing high-sector Schur complement is $O(R^2/J^2)$; the double buffer places the Galerkin edge beyond the active block. \\
Broadband wrapped-mask turning-scale test & \statusproved & Increasing pulse area gives uniform high-band mass, but vanishing charged cost and clean transport force the mask far outside the generated modes' turning radius. \\
Outer unitary converter ordering & \statusproved & A reflectionless outer right factor leaves the scalar Pl\"ucker quotient unchanged; suppression requires conversion within or inside the turning dynamics. \\
Zero-background dimension-normalized matrix sum rule & \statusproved & The determinant trace identity is linear in matrix dimension after normalization and matches the scalar $X_3$ cost. \\
Naive insertion of the Bessel background & \statusfailed & The background-subtracted quadratic invariant contains an explicit angular-mean term with coefficient of order $L^2$. \\
Radial-mean Bessel cancellation and signed relative sum rule & \statusproved & Taking $D_L/r^2+q(r)I$ as reference removes every cross term and gives a dimension-uniform signed logarithmic identity. \\
Star--Feshbach charged compressed resolvent & \statusproved & For $N=1$, retaining the full high--high block gives a positive trace-class negative-energy correction bounded only by $\int\|Q_0M_{V^\circ}e_0\|^2$. \\
Star-relative bound-state divisor & \statusproved & Off-diagonal sign-flip convexity makes the relative negative $3/2$-moment nonnegative, so the signed second trace is bounded by twice the charged column cost. \\
Scalar Schur high-minor divisor & \statusopen & Complex zeros of the chosen high minor have indefinite cubic trace contribution; the pole-free radiation column must retain their reciprocal refocusing. \\
Radial-mean two-port and reciprocal refocusing & \statusproved & Exact unitary gauges remove the forward blocks; the physical radiation column changes only through $\Theta'R$, where $|\Theta_n'|=2k/|f_{q,n}|^2$. \\
Bessel-balanced scalar generator & \statusproved & Sandwiching the return coupling by the reciprocal phase velocity gives a unitary conjugate of $W_L^\circ$, exactly preserving the charged low column and angular bandwidth. \\
Signed pole-free work reduction & \statusproved & Rationally weighted spectral mass controls density amplification, so a signed bound on the exact refocusing work already implies the one-sided $N=1$ entropy estimate. \\
Poisson/interior-row shortcut & \statusfailed & A fixed-band mean-zero slab can have vanishing charged cost and vanishing compact real-band forward loss while its forward Jost row at $-\kappa^2$ grows exponentially; the singular mass at infinity is not $L^2$-controlled. \\
Transport-free scalar $N=1$ entropy & \statusproved & After deleting $A_3/r^2$, angular fibers are scalar one-dimensional problems; Jensen and the scalar Deift--Killip transmission inequality give the selected-density entropy with exactly the charged cost. \\
Universal signed centrifugal endpoint & \statusfailed & The interpolation identity and endpoint cancellation are exact, but the universal estimate $\mathrm{SCI}_1$ is false because it would imply the disproved fixed selected-channel entropy bound. \\
Negative-energy centrifugal action & \statusproved & The compressed resolvent decreases in $\lambda$ and its total $D_L/r^2$ action is bounded by the charged star column; a monotone Stieltjes-density example proves that this does not by itself control boundary entropy. \\
Universal radially scaled angular heat endpoint & \statusfailed & The heat-flow identity and budget $\frac12\|V^\circ\|_{X_3}^2$ are exact, but the universal product-Fisher estimate $\mathrm{HFI}_1$ is false by the fixed Riesz--Talbot potential. \\
Raw directional reflection/bilinear density action & \statusdiagnostic & It is sufficient but over-strong: an ideal scalar shell has arbitrarily large evanescent reflection while leaving the observation-plane density unchanged. \\
Cutoff-compatible refocusing-work entropy & \statusfailed & The signed reduction remains exact, but a universal bound is contradicted by the reciprocal Riesz--Talbot cascade. \\
Arbitrary matrix channel rotator/barrier & \statusfailed & The low-column cost tends to zero while the positive Grassmannian entropy diverges; the cutoff-edge scalar analogue is proved below. \\
Exact cutoff-edge scalar rotator/barrier & \statusproved & Every real $2\times2$ block has a reducing scalar lift; a pure degree-$2L$ harmonic is an exact high-only barrier of cost $\asymp L^{-1/2}$. \\
Scalar $\mathrm{PSR}_1$ & \statusfailed & The exact cutoff-edge lift gives bounded charged cost and fixed-band quotient entropy tending to infinity. \\
Full-space scalar converter & \statusproved & A zonal quarter-Talbot Riesz product concentrates the selected state exponentially, a localized reciprocal reflector charges exponentially little, and the reverse train refocuses it. \\
Physically scaled scalar nonlinear Plancherel $\mathrm{SSNP}_\ell$ & \statusopen & A uniform fixed-band bound would exclude the frozen converter before restoring radial second-order errors. \\
Pure-detuning nonlinear Plancherel budget & \statusproved & Fixed-band transition probability has an exact nonlinear $L^2$ cost. \\
Abstract scaled-ensemble extension & \statusfailed & A first-order matrix rotator gives uniform band conversion at vanishing charged cost. \\
Scalar ground-state adiabatic swap & \statusproved & Positivity fixes the ground branch, and a small endpoint potential at $R=O(\ell)$ has vanishing overlap with $h_\ell$. \\
Monotone chirped scalar passage & \statusproved & Uniform adiabaticity across a fixed energy band forces a divergent charged cost. \\
Rotating-frame ground ladder & \statusproved & At a free endpoint its ground angular momentum depends on $s$; no fixed rotation rate selects one common $\ell$ across a nondegenerate band. \\
Gap-cancelled finite-depth composite passage & \statusproved & Vanishing cost forces angular frequency $\gg\ell^{3/2}$; fixed-depth and $o(\sqrt\ell)$-depth bounded-gradient designs cannot convert. \\
Exact scalar star-frame reduction & \statusproved & The charged column becomes an off-diagonal vector of exactly the same norm; all uncharged danger is isolated in its energy-dependent high-sector rotation. \\
Scalar two-plane mixer--barrier duality & \statusproved & The mixer leakage is $S b^2G_h$, the cheapest high-only barrier costs $Q^2/(S G_h)$, and their product is exactly $b^2Q^2$. \\
Charged pulse-area and scalar-roughness toll & \statusproved & Gap cancellation forces $\mathcal E_\ell\gtrsim\ell^{-1/4}$ and average angular gradient $\gtrsim\ell/\mathcal E_\ell^2$; target-scale Lipschitz controls have a positive cost floor. \\
Born high-band recoil & \statusproved & A direct fixed-band launch to degrees $\geq L$ costs at least order $L^2/R^2$ at first order; only a nonlinear roughness bootstrap remains. \\
Nonlinear charged spectral gate & \statusproved & Even with arbitrary high--high amplification retained, first escape from degrees $\leq N$ has angle at most $s(N+1)\int\|M_ve_0\|$; a small charged component acts strongly only on a preconcentrated state. \\
Density-concentration action & \statusproved & Conversion to $h_\ell$ forces $\int|\psi|^2|\nabla\psi|^2\gtrsim R^2\sqrt\ell$; scalar phase masks leave density unchanged and cannot supply this action directly. \\
Wrapped-phase gate additivity & \statusfailed & A bounded azimuthal sawtooth injects uniform fixed-band mass into degrees $\geq L$ at vanishing charged quadratic cost when centrifugal transport is deleted. \\
Radial energy/BV obstruction & \statusproved & With centrifugal transport restored, angular current is at most $\sqrt{\Gamma_v}\log(R_2/R_1)$; stationary and radially low-variation weak slabs cannot convert. \\
Radial superoscillation toll & \statusproved & Gap cancellation forces radial $L^\infty$ variation, and hence bounded-amplitude layer count, $\gtrsim\ell/\mathcal E_\ell^2$. \\
Raw radial variation as a cost proxy & \statusfailed & Tiny rapid radial oscillations make $\Gamma_v\to\infty$ while their charged cost, angular gradient norm, and full propagator effect vanish. \\
Effective radial-work bridge & \statusproved & Conversion forces state-correlated work of order one and $\mathfrak R_v\mathfrak D_s\gtrsim R^2$; for angular bandwidth $L$, $\mathcal E\mathfrak D_s\gtrsim R^2/L^2$. \\
Moving-support effective-work bridge & \statusopen & The remaining profile is doubly super-rough; no charged-cost estimate yet controls its state-correlated work with broadband reciprocal refocusing. \\
Frozen resonant-drive test & \statusdiagnostic & Single-energy transfer is cheap, but the Born and pure-detuning nonlinear models impose fixed-band costs. \\
Exterior/full-space bridge after inserting a free collar & \statusknown & For $q_d=\lceil d/2\rceil$ the resolvent-power difference is trace class, so a.c. multiplicity is preserved. \\
\bottomrule
\end{longtable}

The table is only an index.  The hypotheses, equations, and logical content of every row are recorded next so that none of the status labels has to be accepted on trust.

\subsection{Controlled truncations: strong resolvent and measure convergence}

On the exterior space set
\begin{equation}
 H_0=-\partial_r^2+r^{-2}A_d,
 \qquad
 \widehat H_{R,L}=H_0+P_LV_RP_L,
 \qquad
 V_R=\mathbf1_{[1,R]}M_V .
 \label{eq:audit-truncated-H}
\end{equation}
Here $P_L\to\Id$ strongly, $V_R\to M_V$ strongly, and all three multiplication operators have norm at most $\|V\|_\infty$.  Hence, for every $f$,
\begin{align}
 (P_LV_RP_L-M_V)f
 &=P_LV_R(P_L-\Id)f+P_L(V_R-M_V)f+(P_L-\Id)M_Vf
 \longrightarrow0.
 \label{eq:audit-potential-strong}
\end{align}
If $H_{\rm ext}=H_0+M_V$ and $z\in\C\setminus\R$, the second resolvent identity gives
\begin{align}
 &(\widehat H_{R,L}-z)^{-1}-(H_{\rm ext}-z)^{-1} \notag\\
 &\quad=-(\widehat H_{R,L}-z)^{-1}
 (P_LV_RP_L-M_V)(H_{\rm ext}-z)^{-1},
 \label{eq:audit-resolvent-identity}
\end{align}
and the two resolvents have norm at most $|\Im z|^{-1}$.  Equations~\eqref{eq:audit-potential-strong}--\eqref{eq:audit-resolvent-identity} prove strong resolvent convergence.  For a fixed finite probe $\Phi_N$ this yields
\begin{equation}
 M_N^{R,L}(z)
 :=\Phi_N^*(\widehat H_{R,L}-z)^{-1}\Phi_N
 \longrightarrow
 \Phi_N^*(H_{\rm ext}-z)^{-1}\Phi_N=:M_N(z)
 \label{eq:audit-herglotz-convergence}
\end{equation}
locally uniformly on $\C_+$.  Equivalently, the representing matrix measures converge vaguely:
\begin{equation}
 \int f(E)\dd\mu_N^{R,L}(E)
 \longrightarrow
 \int f(E)\dd\mu_N(E),
 \qquad f\in C_c(\R).
 \label{eq:audit-vague}
\end{equation}
This step alone says nothing about preservation of absolutely continuous density; that requires the entropy bound below.

\subsection{Finite-model Jost/Gram formula and local entropy}

For fixed $R,L$, let $m=m_L$ and let $S$ and $F$ solve
\begin{align}
 -S''+Q_{R,L}S&=k^2S,
 &S(1,k)&=0,&S'(1,k)&=\Id_m,\notag\\
 -F''+Q_{R,L}F&=k^2F,
 &F(r,k)&\sim F_{0,L}(r,k)\quad(r\to\infty),
 \label{eq:audit-fundamental-solutions}
\end{align}
where $Q_{R,L}=D_L/r^2+P_LV_RP_L$, and put $J(k)=F(1,k)$.  With the Wronskian of $F_{0,L}$ normalized to $2ik\Id_m$, the Dirichlet Weyl matrix satisfies
\begin{equation}
 \Im m_{R,L}(k^2+i0)
 =kJ(k)^{-*}J(k)^{-1},
 \qquad
 h_{R,L}(k^2)=\frac{k}{\pi}J(k)^{-*}J(k)^{-1}.
 \label{eq:audit-jost-density}
\end{equation}
For a compact probe $\Phi_N$ define, with a transpose rather than an adjoint in the analytic variable,
\begin{equation}
 T_N(k)=\int_1^\infty S(r,k)^T\Phi_N(r)\dd r,
 \qquad
 \mathcal A_N(k)=J(k)^{-1}T_N(k).
 \label{eq:audit-regular-transform}
\end{equation}
Then the Gram formula and Cauchy--Binet identity give
\begin{align}
 W_N^{R,L}(k^2)
 &=\frac{k}{\pi}\mathcal A_N(k)^*\mathcal A_N(k),
 \label{eq:audit-probe-density}\\
 \det W_N^{R,L}(k^2)
 &=\left(\frac{k}{\pi}\right)^N
   \left\|\wedgeN\mathcal A_N(k)\right\|^2
 =\left(\frac{k}{\pi}\right)^N
   \sum_{|S|=N}|\det\mathcal A_{N,S}(k)|^2.
 \label{eq:audit-gram-determinant}
\end{align}
The large-imaginary-$k$ Volterra asymptotic makes $\wedgeN\mathcal A_N$ nonzero as an analytic vector.  Its real zeros on a compact positive interval are therefore isolated and of finite order.  Thus
\begin{equation}
 \det W_N^{R,L}(E)>0\quad\text{for a.e. }E>0,
 \qquad
 \int_I\log^-\det W_N^{R,L}(E)\dd E<\infty
 \label{eq:audit-finite-entropy}
\end{equation}
for each fixed model.  No constant uniform in $R$ or $L$ follows from this argument.

\subsection{Entropy lower semicontinuity}

Write $\mu_n=W_n\dd E+\mu_n^s$, $\mu=W\dd E+\mu^s$, and suppose $\mu_n\to\mu$ vaguely near $I=[a,b]$.  On Hermitian matrices define
\begin{equation}
 \varphi(A)=
 \begin{cases}
  \max\{0,-\log\det A\},&A>0,\\
  +\infty,&A\not>0.
 \end{cases}
 \label{eq:audit-entropy-integrand}
\end{equation}
The map $\varphi$ is convex, lower semicontinuous, and decreasing in the Loewner order.  If $\rho_\eps$ is a positive approximate identity and $I_\delta=[a+\delta,b-\delta]$, positivity of $\rho_\eps*\mu_n^s$ and Jensen's inequality imply
\begin{align}
 \int_{I_\delta}\varphi((\rho_\eps*\mu_n)(x))\dd x
 &\leq
 \int_I\varphi(W_n(E))\dd E.
 \label{eq:audit-mollified-entropy}
\end{align}
Vague convergence passes the convolved matrices to the limit.  Fatou's lemma, differentiation of finite matrix measures as $\eps\downarrow0$, and then $\delta\downarrow0$ give
\begin{equation}
 \int_I\log^-\det W(E)\dd E
 \leq\liminf_{n\to\infty}
 \int_I\log^-\det W_n(E)\dd E.
 \label{eq:audit-lsc}
\end{equation}
The inner interval is indispensable: vague convergence does not prevent mass from approaching an endpoint.

\subsection{The rank and countability finish}

If the limiting entropy is finite, then a positive-measure set on which $\det W_N(E)=0$ is impossible.  Hence
\begin{equation}
 \int_I\log^-\det W_N<\infty
 \Longrightarrow
 \rank W_N(E)=N\quad\text{for a.e. }E\in I.
 \label{eq:audit-rank-density}
\end{equation}
In the direct-integral representation of the absolutely continuous part, a density generated by $N$ probes has rank no larger than the fiber multiplicity.  Therefore
\begin{equation}
 \rank W_N(E)=N
 \Longrightarrow
 m_{\ac}(E;H)\geq N
 \quad\text{for a.e. }E\in I.
 \label{eq:audit-rank-multiplicity}
\end{equation}
Intersecting the full-measure sets over $N\in\N$ and compact positive intervals with rational endpoints yields
\begin{equation}
 m_{\ac}(E;H)=\infty
 \quad\text{for a.e. }E>0.
 \label{eq:audit-countability}
\end{equation}
This finish is rigorous but conditional on a uniform finite-model entropy estimate and a valid exterior/full-space bridge.

\subsection{Why leakage and Herglotz positivity alone are false}

The first exact counterexample is the fiberwise projection
\begin{equation}
 \mathcal H(x)=
 \begin{pmatrix}
 x&\sqrt{x(1-x)}\\
 \sqrt{x(1-x)}&1-x
 \end{pmatrix},
 \qquad \mathcal H(x)^2=\mathcal H(x),
 \label{eq:fiber-counterexample}
\end{equation}
on $L^2(0,1)\oplus L^2(0,1)$.  Its off-diagonal leakage has finite cost
\begin{equation}
 \int_0^1x(1-x)\dd x=\frac16,
 \end{equation}
but the first coordinate has spectral measure
\begin{equation}
 d\mu_{e_1}=\tfrac12\delta_0+\tfrac12\delta_1.
\end{equation}
Thus its density vanishes a.e. on $(0,1)$.  Nevertheless the block self-energy obeys
\begin{equation}
 \Im\!\left[B^*(H_{QQ}-z)^{-1}B\right]\succeq0,
 \qquad z\in\C_+.
 \label{eq:audit-herglotz-insufficient}
\end{equation}
Consequently neither finite Hilbert--Schmidt leakage nor the Herglotz sign implies finite-channel entropy in abstract block theory.  The family obtained by multiplying the off-diagonal entries by $\alpha$ has varying eigenbranches for generic $\alpha$ but collapses at $\alpha=1$.  Thus almost-every-coupling spectral averaging cannot simply be evaluated at the physical coupling.

\subsection{Why truncated leakage is unsafe}

On $\mathbb S^1$, with
\begin{equation}
 v=a+c\cos(\ell\omega)+d\cos(2\ell\omega),
 \label{eq:alias-v}
\end{equation}
the retained block on $\operatorname{span}\{Y_0,C_\ell\}$ is
\begin{equation}
 \begin{pmatrix}
 a&c/\sqrt2\\ c/\sqrt2&a+d/2
 \end{pmatrix}.
 \label{eq:alias-block}
\end{equation}
Writing $b=c/\sqrt2$, the truncated column sees $a^2+b^2$, whereas the full column sees
\begin{equation}
 \|Q M_vY_0\|^2=b^2+\frac{d^2}{2}.
 \label{eq:alias-cost}
\end{equation}
The choice $a=c=0$, $d=2Q$ therefore creates a $+Q$ retained-channel barrier at zero truncated-column cost.  Every proposed estimate below uses the full complement or the full scalar cost.

\subsection{The corrected full-leakage endpoint}

With $Q_N=\Id-P_N$, orthogonality gives the exact identity
\begin{align}
 \|P_NM_VP_N\|_{\HS}^2+\|Q_NM_VP_N\|_{\HS}^2
 &=\|M_VP_N\|_{\HS}^2
 =\sum_{j=1}^N\|VY_j\|_2^2.
 \label{eq:audit-full-leakage}
\end{align}
If $Y_1=|\Sd|^{-1/2}$, then
\begin{equation}
 |\Sd|^{-1}\|V(r,\cdot)\|_2^2
 \leq\|M_{V(r)}P_N\|_{\HS}^2
 \leq\left(\sum_{j=1}^N\|Y_j\|_\infty^2\right)
 \|V(r,\cdot)\|_2^2.
 \label{eq:audit-full-scalar-equivalence}
\end{equation}
The safe, still open endpoint is therefore
\begin{equation}
 \sup_{R,L}\int_I\log^-\det W_N^{R,L}(E)\dd E
 \leq C_{I,N,\Phi}
 \left(1+\int_1^\infty\|M_{V(r)}P_N\|_{\HS}^2\dd r\right).
 \label{eq:audit-corrected-fcen}
\end{equation}
Equation~\eqref{eq:audit-full-scalar-equivalence} proves that the corrected cost is exactly of the strength of Simon's hypothesis.  It does not prove~\eqref{eq:audit-corrected-fcen}.

\subsection{The arbitrary matrix rotator/barrier obstruction is exact}

In the two-channel matrix model, let
\begin{equation}
 V_{\rm rot}(r)=\delta_\Lambda u_2(r)u_2(r)^*,
 \qquad
 u_2=(-\sin\theta,\cos\theta),
 \qquad
 \delta_\Lambda=\Lambda^{-2/3},
 \label{eq:audit-rotator}
\end{equation}
on an interval of length $\Lambda$, with
$\theta(r)=\pi r/(2\Lambda)$.  Its low-column cost is exactly
\begin{equation}
 \int_{\rm rot}\|V_{\rm rot}(r)e_1\|^2\dd r
 =\frac12\Lambda^{-1/3}.
 \label{eq:audit-rotator-cost}
\end{equation}
The rotating-frame equation has constant coefficients and gives a full modal
quarter-turn with error $O_K(\Lambda^{-1/3})$.  Follow it by
$\diag(0,Q_0)$, which has zero low column, and choose a shorter logarithmic
barrier so that its amplification grows like $\Lambda^\beta$ with
$0<\beta<1/3$ small enough.  Lemma~\ref{lem:rapid-phase-good-set} shows that
forward/backward cancellation fails on a fixed positive fraction of every
compact energy interval.  Theorem~\ref{thm:matrix-rotator-counterexample}
therefore proves
\begin{equation}
 \int_K\log^+\Gamma_P(F_\Lambda(0,k))\dd k
 \geq c_K\log\Lambda-C_K,
 \qquad
 \int\|\mathcal V_\Lambda(r)P\|_{\HS}^2\dd r\longrightarrow0.
 \label{eq:audit-matrix-rotator-failure}
\end{equation}

This rejects the arbitrary matrix low-column endpoint, not yet the scalar
angular endpoint~\eqref{eq:audit-corrected-fcen}.  However, scalar compression
does not rescue the estimate at the algebraic level.  Proposition
~\ref{prop:minimum-scalar-block-lift} gives an exact scalar lift of every real
$2\times2$ compressed block.  For concentrated highest-weight harmonics on
$\mathbb S^2$, the extra full constant-column cost tends to zero.  More
decisively, Proposition~\ref{prop:thin-belt-barrier} constructs a uniformly
bounded scalar high-only barrier with
\begin{equation}
 \|M_{b_{\ell,Q}}e_0\|_2^2=O_Q(\ell^{-\alpha}),
 \qquad
 \|(I-P_{\{e_0,h_\ell\}})M_{b_{\ell,Q}}h_\ell\|_2^2
 \leq C_Qe^{-c\ell^{1-2\alpha}}.
 \label{eq:audit-scalar-belt-barrier}
\end{equation}
Thus the barrier stage can be scalarized at vanishing full low-column cost.
The unresolved obstruction has narrowed to the converter: can a bounded real
scalar evolution rotate the constant harmonic into $h_\ell$ with sufficiently
small error and quadratic cost when the entire angular space and the
centrifugal flow are retained?

\subsection{The exterior/full-space bridge as a separate theorem}

Let $\Sigma=\{|x|=a\}$ lie inside a collar on which the tail potential vanishes, and let
\begin{equation}
 H_{\rm dec}=H_{\rm int}^D\oplus H_{\rm ext}^D.
 \label{eq:audit-decoupled}
\end{equation}
The exact sufficient bridge is the existence of an integer $q$ and $c>\|V\|_\infty$ such that
\begin{equation}
 (H_{\rm tail}+c)^{-q}-(H_{\rm dec}+c)^{-q}\in\TC.
 \label{eq:audit-resolvent-power-bridge}
\end{equation}
Birman--Kuroda then gives
\begin{equation}
 (H_{\rm tail})_{\ac}\simeq(H_{\rm dec})_{\ac}
 \simeq(H_{\rm ext}^D)_{\ac},
 \label{eq:audit-ac-bridge}
\end{equation}
because $H_{\rm int}^D$ has compact resolvent.  The needed theorem is available with exactly local, rather than global, regularity: Lotoreichik~\cite[Theorem~4.17 and Corollary~4.19]{Lotoreichik2013} proves the weak Schatten estimate for the coupled/Dirichlet-decoupled pair when the bounded potential has finitely many bounded derivatives only near the cutting hypersurface.  A free collar supplies derivatives of every order there, regardless of roughness farther out.  The precise power and implication are recorded in Proposition~\ref{prop:free-collar} below.

\section{Three exact reductions}

\subsection{A compact core can be discarded}

\begin{proposition}[Small-tail reduction]\label{prop:small-tail}
Fix a compact energy interval $I\Subset(0,\infty)$ and $N\in\N$.  Suppose there is $\eps_{I,N}>0$ with the following property: every bounded real exterior potential $W$ satisfying
\begin{equation}
 \int_R^\infty\|W(r,\cdot)\|_2^2\dd r<\eps_{I,N}^2
 \label{eq:small-tail-hyp}
\end{equation}
has absolutely continuous multiplicity at least $N$ for a.e. $E\in I$.  Then the same conclusion holds for every bounded real $V$ satisfying~\eqref{eq:polar-cost}.
\end{proposition}

\begin{proof}
Choose $R$ so large that $V_{\mathrm{tail}}=\mathbf1_{\{r>R\}}V$ satisfies~\eqref{eq:small-tail-hyp}.  The difference $V-V_{\mathrm{tail}}$ is bounded and compactly supported.  Standard short-range scattering, equivalently the Birman--Kuroda theorem applied to a sufficiently high resolvent power, gives unitary equivalence of the absolutely continuous parts of $-\Delta+V$ and $-\Delta+V_{\mathrm{tail}}$.  Absolutely continuous multiplicity is therefore preserved.  Applying the assumed small-tail statement to $V_{\mathrm{tail}}$ proves the claim.
\end{proof}

\begin{remark}
Proposition~\ref{prop:small-tail} removes a major distraction.  The endpoint $\alpha=1$ need not be reached by continuation from small coupling; after removing a compact core, the physical potential itself has arbitrarily small $X_d$ tail norm.  A proof may therefore focus on a perturbative theorem in the correct norm while allowing $\|V\|_\infty$ to remain large on sparse sets.
\end{remark}

\begin{proposition}[Free-collar Dirichlet decoupling]\label{prop:free-collar}
Let $V_{\mathrm{tail}}\in L^\infty(\R^d;\R)$ vanish on a neighbourhood of the sphere $\Sigma=\{|x|=a\}$.  Let
\begin{equation}
 H_{\mathrm{tail}}=-\Delta+V_{\mathrm{tail}},
 \qquad
 H_{\rm dec}=H_{\rm int}^D\oplus H_{\rm ext}^D
 \label{eq:free-collar-operators}
\end{equation}
be the coupled and Dirichlet-decoupled realizations of the same differential expression.  Set
\begin{equation}
 q_d=\left\lceil\frac d2\right\rceil,
 \label{eq:free-collar-power}
\end{equation}
then, for every $c>\|V_{\mathrm{tail}}\|_\infty$,
\begin{equation}
 (H_{\mathrm{tail}}+c)^{-q_d}-(H_{\rm dec}+c)^{-q_d}
 \in\mathfrak S_{\frac{d-1}{2q_d},\infty}\subset\TC.
 \label{eq:free-collar-bridge}
\end{equation}
Consequently the absolutely continuous parts of $H_{\mathrm{tail}}$ and $H_{\rm ext}^D$ are unitarily equivalent, including multiplicity.
\end{proposition}

\begin{proof}
The collar hypothesis implies
\begin{equation}
 V_{\mathrm{tail}}\in W_\Sigma^{m,\infty}(\R^d)
 \qquad\text{for every }m\in\N_0,
 \label{eq:collar-local-regularity}
\end{equation}
where $W_\Sigma^{m,\infty}$ means bounded Sobolev regularity in some neighbourhood of $\Sigma$ only.  Apply Lotoreichik~\cite[Theorem~4.17]{Lotoreichik2013} with $\ell=m=q_d$ and $z=-c$.  The choice $c>\|V_{\mathrm{tail}}\|_\infty$ puts $-c$ in the resolvent sets of the coupled, Dirichlet-decoupled, and auxiliary Neumann-decoupled operators.  The theorem gives
\begin{equation}
 (H_{\mathrm{tail}}+c)^{-q_d}-(H_{\rm dec}+c)^{-q_d}
 \in\mathfrak S_{\frac{d-1}{2q_d},\infty}.
 \label{eq:collar-weak-schatten}
\end{equation}
Because $(d-1)/(2q_d)<1$, its singular values are
\begin{equation}
 O\!\left(j^{-2q_d/(d-1)}\right),
 \end{equation}
with summable exponent.  Hence~\eqref{eq:collar-weak-schatten} is trace class and proves~\eqref{eq:free-collar-bridge}.

Kato--Rosenblum applied to the bounded self-adjoint resolvent powers makes their absolutely continuous parts unitarily equivalent.  The function $t\mapsto(t+c)^{-q_d}$ is one-to-one on the spectra of these semibounded operators, so the invariance principle transfers the equivalence back to $H_{\mathrm{tail}}$ and $H_{\rm dec}$.  The interior Dirichlet summand has compact resolvent, leaving precisely $(H_{\rm ext}^D)_{\ac}$.
\end{proof}

The free collar also permits the probes of Section~\ref{sec:probes} to be placed where their regular transform is an explicit low-channel diagonal factor.

\subsection{The full column is the scalar cost}

Let $P_N$ be a finite-rank projection commuting with $A_d$, with real orthonormal eigenharmonics $Y_1,\ldots,Y_N$, and assume
\begin{equation}
 Y_1=|\Sd|^{-1/2}.
 \label{eq:constant-harmonic}
\end{equation}
Put $Q_N=\Id-P_N$.

\begin{proposition}[Exact full-column identity]\label{prop:full-column}
For a.e. $r$,
\begin{align}
 &\|P_NM_VP_N\|_{\HS}^2+\|Q_NM_VP_N\|_{\HS}^2
   =\|M_VP_N\|_{\HS}^2 \notag\\
 &\hspace{35mm}=\sum_{j=1}^N\|V(r,\cdot)Y_j\|_{L^2(\Sd)}^2.
 \label{eq:column-exact}
\end{align}
Consequently,
\begin{equation}
 \frac{1}{|\Sd|}\|V(r,\cdot)\|_2^2
 \leq \|M_VP_N\|_{\HS}^2
 \leq C_N\|V(r,\cdot)\|_2^2,
 \qquad
 C_N=\sum_{j=1}^N\|Y_j\|_\infty^2.
 \label{eq:column-comparable}
\end{equation}
\end{proposition}

\begin{proof}
The first equality is Pythagoras in the orthogonal decomposition $\K=P_N\K\oplus Q_N\K$.  Expanding the Hilbert--Schmidt norm in the basis $Y_j$ gives the second equality.  The upper bound follows from $|Y_j(\omega)|\leq\|Y_j\|_\infty$.  The lower bound is the $j=1$ term and~\eqref{eq:constant-harmonic}.
\end{proof}

Define the finite-channel energy
\begin{equation}
 \cE_N(V)
 :=\int_1^\infty\|M_{V(r)}P_N\|_{\HS}^2\dd r.
 \label{eq:EN}
\end{equation}
For fixed $N$, Proposition~\ref{prop:full-column} says
\begin{equation}
 |\Sd|^{-1}\|V\|_{X_d,\mathrm{ext}}^2
 \leq\cE_N(V)\leq C_N\|V\|_{X_d,\mathrm{ext}}^2.
 \label{eq:EN-equivalence}
\end{equation}
Thus the corrected full leakage does not introduce a stronger hypothesis.  It is simply a finite-channel representation of the original scalar angular cost.

\subsection{Finite models and normalization}

Let $P_L$ be the complete spectral projection of $A_d$ through angular degree $L$, let $m_L=\rank P_L$, and assume $P_N\leq P_L$.  With compact radial truncation $V_R=\mathbf1_{[1,R]}V$, define
\begin{equation}
 H_{R,L}=-\partial_r^2+\frac{D_L}{r^2}+P_LM_{V_R}P_L
 \label{eq:finite-H}
\end{equation}
on $L^2((1,\infty);\C^{m_L})$, with Dirichlet boundary condition at $r=1$.  The Bessel diagonal $D_L/r^2$ is treated as background, not as perturbation.

Let $F(r,k)$ be the outgoing Jost matrix with the unambiguous normalization
\begin{equation}
 F(r,k)=F_{0,L}(r,k)(\Id+o(1)),
 \qquad r\to\infty,
 \label{eq:jost-normalization}
\end{equation}
where $F_{0,L}$ is the outgoing diagonal Bessel matrix and is normalized so that its real-axis Wronskian is $2ik\Id$.  Write
\begin{equation}
 J_{R,L}(k)=F(1,k).
 \label{eq:J}
\end{equation}
After the harmless channelwise Bessel normalization, the canonical Dirichlet density is
\begin{equation}
 h_{R,L}(k^2)=\frac{k}{\pi}J_{R,L}(k)^{-*}J_{R,L}(k)^{-1},
 \qquad k>0.
 \label{eq:canonical-density}
\end{equation}
All later formulas use this normalization.

\section{The Pl\"ucker--Schur reduction}

This section gives the main new algebraic reduction.

\begin{theorem}[Exact low-channel Pl\"ucker formula]\label{thm:plucker}
Let $J\in\C^{m\times m}$ be invertible, let $P$ be an orthogonal projection of rank $N$, and let $Q=\Id-P$.  Put
\begin{equation}
 h=\frac{k}{\pi}J^{-*}J^{-1},
 \qquad h_P=PhP\big|_{P\C^m}.
\end{equation}
Then
\begin{align}
 \det h_P
 &=\left(\frac{k}{\pi}\right)^N
   \frac{\det\!\left(QJJ^*Q\big|_{Q\C^m}\right)}{|\det J|^2}
 \label{eq:jacobi-ratio}\\
 &=\left(\frac{k}{\pi}\right)^N
   \frac{\left\|\bigwedge^{m-N}(QJ)\right\|^2}{|\det J|^2}.
 \label{eq:plucker-ratio}
\end{align}
Here the exterior norm is the Euclidean norm of the vector of all $(m-N)\times(m-N)$ minors of $QJ$.
\end{theorem}

\begin{proof}
Set $G=JJ^*>0$.  Since $J^{-*}J^{-1}=G^{-1}$, Jacobi's complementary-minor identity gives
\begin{equation}
 \det\!\left(PG^{-1}P\big|_{P\C^m}\right)
 =\frac{\det\!\left(QGQ\big|_{Q\C^m}\right)}{\det G}.
 \label{eq:jacobi-proof}
\end{equation}
Now $\det G=|\det J|^2$.  Moreover, $QGQ=(QJ)(QJ)^*$ on $Q\C^m$.  The Cauchy--Binet identity gives
\begin{equation}
 \det((QJ)(QJ)^*)
 =\sum_{|S|=m-N}|\det (QJ)_S|^2
 =\left\|\bigwedge^{m-N}(QJ)\right\|^2.
\end{equation}
Multiplication by $(k/\pi)^N$ proves both formulas.
\end{proof}

Define the Grassmannian amplification factor
\begin{equation}
 \Gamma_{P;R,L}(k)
 :=\frac{|\det J_{R,L}(k)|}
 {\|\bigwedge^{m_L-N}(QJ_{R,L}(k))\|}.
 \label{eq:Gamma}
\end{equation}

\begin{proposition}[Dual inverse-frame Pl\"ucker identity]\label{prop:dual-plucker}
For every $J\in\mathrm{GL}(m,\C)$ and every coordinate projection $P$ of
rank $N$,
\begin{equation}
 \frac{\|\bigwedge^{m-N}(QJ)\|}{|\det J|}
 =\|\wedgeN(J^{-1}P)\|.
 \label{eq:dual-plucker}
\end{equation}
Consequently
\begin{equation}
 \Gamma_P(J)
 =\|\wedgeN(J^{-1}P)\|^{-1}
 =\det\!\left(PJ^{-*}J^{-1}P\big|_{P\C^m}\right)^{-1/2}.
 \label{eq:Gamma-inverse-frame}
\end{equation}
\end{proposition}

\begin{proof}
For every $N$-element row set $T$, Jacobi's complementary-minor identity gives
\begin{equation}
 \det(J^{-1})_{T,P}
 =\pm\frac{\det J_{Q,T^c}}{\det J}.
\end{equation}
Square the moduli and sum over $T$.  Cauchy--Binet identifies the left
collection with the coordinates of $\wedgeN(J^{-1}P)$ and the right
collection with those of $\bigwedge^{m-N}(QJ)/\det J$.  The Gram determinant
formula gives~\eqref{eq:Gamma-inverse-frame}.
\end{proof}

Thus the natural analytic exterior vector is
\begin{equation}
 \mathfrak g_P(k)=\wedgeN\!\left(J(k)^{-1}P\right),
 \qquad
 \log\Gamma_P(k)=-\log\|\mathfrak g_P(k)\|.
 \label{eq:dual-exterior-frame}
\end{equation}
It has no distinguished-minor denominator and works for every fixed $N$.

\begin{proposition}[Exact differential Grassmannian flux]\label{prop:differential-flux}
Let $\mathcal J(r)\in\mathrm{GL}(m,\C)$ solve
\begin{equation}
 \mathcal J'=\mathcal A\mathcal J,
 \qquad
 \mathcal A=
 \begin{pmatrix}
  \mathcal A_{PP}&\mathcal A_{PQ}\\
  \mathcal A_{QP}&\mathcal A_{QQ}
 \end{pmatrix}.
 \label{eq:left-multiplicative-cocycle}
\end{equation}
Assume $N<m$; when $N=m$, omit the $Q$ rows and take $C=0$.
Put
\begin{align}
 X&=P\mathcal J,&Y&=Q\mathcal J,&S&=YY^*,\label{eq:quotient-rows}\\
 C&=XY^*S^{-1},&E&=X-CY,&\Sigma&=EE^*.
 \label{eq:quotient-regression}
\end{align}
Then
\begin{align}
 EY^*&=0,
 &\Gamma_P(\mathcal J)^2&=\det\Sigma,
 \label{eq:residual-Gram-volume}\\
 \Sigma'
 &=\mathcal B\Sigma+\Sigma\mathcal B^*,
 &\mathcal B&=\mathcal A_{PP}-C\mathcal A_{QP},
 \label{eq:residual-Gram-evolution}\\
 \frac{\dd}{\dd r}\log\Gamma_P(\mathcal J(r))
 &=\Re\tr\!\left(\mathcal A_{PP}-C\mathcal A_{QP}\right).
 \label{eq:differential-Grassmannian-flux}
\end{align}
In particular, neither $\mathcal A_{QQ}$ nor $\mathcal A_{PQ}$ occurs
directly in the logarithmic increment.  The displayed generator contains only
the charged column $\mathcal AP$; the high--high evolution is absorbed into
the dressed Grassmannian state $C$.
\end{proposition}

\begin{proof}
The definition of $C$ gives $EY^*=0$.  The Schur complement of $S=YY^*$ in
$\mathcal J\mathcal J^*$ is $EE^*=\Sigma$, so
\begin{equation}
 |\det\mathcal J|^2=\det S\det\Sigma.
\end{equation}
This proves~\eqref{eq:residual-Gram-volume}.  From
\eqref{eq:left-multiplicative-cocycle},
\begin{equation}
 X'=\mathcal A_{PP}X+\mathcal A_{PQ}Y,
 \qquad
 Y'=\mathcal A_{QP}X+\mathcal A_{QQ}Y.
\end{equation}
Substitution of $X=CY+E$ into $E'=X'-C'Y-CY'$ gives
\begin{equation}
 E'=(\mathcal A_{PP}-C\mathcal A_{QP})E+\mathcal K Y
\end{equation}
for a matrix $\mathcal K$.  Since $EY^*=0$, the $\mathcal K Y$ terms vanish
in $E'E^*+EE'^*$, proving~\eqref{eq:residual-Gram-evolution}.  Taking one
half of the logarithmic determinant derivative proves
\eqref{eq:differential-Grassmannian-flux}.
\end{proof}

\begin{proposition}[No low-column estimate for unrestricted cocycles]
\label{prop:cocycle-low-column-no-go}
Pointwise cancellation of the high--high generator in
\eqref{eq:differential-Grassmannian-flux} does not imply a bound for the
Grassmannian increment in terms of $\int\|\mathcal AP\|_{\HS}^2$ alone.  More
precisely, for $m=2$, $N=1$, and $P=e_1e_1^*$ there are piecewise constant
left cocycles with $\|\mathcal A(r)\|\leq1$ such that
\begin{equation}
 \int_0^{H+L}\|\mathcal A(r)P\|_{\HS}^2\dd r
 =\frac{\pi^2}{4L},
 \qquad
 \log\Gamma_P(\mathcal J(H+L))=H.
 \label{eq:cocycle-no-go-cost-growth}
\end{equation}
Consequently there is no uniform estimate of the terminal increment by a
constant plus this low-column cost.
\end{proposition}

\begin{proof}
Let
\begin{equation}
 \mathsf K=\begin{pmatrix}0&-1\\1&0\end{pmatrix},
 \qquad
 \omega_L=\frac{\pi}{2L},
\end{equation}
and, with $L\geq\pi/2$, set
\begin{equation}
 \mathcal A(r)=
 \begin{cases}
  \diag(0,1),&0<r<H,\\
  \omega_L\mathsf K,&H<r<H+L.
 \end{cases}
 \label{eq:cocycle-no-go-generator}
\end{equation}
For $\mathcal J(0)=\Id$, the first interval gives
$\mathcal J(H)=D_H:=\diag(1,e^H)$.  The second gives the exact quarter-turn
\begin{equation}
 \mathcal J(H+L)=e^{(\pi/2)\mathsf K}D_H
 =\begin{pmatrix}0&-e^H\\1&0\end{pmatrix}.
\end{equation}
The first generator has zero $P$-column, while the second has $P$-column norm
$\omega_L$.  This proves the cost identity in
\eqref{eq:cocycle-no-go-cost-growth}.  Moreover
$\Gamma_P(\Id)=\Gamma_P(D_H)=1$, whereas the final determinant has modulus
$e^H$ and its $Q$ row has norm one.  Hence the entire logarithmic growth $H$
occurs during a rotation whose charged cost tends to zero as $L\to\infty$.
\end{proof}

\begin{remark}[What the counterexample does and does not reject]
The dressed state $C$ in~\eqref{eq:differential-Grassmannian-flux} remembers
the uncharged amplification.  A later charged rotation can therefore make the
product $C\mathcal A_{QP}$ large even when $\mathcal A_{QP}$ is small.  Thus
any successful flux correction must retain quantitative control of the dressed
state.  Proposition~\ref{prop:cocycle-low-column-no-go} concerns arbitrary
matrix cocycles; it does not yet reject the Schr\"odinger endpoint, where
$V_m=P_LM_VP_L$ is a compression of one scalar multiplication operator and the
cost includes its leakage into every omitted angular channel.
\end{remark}

Theorem~\ref{thm:plucker} is equivalently
\begin{equation}
 -\log\det h_P(k^2)
 =-N\log\frac{k}{\pi}+2\log\Gamma_{P;R,L}(k).
 \label{eq:entropy-Gamma}
\end{equation}
Thus high-channel volume is not to be bounded separately: it is divided out exactly.

\begin{corollary}[Sharp sufficient sum rule]\label{cor:PSR}
Let $K=[k_0,k_1]\Subset(0,\infty)$.  If
\begin{equation}
 \sup_{R,L:\,P_N\leq P_L}
 \int_K\log^+\Gamma_{P_N;R,L}(k)\dd k
 \leq C_{K,N}\bigl(1+\cE_N(V)\bigr)
 \label{eq:PSR}
\end{equation}
then the boundary $N$-channel entropy is uniformly bounded on $K^2$.
\end{corollary}

\begin{proof}
On $K$, the term $-N\log(k/\pi)$ is bounded.  Equation~\eqref{eq:entropy-Gamma} and $x^+\geq x$ yield the desired estimate.
\end{proof}

\begin{remark}[Why this is sharper than a distinguished minor]
Replacing the denominator in~\eqref{eq:Gamma} by one chosen minor produces artificial poles whenever that minor vanishes.  The Pl\"ucker norm sums all minors and cannot vanish on the real axis when $J$ is invertible.  The appropriate analytic object is therefore an exterior vector or its scalar outer factor, not the quotient by a preselected high-channel determinant.
\end{remark}

\begin{corollary}[Schur--Jost sufficient target]\label{cor:schur-jost}
Write, at points where the $Q$ block is invertible,
\begin{equation}
 J=\begin{pmatrix}A&B\\ C&D\end{pmatrix},
 \qquad S_P=A-BD^{-1}C.
 \label{eq:J-blocks}
\end{equation}
Then
\begin{equation}
 \det h_P(k^2)
 =\left(\frac{k}{\pi}\right)^N
 \frac{\det\!\left(\Id+C^*D^{-*}D^{-1}C\right)}{|\det S_P|^2}
 \geq\left(\frac{k}{\pi}\right)^N|\det S_P|^{-2}.
 \label{eq:schur-jost-bound}
\end{equation}
Consequently it is sufficient to prove the meromorphic relative estimate
\begin{equation}
 \sup_{R,L}\int_K\log^+|\det S_{P;R,L}(k)|\dd k
 \leq C_{K,N}\bigl(1+\cE_N(V)\bigr),
 \label{eq:schur-jost-target}
\end{equation}
provided the zeros of $D$ are included in a complete pole/Blaschke audit.
\end{corollary}

\begin{proof}
Block inversion gives
\begin{equation}
 J^{-1}P=
 \begin{pmatrix}\Id\\-D^{-1}C\end{pmatrix}S_P^{-1}.
\end{equation}
Taking the Gram determinant of these $N$ columns and using~\eqref{eq:canonical-density} proves~\eqref{eq:schur-jost-bound}.  The last assertion follows on compact $K$ from the bounded free factor.
\end{proof}

\subsection{Stage 1 scalar test: a single normalized radiation column}

Take
\begin{equation}
 P_0=|Y_0\rangle\langle Y_0|,
 \qquad Y_0=|\Sd|^{-1/2},
 \qquad e_0=(1,0,\ldots,0)^T.
 \label{eq:scalar-P0}
\end{equation}
The Pl\"ucker quotient has a particularly simple form which removes every irrelevant high-channel volume:
\begin{proposition}[Exact scalar Pl\"ucker identities]\label{prop:scalar-plucker}
For every finite model and every $k>0$,
\begin{align}
 h_0(k^2)
 &=\frac{k}{\pi}\|J(k)^{-1}e_0\|^2,
 \label{eq:scalar-density-column}\\
 \Gamma_0(k)
 &=\|J(k)^{-1}e_0\|^{-1}.
 \label{eq:scalar-gamma-column}
\end{align}
If $J=\left(\begin{smallmatrix}a&b\\c&D\end{smallmatrix}\right)$ and $D$ is invertible, then
\begin{equation}
 \Gamma_0(k)
 =\frac{|a-bD^{-1}c|}
 {\sqrt{1+\|D^{-1}c\|^2}}
 \leq |a-bD^{-1}c|.
 \label{eq:scalar-schur-exact}
\end{equation}
Moreover, the charged column is exactly
\begin{equation}
 \int_1^\infty\|M_{V(r)}Y_0\|_2^2\dd r
 =\frac1{|\Sd|}\int_1^\infty\|V(r,\cdot)\|_2^2\dd r.
 \label{eq:scalar-exact-cost}
\end{equation}
\end{proposition}

\begin{proof}
Equation~\eqref{eq:scalar-density-column} is the $(0,0)$ entry of~\eqref{eq:canonical-density}.  Comparing it with~\eqref{eq:plucker-ratio} for $N=1$ proves~\eqref{eq:scalar-gamma-column}.  Block inversion gives
\begin{equation}
 J^{-1}e_0
 =(a-bD^{-1}c)^{-1}
 \binom{1}{-D^{-1}c},
\end{equation}
which proves~\eqref{eq:scalar-schur-exact}.  Finally, $M_VY_0=|\Sd|^{-1/2}V$ proves~\eqref{eq:scalar-exact-cost}.
\end{proof}

Thus the compactly supported small-tail problem for $N=1$ is precisely
\begin{equation}
 \sup_{R,L}\int_K
 \log^+\frac1{\|J_{R,L}(k)^{-1}e_0\|}\dd k
 \leq C_K\left(1+\frac1{|\Sd|}\|V\|_{X_d,\rm ext}^2\right).
 \label{eq:scalar-PSR-target}
\end{equation}
Theorem~\ref{thm:scalar-PSR1-counterexample} below disproves this target under
its stated uniform finite-cutoff quantifiers.  The denominator in
~\eqref{eq:scalar-schur-exact} is beneficial and must not be discarded in any
cutoff-compatible replacement.  The following scalar extraction lemma records
a rigorous sufficient route, although its hypotheses cannot hold uniformly for
the counterexample family.

\begin{proposition}[Scalar Hardy extraction]\label{prop:scalar-hardy-extraction}
Let $K\Subset(0,\infty)$ and choose a bounded piecewise $C^{1,\alpha}$ Jordan domain $\Omega_K$ in the first quadrant whose boundary contains $K$ as a straight segment away from the corners.  Let $k_*\in\Omega_K$ and let $\omega_*$ be harmonic measure at $k_*$.  Its density with respect to arclength is then bounded below on $K$.  Let $j_0(k)$ be the free constant-channel Jost factor and set
\begin{equation}
 f_{R,L}(k)=j_0(k)\langle e_0,J_{R,L}(k)^{-1}e_0\rangle.
 \label{eq:scalar-f}
\end{equation}
Suppose uniformly in $R,L$ that
\begin{equation}
 |f_{R,L}(k_*)|\geq c_*>0,
 \qquad
 \int_{\partial\Omega_K}|f_{R,L}(k)|^2\dd\omega_*(k)
 \leq C_*^2.
 \label{eq:scalar-Hardy-data}
\end{equation}
Then
\begin{equation}
 \sup_{R,L}\int_K\log^+\Gamma_{0;R,L}(k)\dd k
 \leq C(K,j_0,c_*,C_*).
 \label{eq:scalar-Hardy-conclusion}
\end{equation}
\end{proposition}

\begin{proof}
For $k\in\Omega_K$ one has $\Im(k^2)=2\Re k\Im k>0$, so $k^2$ lies in the resolvent set of the self-adjoint finite model and $J(k)$ is invertible.  Thus~\eqref{eq:scalar-f} is analytic.  The $H^2(\Omega_K)$ norm below is understood through nontangential boundary values in $L^2(\partial\Omega_K,\omega_*)$.  Since
\begin{equation}
 \|J^{-1}e_0\|
 \geq|\langle e_0,J^{-1}e_0\rangle|
 =|j_0|^{-1}|f_{R,L}|,
\end{equation}
the free factor, which is bounded above and below on $K$, gives
\begin{equation}
 \log^+\Gamma_0
 \leq C_K+\log^-|f_{R,L}|.
 \label{eq:scalar-log-reduction}
\end{equation}
Subharmonicity of $\log|f|$ and~\eqref{eq:scalar-Hardy-data} imply
\begin{align}
 \int_{\partial\Omega_K}\log^-|f|\dd\omega_*
 &\leq
 \int_{\partial\Omega_K}\log^+|f|\dd\omega_*
 -\log|f(k_*)|\notag\\
 &\leq \frac12C_*^2+\log(c_*^{-1}).
 \label{eq:scalar-harmonic-measure}
\end{align}
The lower bound for the harmonic-measure density on $K$ converts~\eqref{eq:scalar-harmonic-measure} to Lebesgue measure there.  Combining with~\eqref{eq:scalar-log-reduction} proves the claim.
\end{proof}

Proposition~\ref{prop:scalar-hardy-extraction} gives a sufficient scalar-coordinate criterion.  For a compactly supported tail with
\begin{equation}
 \eps^2=|\Sd|^{-1}\int_R^\infty\|V(r,\cdot)\|_2^2\dd r,
 \label{eq:scalar-epsilon}
\end{equation}
the first attempted Stage~1 obligation was to prove, uniformly in the angular cutoff,
\begin{equation}
 \|f_{R,L}\|_{H^2(\Omega_K)}\leq C_K,
 \qquad
 |f_{R,L}(k_*)-1|\leq C_K\eps,
 \qquad C_K\eps<1.
 \label{eq:RCH1}
\end{equation}
The second inequality would supply $c_*=1-C_K\eps$ and the first would supply $C_*$.  Thus $\mathrm{RCH}_1$ would imply the scalar Pl\"ucker estimate at the physical coupling without spectral averaging.  The following exact slab test shows, however, that this unrenormalized condition is false.

\begin{proposition}[$L^2$-small tails can destroy the interior normalization]\label{prop:RCH1-false}
Fix $k_*\in\C$ with $\Re k_*,\Im k_*>0$.  In dimension $d=3$, already in the radial constant channel there are bounded compactly supported potentials $q_\ell$ such that
\begin{equation}
 \|q_\ell\|_2\longrightarrow0,
 \qquad
 |f_\ell(k_*)|\longrightarrow0.
 \label{eq:RCH1-counterexample-conclusion}
\end{equation}
Consequently neither the interior lower bound in~\eqref{eq:scalar-Hardy-data} nor the closeness assertion in~\eqref{eq:RCH1} can hold under the $L^2$ tail hypothesis alone.
\end{proposition}

\begin{proof}
Let
\begin{equation}
 A_\ell=\ell^2,
 \qquad
 a_\ell=\ell^{-3/4},
 \qquad
 q_\ell(r)=a_\ell\mathbf1_{[A_\ell,A_\ell+\ell]}(r).
 \label{eq:RCH1-slab}
\end{equation}
Then
\begin{equation}
 \|q_\ell\|_2^2=a_\ell^2\ell=\ell^{-1/2}\longrightarrow0,
 \qquad
 M_\ell:=\int q_\ell=a_\ell\ell=\ell^{1/4}\longrightarrow\infty.
 \label{eq:RCH1-mass-cost}
\end{equation}
For $d=3$ the flattened constant channel has no centrifugal term.  Put
\begin{equation}
 p_\ell(k)=\sqrt{k^2-a_\ell},
 \qquad p_\ell(k)\longrightarrow k
 \quad(a_\ell\downarrow0),
 \label{eq:slab-momentum}
\end{equation}
with the analytic branch near $k_*$.  Matching the solution and its derivative at the two edges of the constant slab gives, after division by the free Jost factor,
\begin{align}
 \frac{J_\ell(k)}{J_0(k)}
 &=\frac{(k+p_\ell)^2}{4kp_\ell}
   e^{i(k-p_\ell)\ell}
   -\frac{(k-p_\ell)^2}{4kp_\ell}
   e^{i(k+p_\ell)\ell}
   +O_{k_*}(e^{-2\Im k_*A_\ell})\notag\\
 &=\exp\!\left(\frac{iM_\ell}{2k}
       +O_{k_*}(a_\ell^2\ell+a_\ell)\right)(1+o(1)).
 \label{eq:slab-transfer}
\end{align}
Here the second transmitted term and the reflected wave are exponentially small at $k_*$, while
$k-p_\ell=a_\ell/(2k)+O_{k_*}(a_\ell^2)$.  Since
$f_\ell=J_0/J_\ell$ in the scalar model,
\begin{equation}
 f_\ell(k_*)
 =\exp\!\left(\frac{M_\ell}{2ik_*}+o(1)\right),
 \qquad
 \Re\frac1{2ik_*}<0.
 \label{eq:slab-f-decay}
\end{equation}
Equation~\eqref{eq:RCH1-counterexample-conclusion} follows.
\end{proof}

The failed factor is a WKB phase: it has unit modulus on the physical axis but exponential size inside the quadrant.  With angular mixing it becomes an operator evolution rather than a scalar exponential, exactly the phenomenon identified in Denisov's WKB analysis~\cite{DenisovWKB2009}.  For a tail supported in $[R_0,T]$ and a finite angular cutoff, define
\begin{align}
 2ik\,\partial_rU_{V,L}(k;r)&=
 \left(\frac{D_L}{r^2}+P_LM_{V(r)}P_L\right)U_{V,L}(k;r),
 &U_{V,L}(k;R_0)&=\Id,\notag\\
 2ik\,\partial_rU_{0,L}(k;r)&=\frac{D_L}{r^2}U_{0,L}(k;r),
 &U_{0,L}(k;R_0)&=\Id,
 \label{eq:forward-wkb-evolution}\\
 \mathcal U_{T,L}(k)&=
 U_{0,L}(k;T)^{-1}U_{V,L}(k;T).
 \label{eq:relative-wkb-gauge}
\end{align}
For real $k\ne0$ both evolutions are unitary, so $\mathcal U_{T,L}(k)$ is unitary.  For $k$ in the first quadrant they are analytic and invertible.  In the radial $d=3$ model,
\begin{equation}
 \mathcal U_T(k)=
 \exp\!\left(\frac1{2ik}\int_{R_0}^Tq(r)\dd r\right),
 \label{eq:scalar-wkb-gauge}
\end{equation}
which is precisely the factor exposed by~\eqref{eq:slab-f-decay}.

Let $\mathcal J_{T,L}=J_{0,L}^{-1}J_{T,L}$, put
\begin{equation}
 w_{T,L}(k)=\mathcal J_{T,L}(k)^{-1}e_0,
 \qquad
 g_{T,L}(k)=\mathcal U_{T,L}(k)^{-1}w_{T,L}(k).
 \label{eq:gauged-radiation-column}
\end{equation}
The following vector criterion is invariant under the harmless real-axis rotation and avoids zeros of a single coordinate.

\begin{proposition}[Gauge-invariant vector Hardy extraction]\label{prop:gauged-vector-hardy}
Suppose, uniformly in $T,L$, that
\begin{equation}
 \|g_{T,L}(k_*)\|\geq c_*>0,
 \qquad
 \int_{\partial\Omega_K}\|g_{T,L}(k)\|^2\dd\omega_*(k)
 \leq C_*^2.
 \label{eq:GVH1}
\end{equation}
Then
\begin{equation}
 \sup_{T,L}\int_K\log^+\widetilde\Gamma_{0;T,L}(k)\dd k
 \leq C(K,c_*,C_*).
 \label{eq:gauged-vector-entropy}
\end{equation}
\end{proposition}

\begin{proof}
On $K$, unitarity of~\eqref{eq:relative-wkb-gauge} gives
\begin{equation}
 \|g_{T,L}(k)\|=\|w_{T,L}(k)\|,
 \qquad
 \log^+\widetilde\Gamma_{0;T,L}(k)
 =\log^-\|g_{T,L}(k)\|.
 \label{eq:gauge-boundary-invariance}
\end{equation}
The function $u(k)=\log\|g_{T,L}(k)\|$ is subharmonic.  Hence
\begin{align}
 \int_{\partial\Omega_K}u^-(k)\dd\omega_*(k)
 &\leq
 \int_{\partial\Omega_K}u^+(k)\dd\omega_*(k)-u(k_*)\notag\\
 &\leq\frac12C_*^2+\log(c_*^{-1}).
 \label{eq:vector-harmonic-measure}
\end{align}
The positive lower bound for the harmonic-measure density on $K$ and~\eqref{eq:gauge-boundary-invariance} prove the result.
\end{proof}

The replacement for $\mathrm{RCH}_1$ is therefore the gauged vector condition~\eqref{eq:GVH1}, denoted $\mathrm{GVH}_1$.  Its linear term is genuinely controlled.  In the radial $d=3$ model, writing $g_T=1+g_T^{(1)}+O(V^2)$ gives the exact cancellation
\begin{equation}
 g_T^{(1)}(k)
 =-\frac1{2ik}\int_{R_0}^T
   e^{2ik(r-R_0)}q(r)\dd r.
 \label{eq:gauged-first-born}
\end{equation}

\begin{proposition}[First-order gauged Hardy bound]\label{prop:gauged-first-order}
For a rectangular $\Omega_K$ with its lower side $K\Subset(0,\infty)$, rounded at the corners if desired,
\begin{equation}
 \sup_T\|g_T^{(1)}\|_{H^2(\Omega_K)}^2
 +\int_K\sup_{T>R_0}|g_T^{(1)}(k)|^2\dd k
 \leq C_K\int_{R_0}^\infty|q(r)|^2\dd r.
 \label{eq:gauged-first-order-bound}
\end{equation}
The same estimate holds for any fixed Bessel order; the $O_K(r^{-1})$ amplitude correction is absorbed by Cauchy--Schwarz.
\end{proposition}

\begin{proof}
On the lower side of the rectangle,~\eqref{eq:gauged-first-born} is the maximal Fourier integral controlled by Carleson--Hunt.  On the upper side, Plancherel applies after the damping factor $e^{-2r\Im k}$.  On a vertical side $k=k_0+iy$,
\begin{align}
 &\int_0^\eta
 \left|\int_{R_0}^\infty e^{2ik_0r}e^{-2yr}q(r)\dd r\right|^2\dd y\notag\\
 &\qquad=\int_{R_0}^\infty\!\int_{R_0}^\infty
 e^{2ik_0(r-s)}q(r)\overline{q(s)}
 \frac{1-e^{-2\eta(r+s)}}{2(r+s)}\dd r\dd s
 \leq C\|q\|_2^2,
 \label{eq:vertical-laplace-bound}
\end{align}
by Hilbert's integral inequality.  Since $|k|$ is bounded below on $\Omega_K$, the factor $(2ik)^{-1}$ is harmless.  The fixed-order Bessel remainder is bounded by
$\|q\|_2\|r^{-1}\|_{L^2(R_0,\infty)}$.
\end{proof}

The classical Christ--Kiselev multilinear bounds~\cite{ChristKiselev1998} require an input exponent strictly below $2$ and therefore do not settle the endpoint needed here.  Equations~\eqref{eq:forward-wkb-evolution}--\eqref{eq:gauged-first-order-bound} remove the zero-phase evolution before that endpoint is addressed.  The first nontrivial perturbative diagnostic is the quadratic coefficient of the gauged vector norm, whose real-axis value is still the quotient combination below because the gauge is unitary there.  Theorem~\ref{thm:QBC-two-slab} will show that this diagnostic is not uniformly bounded coefficientwise, even though the exact physical evolution remains bounded.

There is an exact perturbative diagnostic for where $\mathrm{GVH}_1$ first becomes nontrivial.  Let $J_0(k)$ be the diagonal free Bessel Jost matrix and use the \emph{left} normalization
\begin{equation}
 \begin{aligned}
 \mathcal J_\alpha(k)&=J_0(k)^{-1}J_\alpha(k)
 =\Id+\alpha A(k)+\alpha^2B(k)+O(\alpha^3),\\
 w_\alpha&=\mathcal J_\alpha^{-1}e_0,
 &\widetilde\Gamma_\alpha
 &:=\frac{\Gamma_\alpha}{|j_0|}
 =\|w_\alpha\|^{-1}.
 \end{aligned}
 \label{eq:scalar-coupling-expansion}
\end{equation}
If $\beta=\langle e_0,Ae_0\rangle$, direct inversion gives
\begin{align}
 \|w_\alpha\|^2
 &=1-2\alpha\Re\beta
 +\alpha^2\left(
   \|Ae_0\|^2
   +2\Re\langle e_0,(A^2-B)e_0\rangle
  \right)+O(\alpha^3),
 \label{eq:scalar-column-expansion}\\
 \log\widetilde\Gamma_\alpha
 &=\alpha\Re\beta
 +\alpha^2\left(
   (\Re\beta)^2-\frac12\|Ae_0\|^2
   -\Re\langle e_0,(A^2-B)e_0\rangle
  \right)+O(\alpha^3).
 \label{eq:scalar-log-expansion}
\end{align}
At first order only the low-low matrix element occurs.  At second order the full outgoing column appears through
\begin{equation}
 \|Ae_0\|^2=|\beta|^2+\|Q_0Ae_0\|^2.
 \label{eq:scalar-first-column}
\end{equation}
The summands in the quadratic coefficient must not be estimated separately.  Long weak radial potentials can make $\|Ae_0\|^2$ and $\Re\langle e_0,(A^2-B)e_0\rangle$ individually large while they cancel in $\log\widetilde\Gamma$.  The correct bilinear object is the full quotient-normalized combination
\begin{equation}
 \mathcal Q_2(k)
 :=(\Re\beta)^2-\frac12\|Ae_0\|^2
   -\Re\langle e_0,(A^2-B)e_0\rangle.
 \label{eq:scalar-Q2}
\end{equation}
Indeed, if
\begin{equation}
 \mathcal J_\alpha=
 \begin{pmatrix}
  1+\alpha a_1+\alpha^2a_2&\alpha b_1+O(\alpha^2)\\
  \alpha c_1+O(\alpha^2)&\Id+\alpha D_1+O(\alpha^2)
 \end{pmatrix},
 \label{eq:scalar-block-expansion}
\end{equation}
then the exact Schur quotient~\eqref{eq:scalar-schur-exact} gives
\begin{equation}
 \mathcal Q_2
 =\Re\!\left(a_2-b_1c_1-\frac12a_1^2\right)
  -\frac12\|c_1\|^2.
 \label{eq:scalar-Q2-block}
\end{equation}
The high-high coefficient $D_1$ cancels completely at this order.  This made a gauged $H^2(\Omega_K)$ bound for the \emph{combined positive part} of~\eqref{eq:scalar-Q2} the natural next candidate.  The radial two-slab test in Theorem~\ref{thm:QBC-two-slab} disproves that coefficientwise candidate.  Thus~\eqref{eq:scalar-Q2} remains a useful cancellation diagnostic, but it cannot be promoted to the endpoint lemma.

\subsection{Ordinary compact probes}\label{sec:probes}

The boundary density may be converted into an $L^2$ probe density without changing the open estimate.  Remove a compact core so that $V=0$ on a collar $(1,c)$.  Let $g\in C_c^\infty((1,c))$ be real and nonzero, and in flattened coordinates take
\begin{equation}
 \Phi_Ne_j=g(r)Y_j(\omega).
 \label{eq:flattened-probe}
\end{equation}
In physical coordinates the same probe is
\begin{equation}
 \phi_j(x)=|x|^{-(d-1)/2}g(|x|)Y_j(x/|x|).
 \label{eq:physical-probe}
\end{equation}
Because the collar dynamics is diagonal in angular momentum, its regular transform is an $N\times N$ diagonal matrix $D_g(k)$ independent of $R$ and $L$.  Hence
\begin{equation}
 W_{N}^{R,L}(k^2)=D_g(k)^*h_P(k^2)D_g(k),
 \qquad
 \det W_N^{R,L}=|\det D_g(k)|^2\det h_P(k^2).
 \label{eq:probe-boundary}
\end{equation}
The zeros of $\det D_g$ are isolated and contribute a fixed, locally integrable logarithm.  Consequently~\eqref{eq:PSR} implies the finite-channel entropy theorem for genuine $L^2$ probes.

\section{A second route: normalized angular entropy}

Let $P_L$ contain every spherical harmonic of degree at most $L$.  The addition theorem gives
\begin{equation}
 \sum_{j=1}^{m_L}|Y_j(\omega)|^2=\frac{m_L}{|\Sd|}.
 \label{eq:addition}
\end{equation}

\begin{proposition}[Exact normalized trace identity]\label{prop:normalized-cost}
For every $v\in L^2(\Sd)$,
\begin{align}
 \frac1{m_L}\|M_vP_L\|_{\HS}^2
 &=\frac1{|\Sd|}\|v\|_2^2,
 \label{eq:normalized-full-column}\\
 \frac1{m_L}\|P_LM_vP_L\|_{\HS}^2
 &\leq\frac1{|\Sd|}\|v\|_2^2,
 \label{eq:normalized-compression}\\
 \frac1{m_L}\|Q_LM_vP_L\|_{\HS}^2
 &\leq\frac1{|\Sd|}\|v\|_2^2.
 \label{eq:normalized-leakage}
\end{align}
\end{proposition}

\begin{proof}
Expand the Hilbert--Schmidt norm and apply~\eqref{eq:addition}:
\begin{equation}
 \|M_vP_L\|_{\HS}^2
 =\sum_{j=1}^{m_L}\int_{\Sd}|v(\omega)|^2|Y_j(\omega)|^2\dd\omega
 =\frac{m_L}{|\Sd|}\|v\|_2^2.
\end{equation}
The last two claims follow by orthogonal compression.
\end{proof}

This identity suggests auditing the finite-width matrix $L^2$ sum rule of Molchanov--Vainberg~\cite{MolchanovVainberg2004} with constants explicit in $m_L$.  With $h_{0,L}$ denoting the diagonal Bessel reference density, the desired dimension-linear statement is
\begin{equation}
 \frac1{m_L}\int_I
 \log^-\det\!\left(h_{0,L}(E)^{-1/2}h_{R,L}(E)h_{0,L}(E)^{-1/2}\right)\dd E
 \leq C_I\left(1+\|V\|_{X_d}^2\right)
 \label{eq:NMSR}
\end{equation}
The perturbative cost in~\eqref{eq:NMSR} is compatible with Proposition~\ref{prop:normalized-cost}; the Bessel normalization prevents high angular momentum from being charged as a potential.

\subsection{Why normalized entropy is not yet the theorem}

Even a proof of~\eqref{eq:NMSR} would not, by itself, justify the passage $L\to\infty$.  A positive fraction of good spectral directions can drift to higher angular modes while every fixed direction degenerates.  The exact matrix example
\begin{equation}
 W_m=e^{-\sqrt m}P_{\lfloor\sqrt m\rfloor}
 +(P_m-P_{\lfloor\sqrt m\rfloor})
\label{eq:dimension-escape}
\end{equation}
has
\begin{equation}
 \frac1m\log\det W_m\longrightarrow-1,
\end{equation}
yet the density in every fixed coordinate tends to zero.  This is the \emph{dimension-escape obstruction}.

The normalized route therefore requires the following separate statement.

\begin{problem}[Angular tightness/no escape]\label{prob:tightness}
Derive from the scalar multiplication structure, the centrifugal term, and a buffered version of~\eqref{eq:NMSR} that for each fixed $N$ there are a finite-dimensional subspace $\cE_N\subset\K$ and $L_0$ such that
\begin{equation}
 \sup_{R,L\geq L_0}
 \int_I\log^-\det W_{\cE_N}^{R,L}(E)\dd E<\infty
 \label{eq:TIGHT}
\end{equation}
The subspace $\cE_N$ must be fixed before the limit $L\to\infty$.
\end{problem}

If both~\eqref{eq:NMSR} and~\eqref{eq:TIGHT} are proved, the established entropy lower-semicontinuity theorem finishes the limit.  The normalized route is therefore viable, but the no-escape lemma is a genuine theorem rather than a compactness slogan.

\section{The full Feshbach self-energy and its exact energy average}

For the $P_N\oplus Q_N$ decomposition, write
\begin{equation}
 H=
 \begin{pmatrix}
 H_{PP}&B^*\\ B&H_{QQ}
 \end{pmatrix},
 \qquad B=Q_NM_VP_N.
 \label{eq:block-H}
\end{equation}
For $z\in\C_+$,
\begin{equation}
 P_N(H-z)^{-1}P_N
 =\left(H_{PP}-z-\Sigma_N(z)\right)^{-1},
 \qquad
 \Sigma_N(z)=B^*(H_{QQ}-z)^{-1}B.
 \label{eq:feshbach}
\end{equation}
The familiar sign is exact:
\begin{equation}
 \Im\Sigma_N(E+i\eta)
 =\eta B^*\bigl((H_{QQ}-E)^2+\eta^2\bigr)^{-1}B\succeq0.
 \label{eq:selfenergy-positive}
\end{equation}

\begin{proposition}[Poisson average of the full self-energy]\label{prop:poisson}
Let $A$ be self-adjoint and $B$ bounded.  Then, in the strong operator topology,
\begin{equation}
 \int_{\R}\Im\!\left[B^*(A-E-i\eta)^{-1}B\right]\dd E
 =\pi B^*B,
 \qquad \eta>0.
 \label{eq:poisson-identity}
\end{equation}
\end{proposition}

\begin{proof}
The spectral theorem reduces the claim to
\begin{equation}
 \int_{\R}\frac{\eta}{(\lambda-E)^2+\eta^2}\dd E=\pi.
\end{equation}
Sandwiching by $B^*$ and $B$ proves~\eqref{eq:poisson-identity}.
\end{proof}

\subsection{The exact $2m$-dimensional phase-space quotient}

The differential quotient identity has a direct phase-space lift.  In a
finite angular model write
\begin{equation}
 W(r)=\frac{D}{r^2}+V_m(r),
 \qquad
 -F''+WF=k^2F.
 \label{eq:finite-W}
\end{equation}
Order the phase coordinates as
$(u_P,u_P',u_Q,u_Q')$.  The first-order system is
\begin{equation}
 \Psi'=\mathbb A\Psi,
 \qquad
 \mathbb A=
 \begin{pmatrix}
  \mathbb A_{PP}&\mathbb A_{PQ}\\
  \mathbb A_{QP}&\mathbb A_{QQ}
 \end{pmatrix},
 \label{eq:phase-system}
\end{equation}
where
\begin{align}
 \mathbb A_{PP}
 &=\begin{pmatrix}0&\Id_P\\ W_{PP}-k^2\Id_P&0\end{pmatrix},
 &\mathbb A_{PQ}
 &=\begin{pmatrix}0&0\\W_{PQ}&0\end{pmatrix},
 \label{eq:phase-blocks-1}\\
 \mathbb A_{QP}
 &=\begin{pmatrix}0&0\\W_{QP}&0\end{pmatrix},
 &\mathbb A_{QQ}
 &=\begin{pmatrix}0&\Id_Q\\W_{QQ}-k^2\Id_Q&0\end{pmatrix}.
 \label{eq:phase-blocks-2}
\end{align}
Let $\mathbb P$ be the rank-$2N$ low-channel phase projection and let
$\mathbb T' =\mathbb A\mathbb T$ be an invertible fundamental transfer.

\begin{corollary}[Phase-space high--high cancellation]\label{cor:phase-flux}
Define the phase-space regression state $\mathbb C$ from $\mathbb T$ exactly
as $C$ was defined from $\mathcal J$ in
\eqref{eq:quotient-rows}--\eqref{eq:quotient-regression}.  Then
\begin{equation}
 \frac{\dd}{\dd r}\log\Gamma_{\mathbb P}(\mathbb T(r,k))
 =\Re\tr\!\left(
   \mathbb A_{PP}-\mathbb C\mathbb A_{QP}
  \right).
 \label{eq:phase-Grassmannian-flux}
\end{equation}
The full block $W_{QQ}$ cancels.  After conjugating the known free and
centrifugal low evolution, the variable part of the displayed charged column
contains only $V_{PP}$ and $V_{QP}$, hence only $M_VP_N$.
\end{corollary}

\begin{proof}
Apply Proposition~\ref{prop:differential-flux} to the decomposition
$\mathbb P\oplus(\Id-\mathbb P)$ and use
\eqref{eq:phase-blocks-1}--\eqref{eq:phase-blocks-2}.  The only potential
entries in $\mathbb A_{PP}$ and $\mathbb A_{QP}$ are $V_{PP}$ and $V_{QP}$.
\end{proof}

For $m=2$ and $N=1$, Corollary~\ref{cor:phase-flux} is the promised
$4\times4$ low/high test.  It passes algebraically even when the high-channel
potential is arbitrary: $V_{22}$ changes the dressed state $\mathbb C$ but
does not appear as a separately charged term.

The outgoing solution is not the full transfer matrix.  It is the
$2m\times m$ flux-normalized Lagrangian frame
\begin{equation}
 \mathbb L(r,k)=\binom{F(r,k)}{F'(r,k)},
 \qquad
 \mathbb L(r,k)^*
 \begin{pmatrix}0&\Id\\-\Id&0\end{pmatrix}
 \mathbb L(r,k)=2ik\Id
 \label{eq:outgoing-Lagrangian-frame}
\end{equation}
on the positive axis, while the physical Jost matrix is the configuration
block $J(k)=F(1,k)$.  The following identity removes this apparent mismatch.

\begin{proposition}[Physical Riccati descent]\label{prop:physical-Riccati-descent}
For every real $k>0$, the outgoing configuration matrix $F(r,k)$ is
invertible.  Assume $N<m$; the case $N=m$ is obtained by omitting the $Q$
rows and setting $C_F=0$.  Define its logarithmic derivative and its low-row regression by
\begin{align}
 \mathcal R&=F'F^{-1},
 &C_F&=(PF)(QF)^*\bigl((QF)(QF)^*\bigr)^{-1}.
 \label{eq:physical-Riccati-variables}
\end{align}
Then
\begin{align}
 \mathcal R'&=W-k^2\Id-\mathcal R^2,
 \label{eq:physical-Riccati-equation}\\
 \Im\mathcal R&=kF^{-*}F^{-1}>0,
 \label{eq:physical-Riccati-flux}\\
 \det\!\left(P\Im\mathcal R P\big|_{P\C^m}\right)
 &=k^N\Gamma_P(F)^{-2},
 \label{eq:physical-Riccati-Gamma}\\
 \frac{\dd}{\dd r}\log\Gamma_P(F(r,k))
 &=\Re\tr\!\left(\mathcal R_{PP}-C_F\mathcal R_{QP}\right).
 \label{eq:physical-Riccati-quotient-flux}
\end{align}
Moreover, the charged Riccati column obeys
\begin{equation}
 (\mathcal RP)'=\left(\frac{D}{r^2}-k^2\Id\right)P
 +V_mP-\mathcal R^2P.
 \label{eq:charged-Riccati-column}
\end{equation}
Thus $V_mP$ is the only variable inhomogeneous source.  The high--high block
$V_{QQ}$ changes the dressed state $\mathcal R$ but is not charged separately.
\end{proposition}

\begin{proof}
If $F(r,k)c=0$, the flux identity
\eqref{eq:outgoing-Lagrangian-frame}, tested on $c$, gives
$0=2ik\|c\|^2$; hence $F$ is invertible.  Differentiating $F'F^{-1}$ and using
$F''=(W-k^2\Id)F$ proves~\eqref{eq:physical-Riccati-equation}.  Multiplication
of the flux identity on the left by $F^{-*}$ and on the right by $F^{-1}$
gives~\eqref{eq:physical-Riccati-flux}.  Proposition~\ref{prop:dual-plucker}
then gives~\eqref{eq:physical-Riccati-Gamma}.  Finally apply
Proposition~\ref{prop:differential-flux} to the exact left cocycle
$F'=\mathcal RF$; this proves~\eqref{eq:physical-Riccati-quotient-flux}.
Right multiplication of~\eqref{eq:physical-Riccati-equation} by $P$ gives
\eqref{eq:charged-Riccati-column}.
\end{proof}

This proves the algebraic descent, including the $4\times4$ one-low/one-high
test, without Airy asymptotics.  It also identifies the genuinely analytic
remainder.

\begin{problem}[Riccati/Feshbach energy square function]
\label{prob:Riccati-square-function}
After conjugating the known free centrifugal flow and the real-axis WKB phase,
prove for each fixed $N$ that the energy-positive shell increment generated by
\eqref{eq:physical-Riccati-quotient-flux} satisfies
\begin{equation}
 \int_K\bigl[\log\Gamma_P(F(R_1,k))
             -\log\Gamma_P(F(R_2,k))\bigr]_+\dd k
 \leq C_{K,N}\int_{R_1}^{R_2}\|M_{V(r)}P_N\|_{\HS}^2\dd r
      +\mathcal F(R_1)-\mathcal F(R_2),
 \label{eq:Riccati-square-function}
\end{equation}
with $\mathcal F\geq0$, uniformly in the angular cutoff and the support radius.
The proof must identify elimination of buffered elliptic channels with the
Feshbach complement~\eqref{eq:feshbach} and keep the energy-dependent dressed
 state inside the estimate.  The unrestricted one-low/one-high matrix estimate
 fails by Theorem~\ref{thm:matrix-rotator-counterexample}, and the stronger
 uniform scalar statement in~\eqref{eq:Riccati-square-function} fails by
 Theorem~\ref{thm:scalar-PSR1-counterexample}.  The cutoff-edge scalar lift
 shows that a degree-$2L$ multiplier can create the exact rotator and barrier
 inside the degree-$L$ model at bounded full constant-column cost.  Any viable
 replacement must therefore impose a cutoff-compatible bandwidth buffer or a
 justified diagonal order of limits.  The full-space converter in
 Problem~\ref{prob:scalar-converter-dichotomy} remains a stronger diagnostic,
 not the unresolved test for the displayed uniform estimate.
\end{problem}

The orientation in~\eqref{eq:Riccati-square-function} is essential.  The outgoing solution is normalized outside
the support, so attaching the shell $[R_1,R_2]$ changes the boundary value by
propagating from $R_2$ inward to $R_1$.  In the scalar case
$\Gamma(F(R_2,k))=1$ outside the support, while the inner Jost modulus can be
larger than one.  The opposite orientation would therefore miss precisely the
entropy loss that must be controlled.

\paragraph{The finite laboratory.}
For $m=2$ and $N=1$, write
\begin{equation}
 \mathcal R=\begin{pmatrix}a&b\\c&d\end{pmatrix},
 \qquad
 V_m=\begin{pmatrix}v_{11}&\overline{v_{21}}\\v_{21}&v_{22}\end{pmatrix}.
\end{equation}
After the known diagonal centrifugal term is separated, the charged column of
\eqref{eq:physical-Riccati-equation} is exactly
\begin{align}
 a'&=v_{11}+\text{known free term}-a^2-bc,
 \label{eq:finite-Riccati-a}\\
 c'&=v_{21}-ca-dc,
 \label{eq:finite-Riccati-c}\\
 \frac{\dd}{\dd r}\log\Gamma_P(F)
 &=\Re(a-C_Fc).
 \label{eq:finite-Riccati-flux}
\end{align}
The truncated variable source cost is $|v_{11}|^2+|v_{21}|^2$.
The coefficient $v_{22}$ acts only through the dressed variables $d,c,C_F$,
but Proposition~\ref{prop:cocycle-low-column-no-go} shows that this algebraic
fact alone gives no estimate: a high block may preload the state and a later
small rotation may expose it.  Allowing an arbitrary bounded $v_{22}$ is thus
a deliberately stronger matrix stress test, not the actual scalar endpoint.
For the physical compression $V_m=P_LM_VP_L$, the charged norm is
$\|M_VP_N\|_{\HS}$ in the full angular space.  It is tempting to conclude that
creating $v_{22}$ must therefore be expensive.  Propositions
~\ref{prop:minimum-scalar-block-lift} and~\ref{prop:thin-belt-barrier} show that
this is false: high-harmonic concentration produces an exact $v_{22}=Q$ block
with vanishing full constant-column cost, and the high state can be made almost
invariant during the barrier.  What remains genuinely non-finite is the
converter, because the scalar lift of its $2\times2$ compression also couples
the selected high harmonic to the rest of the high sector.

\begin{problem}[Stronger full-space scalar converter dichotomy]
\label{prob:scalar-converter-dichotomy}
Let $e_0=(4\pi)^{-1/2}$ and let $h_\ell$ be the real highest-weight harmonic
~\eqref{eq:real-highest-weight}.  Decide between the following two alternatives
for the full angular Schr\"odinger evolution, including $A_3/r^2$:
\begin{enumerate}[label=(\roman*)]
 \item prove a quantitative lower bound excluding any uniformly bounded real
 scalar converter which maps both directional modes $e_0^\pm$ to
 $h_\ell^\pm$ with error $o(1)$ while
 $\int\|M_{v(r)}e_0\|_2^2\dd r=o(1)$; or
 \item construct such converters with modal error
 $o(e^{-\kappa T_\ell})$ and concatenate them with the thin-belt barrier of
 Proposition~\ref{prop:thin-belt-barrier}.
\end{enumerate}
A proof of (i) would supply structural input for a cutoff-stable relative sum
rule.  A construction in (ii) would give a full-space version of the
rotator/barrier obstruction.  Theorem~\ref{thm:scalar-PSR1-counterexample}
already disproves $\mathrm{PSR}_1$ by the weaker finite-cutoff mechanism, so
this dichotomy is retained as a stronger question rather than as the logical
gate to $\mathrm{PSR}_1$.
Propositions~\ref{prop:decomposable-converter-no-go} and
~\ref{prop:angular-current-requirement} prove the first part of (i) in the
centrifugal-perturbative regime: any successful converter must generate
order-one angular current and cannot be close to the decomposable evolution
obtained by deleting $A_3/r^2$.  What remains is to control or construct that
nonperturbative angular transport using the charged scalar column.
Propositions~\ref{prop:minimum-lift-complement-leakage} and
~\ref{prop:nodal-set-lift-obstruction} also exclude a near-invariant
 two-state scalar realization of the matrix rotator.  The sole constructive
 escape is therefore a genuinely multimode excursion followed by uniform
 refocusing into $h_\ell$.  Proposition~\ref{prop:cheap-scalar-phase-mask}
 shows that high angular frequency can itself be generated at vanishing cost;
 Proposition~\ref{prop:phase-mask-transport-incompatibility} then shows that a
 clean cheap mask cannot be placed at the same radial scale as degree-$\ell$
 centrifugal transport.  The remaining core is therefore the coupled
 weak-slab version of the achromatic bidirectional refocusing problem
 ~\ref{prob:achromatic-phase-mask-refocusing}, not the ideal phase-mask product
 by itself.  The resonant-drive diagnostic
 ~\eqref{eq:born-broadband-lower-bound} further shows that the obvious weak
 resonant pulse is cheap only pointwise in energy: at first order, uniform
 transfer over a fixed band already has order-one charged cost.  Proposition
 ~\ref{prop:SU2-nonlinear-Plancherel} makes that budget exact and nonlinear for
 a pure-detuning ensemble.  Proposition
 ~\ref{prop:scaled-matrix-ensemble-no-go} shows why it cannot simply be imported
 into the physically scaled system.  Propositions
 ~\ref{prop:scalar-ground-adiabatic-no-go},
 ~\ref{prop:chirped-ensemble-cost}, and
 ~\ref{prop:rotating-ground-ladder-fanout} exclude respectively the direct
 positive-ground rotator, the standard monotone highest-weight chirp, and the
 adjacent-rung rotating ground ladder.  Corollary
 ~\ref{cor:scalar-pulse-area-roughness} forces any gap-cancelled cheap design
 to use scalar angular bandwidth $\gtrsim\ell/\mathcal E_\ell^2$, while
 Proposition~\ref{prop:born-high-band-recoil} excludes direct broadband
 injection into that rough band at first order.  What remains is a genuinely
 nonlinear bootstrap: lower-mode concentration, rough complementary-sector
 rotation, and reciprocal refocusing in the exact scaled radial evolution.
\end{problem}

The formula sees the full leakage and is immune to cutoff aliasing.  It does not yet prove entropy, because the low-channel boundary solution depends on the same energy being integrated.  If
\begin{equation}
 \mathcal U_{P;R,L}(z):\C^N\longrightarrow L^2((1,\infty);P_N\K)
 \label{eq:dressed-solution-map}
\end{equation}
denotes the low solution generated by the boundary data, the missing nonlinear embedding has the schematic form
\begin{equation}
 \sup_{0<\eta<1}\int_I
 \tr\!\left[
 \mathcal U_P(E+i\eta)^*\Im\Sigma_N(E+i\eta)\mathcal U_P(E+i\eta)
 \right]\dd E
 \leq C_{I,N}\cE_N(V)
 \label{eq:DCE}
\end{equation}
The fixed-vector Poisson identity~\eqref{eq:poisson-identity} is the linear seed of~\eqref{eq:DCE}; the energy dependence of $\mathcal U_P$ is the entire difficulty.

\section{A rigorous free spacetime estimate in the conjectured norm}

The weight in Simon's conjecture is exactly the weight selected by free finite-channel propagation.

\begin{proposition}[Restricted free smoothness]\label{prop:free-smoothing}
Let $I=[E_0,E_1]\Subset(0,\infty)$, let $P_N$ be as above, and let $H_0=-\partial_r^2+A_d/r^2$ on the exterior half-line.  There is $C_{I,N}<\infty$ such that
\begin{equation}
 \int_{\R}\left\|M_Ve^{-itH_0}\mathbf1_I(H_0)P_Nf\right\|_{L^2((1,\infty)\times\Sd)}^2\dd t
 \leq C_{I,N}\cE_N(V)\|f\|^2.
 \label{eq:free-smoothing}
\end{equation}
\end{proposition}

\begin{proof}
Let $k\in K=[\sqrt{E_0},\sqrt{E_1}]$.  In each of the finitely many channels, let $\varphi_j(r,k)$ be the Bessel generalized eigenfunction in flattened coordinates, normalized for the spectral transform.  Since $r\geq1$ and $K\Subset(0,\infty)$,
\begin{equation}
 \sup_{1\leq j\leq N}\sup_{r\geq1,\,k\in K}|\varphi_j(r,k)|\leq C_{I,N}.
 \label{eq:bessel-bound}
\end{equation}
Writing the spectral coefficients as $\widehat f_j(k)$ gives
\begin{equation}
 u(t,r,\omega)
 =\sum_{j=1}^NY_j(\omega)
   \int_K e^{-itk^2}\varphi_j(r,k)\widehat f_j(k)\dd k.
 \label{eq:free-expansion}
\end{equation}
Plancherel in $t$, with $E=k^2$ and $\dd E=2k\dd k$, followed by Cauchy--Schwarz in the finite channel index, yields
\begin{equation}
 \int_{\R}|u(t,r,\omega)|^2\dd t
 \leq C_{I,N}\sum_{j=1}^N|Y_j(\omega)|^2
 \int_K|\widehat f_j(k)|^2\dd k.
 \label{eq:pointwise-time}
\end{equation}
Multiply by $|V(r,\omega)|^2$, integrate in $r$ and $\omega$, and use~\eqref{eq:column-exact}.  The spectral transform is unitary, so the resulting right-hand side is bounded by $C_{I,N}\cE_N(V)\|f\|^2$.
\end{proof}

\begin{remark}
Estimate~\eqref{eq:free-smoothing} is an $L_t^2$ statement, whereas Cook's method asks for $L_t^1$.  Applying Cauchy--Schwarz on an infinite time interval loses a divergent factor.  The correct next step is therefore not to quote Cook's theorem, but to prove a maximal or dressed square-function estimate.  This is precisely the role of~\eqref{eq:DCE} and the nonlinear Carleson module below.
\end{remark}

\section{Auxiliary approaches and their limitations}

\subsection{Module A: a Pl\"ucker outer-factor sum rule}

For the finite model, define the analytic exterior vector
\begin{equation}
 p_Q(z)=\bigwedge^{m_L-N}\!\left(Q_NJ_{R,L}(z)\right).
 \label{eq:pQ}
\end{equation}
On the real axis its norm is the denominator of~\eqref{eq:Gamma}.  A distinguished coordinate of $p_Q$ is unsafe because it may have zeros unrelated to the spectral problem.  The natural scalar replacement is an outer function $\mathcal O_Q$ whose boundary modulus is $\|p_Q\|$:
\begin{equation}
 \log|\mathcal O_Q(z)|
 =\frac1\pi\int_{\R}
   \frac{\Im z}{|t-z|^2}\log\|p_Q(t)\|\dd t,
 \label{eq:outer-Q}
\end{equation}
with the standard integrable regularization at infinity.

After removing bound-state zeros of $\det J$ by a Blaschke product $\mathcal B$, set formally
\begin{equation}
 \mathcal D_P(z)
 =\frac{\det J(z)}{\mathcal B(z)\mathcal O_Q(z)}.
 \label{eq:relative-outer}
\end{equation}
Since $|\mathcal B|=1$ on $\R$,
\begin{equation}
 \log|\mathcal D_P(k)|=\log\Gamma_{P;R,L}(k).
 \label{eq:D-Gamma}
\end{equation}

\begin{problem}[Grassmannian sum rule]\label{prob:grassmann-sum-rule}
Prove that $\mathcal D_P$ is of uniformly bounded characteristic and that its trace identity has quadratic coefficient
\begin{equation}
 \int_1^\infty\tr\!\left(P_NM_{V(r)}^2P_N\right)\dd r
 =\cE_N(V),
 \label{eq:quadratic-coefficient}
\end{equation}
with all zero and pole terms of favorable sign or bounded by the same cost.  Deduce~\eqref{eq:PSR}.
\end{problem}

The cancellation in~\eqref{eq:quadratic-coefficient} is the central advantage of the Pl\"ucker quotient.  A full determinant charges $\tr(P_LM_VP_L)^2$, while the relative $P$-volume should charge only the $P$ corner of $M_V^2$:
\begin{equation}
 \tr(P_NM_V^2P_N)
 =\|P_NM_VP_N\|_{\HS}^2+\|Q_NM_VP_N\|_{\HS}^2.
 \label{eq:desired-trace}
\end{equation}
This is exactly the safe cost.

\subsection{Module B: dressed Carleson and wave-packet control}

Proposition~\ref{prop:free-smoothing} gives the linear estimate.  The desired nonlinear upgrade is a maximal bound for a flux-normalized interaction variable $Z_{R,L}(r,k)$:
\begin{equation}
 \int_K\sup_{\rho\geq1}
 \left\|
  \int_1^\rho e^{2iks}\,M_{V(s)}Z_{R,L}(s,k)P_N\dd s
 \right\|_{\HS}^2\dd k
 \leq C_{K,N}\cE_N(V)
 \label{eq:NC}
\end{equation}
The factor $Z$ is constrained by flux conservation and the common second-order radial kinetic term; it is not an arbitrary bounded matrix.  An estimate of the form~\eqref{eq:NC} would control the logarithmic growth of the exterior volume in~\eqref{eq:Gamma} shell by shell and imply~\eqref{eq:DCE}.

Recent wave-packet work of Denisov~\cite{Denisov2025} supplies a model for exploiting phase separation, but its conclusions are generic in a parameter and require additional decay/frequency hypotheses.  The proposed use here is narrower: establish an energy square function for the exact full leakage, not uniform pointwise free asymptotics.

The first proposed harmonic-analysis milestone was already linear.  Let
\begin{equation}
 v_{\nu j}(r)=\langle Y_\nu,M_{V(r)}Y_j\rangle_{L^2(\Sd)},
 \qquad
 \sum_\nu|v_{\nu j}(r)|^2=\|V(r,\cdot)Y_j\|_2^2,
 \label{eq:parseval-column}
\end{equation}
and let $\phi_\nu^\pm(k,r)$ be flux-normalized incoming and outgoing Bessel waves.  The tempting raw estimate was
\begin{equation}
 \int_K\sup_{T>R}\sum_\nu
 \left|
  \int_R^T
  \phi_\nu^-(k,r)v_{\nu j}(r)\phi_j^+(k,r)\dd r
 \right|^2\dd k
 \leq C_{K,j}\int_R^\infty\|V(r,\cdot)Y_j\|_2^2\dd r.
 \label{eq:Bessel-Carleson}
\end{equation}
The full sum in~\eqref{eq:Bessel-Carleson} avoids angular aliasing, but that fact does not remove a radial forward mode.

\begin{proposition}[The raw on-shell estimate is false]\label{prop:raw-bessel-false}
For every nontrivial compact $K\Subset(0,\infty)$, estimate~\eqref{eq:Bessel-Carleson} fails already for a radial scalar potential in the constant angular channel.
\end{proposition}

\begin{proof}
Uniformly for $k\in K$, the flux-normalized constant-channel waves have the large-$r$ form
\begin{equation}
 \phi_0^\pm(k,r)
 =(2k)^{-1/2}e^{\pm i(kr-\vartheta_0)}
 \bigl(1+O_K(r^{-1})\bigr),
 \qquad
 \phi_0^-(k,r)\phi_0^+(k,r)
 =\frac1{2k}+O_K(r^{-1}).
 \label{eq:forward-bessel-product}
\end{equation}
Let $A_\ell=\ell^2$ and choose the bounded compactly supported radial potential
\begin{equation}
 q_\ell(r)=\eps\ell^{-1/2}\mathbf1_{[A_\ell,A_\ell+\ell]}(r),
 \qquad
 V_\ell(r,\omega)=q_\ell(r).
 \label{eq:forward-counterexample}
\end{equation}
Its charged cost is $\int|q_\ell|^2\dd r=\eps^2$.  Taking $j=\nu=0$ and $T=A_\ell+\ell$, however, gives
\begin{align}
 \int_{A_\ell}^{T}
 \phi_0^-q_\ell\phi_0^+\dd r
 &=\frac{\eps\sqrt\ell}{2k}
   +O_K\!\left(\frac{\eps\sqrt\ell}{A_\ell}\right),
 \label{eq:forward-counterexample-amplitude}\\
 \int_K\sup_{T>A_\ell}
 \left|\int_{A_\ell}^{T}\phi_0^-q_\ell\phi_0^+\dd r\right|^2\dd k
 &\geq c_K\eps^2\ell
 \label{eq:forward-counterexample-growth}
\end{align}
for all sufficiently large $\ell$.  This contradicts~\eqref{eq:Bessel-Carleson} with a constant independent of the support location and length.
\end{proof}

\begin{remark}
Proposition~\ref{prop:raw-bessel-false} does not contradict the free spacetime estimate in Proposition~\ref{prop:free-smoothing}.  Plancherel controls an $L_t^2$ evolution before restricting both waves to the same energy shell; the failure above is precisely the nonoscillatory diagonal created by that restriction.
\end{remark}

The scalar Jost relation identifies the exact subtraction.  Let $f_0^+$ be the free outgoing constant-channel solution, $f_0^-=\overline{f_0^+}$ on the positive axis, $j_0$ its Jost factor, and $s_0$ the regular solution.  Then
\begin{equation}
 s_0(r,k)=\frac{\overline{j_0(k)}f_0^+(r,k)-j_0(k)f_0^-(r,k)}{2ik}.
 \label{eq:scalar-regular-jost-relation}
\end{equation}
For the truncated first Born coefficient
\begin{equation}
 \beta_T(k)=\frac1{j_0(k)}\int_R^T
 s_0(r,k)v_{00}(r)f_0^+(r,k)\dd r,
 \label{eq:scalar-beta-T}
\end{equation}
substitution of~\eqref{eq:scalar-regular-jost-relation} gives
\begin{align}
 \beta_T(k)
 &=\frac{\overline{j_0(k)}}{2ikj_0(k)}
   \int_R^T v_{00}(r)\bigl(f_0^+(r,k)\bigr)^2\dd r
   -\frac1{2ik}\int_R^T v_{00}(r)|f_0^+(r,k)|^2\dd r,
 \label{eq:scalar-forward-splitting}\\
 \Re\beta_T(k)
 &=\Re\!\left[
   \frac{\overline{j_0(k)}}{2ikj_0(k)}
   \int_R^T v_{00}(r)\bigl(f_0^+(r,k)\bigr)^2\dd r
  \right].
 \label{eq:scalar-forward-subtraction}
\end{align}
The omitted term in~\eqref{eq:scalar-forward-subtraction} is purely imaginary.  Since $(f_0^+)^2$ has phase $e^{2ikr}$ and a fixed-order Bessel amplitude, the Carleson--Hunt maximal theorem~\cite{Hunt1968}, together with the $O_K(r^{-1})$ asymptotic remainder, proves the corrected scalar estimate
\begin{equation}
 \int_K\sup_{T>R}|\Re\beta_T(k)|^2\dd k
 \leq C_K\int_R^\infty|v_{00}(r)|^2\dd r
 \leq\frac{C_K}{|\Sd|}
      \int_R^\infty\|V(r,\cdot)\|_2^2\dd r.
 \label{eq:scalar-renormalized-carleson}
\end{equation}
Thus the scalar linear part is controlled after the exact forward phase has been removed; no Airy turning analysis is needed at this order because only the constant channel occurs.

The first proposed nonlinear obligation was quadratic and retained the
Pl\"ucker quotient cancellation.  Let $A_{R,L,T}$ and $B_{R,L,T}$ be the first
two coefficients of the left-normalized Jost--Volterra expansion cut off at
$T$, put $\beta_{R,L,T}=\langle e_0,A_{R,L,T}e_0\rangle$, and define
\begin{equation}
 \mathcal Q_{R,L,T}(k)
 :=(\Re\beta_{R,L,T})^2
 -\frac12\|A_{R,L,T}e_0\|^2
 -\Re\langle e_0,(A_{R,L,T}^2-B_{R,L,T})e_0\rangle.
 \label{eq:QBC1-integrand}
\end{equation}
The proposed quotient Bessel--Carleson estimate $\mathrm{QBC}_1$ asserted
\begin{equation}
 \sup_{L}\int_K\sup_{T>R}
 [\mathcal Q_{R,L,T}(k)]_+\dd k
 \leq\frac{C_K}{|\Sd|}
 \int_R^\infty\|V(r,\cdot)\|_2^2\dd r.
 \label{eq:QBC1}
\end{equation}
It was deliberately imposed on the combined positive part.  Nevertheless, it
is false before any Airy turning channel is present.

In the scalar radial $d=3$ model use
\begin{equation}
 f_0(r,k)=e^{ik(r-1)},
 \qquad
 s_0(r,k)=\frac{\sin(k(r-1))}{k}.
\end{equation}
The exact combined quadratic coefficient reduces to
\begin{equation}
 \mathcal Q_{R,1,T}(k)
 =-\Re\int_{R<r<t<T}
 s_0(r,k)^2e^{2ik(t-1)}q(r)q(t)\dd t\dd r.
 \label{eq:QBC-radial-return}
\end{equation}

\begin{theorem}[Radial two-slab obstruction to $\mathrm{QBC}_1$]
\label{thm:QBC-two-slab}
Fix $K=[k_0,k_1]\Subset(0,\infty)$, $A>R>1$, $b>0$, and
$0<\delta<\pi/k_1$.  Put
\begin{align}
 a_\Lambda&=\Lambda^{-3/4},
 &B_\Lambda&=A+\Lambda+1,\label{eq:QBC-slab-parameters}\\
 q_\Lambda(r)
 &=a_\Lambda\mathbf1_{[A,A+\Lambda]}(r)
 +b\mathbf1_{[B_\Lambda,B_\Lambda+\delta]}(r).
 \label{eq:QBC-slab-potential}
\end{align}
Then
\begin{equation}
 \int_R^\infty|q_\Lambda(r)|^2\dd r
 =\Lambda^{-1/2}+b^2\delta,
 \label{eq:QBC-slab-cost}
\end{equation}
but, with $T_\Lambda=B_\Lambda+\delta$,
\begin{equation}
 \int_K[\mathcal Q_{R,1,T_\Lambda}(k)]_+\dd k
 \longrightarrow\infty.
 \label{eq:QBC-slab-divergence}
\end{equation}
Thus~\eqref{eq:QBC1} is false even for one radial channel and arbitrarily
small fixed $L^2$ tails.
\end{theorem}

\begin{proof}
Decompose~\eqref{eq:QBC-radial-return} into the two self-interactions and the
inner--outer cross term.  The cross term is
\begin{equation}
 \mathcal C_\Lambda(k)
 =-a_\Lambda b I_\Lambda(k)\Re J_\Lambda(k),
\end{equation}
where, uniformly on $K$,
\begin{align}
 I_\Lambda(k)
 &=\int_A^{A+\Lambda}\frac{\sin^2(k(r-1))}{k^2}\dd r
 =\frac{\Lambda}{2k^2}+O_K(1),\\
 J_\Lambda(k)
 &=\int_{B_\Lambda}^{B_\Lambda+\delta}e^{2ik(t-1)}\dd t
 =e^{i(2k(B_\Lambda-1)+k\delta)}\frac{\sin(k\delta)}{k}.
\end{align}
Hence
\begin{equation}
 \mathcal C_\Lambda(k)
 =-\Lambda^{1/4}\frac{b\sin(k\delta)}{2k^3}
   \cos(\lambda_\Lambda k)+O_K(\Lambda^{-3/4}),
 \qquad
 \lambda_\Lambda=2(B_\Lambda-1)+\delta.
\end{equation}
The continuous weight is positive on $K$.  Since the periodic function
$[-\cos\theta]_+$ has mean $1/\pi$,
\begin{equation}
 \int_K[\mathcal C_\Lambda(k)]_+\dd k
 =\frac{b\Lambda^{1/4}}{2\pi}
   \int_K\frac{\sin(k\delta)}{k^3}\dd k+o(\Lambda^{1/4}).
\end{equation}
The inner self-interaction has $L^1(K)$ norm $O_K(\Lambda^{-1/2})$ and the
fixed outer self-interaction has norm $O_{K,b,\delta}(1)$, by integrating the
outer oscillation first in~\eqref{eq:QBC-radial-return}.  Since
$[x+y]_+\geq[x]_+-|y|$, the cross growth proves
\eqref{eq:QBC-slab-divergence}.  Multiplying every $q_\Lambda$ by a fixed
$\varepsilon>0$ multiplies both sides of the proposed quadratic comparison by
$\varepsilon^2$, so small tail norm does not repair the estimate.
\end{proof}

The obstruction is nonuniform perturbation theory, not physical growth.  The
long inner slab accumulates a forward phase of size $\Lambda^{1/4}$; expanding
that phase in the coupling before composing it with the outer reflector creates
the divergent coefficient.

\begin{theorem}[Exact radial scalar entropy at physical coupling]
\label{thm:scalar-physical-entropy}
Let $q\in L^2(1,\infty)$ be real and compactly supported, let $J_q(k)$ be its
half-line Jost function, and let $K=[k_0,k_1]\Subset(0,\infty)$.  Then
\begin{equation}
 \int_K\log^+|J_q(k)|\dd k
 \leq |K|\log2+\frac{\pi}{8k_0^2}
 \int_1^\infty|q(r)|^2\dd r.
 \label{eq:scalar-physical-entropy}
\end{equation}
The constant is independent of the support radius.
\end{theorem}

\begin{proof}
Extend $q$ by zero to the left of $r=1$.  To the left of its support, the
outgoing solution has the form
$a_q(k)e^{ik(r-1)}+b_q(k)e^{-ik(r-1)}$.  Flux conservation gives
$|a_q|^2-|b_q|^2=1$, so
\begin{equation}
 |J_q(k)|\leq |a_q(k)|+|b_q(k)|\leq2|a_q(k)|.
\end{equation}
The Deift--Killip transmission inequality~\cite{DeiftKillip1999} states
\begin{equation}
 \int_\R k^2\log|a_q(k)|\dd k
 \leq\frac{\pi}{8}\int_\R|q(r)|^2\dd r.
\end{equation}
Since $\log|a_q|\geq0$ and $k\geq k_0$ on $K$, the claimed estimate follows.
\end{proof}

In particular, the exact entropy of the two-slab family in
Theorem~\ref{thm:QBC-two-slab} is uniformly bounded although its quadratic
Taylor coefficient diverges.  Therefore no coefficientwise coupling series
whose second-order bound implies~\eqref{eq:QBC1} can be the endpoint proof.

The direct replacement is the physical-coupling exterior frame.  With the
relative Jost matrix $\mathcal J_{T,L}$ and WKB evolution
$\mathcal U_{T,L}$ from~\eqref{eq:relative-wkb-gauge}, define
\begin{equation}
 \widetilde{\mathfrak g}_{P;T,L}(k)
 :=\wedgeN\!\left(
 \mathcal U_{T,L}(k)^{-1}\mathcal J_{T,L}(k)^{-1}P_N
 \right).
 \label{eq:physical-gauged-exterior-frame}
\end{equation}
On the real axis the gauge is unitary, so
\begin{equation}
 \log\widetilde\Gamma_{P;T,L}(k)
 :=-\log\|\wedgeN(\mathcal J_{T,L}(k)^{-1}P_N)\|
 =-\log\|\widetilde{\mathfrak g}_{P;T,L}(k)\|.
\end{equation}
The free diagonal normalization changes this from $\log\Gamma_P$ only by a
fixed bounded term on $K$.

\begin{problem}[Physical-coupling exterior Hardy theorem]
\label{prob:physical-exterior-Hardy}
For every fixed $N$ and compact $K\Subset(0,\infty)$, prove under the small-tail
hypothesis, uniformly in $R,L,T$, an interior lower bound
\begin{equation}
 \|\widetilde{\mathfrak g}_{P;T,L}(k_*)\|\geq c_{K,N}>0
\end{equation}
and the boundary estimate
\begin{equation}
 \int_{\partial\Omega_K}
 \|\widetilde{\mathfrak g}_{P;T,L}(k)\|^2\dd\omega_*(k)
 \leq C_{K,N}\bigl(1+\cE_N(V)\bigr).
 \label{eq:physical-exterior-Hardy}
\end{equation}
Subharmonicity and Proposition~\ref{prop:dual-plucker} would then give
$\mathrm{PSR}_N$.  The first route to~\eqref{eq:physical-exterior-Hardy} is
Proposition~\ref{prop:physical-Riccati-descent} followed by
Problem~\ref{prob:Riccati-square-function} and the dressed estimate
\eqref{eq:DCE}; every estimate must be made directly at $\alpha=1$.
\end{problem}

\subsection{Module C: buffered impact-parameter channels}

Let $K=[k_0,k_1]$ and divide the exterior into dyadic shells
\begin{equation}
 \mathcal S_j=[R_j,2R_j],
 \qquad R_j=2^jR_0.
 \label{eq:shells}
\end{equation}
Choose a buffered spectral projector
\begin{equation}
 P_j=\mathbf1_{\{A_d\leq \Lambda^2R_j^2\}}(A_d),
 \qquad
 \Lambda^2>4(k_1^2+\|V\|_\infty+1).
 \label{eq:moving-proj}
\end{equation}
For $r\in\mathcal S_j$ and $Q_j=\Id-P_j$,
\begin{equation}
 Q_j\left(\frac{A_d}{r^2}+M_V-k^2\right)Q_j
 \succeq Q_j,
 \qquad k\in K,
 \label{eq:elliptic-high}
\end{equation}
as a quadratic form.  Thus modes above the buffer are elliptic on the shell and can be integrated out without a real-axis resonance.  Projectors are held constant on each shell, so no derivative $P_j'(r)$ is created.

Proposition~\ref{prop:physical-Riccati-descent} identifies the exact shell
increment, but it does not estimate its positive energy part.  Let
$\Delta_j(k)$ denote the inward shell increment
\begin{equation}
 \Delta_j(k)
 :=\log\Gamma_P(F(R_j,k))-\log\Gamma_P(F(2R_j,k)).
 \label{eq:oriented-shell-increment}
\end{equation}
The dimension-free target is
\begin{equation}
 \int_K\Delta_j(k)^+\dd k
 \leq C_{K,N}\int_{R_j}^{2R_j}\|V(r,\cdot)\|_2^2\dd r
 +\mathcal F_j-\mathcal F_{j+1}
 \label{eq:SHELL}
\end{equation}
where $\mathcal F_j\geq0$ is a flux correction that telescopes.  This is the
dyadic form of~\eqref{eq:Riccati-square-function}.  Here $\mathcal F_j$ must be
built from the dressed Riccati/Feshbach state; the exact cancellation says that
$V_{QQ}$ may influence this state but may not appear as an additional cost.
Summing~\eqref{eq:SHELL} gives~\eqref{eq:PSR}.

\begin{remark}
The buffer must be applied to the dynamics, but the cost must remain the full scalar norm.  Replacing $Q_N$ by $P_j-P_N$ on the right of~\eqref{eq:SHELL} would reintroduce the exact aliasing failure~\eqref{eq:alias-v}--\eqref{eq:alias-cost}.
\end{remark}

\subsection{Module D: normalized entropy plus angular tightness}

This module begins with the exact identity~\eqref{eq:normalized-full-column}.  Its two steps are deliberately separated:
\begin{align}
 \text{dimension-linear matrix sum rule}
 &\Longrightarrow \mathrm{NMSR},
 \label{eq:norm-route-1}\\
 \mathrm{NMSR}+\text{buffered angular tightness}
 &\Longrightarrow \mathrm{TIGHT}_N
 \Longrightarrow \mathrm{FCEN}_N.
 \label{eq:norm-route-2}
\end{align}
The likely tightness mechanism is an impact-parameter buffer combined with a determinant chain rule.  If $W_L>0$ and $P_a<P_b$, then
\begin{equation}
 \log\det W_{P_b}
 =\log\det W_{P_a}
 +\log\det\!\left(W_{P_b}/W_{P_a}\right),
 \label{eq:schur-chain}
\end{equation}
where the quotient is the positive Schur complement.  One must show that entropy cannot be carried indefinitely by successive high-degree Schur complements without spending the scalar energy in~\eqref{eq:normalized-full-column}.  Equation~\eqref{eq:dimension-escape} shows why a general matrix argument cannot do this; the proof must use the centrifugal ordering and scalar multiplication structure.

\subsection{Module E: the solved raywise operator and one unbounded short-range theorem}

Remove the angular kinetic term but retain the potential exactly:
\begin{equation}
 H_{\mathrm{ray}}
 =-\partial_r^2+M_{V(r,\omega)}
 =\int_{\Sd}^{\oplus}
   \bigl(-\partial_r^2+V(r,\omega)\bigr)\dd\omega,
 \label{eq:H-ray}
\end{equation}
with a Dirichlet boundary condition at $r=R$.  This comparison problem is already solved.

\begin{proposition}[Infinite multiplicity for the raywise comparison]\label{prop:raywise}
If $V$ is real, bounded, and satisfies~\eqref{eq:polar-cost}, then
\begin{equation}
 m_{\ac}(E;H_{\mathrm{ray}})=\infty
 \quad\text{for a.e. }E>0.
 \label{eq:raywise-ac}
\end{equation}
\end{proposition}

\begin{proof}
Fubini gives $V(\cdot,\omega)\in L^2(R,\infty)$ for a.e. $\omega$.  The scalar half-line $L^2$ theorem gives a positive absolutely continuous density for a.e. $E>0$ in almost every such fiber.  A second application of Fubini makes the exceptional set independent of $\omega$: for a.e. $E>0$, the density is positive for a.e. $\omega$.  The multiplicity fiber of~\eqref{eq:H-ray} then contains the infinite-dimensional space $L^2(\Sd,w_\omega(E)\dd\omega)$.
\end{proof}

The original exterior operator is the form sum
\begin{equation}
 H_{\mathrm{ext}}=H_{\mathrm{ray}}+\frac{A_d}{r^2}.
 \label{eq:ray-comparison}
\end{equation}
The coefficient $r^{-2}$ is integrable at infinity, but $A_d$ is unbounded and $M_V$ has no angular regularity.  The entire conjecture would therefore follow from one comparison theorem.

\begin{problem}[Unbounded angular short-range comparison]\label{prob:raywise-wave}
Prove existence and completeness of
\begin{equation}
 W_\pm\!\left(H_{\mathrm{ray}}+r^{-2}A_d,H_{\mathrm{ray}}\right).
 \label{eq:ray-wave}
\end{equation}
For bounded angular cutoffs the perturbation is short range.  The new content is a cutoff-uniform estimate through the domains of $A_d$ and $H_{\mathrm{ray}}$; an evanescent-channel bound of the type~\eqref{eq:elliptic-high} is the natural compactness mechanism.
\end{problem}

\subsection{Module F: stochastic angular paths and quadratic variation}

Let $(\Omega_t)_{t\geq0}$ be spherical Brownian motion started in its invariant uniform distribution, and set
\begin{equation}
 \tau(r)=\int_R^r s^{-2}\dd s=R^{-1}-r^{-1}.
 \label{eq:angular-time}
\end{equation}
Stationarity and Tonelli give the exact path-cost identity
\begin{equation}
 \mathbb E\int_R^\infty
 |V(r,\Omega_{\tau(r)})|^2\dd r
 =\frac1{|\Sd|}
  \int_R^\infty\|V(r,\cdot)\|_2^2\dd r.
 \label{eq:path-cost}
\end{equation}
Thus Simon's norm has exactly the form of an expected quadratic-variation budget along a stationary angular path.  This is the continuum analogue of path averaging in tree sum rules and operates directly at coupling one.

\begin{problem}[Stochastic exterior-power entropy]\label{prob:stochastic}
Construct the antisymmetric $N$-path evolution associated with~\eqref{eq:angular-time} and prove an inequality of the form
\begin{equation}
 \int_I\log\det W_N(E)\dd E
 \geq-C_{I,N}
 -C_I\,\mathbb E\sum_{j=1}^N
  \int_R^\infty|V(r,\Omega_{\tau(r)}^{(j)})|^2\dd r.
 \label{eq:stochastic-entropy}
\end{equation}
The obstruction is conceptual and explicit: Feynman--Kac controls imaginary-time propagation, whereas~\eqref{eq:stochastic-entropy} concerns real-energy boundary values.  A successful martingale or Girsanov argument must retain the exterior-power phase rather than replacing the unitary dynamics by a heat estimate.
\end{problem}

\section{Endpoint device: fixed coupling without spectral averaging}

The fixed-coupling counterexample~\eqref{eq:fiber-counterexample} prevents the inference ``almost every $\alpha$ implies $\alpha=1$.''  A different use of the coupling parameter may nevertheless be viable.  For a Hilbert--Schmidt Birman--Schwinger operator $K(z)$,
\begin{equation}
 D_2(\zeta,z)=\det{}_2(\Id+\zeta K(z))
 \label{eq:det2}
\end{equation}
is entire in $\zeta$ and satisfies
\begin{equation}
 \log|D_2(\zeta,z)|
 \leq\frac{|\zeta|^2}{2}\|K(z)\|_{\HS}^2.
 \label{eq:det2-growth}
\end{equation}
Subharmonic mean-value at the point $\zeta=1$ can therefore turn a complex-neighbourhood estimate into a fixed-coupling estimate; no specialization of an almost-every-real-coupling statement is involved.

For the present problem, a formal relative perturbation is
\begin{equation}
 K_{N,L}(r)=
 \begin{pmatrix}
 V_{PP}(r)&V_{PQ}(r)\\
 V_{QP}(r)&0
 \end{pmatrix},
 \qquad
 \|K_{N,L}(r)\|_{\HS}^2
 =\|V_{PP}\|_{\HS}^2+2\|V_{QP}\|_{\HS}^2.
 \label{eq:relative-K}
\end{equation}
The high-high block is placed in the reference operator.  This bookkeeping is attractive, but a general Hilbert--Schmidt perturbation determinant can fail to have nontangential boundary values~\cite{GPS2007}.  Moreover Theorem~\ref{thm:matrix-rotator-counterexample} proves that an arbitrary high block recreates the channel-rotator obstruction.  The determinant route is therefore conditional on a Schr\"odinger-specific boundary theorem that uses the scalar angular compression and its full leakage.

There is nevertheless an exact reason to test this route.  With
\begin{equation}
 \mathcal V=\begin{pmatrix}V_{PP}&V_{PQ}\\V_{QP}&V_{QQ}\end{pmatrix},
 \qquad
 \mathcal V_0=\begin{pmatrix}0&0\\0&V_{QQ}\end{pmatrix},
\end{equation}
the quadratic trace difference is
\begin{equation}
 \tr(\mathcal V^2)-\tr(\mathcal V_0^2)
 =\|V_{PP}\|_{\HS}^2+2\|V_{QP}\|_{\HS}^2
 =\|K_{N,L}\|_{\HS}^2.
 \label{eq:relative-trace-coefficient}
\end{equation}
Thus the high-energy coefficient of a formal relative matrix sum rule deletes
$V_{QQ}$ exactly and leaves a quantity equivalent, with dimension-independent
constants, to the charged column
$\|\mathcal VP\|_{\HS}^2=\|V_{PP}\|_{\HS}^2+
\|V_{QP}\|_{\HS}^2$.  Existing matrix $L^2$ sum rules charge the full matrix
potential~\cite{MolchanovVainberg2004}; they do not supply this relative
boundary statement with an arbitrary high-channel reference.

Proposition~\ref{prop:minimum-scalar-block-lift} now shows that even imposing
scalar compression does not validate this bookkeeping by itself.  A
concentrated high harmonic permits an arbitrary $V_{QQ}$ entry with an
$o(1)$ addition to the full constant-column cost, and Proposition
~\ref{prop:thin-belt-barrier} makes the high state almost invariant.  Hence the
relative trace coefficient contains no purely algebraic obstruction to the
rotator/barrier.  Any positive relative theorem must use the dynamics of the
converter, including the centrifugal flow and the scalar dilation into the
whole high sector.

\begin{problem}[Relative matrix Schr\"odinger sum rule]
\label{prob:relative-matrix-sum-rule}
Assume specifically that
$\mathcal V=P_LM_vP_L$ is the finite angular compression of one real scalar
multiplication potential, and let $\mathcal V_0$ retain only its $QQ$ block.
Construct the Schr\"odinger-specific relative Jost or perturbation determinant
for $(\mathcal V,\mathcal V_0)$ and prove that its boundary entropy controls
$\log^+\Gamma_P$ while its trace side is bounded by the full scalar column
$\|M_vP_N\|_{\HS}^2$; the finite part is
~\eqref{eq:relative-trace-coefficient}.  The constants must be uniform in
support radius and angular model dimension.  The proof must show that scalar
compression and the full angular dynamics prevent the converter in
Problem~\ref{prob:scalar-converter-dichotomy}; compressed trace algebra and the
barrier cost alone are insufficient.  It must also retain the full
Grassmannian quotient.
\end{problem}

\begin{problem}[Fixed-coupling subharmonic transfer]\label{prob:subharmonic}
Construct a scalar analytic quantity $D_{P}(\zeta,z)$ such that
\begin{enumerate}[label=(\roman*)]
 \item $\log\Gamma_{P;R,L}(k)\leq\log^+|D_P(1,k+i0)|+C$;
 \item $D_P$ has a complex-coupling growth bound controlled by $\cE_N(V)$;
 \item its boundary values and zero terms are uniform in $R,L$.
\end{enumerate}
Then apply the subharmonic mean-value inequality in a disk centred at $\zeta=1$ to obtain the physical endpoint directly.
\end{problem}

\section{Consequences for the selected-channel approach}

The shortest fully rigorous completion route is now
\begin{align}
 \|V\|_{X_d}^2<\infty
 &\overset{\text{Prop.~\ref{prop:small-tail}}}{\Longrightarrow}
 \text{small physical tail},
 \label{eq:logic-1}\\
 \text{small tail}+\mathrm{PSR}_N
 &\overset{\text{Thm.~\ref{thm:plucker}}}{\Longrightarrow}
 \sup_{R,L}\int_I\log^-\det W_N^{R,L}<\infty,
 \label{eq:logic-2}\\
 \text{strong resolvent convergence}+\text{entropy l.s.c.}
 &\Longrightarrow
 \int_I\log^-\det W_N<\infty,
 \label{eq:logic-3}\\
 \det W_N(E)>0\ \text{a.e.}
 &\Longrightarrow m_{\ac}(E;H)\geq N\ \text{a.e.},
 \label{eq:logic-4}\\
 N\in\N\ \text{and rational }I
 &\Longrightarrow m_{\ac}(E;H)=\infty\ \text{a.e. on }(0,\infty).
 \label{eq:logic-5}
\end{align}
The alternative normalized route replaces $\mathrm{PSR}_N$ by $\mathrm{NMSR}+\mathrm{TIGHT}_N$.  The raywise route bypasses entropy entirely if Problem~\ref{prob:raywise-wave} is solved, while Problem~\ref{prob:stochastic} would supply the entropy estimate directly from the exact path cost~\eqref{eq:path-cost}.

Theorem~\ref{thm:fixed-potential-riesz-counterexample}, proved in the final
appendix, now shows that even the buffered fixed-potential selected $N=1$
replacement for the first route is false.  Thus~\eqref{eq:logic-1}--
\eqref{eq:logic-5} remains a correct conditional chain, but its selected-probe
input cannot hold universally.  The normalized, moving-probe, full-operator,
raywise, and stochastic routes are not decided by that counterexample.

\subsection{Remaining open directions}

The following order maximizes falsifiability and minimizes dependence on unproved machinery.

\begin{enumerate}[label=\textbf{Stage \arabic*.},leftmargin=*]
 \item \textbf{Free collar and exact $N=1$ column complete.}  Proposition~\ref{prop:free-collar} proves the multiplicity-preserving bridge with $q_d=\lceil d/2\rceil$, while Proposition~\ref{prop:scalar-plucker} identifies the exact scalar column and cost $|\Sd|^{-1}\|V\|_{X_d}^2$.
 \item \textbf{Zero-phase audit and WKB gauge complete.}  Propositions~\ref{prop:raw-bessel-false} and~\ref{prop:RCH1-false} rule out the raw same-shell estimate and the unrenormalized interior normalization.  Equations~\eqref{eq:forward-wkb-evolution}--\eqref{eq:gauged-radiation-column} give the operator correction, Proposition~\ref{prop:gauged-vector-hardy} gives the invariant entropy extraction, and Proposition~\ref{prop:gauged-first-order} proves its first-order $H^2$ bound.
 \item \textbf{Coefficientwise quadratic route decided.}  Theorem~\ref{thm:QBC-two-slab} disproves $\mathrm{QBC}_1$ even in the scalar radial model, while Theorem~\ref{thm:scalar-physical-entropy} proves that the corresponding exact physical entropy remains uniformly bounded.  Coupling-Taylor control is therefore removed from the endpoint route.
 \item \textbf{Dual geometry, physical descent, and the generic no-go theorem complete.}  Proposition~\ref{prop:dual-plucker} removes the distinguished-minor denominator, and Propositions~\ref{prop:differential-flux} and~\ref{prop:physical-Riccati-descent} identify the exact high--high cancellation in the physical outgoing frame.  Proposition~\ref{prop:cocycle-low-column-no-go} proves that this cancellation alone cannot yield a low-column estimate.
 \item \textbf{Finite matrix rotator/barrier decided.}  Lemma~\ref{lem:exact-rotator-conversion}, the rapid-phase good-set lemma, and Theorem~\ref{thm:matrix-rotator-counterexample} prove that the unrestricted two-channel matrix low-column endpoint is false.
 \item \textbf{Scalar $\mathrm{PSR}_1$ decided negatively.}  Proposition~\ref{prop:exact-cutoff-edge-lift}, Lemma~\ref{lem:exact-cutoff-edge-barrier}, and Theorem~\ref{thm:scalar-PSR1-counterexample} give an exact scalar rotator/barrier in the finite cutoff with bounded full-column cost and divergent quotient entropy.  The proposed uniform $\mathrm{PSR}_N$ route is therefore closed already at $N=1$.
 \item \textbf{Fixed-potential buffering decided negatively.}  Theorem~\ref{thm:fixed-potential-riesz-counterexample} produces one bounded mean-zero $X_3$ potential whose selected entropy diverges along every cofinal positive smoothing, radial truncation, and double-buffered cutoff.  Buffering removes the algebraic cutoff edge but does not prevent a genuine full-space reciprocal cascade.
 \item \textbf{Full-space scalar converter decided positively.}  Lemma~\ref{lem:translated-barrier-hardy}, Propositions~\ref{prop:robust-reciprocal-cascade} and~\ref{prop:talbot-riesz-concentration}, and Proposition~\ref{prop:scalar-reciprocal-riesz-cell} give concentration, localized reflection, reciprocal synchronization, radial realization, and refocusing at summable charged cost.
 \item \textbf{Do not replace the boundary theorem by one imaginary-energy row.}  Proposition~\ref{prop:negative-energy-row-no-go} gives both vanishing-cost single slabs and a fixed finite-cost slab train for which the negative-energy forward row diverges.  The rational second-trace boundary localization is essential; Poisson subharmonicity carries an uncontrolled singular mass at infinity.
 \item \textbf{The transport-free scalar endpoint is complete but cannot be continued uniformly.}  Proposition~\ref{prop:decomposable-scalar-entropy} proves the selected entropy after deleting $A_3/r^2$.  The Riesz--Talbot construction shows that centrifugal angular transport is not a perturbative remainder: it enables the counterexample.
 \item \textbf{Signed centrifugal interpolation is an identity, not an endpoint bound.}  Proposition~\ref{prop:exact-centrifugal-interpolation} and its pure-gauge cancellation remain exact, but the universal estimate $\mathrm{SCI}_1$ is false by Corollary~\ref{cor:fixed-potential-route-closed}.  The absolute-variation replacement was already algebraically false.
 \item \textbf{Keep negative-energy action as a budget, not an entropy theorem.}  Proposition~\ref{prop:negative-energy-centrifugal-action} proves a dimension-free charged bound for the full positive centrifugal resolvent action.  Remark~\ref{rem:resolvent-order-not-entropy} gives a monotone Stieltjes family with divergent entropy, so the missing step must retain real-axis local propagation.
 \item \textbf{The angular heat endpoint is also decided negatively.}  Proposition~\ref{prop:angular-heat-flow-interpolation} and its exact $X_3$ heat budget remain valid, but the universal product-Fisher estimate $\mathrm{HFI}_1$ would imply the false selected entropy theorem.
 \item \textbf{Audit the alternative routes under the new falsification test.}  The normalized entropy plus tightness, raywise comparison, stochastic path, and direct full-operator determinant routes must be formulated without uniform cutoff-edge compression.
 \item \textbf{Move to a basis-free endpoint.}  Entropy lower semicontinuity, Proposition~\ref{prop:small-tail}, and the free-collar Birman--Kuroda bridge remain valid finishing tools, but no theorem controlling one fixed selected probe can supply their hypothesis.  The next viable targets are moving-probe or full-operator entropy, or a raywise/non-entropic mechanism which cannot lose all mass by rotation of the good angular directions.
\end{enumerate}

\section{Obstructions to stronger intermediate estimates}

Every module must pass the following tests before it can be promoted to a lemma.

\begin{longtable}{>{\raggedright\arraybackslash}p{0.26\textwidth}>{\raggedright\arraybackslash}p{0.32\textwidth}>{\raggedright\arraybackslash}p{0.33\textwidth}}
\toprule
Candidate step & Required test & Failure signal \\
\midrule
Pl\"ucker outer factor & Treat the norm of all minors; audit zeros, poles, and high-$k$ coefficients. & Division by one minor or an uncontrolled destructive zero term. \\
Hardy interior normalization & Remove the operator WKB evolution and use the full vector norm as in~\eqref{eq:GVH1}. & The long weak slab~\eqref{eq:RCH1-slab} makes an unrenormalized interior coordinate exponentially small at vanishing $L^2$ cost. \\
Coefficientwise Born expansion & Test the combined quotient coefficient on two separated radial slabs before adding angular or Airy structure. & Theorem~\ref{thm:QBC-two-slab} gives divergent positive quadratic mass at bounded, arbitrarily small $L^2$ cost; this route is rejected. \\
Physical Riccati descent & Recover $\Gamma_P(F)$ from $P\Im\mathcal RP$ and audit the differential generator in the $4\times4$ model. & A direct charge involving $V_{QQ}$, or failure to match the physical Jost quotient. \\
Unrestricted cocycle square function & Test a high-channel amplifier followed by a slow quarter-turn. & Proposition~\ref{prop:cocycle-low-column-no-go} produces growth $H$ at charged cost $\pi^2/(4L)$; this route is rejected. \\
Uniform scalar-compression square function & Test a degree-$2L$ multiplier against the degree-$L$ model. & Proposition~\ref{prop:exact-cutoff-edge-lift} and Lemma~\ref{lem:exact-cutoff-edge-barrier} give the exact scalar rotator/barrier counterexample. \\
Self-energy averaging & Keep the energy-dependent dressed solution inside the estimate. & Applying~\eqref{eq:poisson-identity} as if the test vector were independent of $E$. \\
Moving cutoff & Charge the full scalar potential, use a bandwidth buffer, and justify the order of limits for one fixed potential. & Full-column charging alone does not stop a pure degree-$2L$ multiplier from producing an exact high-only barrier in $P_{\leq L}$. \\
Normalized entropy & Prove $\mathrm{TIGHT}_N$ for a fixed subspace. & Good directions drift with $L$ as in~\eqref{eq:dimension-escape}. \\
Raywise comparison & Control the unbounded $A_d$ uniformly in angular cutoff. & Treating $r^{-2}A_d$ as a bounded short-range potential. \\
Stochastic path sum rule & Preserve real-energy exterior-power phases. & Substituting an imaginary-time Feynman--Kac estimate for the required boundary entropy. \\
Complex coupling & Use a mean-value estimate centred at $1$. & Specializing an almost-every-real-$\alpha$ theorem to $\alpha=1$. \\
Relative determinant & Establish Schr\"odinger boundary values. & Appealing to general $\det_2$ theory, contradicted by known Hilbert--Schmidt examples. \\
Full-space scalar channel rotator & Test the full converter into a concentrated highest-weight packet; verify the nodal leakage, angular-current, radial-action, nonlinear Fourier, pulse-area, roughness, and high-band recoil bounds. & These tests constrain a cutoff-stable converter, but they cannot rescue uniform $\mathrm{PSR}_1$ from the exact cutoff-edge lift. \\
Limit $R,L\to\infty$ & Use entropy lower semicontinuity on inner intervals. & Claiming that strong resolvent convergence preserves a.c. spectrum. \\
\bottomrule
\end{longtable}

\section{A conditional implication for Simon's conjecture}

\begin{theorem}[Completion from the Grassmannian sum rule]\label{thm:completion}
Assume the following for every $N$ and every compact $K\Subset(0,\infty)$:
\begin{enumerate}[label=(\alph*)]
 \item the finite models~\eqref{eq:finite-H} satisfy $\mathrm{PSR}_N$;
 \item their matrix spectral measures converge vaguely to the exterior measure;
 \item $\int\log^-\det$ is lower semicontinuous under that convergence;
 \item the exterior and full operators have unitarily equivalent absolutely continuous parts.
\end{enumerate}
Then Simon's conclusion~\eqref{eq:simon} holds.
\end{theorem}

\begin{proof}
Corollary~\ref{cor:PSR} and~\eqref{eq:probe-boundary} give a uniform finite-model entropy bound.  Assumptions (b)--(c) pass it to the exterior spectral density $W_N$:
\begin{equation}
 \int_I\log^-\det W_N(E)\dd E<\infty.
\end{equation}
Therefore $\det W_N(E)>0$ for a.e. $E\in I$, so $\rank W_N(E)=N$.  In a direct-integral representation, the rank of a density generated by $N$ vectors is bounded by the a.c. fiber multiplicity; hence $m_{\ac}(E)\geq N$ a.e. on $I$.  Assumption (d) transfers the conclusion to the full operator.  Intersect the full-measure sets for all $N$ and all compact intervals with rational endpoints.
\end{proof}

\section{Conclusion}

The audit has moved the project beyond the unsafe truncated-column formulation.  The central open object is no longer an unspecified determinant of a large matrix.  It is the low-channel Grassmannian amplification
\begin{equation}
 \Gamma_P(k)
 =\frac{|\det J(k)|}{\|\wedge^{m-N}(QJ(k))\|},
\end{equation}
whose logarithm is exactly the entropy loss of the low spectral density, up to the free factor.  Its expected quadratic trace coefficient is
\begin{equation}
 \tr(PM_V^2P)=\|M_VP\|_{\HS}^2,
\end{equation}
which is precisely Simon's scalar angular $L^2$ cost when $P$ contains the constant harmonic.

The latest audit changes the principal route.  The raw same-shell estimate,
$\mathrm{RCH}_1$, and now the combined coefficientwise estimate
$\mathrm{QBC}_1$ are all false.  The two-slab obstruction is especially
decisive: its quadratic Taylor coefficient diverges although the exact scalar
physical-coupling entropy is uniformly bounded by
Theorem~\ref{thm:scalar-physical-entropy}.  Therefore the conjecture has not
been disproved; rather, any endpoint proof must resum the forward phase and
work directly at physical coupling.

The exact replacement is now more concrete.  The dual inverse frame
\eqref{eq:dual-exterior-frame} removes zeros and poles of a selected minor,
and the physical Riccati identities
\eqref{eq:physical-Riccati-Gamma}--\eqref{eq:charged-Riccati-column} show that
the high--high potential is absorbed into a dressed state while the only
variable source is the charged column $V_mP_N$.  Proposition
~\ref{prop:cocycle-low-column-no-go} now supplies a crucial negative result:
for an unrestricted left cocycle, a high block can preload amplification and a
later arbitrarily cheap rotation can transfer it into $\Gamma_P$.  Therefore
the differential cancellation is an identity, not by itself an estimate.

Theorem~\ref{thm:matrix-rotator-counterexample} completes the next negative
step: the exact converter and a shorter barrier make the arbitrary two-channel
matrix Grassmannian estimate false, with vanishing low-column cost.
Proposition~\ref{prop:minimum-scalar-block-lift} then removes the hoped-for
purely algebraic rescue: every real $2\times2$ block has an exact scalar lift,
and highest-weight concentration makes the extra full constant-column cost
vanish.  Proposition~\ref{prop:thin-belt-barrier} goes further by scalarizing
the high-only barrier with exponentially small high-state leakage.

The cutoff quantifier supplies the decisive scalar result.  Proposition
~\ref{prop:exact-cutoff-edge-lift} proves that degrees $L+1$ through $2L$ can
cancel every unwanted retained component and make
$\operatorname{span}\{e_0,h_L\}$ exactly reducing.  Lemma
~\ref{lem:exact-cutoff-edge-barrier} identifies the closed-form barrier
$p_{2L}=\sin^{2L}\theta\cos(2L\phi)$: it has zero compressed constant column,
an order-one diagonal entry on $h_L$, and full charged density
$\asymp L^{-1/2}$.  Theorem~\ref{thm:scalar-PSR1-counterexample} combines these
facts with the exact rotator and rapid-phase lemma to disprove scalar
$\mathrm{PSR}_1$ with uniformly bounded scalar amplitude and bounded full
constant-column cost.

This does not touch Simon's conjecture.  It rejects the proposed uniform
sufficient strengthening and identifies the precise failure: multiplier
bandwidth was allowed to track and overtake the Galerkin cutoff.  Any repaired
entropy route must fix one scalar potential first and justify a
bandwidth-buffered or diagonal limit, or work directly in the full operator.
The conditional completion theorem remains correct, but its assumption~(a) is
false as presently formulated.  The full-space converter analysis below is a
stronger diagnostic rather than the gate to $\mathrm{PSR}_1$.

The first part of the converter fork is now decided.  Proposition
~\ref{prop:decomposable-converter-no-go} shows that scalar multiplication
without angular transport has conversion error asymptotic to $\sqrt2$, and
Proposition~\ref{prop:angular-current-requirement} forces every successful
converter to carry an order-one angular flux.  Thus the previously convenient
far-field placement, where $A_3/r^2$ was intended to be a small error, cannot
validate the compressed rotator.  The remaining alternatives are singular but
concrete: prove that the charged scalar column cannot sustain the required
current, including any frequency cascade, or build a converter at the
centrifugal scale where that current is genuinely nonperturbative.

The compressed lift itself is now decided more sharply.  Proposition
~\ref{prop:minimum-lift-complement-leakage} shows that its high-state leakage
exceeds the intended adiabatic gap by a factor of order $\ell^{1/4}$, while
Proposition~\ref{prop:nodal-set-lift-obstruction} proves that every uniformly
bounded alternative scalar lift has leakage comparable to its off-diagonal
coupling.  Thus no perturbation of the two-channel adiabatic argument can close
the construction while keeping the chosen plane approximately invariant.  A
counterexample must instead exploit the entire complementary sector and then
refocus it into $h_\ell$ for both directional modes, uniformly on an energy
set of positive measure.  Conversely, a positive theorem need only exclude
this cascade-and-refocusing mechanism, not arbitrary compressed rotations.

Proposition~\ref{prop:cheap-scalar-phase-mask} prevents one final premature
positive conclusion: the low-column cost does not control angular frequency.
A slab of height $\epsilon$ and length $\epsilon^{-1}$ has vanishing cost but
imprints a fixed bounded angular phase, even when that phase oscillates at an
arbitrarily high degree after centrifugal transport is deleted.  Proposition
~\ref{prop:phase-mask-transport-incompatibility} now supplies the missing radial
check: at $R=O(L)$, making that mask cheap forces unbounded degree-$L$
centrifugal action during the slab, while making it clean forces its scalar
cost to diverge.  Thus a succession of ideal masks and free propagators is only
an auxiliary controllability test, not a radial construction.  The decisive
object is the coupled scalar--centrifugal evolution inside the weak slabs,
uniform in $k$ and reciprocal for the two directional modes.  This is now the
narrowest unresolved converter mechanism.  The resonant-drive calculation
~\eqref{eq:single-energy-resonant-cost} confirms that this mechanism can be
 cheap at a single energy while remaining at $R\asymp\ell$, so no
 pointwise-in-energy converter bound can be sufficient.  But its linearized fixed-band
 amplitude is a Fourier transform, and
 ~\eqref{eq:born-broadband-lower-bound} forces an order-one charged cost.  The
 nonlinear Fourier calculation in Proposition
 ~\ref{prop:SU2-nonlinear-Plancherel} upgrades this to an exact theorem when the
 detuning varies but the coupling potential is fixed.  The physical radial
 parameter scales both terms, and Proposition
 ~\ref{prop:scaled-matrix-ensemble-no-go} shows that no such theorem survives
 for arbitrary matrix controls: the adiabatic rotator is already a uniform-band
 counterexample.  Positivity is now the scalar dividing line.  Proposition
 ~\ref{prop:scalar-ground-adiabatic-no-go} proves that the scalar angular ground
 branch cannot rotate from $e_0$ to $h_\ell$, and even a vanishing endpoint
 matching potential has $o(1)$ overlap with $h_\ell$ at $R=O(\ell)$.  The next
 rotating-frame candidate, a monotone chirped highest-weight pulse, is also
 excluded in its standard uniformly adiabatic regime by Proposition
 ~\ref{prop:chirped-ensemble-cost}: sweeping a continuum of resonances forces
 divergent charged cost.  Proposition
 ~\ref{prop:rotating-ground-ladder-fanout} also excludes the tempting cheap
 adjacent-rung escape: the terminal rotating-frame ground momentum depends on
 $s$ and cannot equal one common $\ell$ on a fixed band.  Proposition
 ~\ref{prop:gap-cancelled-supercascade} rules out finite-depth composites, and
 Corollary~\ref{cor:scalar-pulse-area-roughness} adds the missing star-frame
 pulse area.  A gap-cancelled converter must satisfy
 $\mathcal E_\ell\gtrsim\ell^{-1/4}$ and, if its cost still vanishes, its scalar
 bandwidth must be at least $\ell/\mathcal E_\ell^2\gg\ell$; target-scale
 Lipschitz controls have a positive cost floor.  Proposition
 ~\ref{prop:born-high-band-recoil} then proves that the forced rough scale
 cannot be populated in one broadband Born transition.  Proposition
 ~\ref{prop:scalar-star-frame} isolates what is left: an energy-dependent
 nonlinear bootstrap which first concentrates mass in lower modes, applies the
 rough scalar component only after that concentration, and reverses the entire
 excursion for both directions.  Proposition
 ~\ref{prop:scalar-preconcentration-gate} now proves that the first crossing of
 every fixed spectral cut is charged even after retaining the full nonlinear
 high-sector dynamics, and Proposition
 ~\ref{prop:density-concentration-action} forces centrifugal density action
 $\gtrsim R^2\sqrt\ell$, which scalar phase masks themselves cannot produce.
 Proposition~\ref{prop:wrapped-phase-gate-additivity-no-go} makes the first
 failure exact: one bounded wrapped phase crosses all cuts below $L$ with
 fixed-band high-mode mass and vanishing quadratic cost once centrifugal
 transport is deleted.  Multiscale charged-gate additivity is therefore false.
Proposition~\ref{prop:radial-energy-BV-obstruction} then restores the radial
centrifugal term and excludes all stationary or radially low-variation weak
slabs by monotone angular energy.  Corollary
~\ref{cor:radial-superoscillation-toll} sharpens this to radial variation
$\gtrsim\ell/\mathcal E_\ell^2$ under gap cancellation, the same scale already
forced in angular bandwidth.  Proposition
~\ref{prop:raw-radial-variation-no-go} then shows that raw BV cannot itself be
charged: dynamically inert oscillations inflate it without changing the
propagator.  Proposition~\ref{prop:effective-radial-work-bridge} isolates the
state-correlated radial work and proves the first exact product bridge
$\mathfrak R_v\mathfrak D_s\gtrsim R^2$, or
$\mathcal E\mathfrak D_s\gtrsim R^2/L^2$ at angular bandwidth $L$.  These
statements remain a rigorous analysis of the stronger full-space problem, but
the unresolved doubly super-rough refocusing fork is ancillary after the
finite-cutoff disproof.

The buffered completion theorem is correct as a conditional statement:
positive spherical Fej\'er smoothing and a Galerkin cutoff placed beyond both
bandwidth and support give the correct strong-resolvent limit, and
Corollary~\ref{cor:double-buffer-removes-edge} removes the cutoff edge.  The
final Riesz--Talbot construction nevertheless shows that its selected $N=1$
entropy hypothesis is false for every such diagonal.  It combines an exact
translated-barrier Hardy isometry, exponentially concentrating zonal Talbot
products, fractional-moment reciprocal synchronization, summable-cost radial
realization, and uniform lacunary phase placement.  The resulting single
bounded real mean-zero $V\in X_3$ satisfies
\eqref{eq:fixed-potential-selected-entropy-divergence}.  Thus
$\mathrm{SCI}_1$, $\mathrm{HFI}_1$, and every other universal estimate which
would bound that fixed selected channel are closed routes, although their
interpolation identities remain valid.

This is not a counterexample to Simon's conjecture.  A fixed probe can lose
all of its density while absolutely continuous spectral mass survives in
directions which move with the angular scale.  The theorem therefore changes
the target rather than the conjecture's status: the next endpoint must control
a basis-free portion of the spectral density, permit the good subspace to
move, or avoid selected-channel entropy entirely.  The normalized
full-operator, moving-subspace, raywise, and stochastic routes remain possible;
the fixed selected-channel route does not.

\appendix

\section{Entropy lower semicontinuity used in the limit}

Let $\mu_n\to\mu$ vaguely be positive semidefinite $N\times N$ matrix measures on a neighbourhood of $I=[a,b]$, and write
\begin{equation}
 \mu_n=W_n\dd E+\mu_n^s,
 \qquad
 \mu=W\dd E+\mu^s.
\end{equation}
Define on Hermitian matrices
\begin{equation}
 \varphi(A)=
 \begin{cases}
 \max\{0,-\log\det A\},&A>0,\\
 +\infty,&\text{otherwise}.
 \end{cases}
\end{equation}
The function $\varphi$ is convex, lower semicontinuous, and decreasing in the Loewner order.  For a smooth approximate identity $\rho_\eps$ and an inner interval $I_\delta=[a+\delta,b-\delta]$, Jensen's inequality and positivity of the convolved singular part give
\begin{equation}
 \int_{I_\delta}\varphi((\rho_\eps*\mu_n)(x))\dd x
 \leq\int_I\varphi(W_n(E))\dd E.
\end{equation}
Vague convergence gives pointwise convergence of the mollified matrices.  Fatou's lemma, followed by differentiation of finite matrix measures and then $\delta\downarrow0$, yields
\begin{equation}
 \int_I\log^-\det W(E)\dd E
 \leq\liminf_{n\to\infty}\int_I\log^-\det W_n(E)\dd E.
 \label{eq:lsc}
\end{equation}
The inner interval is essential because vague convergence alone does not control mass arriving at an endpoint.

\section{Exact matrix rotator counterexample and the scalar fork}
\label{app:rotator}

The following calculation closes the channel-rotator mechanism for arbitrary
two-channel matrix potentials.  In a two-channel matrix model, put on
$[0,\Lambda]$
\begin{equation}
 V_{\mathrm{rot}}(r)=\delta_\Lambda
 R(\theta(r))\diag(0,1)R(\theta(r))^*,
 \qquad
 \theta(r)=\omega_\Lambda r,
 \label{eq:exact-rotator-potential}
\end{equation}
where
\begin{equation}
 R(\theta)=\begin{pmatrix}\cos\theta&-\sin\theta\\
                    \sin\theta& \cos\theta\end{pmatrix},
 \qquad
 \omega_\Lambda=\frac{\pi}{2\Lambda},
 \qquad
 \delta_\Lambda=\Lambda^{-2/3}.
\end{equation}
For $P=e_1e_1^*$ its exact truncated-column cost is
\begin{equation}
 \int_0^\Lambda\|V_{\mathrm{rot}}(r)P\|_{\HS}^2\dd r
 =\delta_\Lambda^2\int_0^\Lambda
   \sin^2\!\left(\frac{\pi r}{2\Lambda}\right)\dd r
 =\frac12\Lambda^{-1/3}.
 \label{eq:exact-rotator-cost}
\end{equation}

\begin{lemma}[Uniform two-channel conversion]
\label{lem:exact-rotator-conversion}
Let $K\Subset(0,\infty)$.  For all sufficiently large $\Lambda$, the scattering
matrix of~\eqref{eq:exact-rotator-potential} maps a unit incoming forward wave
in channel one to a transmitted forward wave in channel two, up to a scalar
phase, with total reflected and wrong-channel amplitude
\begin{equation}
 O_K\!\left(\delta_\Lambda+
       \frac{\omega_\Lambda}{\delta_\Lambda}
       +\frac{\Lambda\omega_\Lambda^2}{\delta_\Lambda}\right)
 =O_K(\Lambda^{-1/3}),
 \label{eq:exact-rotator-conversion-error}
\end{equation}
uniformly for $k\in K$.  More strongly, in the normalized free modal basis
\begin{equation}
 e_j^\pm=\binom{e_j}{\pm ik e_j},
 \label{eq:free-directional-modes}
\end{equation}
the full $4\times4$ transfer map differs by the same bound from a
direction-preserving channel swap
$e_j^\pm\mapsto e^{i\phi_{j,\pm}(k)}e_{3-j}^\pm$.
\end{lemma}

\begin{proof}
Let
$\mathsf K=\left(\begin{smallmatrix}0&-1\\1&0\end{smallmatrix}\right)$
and write $u(r)=R(\theta(r))y(r)$.  The equation
$-u''+V_{\mathrm{rot}}u=k^2u$ becomes exactly
\begin{equation}
 -y''-2\omega_\Lambda\mathsf K y'
 +\omega_\Lambda^2y
 +\delta_\Lambda\diag(0,1)y=k^2y.
 \label{eq:exact-rotating-frame}
\end{equation}
Thus the rotating frame has constant coefficients.  Its plane-wave pencil is
\begin{equation}
 p^2\Id-2i\omega_\Lambda p\mathsf K+\omega_\Lambda^2\Id
 +\delta_\Lambda\diag(0,1)-k^2\Id.
 \label{eq:rotator-pencil}
\end{equation}
At $\omega_\Lambda=0$, the two forward wave numbers are $k$ and
$\sqrt{k^2-\delta_\Lambda}$; their separation is comparable to
$\delta_\Lambda$, uniformly on $K$.  Riesz projection perturbation of the
constant $4\times4$ first-order generator therefore changes the forward modal
projections by $O_K(\omega_\Lambda/\delta_\Lambda)$.  The first-order
eigenvalue shift vanishes because $\mathsf K$ is off diagonal, so propagation
over length $\Lambda$ creates phase error
$O_K(\Lambda\omega_\Lambda^2/\delta_\Lambda)$.  Matching the constant
coefficient modes to the free leads contributes $O_K(\delta_\Lambda+
\omega_\Lambda)$.  Finally, the zero-eigenvalue potential direction is
$R(\theta)e_1$; it starts at $e_1$ and ends at $e_2$.  The same projection
calculation applies to all four signs in~\eqref{eq:free-directional-modes}, and
the forward/backward spectral clusters remain separated by a constant
depending only on $K$.  Combining these bounds gives
~\eqref{eq:exact-rotator-conversion-error} and its full modal version.
\end{proof}

\begin{lemma}[Rapid-phase good set]\label{lem:rapid-phase-good-set}
Let $K=[k_0,k_1]$.  Suppose $A_\Lambda,B_\Lambda\in C^1(K)$,
$M_\Lambda>0$, and $\psi_\Lambda\in C^2(K)$ obey
\begin{align}
 cM_\Lambda&\leq\max\{|A_\Lambda|,|B_\Lambda|\}
 \leq CM_\Lambda,\qquad
 |(\log M_\Lambda)'|\leq C\log\Lambda, \label{eq:phase-lemma-size}\\
 |A_\Lambda'|+|B_\Lambda'|&\leq
 C(\log\Lambda)M_\Lambda,\qquad
 c\Lambda\leq|\psi_\Lambda'|\leq C\Lambda,\qquad
 |\psi_\Lambda''|\leq C\Lambda. \label{eq:phase-lemma-variation}
\end{align}
Then there are $c_0,c_1>0$, independent of $\Lambda$, such that for all large
$\Lambda$,
\begin{equation}
 \left|\left\{k\in K:
 |A_\Lambda(k)e^{i\psi_\Lambda(k)}+B_\Lambda(k)|
 \geq c_0M_\Lambda(k)\right\}\right|
 \geq c_1|K|.
 \label{eq:rapid-phase-good-set}
\end{equation}
\end{lemma}

\begin{proof}
Partition $K$, apart from two endpoint pieces, into intervals on which
$\psi_\Lambda$ traverses $2\pi$.  Their lengths are comparable to
$\Lambda^{-1}$.  On each such interval the three coefficients vary by
$O((\log\Lambda)/\Lambda)$ relative to $M_\Lambda$.  If
$\max\{|A|,|B|\}\geq cM$, elementary circle geometry shows that
$|Ae^{i\vartheta}+B|\geq c_0M$ on a fixed positive proportion of
$\vartheta\in[0,2\pi]$, uniformly in the ratio $A/B$.  The change of variables
$\vartheta=\psi_\Lambda(k)$ and
~\eqref{eq:phase-lemma-variation} preserve a fixed proportion.  Summing the
complete intervals proves~\eqref{eq:rapid-phase-good-set}; the two endpoint
pieces have total length $O(\Lambda^{-1})$.
\end{proof}

\begin{theorem}[Failure of the arbitrary matrix low-column bound]
\label{thm:matrix-rotator-counterexample}
Let $K=[k_0,k_1]\Subset(0,\infty)$ and $P=e_1e_1^*$.  There is a uniformly
bounded sequence of compactly supported real symmetric $2\times2$ matrix
potentials $\mathcal V_\Lambda$ such that
\begin{equation}
 \int_0^\infty\|\mathcal V_\Lambda(r)P\|_{\HS}^2\dd r
 =\frac12\Lambda^{-1/3}\longrightarrow0,
 \label{eq:matrix-rotator-total-cost}
\end{equation}
but for their outgoing Jost matrices,
\begin{equation}
 \int_K\log^+\Gamma_P(F_\Lambda(0,k))\dd k
 \geq c_K\log\Lambda-C_K\longrightarrow\infty.
 \label{eq:matrix-rotator-entropy-growth}
\end{equation}
Thus no uniform Grassmannian entropy estimate for arbitrary matrix
Schr\"odinger potentials can be charged only to the low matrix column.
\end{theorem}

\begin{proof}
Use~\eqref{eq:exact-rotator-potential} on $[0,\Lambda]$ and follow it by
$\diag(0,Q_0)$, where $Q_0>k_1^2$.  Write
\begin{equation}
 \kappa_{\min}=\sqrt{Q_0-k_1^2},\qquad
 \kappa_{\max}=\sqrt{Q_0-k_0^2},
\end{equation}
choose
\begin{equation}
 0<\beta<\frac{\kappa_{\min}}{3\kappa_{\max}},
 \qquad
 T_\Lambda=\frac{\beta}{\kappa_{\min}}\log\Lambda,
 \label{eq:shorter-barrier-choice}
\end{equation}
and put the barrier on $[\Lambda,\Lambda+T_\Lambda]$.  The barrier has zero
$P$-column, so~\eqref{eq:matrix-rotator-total-cost} follows from
~\eqref{eq:exact-rotator-cost}.

For $\kappa(k)=\sqrt{Q_0-k^2}$, backward propagation of the outgoing high
channel through the barrier gives, in the free directional basis at its left
edge, coefficients
\begin{align}
 A_\Lambda(k)
 &=e^{ikT_\Lambda}\left[
 \cosh(\kappa T_\Lambda)
 +\frac{i}{2}\left(\frac{\kappa}{k}-\frac{k}{\kappa}\right)
 \sinh(\kappa T_\Lambda)\right],\label{eq:barrier-A}\\
 B_\Lambda(k)
 &=-\frac{i}{2}e^{ikT_\Lambda}
 \left(\frac{\kappa}{k}+\frac{k}{\kappa}\right)
 \sinh(\kappa T_\Lambda).\label{eq:barrier-B}
\end{align}
In particular,
\begin{equation}
 c_K e^{\kappa(k)T_\Lambda}
 \leq\max\{|A_\Lambda(k)|,|B_\Lambda(k)|\}
 \leq C_K e^{\kappa(k)T_\Lambda}.
 \label{eq:barrier-coefficient-size}
\end{equation}

The full modal statement in Lemma~\ref{lem:exact-rotator-conversion} transforms
the high Jost column at $r=\Lambda$ into the low configuration entry
\begin{equation}
 G_\Lambda(k)
 =A_\Lambda(k)e^{i\phi_{+,\Lambda}(k)}
  +B_\Lambda(k)e^{i\phi_{-,\Lambda}(k)}
  +O_K\!\left(\Lambda^{-1/3}e^{\kappa(k)T_\Lambda}\right).
 \label{eq:rotator-G}
\end{equation}
The constant-coefficient pencil~\eqref{eq:rotator-pencil} gives
\begin{equation}
 \left|\frac{\dd}{\dd k}
  (\phi_{+,\Lambda}-\phi_{-,\Lambda})\right|\asymp_K\Lambda,
 \qquad
 \left|\frac{\dd^2}{\dd k^2}
  (\phi_{+,\Lambda}-\phi_{-,\Lambda})\right|=O_K(\Lambda).
 \label{eq:rotator-relative-phase}
\end{equation}
Equations~\eqref{eq:barrier-A}--\eqref{eq:barrier-B} have relative derivatives
$O_K(T_\Lambda)$, so Lemma~\ref{lem:rapid-phase-good-set} applies with
$M_\Lambda=e^{\kappa T_\Lambda}$.  Moreover the choice of $\beta$ gives
\begin{equation}
 \Lambda^{-1/3}e^{\kappa(k)T_\Lambda}
 \leq\Lambda^{-1/3+\beta\kappa_{\max}/\kappa_{\min}}=o(1).
 \label{eq:rotator-error-after-barrier}
\end{equation}
On the good set from~\eqref{eq:rapid-phase-good-set}, the configuration Jost
matrix is therefore an anti-diagonal channel swap with one entry of size at
least $c_Ke^{\kappa T_\Lambda}$ and the other bounded above and below, plus
$o(1)$ in the complementary entries.  Hence
\begin{equation}
 \Gamma_P(F_\Lambda(0,k))
 \geq c_Ke^{\kappa(k)T_\Lambda}
 \geq c_K\Lambda^\beta
\end{equation}
on a subset of $K$ of measure at least $c_K|K|$.  Integrating its positive
logarithm proves~\eqref{eq:matrix-rotator-entropy-growth}.
\end{proof}

The theorem rejects an arbitrary matrix replacement of the scalar angular
problem.  It does not disprove Simon's conjecture.  A scalar realization
necessarily creates a structured high-sector dilation, but the next section
shows that angular concentration can make both its full low-column cost and
its barrier leakage vanish.  For a full-space, cutoff-stable counterexample one
would still have to lift the rotating converter in the complete scalar angular
dynamics.  Section~\ref{app:cutoff-edge-scalar-counterexample} shows that the
uniform finite-cutoff statement $\mathrm{PSR}_1$ has weaker quantifiers and can
already be disproved without solving that stronger problem.

\section{Scalar lifts of the two-channel block}
\label{app:scalar-block-lift}

Angular concentration can make a large high-channel matrix element very cheap
in scalar $L^2$.  The exact minimum-cost calculation is as follows.

\begin{proposition}[Minimum-cost scalar lift of a real $2\times2$ block]
\label{prop:minimum-scalar-block-lift}
Let $(X,\mu)$ have finite measure $S$, set $e_0=S^{-1/2}$, and let
$h\in L^\infty(X;\mathbb R)$ satisfy
\begin{equation}
 \int_Xh\dd\mu=0,\qquad \|h\|_2=1,\qquad
 \int_Xh^3\dd\mu=0.
\end{equation}
Put
\begin{equation}
 g_h=h^2-S^{-1},\qquad
 G_h=\|g_h\|_2^2=\|h\|_4^4-S^{-1}.
 \label{eq:scalar-lift-G}
\end{equation}
For real $a,b,d$ and $G_h>0$, define
\begin{equation}
 v_{a,b,d}
 =a+\sqrt S\,b\,h+\frac{d-a}{G_h}g_h.
 \label{eq:minimum-scalar-lift}
\end{equation}
Then
\begin{equation}
 P_{\{e_0,h\}}M_{v_{a,b,d}}P_{\{e_0,h\}}
 =\begin{pmatrix}a&b\\b&d\end{pmatrix}
 \label{eq:exact-scalar-compressed-block}
\end{equation}
and
\begin{align}
 \|v_{a,b,d}\|_2^2
 &=S(a^2+b^2)+\frac{(d-a)^2}{G_h},\label{eq:minimum-scalar-L2}\\
 \|M_{v_{a,b,d}}e_0\|_2^2
 &=a^2+b^2+\frac{(d-a)^2}{SG_h}.
 \label{eq:minimum-scalar-column-cost}
\end{align}
Moreover,~\eqref{eq:minimum-scalar-lift} is the unique minimum-$L^2$ real
scalar potential having the compressed block
~\eqref{eq:exact-scalar-compressed-block}.
\end{proposition}

\begin{proof}
The functions $1,h,g_h$ are pairwise orthogonal; the third-moment condition is
used for $h\perp g_h$.  The three prescribed matrix entries fix respectively
the coefficients of these three functions, giving
~\eqref{eq:minimum-scalar-lift}.  Equations
~\eqref{eq:exact-scalar-compressed-block}--\eqref{eq:minimum-scalar-column-cost}
follow by direct substitution.  Any other lift differs by a function
orthogonal to $1,h,g_h$, so Pythagoras proves minimality and uniqueness.
\end{proof}

\begin{proposition}[Exact scalar mixer--barrier duality]
\label{prop:scalar-mixer-barrier-duality}
Under the hypotheses of Proposition~\ref{prop:minimum-scalar-block-lift}, let
\begin{equation}
 m_{h,b}=\sqrt S\,b h,
 \qquad
 d_{h,Q}=\frac{Q}{G_h}\left(h^2-S^{-1}\right).
 \label{eq:scalar-mixer-barrier-pair}
\end{equation}
The first multiplier has the exact compressed mixer block
\begin{equation}
 P_{\{e_0,h\}}M_{m_{h,b}}P_{\{e_0,h\}}
 =\begin{pmatrix}0&b\\b&0\end{pmatrix},
 \label{eq:pure-scalar-mixer-block}
\end{equation}
and its leakage from the selected high state is
\begin{equation}
 \left\|(I-P_{\{e_0,h\}})M_{m_{h,b}}h\right\|_2^2
 =S b^2G_h.
 \label{eq:pure-scalar-mixer-leakage}
\end{equation}
The second is the unique minimum-$L^2$ high-only barrier with compressed
block $\diag(0,Q)$, and its charged constant-column cost is
\begin{equation}
 \|M_{d_{h,Q}}e_0\|_2^2=\frac{Q^2}{S G_h}.
 \label{eq:pure-scalar-barrier-cost}
\end{equation}
Consequently the exact product law is
\begin{equation}
 \left\|(I-P_{\{e_0,h\}})M_{m_{h,b}}h\right\|_2^2
 \|M_{d_{h,Q}}e_0\|_2^2=b^2Q^2.
 \label{eq:scalar-mixer-barrier-product}
\end{equation}
Thus concentration cannot make a high-only scalar barrier cheap while also
making the corresponding pure scalar two-state mixer invariant.  This is a
two-plane obstruction; it does not exclude a genuinely multimode cascade.
\end{proposition}

\begin{proof}
The functions $1,h,g_h$ are pairwise orthogonal and
$M_{m_{h,b}}e_0=bh$.  Also
\begin{equation}
 M_{m_{h,b}}h=\sqrt S\,b h^2
 =b e_0+\sqrt S\,b g_h.
\end{equation}
Since $g_h\perp\{e_0,h\}$, this proves
~\eqref{eq:pure-scalar-mixer-block}--
~\eqref{eq:pure-scalar-mixer-leakage}.  Proposition
~\ref{prop:minimum-scalar-block-lift} with $(a,b,d)=(0,0,Q)$ proves
~\eqref{eq:pure-scalar-barrier-cost}; multiplication gives
~\eqref{eq:scalar-mixer-barrier-product}.
\end{proof}

\begin{proposition}[Highest-weight concentration on $\mathbb S^2$]
\label{prop:highest-weight-scalar-lift}
Let $S=4\pi$ and, for $\ell\geq1$, let
\begin{equation}
 h_\ell(\theta,\phi)
 =C_\ell(\sin\theta)^\ell\cos(\ell\phi),
 \qquad \|h_\ell\|_2=1.
 \label{eq:real-highest-weight}
\end{equation}
Define
\begin{equation}
 J_n=\sqrt\pi\,\frac{\Gamma(n+1)}{\Gamma(n+3/2)}.
\end{equation}
Then $C_\ell^2=(\pi J_\ell)^{-1}$ and
\begin{align}
 H_{4,\ell}:=\int_{\mathbb S^2}h_\ell^4\dd\omega
 &=\frac{3}{4\pi}\frac{J_{2\ell}}{J_\ell^2}
 \sim\frac{3}{4\pi\sqrt{2\pi}}\sqrt\ell,\label{eq:highest-H4}\\
 H_{6,\ell}:=\int_{\mathbb S^2}h_\ell^6\dd\omega
 &=\frac{5}{8\pi^2}\frac{J_{3\ell}}{J_\ell^3}
 \sim\frac{5}{8\pi^3\sqrt3}\ell.
 \label{eq:highest-H6}
\end{align}
In particular $G_\ell=H_{4,\ell}-(4\pi)^{-1}\asymp\sqrt\ell$.  The exact
high-only lift
\begin{equation}
 v_{\ell,Q}=\frac{Q}{G_\ell}
 \left(h_\ell^2-\frac1{4\pi}\right)
 \label{eq:highest-exact-barrier}
\end{equation}
satisfies
\begin{align}
 P_{\{e_0,h_\ell\}}M_{v_{\ell,Q}}P_{\{e_0,h_\ell\}}
 &=\diag(0,Q),\label{eq:highest-exact-block}\\
 \|M_{v_{\ell,Q}}e_0\|_2^2
 &=\frac{Q^2}{4\pi G_\ell}
 \sim\frac{\sqrt{2\pi}}{3}\frac{Q^2}{\sqrt\ell},
 \label{eq:highest-column-cost}\\
 \limsup_{\ell\to\infty}\|v_{\ell,Q}\|_\infty
 &\leq\frac{4\sqrt2}{3}|Q|.
 \label{eq:highest-uniform-bound}
\end{align}
The selected two-dimensional space is not invariant.  Its high-state leakage
has the nonzero limit
\begin{equation}
 \lim_{\ell\to\infty}
 \frac{\|(I-P_{\{e_0,h_\ell\}})M_{v_{\ell,Q}}h_\ell\|_2^2}{Q^2}
 =\frac{20}{9\sqrt3}-1.
 \label{eq:highest-minimizer-leakage}
\end{equation}
All harmonics in this leakage have degree between $\ell$ and $3\ell$.
\end{proposition}

\begin{proof}
The $\phi$ integrals are
$\int_0^{2\pi}\cos^2(\ell\phi)\dd\phi=\pi$,
$\int_0^{2\pi}\cos^4(\ell\phi)\dd\phi=3\pi/4$, and
$\int_0^{2\pi}\cos^6(\ell\phi)\dd\phi=5\pi/8$.
The remaining integrals are $J_\ell,J_{2\ell},J_{3\ell}$.
The gamma-ratio asymptotic gives
$J_n\sim\sqrt{\pi/n}$ and proves
~\eqref{eq:highest-H4}--\eqref{eq:highest-H6}.
Since $\int h_\ell^3=0$, Proposition
~\ref{prop:minimum-scalar-block-lift} gives
~\eqref{eq:highest-exact-block}--\eqref{eq:highest-column-cost}.
Also
\begin{equation}
 \frac{\|h_\ell\|_\infty^2}{H_{4,\ell}}\longrightarrow
 \frac{4\sqrt2}{3},
\end{equation}
which proves~\eqref{eq:highest-uniform-bound}.  Finally,
\begin{equation}
 \|(I-P_{\{e_0,h_\ell\}})M_{v_{\ell,Q}}h_\ell\|_2^2
 =Q^2\left[
 \frac{H_{6,\ell}-H_{4,\ell}/(2\pi)+1/(16\pi^2)}
 {G_\ell^2}-1\right].
\end{equation}
The displayed asymptotics give~\eqref{eq:highest-minimizer-leakage}.
The identity
$\cos^3(\ell\phi)=(3\cos(\ell\phi)+\cos(3\ell\phi))/4$
shows that $h_\ell^3-(4\pi)^{-1}h_\ell$ has azimuthal frequencies
$\ell$ and $3\ell$; its spherical-harmonic degrees therefore lie between
$\ell$ and $3\ell$.
\end{proof}

The fixed leakage in~\eqref{eq:highest-minimizer-leakage} is not unavoidable.
A slightly non-minimal lift makes the high state almost invariant.

\begin{proposition}[Thin-belt high-only scalar barrier]
\label{prop:thin-belt-barrier}
Fix $0<\alpha<1/2$ and let
\begin{equation}
 E_\ell=\{|\theta-\pi/2|\leq\ell^{-\alpha}\},
 \quad
 \mu_\ell=|E_\ell|,
 \quad
 \epsilon_\ell=\int_{E_\ell^c}h_\ell^2\dd\omega.
\end{equation}
There are real axisymmetric $b_{\ell,Q}$ such that
\begin{align}
 P_{\{e_0,h_\ell\}}M_{b_{\ell,Q}}P_{\{e_0,h_\ell\}}
 &=\diag(0,Q),\label{eq:belt-exact-block}\\
 \|M_{b_{\ell,Q}}e_0\|_2^2
 &=O_Q(\ell^{-\alpha}),&
 \|b_{\ell,Q}\|_\infty&=O_Q(1),\label{eq:belt-cost-bound}\\
 \|(I-P_{\{e_0,h_\ell\}})M_{b_{\ell,Q}}h_\ell\|_2^2
 &\leq C_Qe^{-c\ell^{1-2\alpha}}.
 \label{eq:belt-leakage-bound}
\end{align}
\end{proposition}

\begin{proof}
Put $\rho_\ell=\mu_\ell/(4\pi-\mu_\ell)$ and
\begin{equation}
 b_{\ell,Q}
 =q_{\ell,Q}\bigl(\mathbf1_{E_\ell}
 -\rho_\ell\mathbf1_{E_\ell^c}\bigr),
 \qquad
 q_{\ell,Q}
 =\frac{Q}{(1-\epsilon_\ell)-\rho_\ell\epsilon_\ell}.
 \label{eq:belt-potential}
\end{equation}
It has mean zero, and its inner product with $h_\ell$ vanishes because it is
axisymmetric while $h_\ell$ has azimuthal frequency $\ell$.  Its expectation
in $h_\ell$ is $Q$, proving~\eqref{eq:belt-exact-block}.  Since
$\mu_\ell=4\pi\sin(\ell^{-\alpha})=O(\ell^{-\alpha})$ and the Gaussian-beam
tail obeys
\begin{equation}
 \epsilon_\ell\leq Ce^{-c\ell^{1-2\alpha}},
\end{equation}
one has $q_{\ell,Q}=Q(1+o(1))$ and
\begin{equation}
 \|M_{b_{\ell,Q}}e_0\|_2^2
 =q_{\ell,Q}^2\frac{\mu_\ell}{4\pi-\mu_\ell}.
\end{equation}
This proves~\eqref{eq:belt-cost-bound}.  The leakage is exactly the variance
of this two-valued function under the probability density $h_\ell^2$:
\begin{equation}
 \|(I-P_{\{e_0,h_\ell\}})M_{b_{\ell,Q}}h_\ell\|_2^2
 =q_{\ell,Q}^2(1+\rho_\ell)^2
 \epsilon_\ell(1-\epsilon_\ell),
\end{equation}
which gives~\eqref{eq:belt-leakage-bound}.
\end{proof}

\begin{corollary}[Algebraic scalar lift of the rotator/barrier profile]
\label{cor:scalar-rotator-barrier-lift}
Let the $2\times2$ rotator block have entries
\begin{equation}
 a=\delta_\Lambda\sin^2\theta,\qquad
 b=-\delta_\Lambda\sin\theta\cos\theta,\qquad
 d=\delta_\Lambda\cos^2\theta
\end{equation}
on an interval of length $\Lambda$, with
$\delta_\Lambda=\Lambda^{-2/3}$ and one quarter-turn.  If
$\ell_\Lambda\to\infty$, Proposition
~\ref{prop:minimum-scalar-block-lift} gives an exact scalar compressed lift
whose full constant-column cost is
\begin{equation}
 \frac12\delta_\Lambda^2\Lambda
 \left(1+\frac1{4\pi G_{\ell_\Lambda}}\right).
 \label{eq:scalar-lifted-rotator-cost}
\end{equation}
Follow it by the belt barrier of Proposition~\ref{prop:thin-belt-barrier} for
time $T_\Lambda=O(\log\Lambda)$.  Choosing, for example,
$\ell_\Lambda=(\log\Lambda)^p$ with $p\alpha>1$ makes the barrier cost and its
cumulative high-state leakage tend to zero.  Since
$\|h_\ell\|_\infty=O(\ell^{1/4})$ for~\eqref{eq:real-highest-weight}, every
fixed polylogarithmic choice also keeps the rotator lift uniformly bounded:
$\delta_\Lambda\|h_{\ell_\Lambda}\|_\infty\to0$.

This is an exact scalar lift of the desired compressed matrices, not yet a
scalar dynamical counterexample: multiplication during the converter does not
preserve $\operatorname{span}\{e_0,h_{\ell_\Lambda}\}$.
\end{corollary}

\section{Cutoff-edge scalar counterexample to $\mathrm{PSR}_1$}
\label{app:cutoff-edge-scalar-counterexample}

The preceding full-space converter question is stronger than the quantifiers in
$\mathrm{PSR}_1$.  That estimate is required uniformly over the finite angular
cutoff.  A scalar multiplier may therefore contain degrees above the retained
cutoff, provided its \emph{full} constant-column cost is charged.  The omitted
degrees can cancel every unwanted retained component of the converter.  This
cutoff-edge freedom decides scalar $\mathrm{PSR}_1$ negatively.

Write
\begin{equation}
 \mathcal H_L=\bigoplus_{n=0}^{L}\mathcal Y_n,
 \qquad P_{\leq L}:L^2(\mathbb S^2)\longrightarrow\mathcal H_L,
 \qquad \mathcal S_L=\operatorname{span}\{e_0,h_L\},
 \label{eq:cutoff-edge-spaces}
\end{equation}
where $h_L$ is the normalized real highest-weight harmonic.  All spaces and
maps in the next proposition are taken over the reals.

\begin{proposition}[Exact cutoff-edge scalar lift]
\label{prop:exact-cutoff-edge-lift}
For every $L\geq1$ and every real symmetric matrix
\begin{equation}
 B=\begin{pmatrix}a&b\\b&d\end{pmatrix},
\end{equation}
there is a real spherical polynomial $v_{L,B}$ of degree at most $2L$ such
that $\mathcal S_L$ reduces $P_{\leq L}M_{v_{L,B}}P_{\leq L}$ and
\begin{equation}
 \left.P_{\leq L}M_{v_{L,B}}P_{\leq L}\right|_{\mathcal S_L}=B
 \label{eq:exact-cutoff-edge-block}
\end{equation}
in the ordered basis $(e_0,h_L)$.  The lift $B\mapsto v_{L,B}$ may be chosen
linear.  Consequently, for each fixed $L$ there are finite constants
$C_{2,L},C_{\infty,L}$ such that
\begin{equation}
 \|v_{L,B}\|_2\leq C_{2,L}\|B\|,
 \qquad
 \|v_{L,B}\|_\infty\leq C_{\infty,L}\|B\|.
 \label{eq:cutoff-edge-lift-bounds}
\end{equation}
No uniform estimate on these constants is needed below.
\end{proposition}

\begin{proof}
Put
\begin{equation}
 \mathcal W_L=\bigoplus_{n=L+1}^{2L}\mathcal Y_n,
 \qquad
 T_L:\mathcal W_L\longrightarrow\mathcal H_L,
 \qquad T_Lw=P_{\leq L}(h_Lw).
 \label{eq:cutoff-edge-moment-map}
\end{equation}
We first prove
\begin{equation}
 \operatorname{Ran}T_L=e_0^\perp\cap\mathcal H_L.
 \label{eq:cutoff-edge-moment-range}
\end{equation}
Certainly $T_L^*e_0=0$, because $h_Le_0$ has pure degree $L$.  Conversely,
let $f\in\ker T_L^*$ and let $f_n\neq0$, $n\geq1$, be its highest nonconstant
harmonic component.  If $H_L$ and $F_n$ are the associated homogeneous
harmonic polynomials, the degree-$(L+n)$ harmonic component of $H_LF_n$ is
nonzero.  Indeed, otherwise the Fischer decomposition would make $H_LF_n$
divisible by $|x|^2$; irreducibility of $|x|^2$ and uniqueness of the harmonic
decomposition would force one of the nonzero harmonic factors to be divisible
by $|x|^2$, which is impossible.  Lower components of $f$ cannot cancel this
top degree.  Thus $T_L^*f\neq0$.  Hence $\ker T_L^*=\mathbb Re_0$, proving
~\eqref{eq:cutoff-edge-moment-range}.

Set
\begin{equation}
 v^{\rm low}_{L,B}=a\sqrt{4\pi}\,e_0+b\sqrt{4\pi}\,h_L.
\end{equation}
Then
$P_{\leq L}(v^{\rm low}_{L,B}e_0)=ae_0+bh_L$.  The remaining required column
is
\begin{equation}
 g_{L,B}=be_0+dh_L-P_{\leq L}(v^{\rm low}_{L,B}h_L).
 \label{eq:cutoff-edge-compatible-rhs}
\end{equation}
Self-adjoint symmetry gives $\langle e_0,g_{L,B}\rangle=0$.  By
~\eqref{eq:cutoff-edge-moment-range}, choose $w_{L,B}\in\mathcal W_L$ with
$T_Lw_{L,B}=g_{L,B}$, using once and for all a linear right inverse of $T_L$
on $e_0^\perp$.  Then
$v_{L,B}=v^{\rm low}_{L,B}+w_{L,B}$ has the two required columns.  The
compressed multiplier is self-adjoint, so invariance of $\mathcal S_L$ makes
it reducing.  Finite dimensionality gives~\eqref{eq:cutoff-edge-lift-bounds}.
\end{proof}

The high-only barrier has a simpler closed form.  Define
\begin{equation}
 p_{2L}(\theta,\phi)=\sin^{2L}\theta\cos(2L\phi),
 \qquad
 c_L=\langle h_L,p_{2L}h_L\rangle.
 \label{eq:cutoff-edge-barrier-harmonic}
\end{equation}
The function $p_{2L}$ is a pure real harmonic of degree $2L$ and
$\|p_{2L}\|_\infty=1$.

\begin{lemma}[Exact high-only scalar barrier at the cutoff]
\label{lem:exact-cutoff-edge-barrier}
One has
\begin{align}
 P_{\leq L}M_{p_{2L}}e_0&=0,
 &P_{\leq L}M_{p_{2L}}h_L&=c_Lh_L,
 \label{eq:cutoff-edge-barrier-block}\\
 c_L&=\frac{I_{4L+1}}{2I_{2L+1}}\longrightarrow\frac1{2\sqrt2},
 &\|p_{2L}\|_2^2&=\pi I_{4L+1}\asymp L^{-1/2},
 \label{eq:cutoff-edge-barrier-cost}
\end{align}
where $I_m=\int_0^\pi\sin^m\theta\,\dd\theta$.  Hence, for every fixed
$Q>0$, the uniformly bounded scalar multiplier
\begin{equation}
 b_{L,Q}=\frac{Q}{c_L}p_{2L}
 \label{eq:exact-cutoff-edge-barrier-potential}
\end{equation}
compresses on $\mathcal S_L$ to $\diag(0,Q)$, while
\begin{equation}
 \|M_{b_{L,Q}}e_0\|_2^2\asymp_Q L^{-1/2}.
 \label{eq:exact-cutoff-edge-barrier-charged-cost}
\end{equation}
\end{lemma}

\begin{proof}
The first identity in~\eqref{eq:cutoff-edge-barrier-block} is degree
orthogonality.  Moreover
\begin{equation}
 p_{2L}h_L
 =\frac{C_L}{2}\sin^{3L}\theta
   \bigl(\cos(3L\phi)+\cos(L\phi)\bigr).
\end{equation}
The first summand has azimuthal order $3L$; the second has order $L$, so its
only component of degree at most $L$ is a multiple of $h_L$.  This proves the
second block identity.  Since
$h_L=C_L\sin^L\theta\cos(L\phi)$ and
$C_L^{-2}=\pi I_{2L+1}$, direct integration gives
~\eqref{eq:cutoff-edge-barrier-cost}.  The Wallis asymptotic for $I_m$ and
$e_0=(4\pi)^{-1/2}$ finish the proof.
\end{proof}

\begin{theorem}[Scalar $\mathrm{PSR}_1$ is false]
\label{thm:scalar-PSR1-counterexample}
Let $K=[k_0,k_1]\Subset(0,\infty)$.  There are angular cutoffs $L\to\infty$
and uniformly bounded, compactly supported real scalar potentials $V_L$ on
$[1,\infty)\times\mathbb S^2$ such that
\begin{equation}
 \sup_L\int_1^\infty\|M_{V_L(r)}e_0\|_2^2\dd r<\infty,
 \label{eq:scalar-PSR1-counterexample-cost}
\end{equation}
but for the finite angular models on $\mathcal H_L$,
\begin{equation}
 \int_K\log^+\Gamma_{0;V_L}(k)\dd k\longrightarrow\infty.
 \label{eq:scalar-PSR1-counterexample-entropy}
\end{equation}
Thus~\eqref{eq:scalar-PSR-target}, and hence scalar $\mathrm{PSR}_1$, fails.
This does not disprove Simon's conjecture.
\end{theorem}

\begin{proof}
Fix $Q_0>k_1^2$, put
$\kappa_{\min}=\sqrt{Q_0-k_1^2}$ and
$\kappa_{\max}=\sqrt{Q_0-k_0^2}$, and choose
$0<\gamma<1/2$.  Let
\begin{equation}
 T_L=L^\gamma.
 \label{eq:cutoff-edge-barrier-length}
\end{equation}
Use Proposition~\ref{prop:exact-cutoff-edge-lift} to lift the rotating block
~\eqref{eq:exact-rotator-potential} with
$\delta_\Lambda=\Lambda^{-2/3}$ and
$\omega_\Lambda=\pi/(2\Lambda)$.  For each $L$, choose $\Lambda=\Lambda_L$
so large that
\begin{equation}
 C_{2,L}^2\Lambda_L^{-1/3}\leq1,
 \qquad
 C_{\infty,L}\Lambda_L^{-2/3}\leq1,
 \qquad
 \left(\Lambda_L^{-1/3}+\frac{L^2}{\Lambda_L}
       +\frac{L^{7/2}}{\Lambda_L^3}\right)
 e^{2\kappa_{\max}T_L}\leq L^{-1}.
 \label{eq:cutoff-edge-rotator-scales}
\end{equation}
This is possible with no information about the growth of the lift constants.

Choose a fixed $A_K\geq3$ larger than the uniform constant in the upper bound
of~\eqref{eq:rotator-relative-phase}, and place the rotator on a shell beginning
at $R_L=A_K\Lambda_L$, replacing $\theta(r)$ by
$\omega_{\Lambda_L}(r-R_L)$.  On its
two-dimensional reducing space, cancel the centrifugal entry
$d_L(r)=L(L+1)/r^2$ by adding
$-d_L(r)c_L^{-1}p_{2L}$.  On the following shell of length $T_L$, use
\begin{equation}
 \frac{Q_0-d_L(r)}{c_L}p_{2L}.
 \label{eq:cutoff-edge-physical-barrier}
\end{equation}
Lemma~\ref{lem:exact-cutoff-edge-barrier} then makes the \emph{total}
compressed radial operator on $\mathcal S_L$ exactly the matrix rotator
followed by $\diag(0,Q_0)$.  At radii comparable to $\Lambda_L$, the normalized
free Bessel directional frame differs from the plane-wave frame, together with
one $k$ derivative, by $O_K(L^2/\Lambda_L)$.  Moreover the charged cost of the
centrifugal cancellation on the rotator is
\begin{equation}
 O\!\left(\Lambda_L\frac{L^4}{\Lambda_L^4}
                 \|p_{2L}e_0\|_2^2\right)
 =O\!\left(\frac{L^{7/2}}{\Lambda_L^3}\right).
 \label{eq:cutoff-edge-centrifugal-cost}
\end{equation}
Thus all matching, centrifugal, and converted-mode errors remain negligible
after barrier amplification by~\eqref{eq:cutoff-edge-rotator-scales}.

The converter cost is bounded by
\begin{equation}
 C C_{2,L}^2\delta_{\Lambda_L}^2\Lambda_L
 =C C_{2,L}^2\Lambda_L^{-1/3}\leq C.
 \label{eq:cutoff-edge-converter-cost}
\end{equation}
The barrier cost tends to zero:
\begin{equation}
 \int_{\rm barrier}\|M_{V_L(r)}e_0\|_2^2\dd r
 \leq C_QT_LL^{-1/2}
 =C_QL^{\gamma-1/2}\longrightarrow0.
 \label{eq:cutoff-edge-total-barrier-cost}
\end{equation}
The choice of $\Lambda_L$, together with
~\eqref{eq:cutoff-edge-barrier-cost}, also makes the full scalar amplitudes
uniformly bounded.  This proves~\eqref{eq:scalar-PSR1-counterexample-cost}.

Because $\mathcal S_L$ reduces every compressed scalar multiplier used above,
the full finite Jost matrix is block diagonal relative to
$\mathcal S_L\oplus\mathcal S_L^\perp$, and
\begin{equation}
 \Gamma_0(k)=\|J(k)^{-1}e_0\|^{-1}
\end{equation}
is exactly the two-channel quotient.  Propagation from $R_L$ through the inner
free collar is diagonal in the angular channels.  For real $k$, the incoming
and outgoing degree-$L$ Bessel factors at $r=1$ are conjugates and hence have
equal modulus; their common row size cancels from $\Gamma_0$.  In the degree
zero row the collar merely appends the phases $e^{\pm ik(R_L-1)}$ to the two
directional terms.  Consequently the proof of Theorem
~\ref{thm:matrix-rotator-counterexample} applies with the present $T_L$:
~\eqref{eq:cutoff-edge-rotator-scales} makes the converted-mode and Bessel
matching errors negligible after barrier amplification, while the total
relative phase has first derivative comparable to $\Lambda_L$ and second
derivative $O_K(\Lambda_L)$.  Lemma
~\ref{lem:rapid-phase-good-set} therefore supplies a subset of $K$ of fixed
positive measure on which
\begin{equation}
 \Gamma_0(k)\geq c_Ke^{\kappa_{\min}T_L}.
\end{equation}
Consequently
$\int_K\log^+\Gamma_0\geq c_KT_L-C_K\to\infty$, proving
~\eqref{eq:scalar-PSR1-counterexample-entropy}.  Taking the radial truncation
parameter beyond the support of $V_L$ places this model among the terms in the
supremum in~\eqref{eq:scalar-PSR-target}, whereas its right-hand side stays
bounded by~\eqref{eq:scalar-PSR1-counterexample-cost}.  This completes the
quantifier-level contradiction.
\end{proof}

\begin{remark}[What has been disproved]
The counterexample uses a degree-$2L$ scalar multiplier against a degree-$L$
Galerkin model.  Its full constant-column norm is included in
~\eqref{eq:scalar-PSR1-counterexample-cost}; no truncated-column aliasing is
hidden.  What fails is uniformity of the proposed sufficient estimate under
cutoff-edge scalar compression.  A theorem which first fixes one scalar
potential and passes to all sufficiently large angular cutoffs, or which
imposes an independently justified relation between multiplier bandwidth and
model cutoff, is a different statement.  The full-space converter problem
below remains mathematically open, but it is no longer needed to decide
$\mathrm{PSR}_1$.
\end{remark}

\section{A cutoff-compatible fixed-potential diagonal}
\label{app:buffered-diagonal}

The failure of uniform $\mathrm{PSR}_1$ does not force the same quantifiers on
an approximation theorem.  For one fixed scalar potential it is enough to
choose a single cofinal sequence of angularly smoothed, radially truncated
models, put the Galerkin boundary beyond the smoothing bandwidth, and retain a
bounded entropy subsequence.  This section makes that replacement exact and
separates two issues: approximation is harmless, while the required entropy
estimate remains substantive.

We first record an $L^\infty$-safe finite-band approximation.  Let
\begin{equation}
 \mathcal D_B(\omega,\eta)
 =\sum_{n=0}^{B}\sum_{m=-n}^{n}
   Y_{nm}(\omega)\overline{Y_{nm}(\eta)},
 \qquad m_B=(B+1)^2,
\end{equation}
and define the positive zonal kernel
\begin{equation}
 \mathcal K_B(\omega,\eta)
 =\frac{|\mathcal D_B(\omega,\eta)|^2}{m_B/(4\pi)},
 \qquad
 \mathsf S_Bv(\omega)=\int_{\mathbb S^2}
   \mathcal K_B(\omega,\eta)v(\eta)\dd\eta.
 \label{eq:positive-angular-smoother}
\end{equation}

\begin{proposition}[Positive finite-band approximation]
\label{prop:positive-finite-band-approximation}
The operators $\mathsf S_B$ preserve real functions and constants, map into
spherical polynomials of degree at most $2B$, and satisfy
\begin{equation}
 \|\mathsf S_Bv\|_p\leq\|v\|_p,
 \qquad 1\leq p\leq\infty,
 \qquad
 \mathsf S_Bv\longrightarrow v\quad\hbox{in }L^2(\mathbb S^2).
 \label{eq:positive-angular-smoother-bounds}
\end{equation}
Consequently, if $V\in L^\infty\cap X_3$ and
\begin{equation}
 V_j(r,\omega)=\mathbf1_{[1,R_j]}(r)\,
                   \mathsf S_{B_j}V(r,\cdot)(\omega),
 \qquad R_j,B_j\longrightarrow\infty,
 \label{eq:buffered-potential-approximants}
\end{equation}
then $V_j$ is real, compactly supported, and of angular degree at most $2B_j$,
while
\begin{equation}
 \|V_j\|_\infty\leq\|V\|_\infty,
 \qquad
 \|V_j-V\|_{X_3}\longrightarrow0,
 \qquad
 \|V_j\|_{X_3}\leq\|V\|_{X_3}.
 \label{eq:buffered-potential-convergence}
\end{equation}
\end{proposition}

\begin{proof}
The reproducing identity gives
$\int|\mathcal D_B(\omega,\eta)|^2\dd\eta=m_B/(4\pi)$.
Hence $\mathcal K_B$ is a symmetric Markov kernel.  Positivity, Jensen's
inequality, and duality prove all $L^p$ contractions.  The product of two
degree-$B$ spherical polynomials has degree at most $2B$.  Finally, the
normalized squared reproducing kernels form an approximate identity: the
Christoffel--Darboux formula gives concentration of their mass in every fixed
neighbourhood of the diagonal.  Thus $\mathsf S_Bv\to v$ first for continuous
functions and then in $L^2$ by density and contraction.  Apply this fiberwise
in $r$ and use dominated convergence, followed by $R_j\to\infty$, to obtain
~\eqref{eq:buffered-potential-convergence}.
\end{proof}

Let $H_0=-\partial_r^2+A_3/r^2$ on the exterior half-line.  Extend each finite
model to the full angular space by leaving the complementary channels free:
\begin{equation}
 \widehat H_j
 =H_0+P_{L_j}M_{V_j}P_{L_j}
 =H_{R_j,L_j}(V_j)\oplus H_0\big|_{P_{L_j}^\perp}.
 \label{eq:extended-buffered-model}
\end{equation}

\begin{proposition}[Strong resolvent convergence along every buffered diagonal]
\label{prop:buffered-strong-resolvent}
If $L_j\to\infty$, then
\begin{equation}
 \widehat H_j\longrightarrow H_0+M_V
 \qquad\hbox{in the strong resolvent sense}.
 \label{eq:buffered-strong-resolvent}
\end{equation}
In particular one may impose, without changing the limit, the double buffer
\begin{equation}
 \frac{L_j}{B_j+R_j+1}\longrightarrow\infty.
 \label{eq:genuine-angular-buffer}
\end{equation}
The matrix spectral measures of every fixed finite collection of compact
low-channel probes then converge vaguely.
\end{proposition}

\begin{proof}
The uniform $L^\infty$ bound and~\eqref{eq:buffered-potential-convergence}
imply $M_{V_j}\to M_V$ strongly.  Since $P_{L_j}\to\Id$ strongly and all
three operator families are uniformly bounded,
$P_{L_j}M_{V_j}P_{L_j}\to M_V$ strongly.  The second resolvent identity for
bounded self-adjoint perturbations of $H_0$ proves
~\eqref{eq:buffered-strong-resolvent}.  Condition
~\eqref{eq:genuine-angular-buffer} can be enforced by choosing $L_j$ after
$B_j$ and $R_j$.  Vague convergence of fixed-probe spectral measures follows
from the spectral theorem.
\end{proof}

The buffer excludes the exact mechanism of
Theorem~\ref{thm:scalar-PSR1-counterexample} for a purely algebraic reason.

\begin{lemma}[A nonconstant reducing cutoff-edge block consumes the buffer]
\label{lem:buffer-excludes-cutoff-edge-lift}
Let $v$ be a real spherical polynomial of degree at most $B<L$, and let $h_L$
have pure degree $L$.  If $\operatorname{span}\{e_0,h_L\}$ is invariant under
$P_{\leq L}M_vP_{\leq L}$, then $v$ is constant and its compressed block on
that plane is a scalar multiple of the identity.  In particular, neither the
rotator nor the high-only barrier used in
Theorem~\ref{thm:scalar-PSR1-counterexample} can occur under
~\eqref{eq:genuine-angular-buffer}.
\end{lemma}

\begin{proof}
Because $B<L$,
$P_{\leq L}(ve_0)=ve_0$ has no degree-$L$ component.  Invariance forces it to
be a multiple of $e_0$, hence $v$ is constant.  The conclusion follows.
\end{proof}

\begin{remark}[Why shell stacking does not restore the edge counterexample]
The scalar $\mathrm{PSR}_1$ examples may be placed on disjoint shells and their
charged costs may be made small by lengthening the rotator.  This observation
can produce bad \emph{unbuffered} diagonals for a shell-stacked potential, but
it does not negate the quantifiers in Theorem
~\ref{thm:buffered-diagonal-completion}.  At stage $j$ the angular smoothing
fixes the multiplier bandwidth and the radial truncation fixes the support
radius before $L_j$ is chosen.  Condition~\eqref{eq:genuine-angular-buffer}
then puts both the multiplier edge and the whole interaction region far below
the Galerkin edge.  Any counterexample surviving this order of limits must be
a genuine full-space scalar converter, not the compression identity in
Theorem~\ref{thm:scalar-PSR1-counterexample}.
\end{remark}

The double buffer also gives a quantitative real-energy statement.  It does
not prove the entropy estimate, but it removes the Galerkin boundary from the
remaining converter problem.

\begin{proposition}[Outgoing high-channel Schur complement is coercive]
\label{prop:outgoing-high-channel-coercivity}
Let $K=[k_0,k_1]\Subset(0,\infty)$, let $V$ be real, supported in
$1\leq r\leq R$, and put $M=\|V\|_\infty$.  For $J<L$, set
$Q_{J,L}=P_L-P_J$ and consider on
$L^2((1,R);Q_{J,L}\mathcal K)$
\begin{equation}
 \mathcal A_{J,L}(k)
 =-\partial_r^2+\frac{D_L}{r^2}
   +Q_{J,L}M_VQ_{J,L}-k^2,
 \label{eq:outgoing-high-operator}
\end{equation}
with Dirichlet condition at $r=1$ and the free outgoing Bessel
Dirichlet-to-Neumann condition at $r=R$.  If
\begin{equation}
 g_{J,R}:=\frac{(J+1)(J+2)}{R^2}-M-k_1^2>0,
 \label{eq:high-channel-gap}
\end{equation}
then $\mathcal A_{J,L}(k)$ is invertible for every $k\in K$ and
\begin{equation}
 \sup_{k\in K}\|\mathcal A_{J,L}(k)^{-1}\|
 \leq g_{J,R}^{-1},
 \label{eq:high-channel-resolvent-bound}
\end{equation}
uniformly in $L$.  If
$\mathcal B_{J,L}=Q_{J,L}M_VP_J$, the exact outgoing Feshbach correction obeys
\begin{equation}
 \sup_{k\in K}
 \|\mathcal B_{J,L}^*\mathcal A_{J,L}(k)^{-1}
       \mathcal B_{J,L}\|
 \leq \frac{M^2}{g_{J,R}}.
 \label{eq:high-channel-self-energy-bound}
\end{equation}
\end{proposition}

\begin{proof}
For angular degree $n$, let
$f_n(r,k)=krh_n^{(1)}(kr)$ be the outgoing flattened free solution and set
$\mathfrak m_n(k,R)=f_n'(R,k)/f_n(R,k)$.  The standard half-integer Hankel
expansion gives
\begin{equation}
 |xh_n^{(1)}(x)|^2=\sum_{q=0}^n c_{nq}x^{-2q},
 \qquad c_{nq}>0.
 \label{eq:spherical-Hankel-amplitude}
\end{equation}
Hence $\Re\mathfrak m_n(k,R)=\partial_R\log|f_n(R,k)|\leq0$.
If $u$ satisfies the outgoing boundary condition, integration by parts gives
\begin{align}
 \Re\langle\mathcal A_{J,L}(k)u,u\rangle
 &=\int_1^R\!\left(\|u'\|^2
   +\left\langle u,
    \left(\frac{D_L}{r^2}+Q_{J,L}M_VQ_{J,L}-k^2\right)u
    \right\rangle\right)\dd r\notag\\
 &\quad-\Re\langle u(R),\mathfrak m_{J,L}(k,R)u(R)\rangle
 \geq g_{J,R}\|u\|_2^2.
 \label{eq:outgoing-high-coercive-form}
\end{align}
The complex coercive form theorem proves invertibility and
~\eqref{eq:high-channel-resolvent-bound}.  Since
$\|\mathcal B_{J,L}\|\leq M$, the Feshbach bound follows.
\end{proof}

\begin{corollary}[The double buffer removes the cutoff edge]
\label{cor:double-buffer-removes-edge}
For the approximants~\eqref{eq:buffered-potential-approximants} satisfying
~\eqref{eq:genuine-angular-buffer}, put
\begin{equation}
 J_j=\left\lfloor\sqrt{R_jL_j}\right\rfloor.
 \label{eq:active-angular-cutoff}
\end{equation}
Then $J_j/R_j\to\infty$, $J_j/L_j\to0$, and, for every fixed compact
$K\Subset(0,\infty)$,
\begin{equation}
 \sup_{k\in K}
 \|\mathcal B_{J_j,L_j}^*\mathcal A_{J_j,L_j}(k)^{-1}
       \mathcal B_{J_j,L_j}\|\longrightarrow0.
 \label{eq:buffered-edge-self-energy-vanishes}
\end{equation}
Moreover $J_j+2B_j<L_j$ eventually, so every one-step coupling from the active
block generated by $V_j$ is retained.  Thus any failure of
~\eqref{eq:buffered-diagonal-PSR} must occur inside a growing active block far
from the Galerkin edge.
\end{corollary}

\begin{proof}
Condition~\eqref{eq:genuine-angular-buffer} gives
$R_j/L_j\to0$ and $B_j/L_j\to0$.  Therefore the three elementary assertions
about $J_j$ and $B_j$ hold.  Since
$\|V_j\|_\infty\leq\|V\|_\infty$ and
$g_{J_j,R_j}\sim L_j/R_j\to\infty$, apply
Proposition~\ref{prop:outgoing-high-channel-coercivity}.
\end{proof}

The most dangerous explicit way to enter the active block is the wrapped
azimuthal phase from Proposition
~\ref{prop:wrapped-phase-gate-additivity-no-go}.  Allowing its pulse area to
grow makes the conversion arbitrarily broadband, but the exact scale ledger
shows that it cannot be both cheap and clean at its own centrifugal turning
radius.

\begin{proposition}[Broadband wrapped masks miss the evanescent scale]
\label{prop:wrapped-mask-evanescent-scale-no-go}
Fix $s_*>0$ and $0<\delta<1/2$, put
$S_\delta=[s_*(1-\delta),s_*(1+\delta)]$, and let
$\Theta_L=\operatorname{Arg}(e^{iL\phi})$.  Given $\tau\geq1$ and
$\epsilon>0$, set
\begin{equation}
 T=\frac{\tau}{s_*\epsilon},
 \qquad
 v(r,\omega)=\epsilon\Theta_L(\omega)
       \mathbf1_{[R,R+T]}(r).
 \label{eq:large-area-wrapped-mask}
\end{equation}
After deleting $A_3/r^2$, the terminal state from $e_0$ is
$\psi_s=e^{-i(s/s_*)\tau\Theta_L}e_0$.  If
\begin{equation}
 M_{\tau,L}=L\left\lfloor(1-2\delta)\tau\right\rfloor,
 \label{eq:wrapped-mask-target-degree}
\end{equation}
then, for all sufficiently large $\tau$,
\begin{equation}
 \sup_{s\in S_\delta}
 \|P_{<M_{\tau,L}}\psi_s\|_2^2
 \leq\frac{C_\delta}{\tau}.
 \label{eq:wrapped-mask-broadband-high-mass}
\end{equation}
The charged cost and the centrifugal phase scale at that target degree are
\begin{align}
 \mathcal C_{\tau,L}
 &=\int_R^{R+T}\|M_{v(r)}e_0\|_2^2\dd r
 =\frac{\pi^2\tau^2}{3s_*^2T},
 \label{eq:large-area-wrapped-cost}\\
 \mathcal D_{\tau,L}
 &=M_{\tau,L}(M_{\tau,L}+1)
   \frac{T}{R(R+T)}.
 \label{eq:large-area-wrapped-transport}
\end{align}
Along any sequence with $\tau\to\infty$ and
$\mathcal C_{\tau,L}\to0$, the turning-scale condition
$R\leq C M_{\tau,L}$ forces
$\mathcal D_{\tau,L}\to\infty$.  Consequently, if both the charged cost and
the centrifugal error tend to zero, then
\begin{equation}
 \frac{R}{M_{\tau,L}}\longrightarrow\infty;
 \label{eq:wrapped-mask-far-propagating}
\end{equation}
the modes created in~\eqref{eq:wrapped-mask-broadband-high-mass} are then far
inside their propagating region and cannot supply an exterior centrifugal
barrier in series with the mask.
\end{proposition}

\begin{proof}
Write $a=s/s_*$ and $\beta=a\tau$.  The azimuthal Fourier coefficients of
$e^{-i\beta\Theta_L}$ occur only at frequencies $qL$ and satisfy
\begin{equation}
 |c_q(\beta)|
 =\left|\frac{\sin\pi(\beta-q)}{\pi(\beta-q)}\right|.
 \label{eq:large-area-wrapped-Fourier-coefficients}
\end{equation}
For $|q|<(1-2\delta)\tau$ and
$\beta\in[(1-\delta)\tau,(1+\delta)\tau]$, summing the squared right-hand side
gives
\begin{equation}
 \sum_{|q|<(1-2\delta)\tau}|c_q(\beta)|^2
 \leq \frac{C_\delta}{\tau}.
\end{equation}
A spherical harmonic with azimuthal frequency $qL$ has degree at least
$|q|L$, proving~\eqref{eq:wrapped-mask-broadband-high-mass}.  Since
$(4\pi)^{-1}\|\Theta_L\|_2^2=\pi^2/3$, substitution of
$\epsilon=\tau/(s_*T)$ proves~\eqref{eq:large-area-wrapped-cost}; the second
identity is direct.

Vanishing cost implies $T/\tau^2\to\infty$.  Suppose
$R\leq C M_{\tau,L}$.  If $T\leq R$, then
\begin{equation}
 \mathcal D_{\tau,L}
 \geq\frac{M_{\tau,L}^2T}{2R^2}
 \geq\frac{T}{2C^2}\longrightarrow\infty.
\end{equation}
If $T\geq R$, then
\begin{equation}
 \mathcal D_{\tau,L}
 \geq\frac{M_{\tau,L}^2}{2R}
 \geq\frac{M_{\tau,L}}{2C}\longrightarrow\infty.
\end{equation}
This proves the turning-scale assertion.  If
~\eqref{eq:wrapped-mask-far-propagating} failed, a bounded-ratio subsequence
would contradict that assertion.
\end{proof}

The last sentence has an exact algebraic counterpart.  A clean converter
placed outside an inner amplifier is a right unitary factor of the boundary
Jost matrix, and the scalar Pl\"ucker quotient does not see it at all.

\begin{proposition}[Outer unitary converters are Pl\"ucker-invisible]
\label{prop:outer-unitary-plucker-invisible}
Let $P_0=e_0e_0^*$ and
$\Gamma_0(J)=\|J^{-1}e_0\|^{-1}$.  If a radially inner system has boundary
matrix $J_{\rm in}$ and a reflectionless outer converter contributes a right
factor $U$, so that
\begin{equation}
 J=J_{\rm in}U,
 \label{eq:inner-outer-Jost-factorization}
\end{equation}
then
\begin{equation}
 U^*U=\Id
 \quad\Longrightarrow\quad
 \Gamma_0(J)=\Gamma_0(J_{\rm in}).
 \label{eq:outer-unitary-exact-invariance}
\end{equation}
More generally, if
$\|U^*U-\Id\|\leq\eta<1$, then
\begin{equation}
 \left|\log\Gamma_0(J)-\log\Gamma_0(J_{\rm in})\right|
 \leq-\frac12\log(1-\eta).
 \label{eq:outer-near-unitary-invariance}
\end{equation}
Thus a far-field phase mask may populate arbitrarily high propagating modes,
but it cannot exploit the centrifugal amplification already accumulated to
its inside.  To increase $\Gamma_0$, conversion must overlap the nonunitary
turning dynamics or occur on its inner side.
\end{proposition}

\begin{proof}
Put $y=J_{\rm in}^{-1}e_0$.  Equation
~\eqref{eq:inner-outer-Jost-factorization} gives
$J^{-1}e_0=U^{-1}y$.  This proves
~\eqref{eq:outer-unitary-exact-invariance}.  Under the approximate hypothesis,
the singular values of $U$ lie between $\sqrt{1-\eta}$ and
$\sqrt{1+\eta}$.  Hence
\begin{equation}
 \frac{\|y\|}{\sqrt{1+\eta}}
 \leq\|U^{-1}y\|
 \leq\frac{\|y\|}{\sqrt{1-\eta}},
\end{equation}
and~\eqref{eq:outer-near-unitary-invariance} follows.  The final statement is
the radial ordering encoded by the right factor in
~\eqref{eq:inner-outer-Jost-factorization}.
\end{proof}

The weakest entropy hypothesis needed on this diagonal uses a liminf, not a
supremum over all cutoffs.  Let $\Gamma_{P_N;j}$ be the finite-block quotient
for $H_{R_j,L_j}(V_j)$.

\begin{theorem}[Completion from buffered diagonal entropy]
\label{thm:buffered-diagonal-completion}
Suppose that, for every bounded real $V\in X_3$, every fixed $N$, and every
$K\Subset(0,\infty)$, one can choose~\eqref{eq:buffered-potential-approximants}
and~\eqref{eq:genuine-angular-buffer} so that
\begin{equation}
 \liminf_{j\to\infty}\int_K
       \log^+\Gamma_{P_N;j}(k)\dd k
 \leq C_{K,N}\bigl(1+\mathcal E_N(V)\bigr).
 \label{eq:buffered-diagonal-PSR}
\end{equation}
Then the exterior spectral density has rank at least $N$ for almost every
positive energy.  If~\eqref{eq:buffered-diagonal-PSR} holds for every $N$, then
Simon's conclusion~\eqref{eq:simon} follows.
\end{theorem}

\begin{proof}
Choose a subsequence realizing the finite liminf.  The Pl\"ucker identity and
the compact-probe boundary factor convert~\eqref{eq:buffered-diagonal-PSR} into
a uniform upper bound for the negative logarithmic determinant of the fixed
$N$-probe densities.  Proposition~\ref{prop:buffered-strong-resolvent} gives
 vague convergence of their measures, and the entropy lower-semicontinuity
 estimate~\eqref{eq:audit-lsc} passes the bound to the limiting exterior
 density.  Its determinant is
therefore positive almost everywhere.  The rank/countability argument and the
free-collar equivalence in Proposition~\ref{prop:free-collar} finish exactly as
in Theorem~\ref{thm:completion}.
\end{proof}

There is also a positive result for the dimension-normalized route.  It shows
that matrix size itself is not the obstruction.

\begin{proposition}[Zero-background dimension-normalized matrix sum rule]
\label{prop:zero-background-NMSR}
Let
$\mathsf H_m=-\partial_r^2+W(r)$ be a half-line Schr\"odinger operator with
componentwise Neumann or Dirichlet boundary condition and a compactly supported
real symmetric $m\times m$ potential.  If $h_m$ and $h_{0,m}$ are its matrix
density and the matching free density, then for every
$I\Subset(0,\infty)$,
\begin{equation}
 \frac1m\int_I\log^-
 \det\!\left(h_{0,m}^{-1/2}h_mh_{0,m}^{-1/2}\right)\dd E
 \leq C_I\left(1+\frac1m\int_0^\infty\tr W(r)^2\dd r\right),
 \label{eq:zero-background-NMSR}
\end{equation}
with $C_I$ independent of $m$.  Consequently, for
$W=P_LM_VP_L$ and $m=m_L$ the right-hand side is bounded by
$C_I(1+(4\pi)^{-1}\|V\|_{X_3}^2)$ uniformly in $L$.
\end{proposition}

\begin{proof}
For smooth compactly supported $W$, extend it by zero to the negative
half-line and let $\mathsf S(E)$ and $\mathsf P(E)$ be the Jost transmission
and reflection matrices.  The direct and reciprocal flux identities are
\begin{equation}
 \mathsf S^*\mathsf S-\mathsf P^*\mathsf P=\Id,
 \qquad
 \mathsf S\mathsf S^*-\overline{\mathsf P}\mathsf P^T=\Id.
 \label{eq:matrix-flux-transmission}
\end{equation}
They follow from Wronskian conservation for the real symmetric equation.  In
particular
\begin{equation}
 \mathsf A=\mathsf S^{-1}\overline{\mathsf P},
 \qquad
 \mathsf A\mathsf A^*
 =\Id-\mathsf S^{-1}\mathsf S^{-*}\preceq\Id.
 \label{eq:reflection-contraction}
\end{equation}
For Neumann boundary condition the incoming coefficient of the regular
solution is
$\mathsf B_N=(\mathsf S+\overline{\mathsf P})/2
=\mathsf S(\Id+\mathsf A)/2$; for Dirichlet boundary condition its ratio to
the free incoming coefficient is
$\mathsf S-\overline{\mathsf P}=\mathsf S(\Id-\mathsf A)$.
Since all singular values of $\mathsf A$ are at most one,
\begin{equation}
 |\det(\Id\mathbin\pm\mathsf A)|\leq2^m,
 \qquad
 |\det\mathsf B_N|\leq|\det\mathsf S|.
 \label{eq:regular-transmission-determinant}
\end{equation}

The $m=1$ trace identity of the matrix $L^2$ sum rule
\cite[Theorem~7]{MolchanovVainberg2004} has quadratic coefficient
$\int\tr W^2$ and a favorable nonpositive bound-state term.  On a compact
positive interval it yields
\begin{equation}
 \int_I\log|\det\mathsf S(E)|\dd E
 \leq C_I\int_0^\infty\tr W(r)^2\dd r.
 \label{eq:matrix-transmission-entropy}
\end{equation}
The density formula
$h_m=c(E)\mathsf B^{-1}\mathsf B^{-*}$, divided by the matching free formula,
and~\eqref{eq:regular-transmission-determinant} give pointwise
\begin{equation}
 \log^-\det(h_{0,m}^{-1/2}h_mh_{0,m}^{-1/2})
 \leq 2m\log2+2\log|\det\mathsf S|.
 \label{eq:zero-background-density-transmission}
\end{equation}
Equations~\eqref{eq:matrix-transmission-entropy} and
\eqref{eq:zero-background-density-transmission} prove
~\eqref{eq:zero-background-NMSR}.  Approximation gives the stated $L^2$
class.  Finally,
\begin{equation}
 \frac1{m_L}\tr(P_LM_VP_L)^2
 =\frac1{m_L}\|P_LM_VP_L\|_{\HS}^2
 \leq\frac1{4\pi}\|V(r,\cdot)\|_2^2
\end{equation}
by Proposition~\ref{prop:normalized-cost}.
\end{proof}

The obstruction to inserting the centrifugal background into this proof is
visible already in the first quadratic trace coefficient.

\begin{proposition}[The naive Bessel subtraction has an $L^2$-uncontrolled
linear term]
\label{prop:naive-Bessel-cross-term}
In dimension three let
$D_LY_{nm}=n(n+1)Y_{nm}$ and $m_L=(L+1)^2$.  For every real
$v\in L^2(\mathbb S^2)$,
\begin{equation}
 \frac1{m_L}\tr\!\left(D_LP_LM_vP_L\right)
 =\frac{L(L+2)}{8\pi}\int_{\mathbb S^2}v(\omega)\dd\omega.
 \label{eq:exact-Bessel-cross-trace}
\end{equation}
Consequently, subtracting the free quadratic invariant from the invariant for
$D_L/r^2+P_LM_VP_L$ leaves the cross term
\begin{equation}
 \frac{2}{m_L}\int_1^\infty\frac1{r^2}
   \tr\!\left(D_LP_LM_{V(r)}P_L\right)\dd r
 =\frac{L(L+2)}{4\pi}
   \int_1^\infty\frac1{r^2}
     \int_{\mathbb S^2}V(r,\omega)\dd\omega\dd r.
 \label{eq:divergent-Bessel-cross-term}
\end{equation}
Its coefficient grows quadratically in $L$ and is not bounded by the scalar
$X_3$ norm uniformly in $L$.  Thus a proof of~\eqref{eq:NMSR} must use a
genuinely Bessel-relative determinant or an explicit channelwise phase
renormalization; subtracting two zero-background trace identities is
insufficient.
\end{proposition}

\begin{proof}
The degree-$n$ addition theorem gives
\begin{align}
 \tr(D_LP_LM_vP_L)
 &=\int_{\mathbb S^2}v(\omega)
   \sum_{n=0}^Ln(n+1)\frac{2n+1}{4\pi}\dd\omega\notag\\
 &=\frac{L(L+1)^2(L+2)}{8\pi}
   \int_{\mathbb S^2}v(\omega)\dd\omega.
\end{align}
Divide by $m_L=(L+1)^2$ and expand
$\tr(D_L/r^2+P_LM_VP_L)^2-\tr(D_L/r^2)^2$.
\end{proof}

The divergent term in Proposition~\ref{prop:naive-Bessel-cross-term} has a
precise origin: it is entirely the radial angular mean.  Splitting that mean
into the reference operator removes the term exactly and leaves the original
quadratic cost.

\begin{proposition}[Radial-mean reference cancels every Bessel cross term]
\label{prop:radial-mean-Bessel-cancellation}
For a real scalar potential on $\mathbb S^2$, write
\begin{equation}
 q(r)=\frac1{4\pi}\int_{\mathbb S^2}V(r,\omega)\dd\omega,
 \qquad
 V^\circ(r,\omega)=V(r,\omega)-q(r),
 \qquad
 W_L^\circ=P_LM_{V^\circ}P_L.
 \label{eq:radial-mean-splitting}
\end{equation}
Then
\begin{align}
 \|V(r,\cdot)\|_2^2
 &=4\pi|q(r)|^2+\|V^\circ(r,\cdot)\|_2^2,
 \label{eq:radial-mean-cost-splitting}\\
 \tr W_L^\circ&=0,
 \qquad
 \tr(D_LW_L^\circ)=0,
 \label{eq:radial-mean-trace-cancellation}
\end{align}
and hence, pointwise in $r$,
\begin{equation}
 \frac1{m_L}\tr\!\left[
 \left(\frac{D_L}{r^2}+q\Id+W_L^\circ\right)^2
 -\left(\frac{D_L}{r^2}+q\Id\right)^2\right]
 =\frac1{m_L}\tr(W_L^\circ)^2
 \leq\frac1{4\pi}\|V^\circ(r,\cdot)\|_2^2.
 \label{eq:radial-mean-relative-quadratic-invariant}
\end{equation}
Thus the correct comparison operator for a Bessel-relative sum rule is
\begin{equation}
 H_{q,L}=-\partial_r^2+\frac{D_L}{r^2}+q(r)\Id,
 \label{eq:radial-mean-reference-operator}
\end{equation}
not the bare Bessel operator.
\end{proposition}

\begin{proof}
The first identity is angular orthogonality.  The addition theorem gives
\begin{equation}
 \tr(P_LM_fP_L)=\frac{m_L}{4\pi}\int_{\mathbb S^2}f,
 \qquad
 \tr(D_LP_LM_fP_L)
 =\frac{L(L+1)^2(L+2)}{8\pi}\int_{\mathbb S^2}f.
\end{equation}
Apply these formulas to $f=V^\circ$.  Expanding the two squares leaves only
$\tr(W_L^\circ)^2$, and Proposition~\ref{prop:normalized-cost} proves the
last inequality.
\end{proof}

\begin{proposition}[Signed radial-mean-relative trace identity]
\label{prop:signed-radial-mean-relative-trace}
Assume first that $q$ and $V^\circ$ are smooth and compactly supported in
$(1,\infty)$.  Put
\begin{equation}
 U_{V,L}=\frac{D_L}{r^2}+q\Id+W_L^\circ,
 \qquad
 U_{q,L}=\frac{D_L}{r^2}+q\Id,
\end{equation}
and let $\mathsf S_{V,L}(E)$ and $\mathsf S_{q,L}(E)$ be their matrix Jost
transmission matrices relative to the flat half-line reference.  If
$\{\lambda_a(V,L)\}$ and $\{\lambda_a(q,L)\}$ are the negative eigenvalues,
then
\begin{align}
 &\frac1{\pi m_L}\int_0^\infty E^{1/2}
 \log\left|\frac{\det\mathsf S_{V,L}(E)}
                    {\det\mathsf S_{q,L}(E)}\right|\dd E
 \notag\\
 &\quad=\frac1{m_L}\int_1^\infty\tr(W_L^\circ(r))^2\dd r
 -\frac{2}{3m_L}\left(
   \sum_a|\lambda_a(V,L)|^{3/2}
  -\sum_a|\lambda_a(q,L)|^{3/2}\right).
 \label{eq:signed-radial-mean-relative-trace}
\end{align}
Moreover, its right-hand side is bounded in absolute value by
\begin{equation}
 C\left(\|V\|_{X_3}^2+\|q\|_{L^2(1,\infty)}^2\right)
 \leq C'\|V\|_{X_3}^2,
 \label{eq:signed-relative-trace-uniform-bound}
\end{equation}
with constants independent of $L$.  The identity and bound extend to bounded
$X_3$ data at the level needed here: every smooth compact approximation obeys
the same bound, uniformly in its support and angular cutoff.  No full-space
transmission determinant is asserted by this statement.
\end{proposition}

\begin{proof}
Apply the $m=1$ trace identity of
Molchanov--Vainberg~\cite[Theorem~7]{MolchanovVainberg2004} separately to
$U_{V,L}$ and $U_{q,L}$ and subtract.  One first truncates the common
$D_L/r^2$ tail; for fixed $L$ it lies in $L^1\cap L^2$, so standard Jost and
negative-spectrum convergence passes the identity to the untruncated tail.
Proposition
~\ref{prop:radial-mean-Bessel-cancellation} gives
\begin{equation}
 \int\tr(U_{V,L}^2-U_{q,L}^2)
 =\int\tr(W_L^\circ)^2,
\end{equation}
which proves~\eqref{eq:signed-radial-mean-relative-trace}.  The matrix
Lieb--Thirring inequality and positivity of $D_L/r^2$ give
\begin{align}
 \sum_a|\lambda_a(V,L)|^{3/2}
 &\leq C\int_1^\infty
       \tr(P_LM_VP_L)_-^2\dd r
 \leq C\int_1^\infty\tr(P_LM_VP_L)^2\dd r,\\
 \sum_a|\lambda_a(q,L)|^{3/2}
 &\leq C m_L\int_1^\infty |q(r)|^2\dd r.
\end{align}
Divide by $m_L$ and use Proposition~\ref{prop:normalized-cost} and
~\eqref{eq:radial-mean-cost-splitting}.
\end{proof}

\begin{remark}[The remaining one-sided gap]
Identity~\eqref{eq:signed-radial-mean-relative-trace} controls a signed
logarithm.  It does not control
\begin{equation}
 \int_I\log^+\left|\frac{\det\mathsf S_{V,L}}
                           {\det\mathsf S_{q,L}}\right|,
 \label{eq:relative-transmission-positive-part-gap}
\end{equation}
because an arbitrarily large positive part could cancel an equally large
negative part.  Thus the divergent Bessel coefficient has been removed, but
the fixed-coupling theorem still requires a one-sided relative determinant
estimate using scalar locality, or the fixed-probe estimate
~\eqref{eq:radial-mean-relative-entropy-target} directly.
\end{remark}

For fixed low channels the radial-mean reference already has the required
entropy.  Indeed, it is a finite direct sum of scalar half-line operators with
potentials $q(r)+n(n+1)/r^2\in L^2(1,\infty)$, so the one-dimensional $L^2$
sum rule applies channel by channel.

\begin{proposition}[Exact reduction to radial-mean-relative entropy]
\label{prop:radial-mean-relative-reduction}
Let $V_j$ be the double-buffered approximants and let $h_{V,j,N}$ and
$h_{q,j,N}$ be the fixed $P_N$ probe densities for $H_{R_j,L_j}(V_j)$ and its
radial-mean reference~\eqref{eq:radial-mean-reference-operator}, respectively.
For every fixed $N$ and $I\Subset(0,\infty)$,
\begin{equation}
 \sup_j\int_I\log^-\det h_{q,j,N}(E)\dd E<\infty,
 \label{eq:radial-mean-base-entropy}
\end{equation}
after any fixed smooth positive normalization of the probe density.  Therefore
the cutoff-compatible estimate
\begin{equation}
 \liminf_{j\to\infty}\int_I\log^-\det\!\left(
 h_{q,j,N}^{-1/2}h_{V,j,N}h_{q,j,N}^{-1/2}\right)\dd E
 \leq C_{I,N}\left(1+\|V^\circ\|_{X_3}^2\right)
 \label{eq:radial-mean-relative-entropy-target}
\end{equation}
implies rank at least $N$ for the limiting exterior density.  If
~\eqref{eq:radial-mean-relative-entropy-target} holds for every $N$, it implies
Simon's conclusion~\eqref{eq:simon}.
\end{proposition}

\begin{proof}
For fixed $N$, the finitely many centrifugal potentials $n(n+1)/r^2$ have
bounded $L^2(1,\infty)$ norm, while
$\|q\|_2^2\leq(4\pi)^{-1}\|V\|_{X_3}^2$.  The scalar physical-coupling entropy
bound, applied to the radial truncations of each channel, proves
~\eqref{eq:radial-mean-base-entropy}.  Determinants give the scalar identity
\begin{equation}
 \det h_{V,j,N}
 =\det h_{q,j,N}\,
  \det(h_{q,j,N}^{-1/2}h_{V,j,N}h_{q,j,N}^{-1/2}).
\end{equation}
Using $(x+y)^+\leq x^++y^+$ for the negative logarithms, equations
~\eqref{eq:radial-mean-base-entropy} and
~\eqref{eq:radial-mean-relative-entropy-target} give a uniform subsequential
entropy bound for $h_{V,j,N}$.  Strong-resolvent convergence, entropy lower
semicontinuity~\eqref{eq:audit-lsc}, and the rank/countability finish prove the
claim.
\end{proof}

For $N=1$ one can remove still more uncharged structure than in the
radial-mean determinant.  Keep the entire high--high block in the reference
operator.  The perturbation is then a star with one charged radial column.

\begin{proposition}[Exact star--Feshbach reduction and charged resolvent bound]
\label{prop:radial-mean-star-Feshbach}
Fix a finite angular model containing $e_0=(4\pi)^{-1/2}$ and decompose it as
$P_0\oplus Q_0$.  With the notation of
~\eqref{eq:radial-mean-splitting}, identify $P_0L^2$ with
$L^2(1,\infty)$ and write
\begin{equation}
 H_{V,L}=
 \begin{pmatrix}
  h_q&B_L^*\\ B_L&H_{Q,L}
 \end{pmatrix},
 \qquad
 h_q=-\partial_r^2+q(r),
 \qquad
 (B_Lf)(r)=b_L(r)f(r),
 \label{eq:star-Feshbach-blocks}
\end{equation}
where
\begin{equation}
 b_L(r)=Q_0P_LM_{V^\circ(r)}e_0.
 \label{eq:star-Feshbach-column}
\end{equation}
Then
\begin{equation}
 \int_1^\infty\|b_L(r)\|^2\dd r
 \leq\frac1{4\pi}\|V^\circ\|_{X_3}^2,
 \label{eq:star-Feshbach-column-cost}
\end{equation}
with equality whenever $P_L(V^\circ e_0)=V^\circ e_0$, as happens eventually
for each buffered finite-band approximant.  For $z\in\C\setminus\R$ the exact
compressed resolvent is
\begin{equation}
 P_0(H_{V,L}-z)^{-1}P_0
 =\left[h_q-z-B_L^*(H_{Q,L}-z)^{-1}B_L\right]^{-1}.
 \label{eq:exact-star-Feshbach-resolvent}
\end{equation}

Assume $\|V\|_\infty\leq M$ and choose $\kappa^2>M+1$.  Put
\begin{equation}
 R_q=(h_q+\kappa^2)^{-1},\qquad
 R_Q=(H_{Q,L}+\kappa^2)^{-1},\qquad
 R_{00}=P_0(H_{V,L}+\kappa^2)^{-1}P_0.
\end{equation}
There is a constant $C_{\kappa,M}$, independent of $L$ and of the radial
support, such that
\begin{align}
 &0\leq R_{00}-R_q\in\mathfrak S_1,
 \label{eq:charged-compressed-resolvent-positive}\\
 &\|R_{00}-R_q\|_{\mathfrak S_1}
 \leq C_{\kappa,M}\int_1^\infty\|b_L(r)\|^2\dd r
 \leq\frac{C_{\kappa,M}}{4\pi}\|V^\circ\|_{X_3}^2.
 \label{eq:charged-compressed-resolvent-trace-bound}
\end{align}
Finally, if $J_{V,L}$ is the physical Jost matrix and $j_q$ is the scalar
Jost factor of $h_q$, then the canonical boundary densities obey
\begin{equation}
 \frac{h_{V,L;0}(k^2)}{h_q(k^2)}
 =\left\|\mathfrak r_L(k)\right\|^2,
 \qquad
 \mathfrak r_L(k):=j_q(k)J_{V,L}(k)^{-1}e_0.
 \label{eq:relative-radiation-column-density}
\end{equation}
Consequently, the $N=1$ case of
~\eqref{eq:radial-mean-relative-entropy-target} is exactly the one-sided
boundary estimate
\begin{equation}
 \liminf_j\int_I
 \log^-\|\mathfrak r_{L_j}(\sqrt E)\|^2\dd E
 \leq C_I\left(1+\|V^\circ\|_{X_3}^2\right).
 \label{eq:relative-radiation-column-entropy-target}
\end{equation}
\end{proposition}

\begin{proof}
The mean-zero identity gives
$P_0M_{V^\circ}P_0=0$, so the block decomposition and
~\eqref{eq:exact-star-Feshbach-resolvent} are the ordinary Schur complement.
Also
\begin{equation}
 \|b_L(r)\|^2\leq\|M_{V^\circ(r)}e_0\|_2^2
 =\frac1{4\pi}\|V^\circ(r,\cdot)\|_2^2,
\end{equation}
which proves~\eqref{eq:star-Feshbach-column-cost} and its equality statement.

At $z=-\kappa^2$, both diagonal blocks in
~\eqref{eq:star-Feshbach-blocks} are positive.  The Schur complement is no
larger than $h_q+\kappa^2$, hence inverse monotonicity gives
~\eqref{eq:charged-compressed-resolvent-positive}.  Moreover,
\begin{equation}
 R_{00}-R_q=R_{00}B_L^*R_QB_LR_q.
 \label{eq:charged-compressed-resolvent-identity}
\end{equation}
The vector-valued heat semigroup is dominated in pointwise norm by
$e^{Mt}$ times the scalar Dirichlet heat semigroup.  Therefore the diagonal
kernels of $R_Q^2$ and $R_q^2$ are bounded by a constant depending only on
$\kappa$ and $M$.  It follows that
\begin{align}
 \|R_QB_L\|_{\mathfrak S_2}^2
 &=\int_1^\infty
   \langle b_L(r),R_Q^2(r,r)b_L(r)\rangle\dd r
 \leq C_{\kappa,M}\int\|b_L(r)\|^2\dd r,\\
 \|B_LR_q\|_{\mathfrak S_2}^2
 &=\int_1^\infty R_q^2(r,r)\|b_L(r)\|^2\dd r
 \leq C_{\kappa,M}\int\|b_L(r)\|^2\dd r.
\end{align}
Thus
$B_L^*R_QB_LR_q=(R_QB_L)^*(B_LR_q)$ is trace class with the
claimed charged bound.  Since $\|R_{00}\|\leq(\kappa^2-M)^{-1}$,
~\eqref{eq:charged-compressed-resolvent-identity} proves
~\eqref{eq:charged-compressed-resolvent-trace-bound}.

Finally,~\eqref{eq:audit-jost-density} and its scalar $q$ analogue give
\begin{equation}
 h_{V,L;0}(k^2)=\frac{k}{\pi}\|J_{V,L}(k)^{-1}e_0\|^2,
 \qquad
 h_q(k^2)=\frac{k}{\pi|j_q(k)|^2},
\end{equation}
which proves~\eqref{eq:relative-radiation-column-density} and the final exact
reformulation.
\end{proof}

\begin{proposition}[The star-relative bound-state divisor has the favorable sign]
\label{prop:star-relative-signed-trace}
Assume first that the finite-model coefficients are smooth and compactly
supported.  In the star decomposition~\eqref{eq:star-Feshbach-blocks}, write
\begin{equation}
 H_{\star,L}=\begin{pmatrix}h_q&0\\0&H_{Q,L}\end{pmatrix},
 \qquad
 \mathsf K_L=\begin{pmatrix}0&B_L^*\\B_L&0\end{pmatrix},
 \qquad H_{V,L}=H_{\star,L}+\mathsf K_L.
 \label{eq:star-sign-flip-pair}
\end{equation}
Let $\mathsf S_{V,L}$ and $\mathsf S_{\star,L}$ be the finite matrix
transmission matrices relative to the same flattened reference, and define
\begin{equation}
 \Delta_{3/2,L}:=\sum_a|\lambda_a(H_{V,L})|^{3/2}
 -\sum_a|\lambda_a(H_{\star,L})|^{3/2}.
 \label{eq:star-relative-negative-moment}
\end{equation}
Then
\begin{equation}
 \Delta_{3/2,L}\geq0
 \label{eq:star-relative-negative-moment-sign}
\end{equation}
and the signed second trace identity has the dimension-free upper bound
\begin{align}
 \frac1\pi\int_0^\infty E^{1/2}
  \log\left|\frac{\det\mathsf S_{V,L}(E)}
                       {\det\mathsf S_{\star,L}(E)}\right|\dd E
 &=2\int_1^\infty\|b_L(r)\|^2\dd r
   -\frac23\Delta_{3/2,L}\notag\\
 &\leq2\int_1^\infty\|b_L(r)\|^2\dd r
 \leq\frac1{2\pi}\|V^\circ\|_{X_3}^2.
 \label{eq:star-relative-signed-trace-bound}
\end{align}
Thus neither the uncharged high--high potential nor the self-adjoint negative
spectrum causes the remaining one-sided $N=1$ gap.
\end{proposition}

\begin{proof}
Let $\mathcal Z=P_0-Q_0$ on the angular channel space.  It commutes with
$H_{\star,L}$ and changes the sign of $\mathsf K_L$, so
\begin{equation}
 \mathcal Z(H_{\star,L}+\mathsf K_L)\mathcal Z
 =H_{\star,L}-\mathsf K_L.
\end{equation}
The trace functional
$A\mapsto\operatorname{Tr}(-A)_+^{3/2}$ is convex.  Since
$H_{\star,L}$ is the midpoint of the two unitarily equivalent operators
$H_{\star,L}\pm\mathsf K_L$, convexity gives
~\eqref{eq:star-relative-negative-moment-sign}.  The assertion follows first
on a finite radial interval and then by monotone exhaustion; the compact
support matrix Lieb--Thirring bound makes all displayed moments finite.

At the potential level the star background is block diagonal and
$\mathsf K_L$ is block off-diagonal.  Hence every cross trace vanishes and
\begin{equation}
 \operatorname{tr}\bigl(U_{V,L}^2-U_{\star,L}^2\bigr)
 =\operatorname{tr}\mathsf K_L^2=2\|b_L(r)\|^2.
\end{equation}
Apply the same finite-matrix second trace identity used in Proposition
~\ref{prop:signed-radial-mean-relative-trace} to the two operators and
subtract.  This proves the equality in
~\eqref{eq:star-relative-signed-trace-bound}; its inequalities follow from
~\eqref{eq:star-relative-negative-moment-sign} and
~\eqref{eq:star-Feshbach-column-cost}.
\end{proof}

\begin{remark}[Why the favorable divisor does not yet prove the entropy bound]
\label{rem:star-trace-versus-schur-divisor}
Proposition~\ref{prop:star-relative-signed-trace} removes a tempting but false
diagnosis: the uncontrolled term is not the relative self-adjoint negative
spectrum.  The scalar Feshbach transmission coefficient is a Schur quotient
of the hybrid Jost matrix.  By the Jacobi identity its denominator is one
chosen high-channel minor.  Zeros of that minor create the artificial
meromorphic pole divisor already warned about after
~\eqref{eq:schur-jost-target}.  In a second trace formula its complex poles
carry cubic terms with no fixed sign; this is the destructive term in the
Feshbach--Szeg\H{o} route of Laptev--Naboko--Safronov~\cite{LNS2005}.

The pole-free Pl\"ucker norm, equivalently the full radiation column
$\|J_{V,L}^{-1}e_0\|$, sums all complementary minors and retains the beneficial
denominator in~\eqref{eq:scalar-schur-exact}.  But its positive boundary part
does not follow from the signed determinant bound
~\eqref{eq:star-relative-signed-trace-bound}.  The reciprocal identity
~\eqref{eq:exact-refocusing-entropy} is precisely a pole-free real-axis
replacement: it identifies the missing term as positive state-correlated
refocusing work rather than as an unbounded negative-eigenvalue contribution.
\end{remark}

\begin{remark}[The boundary-value bottleneck]
The trace-class statement
~\eqref{eq:charged-compressed-resolvent-trace-bound} is genuinely
dimension-free and deletes the whole uncharged high--high block.  It is still
an off-axis estimate.  Trace-class control at one negative spectral point does
not by itself bound the negative logarithm of the boundary density, and the
general Hilbert--Schmidt perturbation determinant can even lack
nontangential boundary values~\cite{GPS2007}.  The remaining theorem is now
precise: use the local radial Schr\"odinger structure and centrifugal ordering
to upgrade~\eqref{eq:charged-compressed-resolvent-trace-bound} to
~\eqref{eq:relative-radiation-column-entropy-target}.  No determinant averaged
over the high block is required for $N=1$.
\end{remark}

There is an exact real-axis factorization of the remaining boundary quantity.
It removes every forward phase, including the long weak phases which destroyed
the ungauged Hardy normalization, and leaves one directional reflection
coefficient.

\begin{proposition}[Exact radial-mean two-port gauge and directional reflection]
\label{prop:radial-mean-directional-reflection}
Assume first that the finite-model coefficients are smooth and compactly
supported.  Let $F_{q,L}(r,k)$ be the diagonal outgoing Jost matrix for
$H_{q,L}$ and let $F_{V,L}(r,k)$ be the outgoing Jost matrix for $H_{V,L}$,
both normalized to have Wronskian $2ik\Id$.  For real $k>0$, define coefficient
matrices $\mathsf T_L(r,k)$ and $\mathsf R_L(r,k)$ by
\begin{equation}
 \binom{F_{V,L}}{F_{V,L}'}
 =\begin{pmatrix}
   F_{q,L}&\overline{F_{q,L}}\\
   F_{q,L}'&\overline{F_{q,L}'}
  \end{pmatrix}
  \binom{\mathsf T_L}{\mathsf R_L},
 \label{eq:radial-mean-two-port-expansion}
\end{equation}
with terminal data $\mathsf T_L=\Id$, $\mathsf R_L=0$ beyond the support.
Then
\begin{equation}
 \mathsf T_L^*\mathsf T_L-\mathsf R_L^*\mathsf R_L=\Id
 \label{eq:radial-mean-two-port-flux}
\end{equation}
at every radius.

Put $F=F_{q,L}$ and $W=W_L^\circ$.  The coefficient evolution is exactly
\begin{equation}
 \frac{\dd}{\dd r}\binom{\mathsf T_L}{\mathsf R_L}
 =\frac1{2ik}
 \begin{pmatrix}
  \overline FWF&\overline FW\overline F\\
  -FWF&-FW\overline F
 \end{pmatrix}
 \binom{\mathsf T_L}{\mathsf R_L}.
 \label{eq:radial-mean-two-port-evolution}
\end{equation}
On the real axis the two diagonal blocks in
~\eqref{eq:radial-mean-two-port-evolution} are anti-Hermitian.  Consequently,
if
\begin{align}
 U_+'&=\frac1{2ik}\overline FWFU_+,
 &U_-'&=-\frac1{2ik}FW\overline FU_-,
 &U_+(R)&=U_-(R)=\Id,
 \label{eq:radial-mean-forward-gauges}\\
 \widetilde{\mathsf T}_L&=U_+^{-1}\mathsf T_L,
 &\widetilde{\mathsf R}_L&=U_-^{-1}\mathsf R_L,
\end{align}
then $U_\pm$ are unitary and
\begin{equation}
 \frac{\dd}{\dd r}
 \binom{\widetilde{\mathsf T}_L}{\widetilde{\mathsf R}_L}
 =\begin{pmatrix}0&\mathsf C_L\\\mathsf C_L^*&0\end{pmatrix}
 \binom{\widetilde{\mathsf T}_L}{\widetilde{\mathsf R}_L},
 \qquad
 \mathsf C_L
 =\frac1{2ik}U_+^{-1}\overline FW\overline FU_-.
 \label{eq:reflection-only-two-port-evolution}
\end{equation}
Thus all real-axis norm growth is carried by the reflection-only off-diagonal
system.

At $r=1$, set
\begin{equation}
 \Theta_{q,L}(k)=J_{q,L}(k)^{-1}\overline{J_{q,L}(k)},
 \qquad J_{q,L}(k)=F_{q,L}(1,k).
 \label{eq:radial-mean-boundary-phase}
\end{equation}
The matrix $\Theta_{q,L}$ is diagonal unitary.  With
\begin{align}
 x_L(k)&=
 \left(\mathsf T_L(k)+\Theta_{q,L}(k)\mathsf R_L(k)\right)^{-1}e_0
 =j_q(k)J_{V,L}(k)^{-1}e_0,
 \label{eq:directional-reflection-radiation-column}\\
 \xi_L(k)&=\frac{x_L(k)}{\|x_L(k)\|},
 &y_L(k)&=\|\mathsf R_L(k)\xi_L(k)\|,
 \label{eq:directional-reflection-state}
\end{align}
one has the pointwise sharp estimate
\begin{align}
 \|x_L(k)\|^{-1}
 &\leq\sqrt{1+y_L(k)^2}+y_L(k),
 \label{eq:directional-reflection-pointwise}\\
 \log^-\|x_L(k)\|^2
 &\leq2\operatorname{arsinh}y_L(k).
 \label{eq:directional-reflection-entropy}
\end{align}
Moreover $y_L=\|\widetilde{\mathsf R}_L\xi_L\|$, so the right-hand side is
unchanged by the exact forward gauges.
\end{proposition}

\begin{proof}
Variation of constants in the reference basis gives
~\eqref{eq:radial-mean-two-port-evolution}; the Wronskian of $F$ and
$\overline F$ is $-2ik\Id$ with the displayed convention.  Conservation of
the full outgoing flux and the vanishing mixed reference Wronskians give
~\eqref{eq:radial-mean-two-port-flux}.  Since $W=W^*$ and $F$ is diagonal,
$\overline FWF$ and $FW\overline F$ are Hermitian.  This proves unitarity of
$U_\pm$.  Direct differentiation gives
~\eqref{eq:reflection-only-two-port-evolution}; its lower-left block is the
adjoint of its upper-right block on the real axis.

At the boundary,
\begin{equation}
 J_{V,L}=J_{q,L}\mathsf T_L+\overline{J_{q,L}}\mathsf R_L
 =J_{q,L}(\mathsf T_L+\Theta_{q,L}\mathsf R_L).
\end{equation}
The reference is a direct sum of scalar channels, hence every diagonal entry
of $\Theta_{q,L}$ has modulus one.  Since $J_{q,L}e_0=j_qe_0$, this proves
~\eqref{eq:directional-reflection-radiation-column}.  Now
\begin{align}
 1
 &=\|(\mathsf T_L+\Theta_{q,L}\mathsf R_L)x_L\|\notag\\
 &\leq\|x_L\|\left(
       \|\mathsf T_L\xi_L\|+\|\mathsf R_L\xi_L\|\right)
 =\|x_L\|\left(\sqrt{1+y_L^2}+y_L\right),
\end{align}
where the last equality is~\eqref{eq:radial-mean-two-port-flux}.  This proves
~\eqref{eq:directional-reflection-pointwise}; take logarithms and use
$\operatorname{arsinh}y=\log(\sqrt{1+y^2}+y)$ to obtain
~\eqref{eq:directional-reflection-entropy}.  Finally $U_-$ is unitary, so it
does not change $y_L$.
\end{proof}

\begin{proposition}[Exact reciprocal-refocusing identity]
\label{prop:exact-reciprocal-refocusing}
In the setting of Proposition~\ref{prop:radial-mean-directional-reflection},
define at every radius
\begin{equation}
 \Theta_{q,L}(r,k)=F_{q,L}(r,k)^{-1}
                    \overline{F_{q,L}(r,k)},
 \qquad
 \mathsf A_L(r,k)=\mathsf T_L(r,k)
                   +\Theta_{q,L}(r,k)\mathsf R_L(r,k).
 \label{eq:radial-refocusing-A}
\end{equation}
Then $\mathsf A_L=F_{q,L}^{-1}F_{V,L}$ and, on the real axis,
\begin{align}
 \Theta_{q,L}'(r,k)&=-2ikF_{q,L}(r,k)^{-2},
 \label{eq:radial-reference-phase-velocity}\\
 \mathsf A_L'(r,k)&=\Theta_{q,L}'(r,k)\mathsf R_L(r,k),
 &
 \mathsf R_L'(r,k)&=-\frac1{2ik}F_{q,L}W_L^\circ F_{q,L}
                         \mathsf A_L.
 \label{eq:closed-refocusing-system}
\end{align}
Thus a large coefficient reflection changes the physical radiation column only
after multiplication by the reciprocal phase velocity $\Theta_{q,L}'$.
The flux law also becomes the exact interference identity
\begin{equation}
 \mathsf A_L^*\mathsf A_L-\Id
 =\mathsf A_L^*\Theta_{q,L}\mathsf R_L
  +\mathsf R_L^*\Theta_{q,L}^*\mathsf A_L.
 \label{eq:exact-reflection-interference}
\end{equation}
In particular, the quadratic size $\mathsf R_L^*\mathsf R_L$ cancels: only its
phase-correlated interference with the physical column can change the density.

Fix the boundary-selected unit vector $\xi_L(k)$ from
~\eqref{eq:directional-reflection-state}, and put
\begin{equation}
 a_L(r,k)=\mathsf A_L(r,k)\xi_L(k),
 \qquad
 b_L(r,k)=\mathsf R_L(r,k)\xi_L(k).
 \label{eq:refocusing-selected-state}
\end{equation}
If $R$ lies beyond the support, then $a_L(R,k)=\xi_L(k)$,
$b_L(R,k)=0$, and the exact state-correlated reciprocal work
\begin{equation}
 \mathfrak W^{\rm ref}_L(k)
 :=-\int_1^R
  \Re\frac{\langle a_L(r,k),
             \Theta_{q,L}'(r,k)b_L(r,k)\rangle}
            {\|a_L(r,k)\|^2}\dd r
 \label{eq:exact-refocusing-work}
\end{equation}
satisfies
\begin{equation}
 \mathfrak W^{\rm ref}_L(k)
 =\log\|\mathsf A_L(1,k)\xi_L(k)\|
 =-\log\|x_L(k)\|,
 \qquad
 \log^-\|x_L(k)\|^2=2[\mathfrak W^{\rm ref}_L(k)]_+.
 \label{eq:exact-refocusing-entropy}
\end{equation}
There is also the exact vector refocusing formula
\begin{equation}
 \frac{e_0}{\|x_L(k)\|}
 =\xi_L(k)-\int_1^R
   \Theta_{q,L}'(r,k)\mathsf R_L(r,k)\xi_L(k)\dd r.
 \label{eq:exact-vector-refocusing}
\end{equation}
Consequently, with
\begin{equation}
 \mathfrak Q_L(k):=\int_1^R
  \|\Theta_{q,L}'(r,k)\mathsf R_L(r,k)\xi_L(k)\|\dd r,
 \label{eq:reciprocal-reflection-action}
\end{equation}
one has the phase-sensitive improvement
\begin{equation}
 \log^-\|x_L(k)\|^2
 \leq2\log\bigl(1+\mathfrak Q_L(k)\bigr).
 \label{eq:reciprocal-reflection-entropy-bound}
\end{equation}
\end{proposition}

\begin{proof}
Since $F=F_{q,L}$ is diagonal, its entries never vanish on the positive real
axis: a zero would contradict the nonzero Wronskian of $f_{q,n}$ and
$\overline{f_{q,n}}$.  The identity $F\mathsf A_L=F_{V,L}$ follows directly
from~\eqref{eq:radial-mean-two-port-expansion}.  The scalar Wronskian convention
$f_{q,n}\overline{f_{q,n}}'-f_{q,n}'\overline{f_{q,n}}=-2ik$ gives
~\eqref{eq:radial-reference-phase-velocity}.  Differentiate
$\mathsf A_L=\mathsf T_L+\Theta_{q,L}\mathsf R_L$ and use
~\eqref{eq:radial-mean-two-port-evolution}.  The two terms multiplying
$\mathsf T_L$ cancel, as do the two terms multiplying $\mathsf R_L$; only
$\Theta_{q,L}'\mathsf R_L$ remains.  The lower row of
~\eqref{eq:radial-mean-two-port-evolution}, together with
$F\mathsf A_L=F\mathsf T_L+\overline F\mathsf R_L$, gives the second equation
in~\eqref{eq:closed-refocusing-system}.

Substitute $\mathsf T_L=\mathsf A_L-\Theta_{q,L}\mathsf R_L$ into
~\eqref{eq:radial-mean-two-port-flux}.  Since $\Theta_{q,L}$ is unitary, the
two copies of $\mathsf R_L^*\mathsf R_L$ cancel, which proves
~\eqref{eq:exact-reflection-interference}.

The vector $a_L(r,k)$ cannot vanish: otherwise $F_{V,L}(r,k)\xi_L=0$, which
would make its outgoing flux zero at that radius, contradicting the positive
outgoing flux of the nonzero coefficient vector $\xi_L$.  Hence
\begin{equation}
 \partial_r\log\|a_L(r,k)\|
 =\Re\frac{\langle a_L(r,k),
             \Theta_{q,L}'(r,k)b_L(r,k)\rangle}{\|a_L(r,k)\|^2}.
\end{equation}
Integrate from $1$ to $R$.  The terminal norm is one, while
~\eqref{eq:directional-reflection-radiation-column} gives
$a_L(1,k)=e_0/\|x_L(k)\|$.  This proves
~\eqref{eq:exact-refocusing-entropy}.  Integrating
$\mathsf A_L'=\Theta_{q,L}'\mathsf R_L$ and using
$\mathsf A_L(R)=\Id$ gives~\eqref{eq:exact-vector-refocusing}.  Its triangle
inequality and the first equality in~\eqref{eq:exact-refocusing-entropy} prove
~\eqref{eq:reciprocal-reflection-entropy-bound}.
\end{proof}

\begin{proposition}[Bessel balancing recovers the scalar charged generator]
\label{prop:bessel-balanced-refocusing-generator}
Keep the notation of Proposition~\ref{prop:exact-reciprocal-refocusing} and put
\begin{equation}
 \Omega_{q,L}(r,k)=2k\bigl(F_{q,L}(r,k)^*F_{q,L}(r,k)\bigr)^{-1},
 \qquad
 \mathcal U_{q,L}(r,k)=F_{q,L}(r,k)
  \bigl(F_{q,L}^*F_{q,L}\bigr)^{-1/2}.
 \label{eq:bessel-balanced-omega-phase}
\end{equation}
Thus $\Omega_{q,L}>0$ is diagonal and $\mathcal U_{q,L}$ is diagonal unitary.
Define
\begin{equation}
 \mathsf C_L^{\rm ph}=\Theta_{q,L}\mathsf R_L,
 \qquad
 \mathsf H_L=\frac1{2k}\overline{F_{q,L}}W_L^\circ F_{q,L}.
 \label{eq:physical-return-and-H}
\end{equation}
Then $\mathsf H_L=\mathsf H_L^*$ and the exact refocusing system becomes
\begin{align}
 \mathsf A_L'&=-i\Omega_{q,L}\mathsf C_L^{\rm ph},
 &
 (\mathsf C_L^{\rm ph})'
 &=-i\Omega_{q,L}\mathsf C_L^{\rm ph}+i\mathsf H_L\mathsf A_L,
 \label{eq:bessel-balanced-refocusing-system}\\
 \mathsf A_L(R)&=\Id,
 &\mathsf C_L^{\rm ph}(R)&=0.
\end{align}
Most importantly, its geometrically balanced interaction is exactly
\begin{equation}
 \boxed{
 \Omega_{q,L}^{1/2}\mathsf H_L\Omega_{q,L}^{1/2}
 =\mathcal U_{q,L}^*W_L^\circ\mathcal U_{q,L}.}
 \label{eq:exact-bessel-balanced-generator}
\end{equation}
Consequently every unitarily invariant norm is preserved.  Since
$\mathcal U_{q,L}$ commutes with every angular-degree projection, for every
fixed $N\leq L$,
\begin{equation}
 \|\Omega_{q,L}^{1/2}\mathsf H_L\Omega_{q,L}^{1/2}P_N\|_{\rm HS}
 =\|W_L^\circ P_N\|_{\rm HS},
 \label{eq:balanced-low-column-isometry}
\end{equation}
and in particular
\begin{equation}
 \|\Omega_{q,L}^{1/2}\mathsf H_L\Omega_{q,L}^{1/2}e_0\|^2
 =\|W_L^\circ e_0\|^2
 =\frac1{4\pi}\|P_LV^\circ(r,\cdot)\|_2^2.
 \label{eq:balanced-constant-column-cost}
\end{equation}
If $V^\circ(r,\cdot)$ has angular bandwidth at most $B$, the balanced
generator has the same degree-coupling bandwidth: its matrix element between
degrees $n$ and $m$ vanishes whenever $|n-m|>B$.
\end{proposition}

\begin{proof}
Equation~\eqref{eq:radial-reference-phase-velocity} gives
$\Theta_{q,L}'\Theta_{q,L}^*=-i\Omega_{q,L}$.  Multiply the second equation in
~\eqref{eq:closed-refocusing-system} by $\Theta_{q,L}$ and differentiate the
product $\Theta_{q,L}\mathsf R_L$ to obtain
~\eqref{eq:bessel-balanced-refocusing-system}.  Diagonality of $F_{q,L}$ and
self-adjointness of $W_L^\circ$ give $\mathsf H_L^*=\mathsf H_L$.  Finally,
\begin{align}
 \Omega_{q,L}^{1/2}\mathsf H_L\Omega_{q,L}^{1/2}
 &=\bigl(F_{q,L}^*F_{q,L}\bigr)^{-1/2}
    F_{q,L}^*W_L^\circ F_{q,L}
    \bigl(F_{q,L}^*F_{q,L}\bigr)^{-1/2}\notag\\
 &=\mathcal U_{q,L}^*W_L^\circ\mathcal U_{q,L},
\end{align}
which proves~\eqref{eq:exact-bessel-balanced-generator}.  The phase matrix is
constant on each spherical-harmonic eigenspace of a fixed degree.  It therefore
commutes with $P_N$ and only changes phases of the matrix elements of
$W_L^\circ$.  This proves~\eqref{eq:balanced-low-column-isometry}, preserves the
degree-coupling zero pattern, and gives
~\eqref{eq:balanced-constant-column-cost} because
$W_L^\circ e_0=(4\pi)^{-1/2}P_LV^\circ$.
\end{proof}

\begin{proposition}[Signed pole-free work is sufficient for the scalar endpoint]
\label{prop:signed-radiation-column-reduction}
Fix $\kappa>0$ and put
\begin{equation}
 \varpi_\kappa(E)=\frac{\sqrt E}{(E+\kappa^2)^3},
 \qquad E>0.
 \label{eq:rational-second-trace-weight}
\end{equation}
For a finite radial-mean model define
\begin{equation}
 u_L(E)=\log\|\mathfrak r_L(\sqrt E)\|^2,
 \qquad
 \mathfrak S_{\kappa,L}
 =-\int_0^\infty\varpi_\kappa(E)u_L(E)\dd E.
 \label{eq:signed-radiation-column-functional}
\end{equation}
Assume $\|V\|_\infty\leq M$ and take $\kappa^2>M+1$.  Along any
double-buffered fixed-potential diagonal, the following statements hold.
At stage $j$, the symbols $q$ and $h_q$ below abbreviate the truncated radial
mean $q_j=\mathbf1_{[1,R_j]}q$ and its density; in particular
$\|q_j\|_2\leq\|q\|_2$, so every displayed scalar bound is uniform in $j$.

First, density amplification has the uniform bound
\begin{align}
 \sup_j\int_0^\infty\varpi_\kappa(E)
       u_{L_j}(E)^+\dd E
 &\leq C_{\kappa,M}
   +C_\kappa\bigl(1+\|q\|_2^2\bigr).
 \label{eq:automatic-radiation-amplification}
\end{align}
Second, if the signed pole-free estimate
\begin{equation}
 \liminf_{j\to\infty}\mathfrak S_{\kappa,L_j}
 \leq C_{\kappa,M}
       \left(1+\|V^\circ\|_{X_3}^2\right)
 \label{eq:signed-radiation-column-endpoint}
\end{equation}
holds, then for every $I\Subset(0,\infty)$,
\begin{equation}
 \liminf_{j\to\infty}\int_I
  \log^-\|\mathfrak r_{L_j}(\sqrt E)\|^2\dd E
 \leq C_{I,\kappa,M}
       \left(1+\|V\|_{X_3}^2\right).
 \label{eq:signed-radiation-column-implies-entropy}
\end{equation}
Consequently~\eqref{eq:signed-radiation-column-endpoint}, which is a
\emph{signed} estimate, is sufficient for the $N=1$ fixed-potential entropy
theorem and hence for rank at least one in the limiting exterior density.

In the notation of Proposition~\ref{prop:exact-reciprocal-refocusing},
\begin{equation}
 \mathfrak S_{\kappa,L}
 =2\int_0^\infty\varpi_\kappa(E)
       \mathfrak W_L^{\rm ref}(\sqrt E)\dd E.
 \label{eq:signed-work-functional}
\end{equation}
Thus the remaining scalar endpoint does not require a direct estimate of
$[\mathfrak W_L^{\rm ref}]_+$: its negative part is already paid for by the
ordinary spectral-mass budget.
\end{proposition}

\begin{proof}
The density identity~\eqref{eq:relative-radiation-column-density} gives
\begin{equation}
 e^{u_L(E)}=\frac{h_{V,L;0}(E)}{h_q(E)}.
\end{equation}
For arbitrary positive numbers $a,b$,
\begin{equation}
 \log^+(a/b)\leq a+\log^-b.
 \label{eq:elementary-density-amplification}
\end{equation}
Hence
\begin{equation}
 u_L(E)^+\leq h_{V,L;0}(E)+\log^-h_q(E).
 \label{eq:relative-density-positive-log}
\end{equation}
The harmless choice of a fixed positive normalization for both densities only
changes the right-hand side by an integrable function depending on $\kappa$.

Let $\rho_{V,L;0}$ be the constant-boundary-channel spectral measure.  The
one-dimensional Dirichlet heat-kernel domination estimate, uniform in the
channel dimension because $D_L/r^2\geq0$ and $\|V\|_\infty\leq M$, gives
\begin{equation}
 \int_\mathbb R\frac{\dd\rho_{V,L;0}(E)}{(E+\kappa^2)^2}
 \leq C_{\kappa,M}.
 \label{eq:uniform-boundary-resolvent-square-mass}
\end{equation}
For smooth coefficients this follows by taking both inward boundary
derivatives in
\begin{equation}
 \|e^{-tH_{V,L}}(r,s)\|
 \leq e^{Mt}e^{t\partial_r^2}_{D}(r,s)
\end{equation}
and using
$\partial_r\partial_s e^{t\partial_r^2}_{D}(1,1)\asymp t^{-3/2}$ in the
Laplace representation of $(H_{V,L}+\kappa^2)^{-2}$; approximation gives the
stated bounded-potential version.  Since
$\varpi_\kappa(E)\leq C_\kappa(E+\kappa^2)^{-2}$ for $E\geq0$, this proves
\begin{equation}
 \sup_L\int_0^\infty\varpi_\kappa(E)h_{V,L;0}(E)\dd E
 \leq C_{\kappa,M}.
 \label{eq:uniform-weighted-boundary-mass}
\end{equation}
The scalar $L^2$ sum rule for $h_q$, in its rationally localized form, gives
\begin{equation}
 \int_0^\infty\varpi_\kappa(E)\log^-h_q(E)\dd E
 \leq C_\kappa(1+\|q\|_2^2).
 \label{eq:scalar-rational-entropy-budget}
\end{equation}
Equations~\eqref{eq:relative-density-positive-log}--
\eqref{eq:scalar-rational-entropy-budget} prove
\eqref{eq:automatic-radiation-amplification}.

Now $u^-_L=u^+_L-u_L$.  If
$c_{I,\kappa}=\min_{E\in I}\varpi_\kappa(E)>0$, then
\begin{align}
 c_{I,\kappa}\int_Iu_L(E)^-\dd E
 &\leq\int_0^\infty\varpi_\kappa(E)u_L(E)^-\dd E\notag\\
 &=\int_0^\infty\varpi_\kappa(E)u_L(E)^+\dd E
   +\mathfrak S_{\kappa,L}.
 \label{eq:signed-plus-amplification-decomposition}
\end{align}
Take a subsequence realizing the liminf in
\eqref{eq:signed-radiation-column-endpoint}, use
\eqref{eq:automatic-radiation-amplification}, and recall
$4\pi\|q\|_2^2+\|V^\circ\|_{X_3}^2=\|V\|_{X_3}^2$.  This proves
\eqref{eq:signed-radiation-column-implies-entropy}.  Finally,
\eqref{eq:exact-refocusing-entropy} gives
$u_L=-2\mathfrak W_L^{\rm ref}$, which is
\eqref{eq:signed-work-functional}.
\end{proof}

\begin{proposition}[A fixed forward row majorizes the signed selected loss]
\label{prop:forward-row-majorizes-signed-work}
Let $J_{\star,L}(k)$ be the Jost matrix of the block-diagonal star reference
$H_{\star,L}=h_q\oplus H_{Q,L}$ from
\eqref{eq:star-sign-flip-pair}, and put
\begin{equation}
 \mathsf M_{\star,L}(k)=J_{\star,L}(k)^{-1}J_{V,L}(k).
 \label{eq:star-relative-jost-matrix}
\end{equation}
Because $J_{\star,L}e_0=j_qe_0$, for real $k>0$,
\begin{equation}
 x=\mathsf M_{\star,L}^{-1}e_0
   =\mathfrak r_L(k).
\end{equation}
Equivalently, if $h_{V,L}(k^2)$ is the full canonical boundary-density
matrix, then the star-normalized relative density is
\begin{equation}
 \mathsf W_{\star,L}^{\rm rel}(k^2)
 :=\frac{\pi}{k}J_{\star,L}(k)^*h_{V,L}(k^2)J_{\star,L}(k)
 =\mathsf M_{\star,L}^{-*}\mathsf M_{\star,L}^{-1},
 \label{eq:star-relative-density-matrix}
\end{equation}
and
\begin{equation}
 \langle e_0,\mathsf W_{\star,L}^{\rm rel}e_0\rangle
 =\|\mathfrak r_L(k)\|^2.
 \label{eq:star-relative-density-compression}
\end{equation}
Then
\begin{equation}
 -\log\|x\|^2
 \leq\langle e_0,
       \log(\mathsf M_{\star,L}\mathsf M_{\star,L}^*)e_0\rangle
 \leq\log\|\mathsf M_{\star,L}^*e_0\|^2.
 \label{eq:harmonic-arithmetic-row-majorant}
\end{equation}
Consequently the fixed-diagonal logarithmic estimate
\begin{equation}
 \liminf_{j\to\infty}\int_0^\infty\varpi_\kappa(E)
  \left\langle e_0,
   \log\!\left(\mathsf M_{\star,L_j}(\sqrt E)
                    \mathsf M_{\star,L_j}(\sqrt E)^*\right)e_0
  \right\rangle\dd E
 \leq C_{\kappa,M}\left(1+\|V^\circ\|_{X_3}^2\right)
 \label{eq:star-fixed-diagonal-log-endpoint}
\end{equation}
implies~\eqref{eq:signed-radiation-column-endpoint} and hence the $N=1$
fixed-potential entropy theorem.  Either of the successively stronger bounds
obtained by replacing the integrand in
\eqref{eq:star-fixed-diagonal-log-endpoint} by
\begin{equation}
 \log\|\mathsf M_{\star,L_j}^*e_0\|^2
 \quad\hbox{or}\quad
 \log^+\|\mathsf M_{\star,L_j}^*e_0\|^2
 \label{eq:star-forward-row-stronger-integrands}
\end{equation}
is also sufficient.  The first estimate is the preferred forward
factorization:
the coefficient difference between $H_{V,L}$ and $H_{\star,L}$ is only
$\mathsf K_L$, and
\begin{equation}
 \int_1^\infty\operatorname{tr}\mathsf K_L(r)^2\dd r
 =2\int_1^\infty\|b_L(r)\|^2\dd r
 \leq\frac1{2\pi}\|V^\circ\|_{X_3}^2.
 \label{eq:star-forward-row-charged-coefficient}
\end{equation}
Thus the unresolved strengthening of
Proposition~\ref{prop:star-relative-signed-trace} is a fixed diagonal entry of
the noncommutative logarithm, not another determinant estimate.  In the
relative-density notation it is
\begin{equation}
 \left\langle e_0,
  \log(\mathsf M_{\star,L}\mathsf M_{\star,L}^*)e_0\right\rangle
 =-\left\langle e_0,
  \log\mathsf W_{\star,L}^{\rm rel}e_0\right\rangle.
 \label{eq:star-diagonal-quantum-entropy}
\end{equation}

The same harmonic--arithmetic majorant also holds with the radial-mean
relative matrix $\mathsf A_L(1,k)=J_{q,L}^{-1}J_{V,L}$ in place of
$\mathsf M_{\star,L}$.  For that choice the forward row obeys an exact
fixed-state system.  Put
\begin{equation}
 p_L(r,k)=\mathsf A_L(r,k)^*e_0,
 \qquad
 z_L(r,k)=(\mathsf C_L^{\rm ph}(r,k))^*e_0,
 \qquad
 \omega_0(r,k)=\langle e_0,\Omega_{q,L}(r,k)e_0\rangle.
 \label{eq:fixed-forward-row-state}
\end{equation}
Then
\begin{align}
 p_L'&=i\omega_0 z_L,
 &p_L(R,k)&=e_0,
 \label{eq:fixed-forward-row-system-a}\\
 z_L'&=i\omega_0z_L-i\mathsf A_L^*\mathsf H_Le_0,
 &z_L(R,k)&=0.
 \label{eq:fixed-forward-row-system-b}
\end{align}
If
\begin{equation}
 \mathsf G_L=\Omega_{q,L}^{1/2}\mathsf H_L
                  \Omega_{q,L}^{1/2},
\end{equation}
then its only forcing column has the exact representation
\begin{equation}
 \mathsf H_Le_0
 =\omega_0^{-1/2}\Omega_{q,L}^{-1/2}\mathsf G_Le_0,
 \qquad
 \|\mathsf G_Le_0\|^2
 =\frac1{4\pi}\|P_LV^\circ(r,\cdot)\|_2^2.
 \label{eq:fixed-forward-row-balanced-forcing}
\end{equation}
Thus the radial-mean version using either row integrand in
\eqref{eq:star-forward-row-stronger-integrands} is a fixed-state nonlinear
Fourier estimate whose bare forcing is exactly the charged scalar column.
The boundary-selected vector has disappeared.  The remaining nonlinearity is
the dressing of that column by the full factor $\mathsf A_L^*$ in
\eqref{eq:fixed-forward-row-system-b}; no estimate of this dressing is claimed
here.
\end{proposition}

\begin{proof}
Apply the following argument first with
$\mathsf C=\mathsf M_{\star,L}$ and then, for the asserted radial-mean
variant, with $\mathsf C=\mathsf A_L(1,k)$.  Set
$\mathsf B=\mathsf C\mathsf C^*>0$.  Since
\begin{equation}
 \|x\|^2
 =\langle e_0,\mathsf B^{-1}e_0\rangle,
\end{equation}
operator concavity of the logarithm, applied first to $\mathsf B^{-1}$ and
then to $\mathsf B$, gives
\begin{align}
 -\log\langle e_0,\mathsf B^{-1}e_0\rangle
 &\leq\langle e_0,\log\mathsf B\,e_0\rangle\\
 &\leq\log\langle e_0,\mathsf B e_0\rangle
 =\log\|\mathsf C^*e_0\|^2.
\end{align}
This proves~\eqref{eq:harmonic-arithmetic-row-majorant}.  Integrating its
upper bound against $\varpi_\kappa$ proves the sufficient implication.

Take adjoints in~\eqref{eq:bessel-balanced-refocusing-system}.  Since
$\Omega_{q,L}e_0=\omega_0e_0$, evaluation on $e_0$ gives
\eqref{eq:fixed-forward-row-system-a}--
\eqref{eq:fixed-forward-row-system-b} and their terminal data.  Finally,
$\Omega_{q,L}^{-1/2}e_0=\omega_0^{-1/2}e_0$, so the definition of
$\mathsf G_L$ gives the first identity in
\eqref{eq:fixed-forward-row-balanced-forcing}; the second is
\eqref{eq:balanced-constant-column-cost}.
\end{proof}

\begin{remark}[The row theorem must use scalar closure]
\label{rem:star-row-scalar-closure}
No estimate of the form
\eqref{eq:star-fixed-diagonal-log-endpoint} can hold for an arbitrary
high-channel background using only the off-diagonal cost
\eqref{eq:star-forward-row-charged-coefficient}: the matrix
rotator--barrier counterexample already disproves that abstract statement.
For a scalar multiplier, however, the high block is not independent of the
charged column.  Identifying $b_L(r)$ with its mean-zero angular function and
assuming the buffered finite-band model contains $V^\circ(r,\cdot)$, one has
the exact closure relation
\begin{equation}
 V^\circ(r,\omega)=\sqrt{4\pi}\,b_L(r,\omega),
 \qquad
 Q_0W_L^\circ(r)Q_0
 =\sqrt{4\pi}\,Q_0M_{b_L(r)}Q_0.
 \label{eq:scalar-star-closure}
\end{equation}
Thus the true remaining theorem is a scalar, cutoff-buffered
\emph{self-background} diagonal-log inequality.  Treating $H_{Q,L}$ as an
arbitrary uncharged bath discards exactly the structure which distinguishes
Simon's conjecture from the false matrix analogue.
\end{remark}

\begin{proposition}[The negative-energy Poisson shortcut is false]
\label{prop:negative-energy-row-no-go}
Delete the angular kinetic term, but retain scalar multiplication on
$L^2(\mathbb S^2)$.  Fix $\kappa>0$ and a nonzero real spherical polynomial
$w$ such that
\begin{equation}
 \int_{\mathbb S^2}w\,\dd\omega=0.
 \label{eq:mean-zero-fixed-band-slab}
\end{equation}
There are compactly supported slabs $v_n$ of this one fixed angular bandwidth
for which
\begin{equation}
 \int_0^\infty\|M_{v_n(r)}e_0\|_2^2\dd r\longrightarrow0,
 \qquad
 \log\|J_{v_n}(i\kappa)^*e_0\|_2^2\longrightarrow+\infty,
 \label{eq:imaginary-row-cost-separation}
\end{equation}
where the Jost row is normalized by the free decaying solution.  On every
$K\Subset(0,\infty)$ the same slabs satisfy
\begin{equation}
 \sup_{k\in K}\log^+\|J_{v_n}(k)^*e_0\|_2^2\longrightarrow0.
 \label{eq:real-row-slab-small}
\end{equation}
Moreover, one can concatenate long weak slabs to obtain a single fixed
mean-zero, fixed-band scalar potential $v$ with
\begin{equation}
 \int_0^\infty\|M_{v(r)}e_0\|_2^2\dd r<\infty
 \label{eq:fixed-slab-train-cost}
\end{equation}
whose compact radial truncations $v^{(N)}$ obey
\begin{equation}
 \log\|J_{v^{(N)}}(i\kappa)^*e_0\|_2^2\longrightarrow+\infty.
 \label{eq:fixed-slab-train-imaginary-growth}
\end{equation}
Consequently no estimate of a negative-energy forward row by the charged
$L^2$ cost, even with an additive constant, can replace the real-axis
rational second-trace estimate
\eqref{eq:star-fixed-diagonal-log-endpoint}.
\end{proposition}

\begin{proof}
Work first with one slab
\begin{equation}
 v_n(r,\omega)=\epsilon_n w(\omega)\mathbf1_{[0,T_n]}(r),
 \qquad \epsilon_n=n^{-3},\qquad T_n=n^4.
 \label{eq:single-long-weak-slab}
\end{equation}
At the angular point $\omega$, put
$\lambda_n(\omega)=(\kappa^2+\epsilon_nw(\omega))^{1/2}$.
For all sufficiently large $n$, the free-normalized value at the left end of
the solution which equals $e^{-\kappa r}$ to the right of the slab is
\begin{equation}
 G_n(\omega)=e^{-\kappa T_n}
 \left(\cosh(\lambda_nT_n)
       +\frac{\kappa}{\lambda_n}\sinh(\lambda_nT_n)\right).
 \label{eq:imaginary-constant-slab-jost}
\end{equation}
Choose $a>0$ so that
$A=\{\omega:w(\omega)\geq a\}$ has positive surface measure.  On $A$,
for all large $n$,
\begin{equation}
 \lambda_n-\kappa
 =\frac{\epsilon_nw}{\lambda_n+\kappa}
 \geq\frac{a\epsilon_n}{3\kappa},
 \qquad
 G_n\geq\frac12
     \exp\!\left(\frac{a\epsilon_nT_n}{3\kappa}\right).
 \label{eq:imaginary-slab-row-lower-bound}
\end{equation}
Since the forward Jost operator in the decomposable model is $M_{G_n}$,
\begin{align}
 \log\|M_{G_n}e_0\|_2^2
 &\geq \frac{2an}{3\kappa}
      +\log\frac{|A|}{16\pi},
 \label{eq:imaginary-slab-row-divergence}\\
 \int_0^{T_n}\|M_{v_n(r)}e_0\|_2^2\dd r
 &=\frac{\epsilon_n^2T_n}{4\pi}\|w\|_2^2
   =\frac{\|w\|_2^2}{4\pi n^2}.
 \label{eq:imaginary-slab-vanishing-cost}
\end{align}
This proves~\eqref{eq:imaginary-row-cost-separation}.

For real $k\in K$, let
$q_n(\omega,k)=(k^2-\epsilon_nw(\omega))^{1/2}$.  Direct propagation through
the constant slab gives
\begin{equation}
 F_n(\omega,k)=e^{ikT_n}
 \left(\cos(q_nT_n)-i\frac{k}{q_n}\sin(q_nT_n)\right).
 \label{eq:real-constant-slab-jost}
\end{equation}
Uniformly on $K\times\mathbb S^2$,
\begin{equation}
 |F_n|^2
 =\cos^2(q_nT_n)+\frac{k^2}{q_n^2}\sin^2(q_nT_n)
 \leq1+C_{K,w}\epsilon_n.
 \label{eq:real-slab-row-modulus}
\end{equation}
The multiplier row therefore proves~\eqref{eq:real-row-slab-small}.  Notice
that the large pulse area $\epsilon_nT_n=n$ is a real-axis phase, but becomes
exponential growth at $i\kappa$.

For the fixed-potential assertion, choose $\epsilon_m=\delta2^{-m}$ and
successive intervals $I_m$ of length $\tau/\epsilon_m$, and set
\begin{equation}
 v(r,\omega)=\epsilon_mw(\omega),\qquad r\in I_m.
 \label{eq:fixed-long-weak-slab-train}
\end{equation}
Then
\begin{equation}
 \int_0^\infty\|M_{v(r)}e_0\|_2^2\dd r
 =\frac{\tau\|w\|_2^2}{4\pi}
   \sum_{m=1}^\infty\epsilon_m<\infty.
 \label{eq:fixed-long-weak-slab-train-cost}
\end{equation}
It remains only to verify that the imaginary-energy growth persists through
the interfaces.  For a scalar nonnegative compactly supported function $u$,
let $f$ be its decaying Jost solution at $-\kappa^2$ and put
$g(r)=e^{\kappa r}f(r)$.  The Volterra equation is
\begin{equation}
 g(r)=1+\int_r^\infty
 \frac{1-e^{-2\kappa(s-r)}}{2\kappa}
 u(s)g(s)\dd s.
 \label{eq:positive-imaginary-volterra}
\end{equation}
Its Picard series is positive.  Partition the half-line into intervals of
length $d=\kappa^{-1}$ and retain the parity class carrying at least half of
$\int u$.  Points chosen from distinct retained intervals are separated by
$d$, so the Volterra kernel is at least
$c_\kappa=(1-e^{-2})/(2\kappa)$.  Summing the terms which use at most one
point from each retained interval gives, whenever
$d\|u\|_\infty$ is bounded by a sufficiently small fixed constant,
\begin{equation}
 g(0)\geq C_\kappa
 \prod_j\left(1+c_\kappa\int_{I_{2j}}u(s)\dd s\right)
 \geq C_\kappa
 \exp\!\left(c_\kappa'\int_0^\infty u(s)\dd s\right).
 \label{eq:positive-imaginary-volterra-growth}
\end{equation}
Changing the retained parity only changes $C_\kappa$.

For $\omega\in A$, the first $N$ slabs give a nonnegative scalar potential
$u_N(r)=v^{(N)}(r,\omega)$ with
\begin{equation}
 \int_0^\infty u_N(r)\dd r
 =\tau Nw(\omega)\geq a\tau N.
\end{equation}
Choose $\delta$ small enough for the hypothesis of
\eqref{eq:positive-imaginary-volterra-growth}.  Integration over $A$ now
gives~\eqref{eq:fixed-slab-train-imaginary-growth}.  The construction has
the angular bandwidth of the single fixed polynomial $w$ throughout.
\end{proof}

\begin{proposition}[The transport-free scalar endpoint is true]
\label{prop:decomposable-scalar-entropy}
Delete $A_3/r^2$ and let
$v\in L^2((0,\infty)\times\mathbb S^2;\mathbb R)$ be bounded, initially with
compact radial support.  For a.e. $\omega$, let $J_\omega(k)$ be the scalar
half-line Jost factor for $v(\,cdot\,,\omega)$, and put
\begin{equation}
 H_v(k^2)=\frac1{4\pi}\int_{\mathbb S^2}|J_\omega(k)|^{-2}\dd\omega.
 \label{eq:decomposable-selected-density}
\end{equation}
This is exactly the free-normalized density of the constant angular vector.
For every $I=[E_0,E_1]\Subset(0,\infty)$,
\begin{equation}
 \int_I\log^-H_v(E)\dd E
 \leq 2\sqrt{E_1}\,|K_I|\log2
 +\frac{\sqrt{E_1}}{16E_0}
       \int_0^\infty\!\int_{\mathbb S^2}|v(r,\omega)|^2
       \dd\omega\dd r,
 \qquad
 K_I=[\sqrt{E_0},\sqrt{E_1}].
 \label{eq:decomposable-scalar-entropy}
\end{equation}
The constant is independent of the radial support and the angular bandwidth.
Consequently the estimate extends by $L^2$ approximation to every bounded
$v$ in the stated space.
\end{proposition}

\begin{proof}
The function $\Phi(t)=\log^-t=\max\{0,-\log t\}$ is convex on
$(0,\infty)$.  Jensen's inequality applied to
\eqref{eq:decomposable-selected-density} gives, for a.e. $k>0$,
\begin{equation}
 \log^-H_v(k^2)
 \leq\frac1{4\pi}\int_{\mathbb S^2}
       \log^+|J_\omega(k)|^2\dd\omega.
 \label{eq:decomposable-density-jensen}
\end{equation}
Apply Theorem~\ref{thm:scalar-physical-entropy} in every angular fiber and
then use Fubini.  Since $k\leq\sqrt{E_1}$ and $\dd E=2k\dd k$,
\begin{align}
 \int_I\log^-H_v(E)\dd E
 &\leq\frac{2\sqrt{E_1}}{4\pi}
   \int_{\mathbb S^2}\int_{K_I}
       \log^+|J_\omega(k)|^2\dd k\dd\omega\notag\\
 &\leq2\sqrt{E_1}\,|K_I|\log2
 +\frac{\sqrt{E_1}}{16E_0}
   \int_0^\infty\!\int_{\mathbb S^2}|v(r,\omega)|^2
       \dd\omega\dd r.
\end{align}
This is~\eqref{eq:decomposable-scalar-entropy}.  The scalar spectral measures
converge under compactly supported $L^2$ approximation, and the entropy
lower-semicontinuity argument in~\eqref{eq:audit-lsc} passes the same bound to
the limit.
\end{proof}

\begin{remark}[The exact location of the remaining obstruction]
\label{rem:only-angular-transport-remains}
Propositions~\ref{prop:negative-energy-row-no-go} and
\ref{prop:decomposable-scalar-entropy} separate two facts which a Poisson
argument conflates.  Long weak scalar phases can make an imaginary-energy row
arbitrarily large at negligible cost, but on the real axis the complete
transport-free selected density still obeys the desired entropy estimate.
Thus neither scalar multiplication, angular concentration by itself, nor the
long radial support is the missing mechanism.  The only term absent from
Proposition~\ref{prop:decomposable-scalar-entropy} is $A_3/r^2$; the remaining
theorem must quantify its noncommutative angular transport together with the
reciprocal return in~\eqref{eq:closed-refocusing-system}.
\end{remark}

\begin{proposition}[Exact signed centrifugal interpolation]
\label{prop:exact-centrifugal-interpolation}
Fix a finite cutoff $L$, a compact radial support $[1,R]$, and real $k>0$.
For $0\leq\lambda\leq1$, replace $A_3/r^2$ by
$\lambda A_3/r^2$ in both the physical operator and the radial-mean reference.
Write
\begin{equation}
 F_{\lambda,L}(r,k)
 =\operatorname{diag}\bigl(f_{\lambda,n}(r,k)\bigr)_{0\leq n\leq L},
 \qquad d_n=n(n+1),
 \label{eq:lambda-reference-jost}
\end{equation}
with the usual multiplicities, and define
\begin{align}
 \Omega_\lambda
 &=2k(F_{\lambda,L}^*F_{\lambda,L})^{-1},
 &
 H_\lambda
 &=\frac1{2k}\overline{F_{\lambda,L}}W_L^\circ F_{\lambda,L},
 \label{eq:lambda-balanced-coefficients}\\
 U_\lambda
 &=F_{\lambda,L}(F_{\lambda,L}^*F_{\lambda,L})^{-1/2},
 &
 G_\lambda
 &=\Omega_\lambda^{1/2}H_\lambda\Omega_\lambda^{1/2}
   =U_\lambda^*W_L^\circ U_\lambda.
 \label{eq:lambda-balanced-generator}
\end{align}
Put
\begin{equation}
 \zeta_\lambda=F_{\lambda,L}^{-1}\partial_\lambda F_{\lambda,L},
 \qquad
 Z_\lambda=-iU_\lambda^*\partial_\lambda U_\lambda
            =\Im\zeta_\lambda.
 \label{eq:centrifugal-phase-generator}
\end{equation}
Then
\begin{align}
 \partial_\lambda\Omega_\lambda
 &=-2\Omega_\lambda\Re\zeta_\lambda,
 \label{eq:centrifugal-omega-variation}\\
 \partial_\lambda H_\lambda
 &=\zeta_\lambda^*H_\lambda+H_\lambda\zeta_\lambda,
 \label{eq:centrifugal-H-variation}\\
 \boxed{\partial_\lambda G_\lambda
 =i[G_\lambda,Z_\lambda]},
 \qquad Z_\lambda e_0=0.
 \label{eq:centrifugal-balanced-commutator}
\end{align}
Thus the Bessel-balanced scalar interaction changes only by a diagonal phase
commutator; the exponentially large variation of evanescent Bessel amplitudes
has disappeared.

For the exact refocusing solution, first set
\begin{equation}
 Y_\lambda=\binom{\mathsf A_\lambda}{\mathsf C_\lambda^{\rm ph}},
 \qquad
 \mathcal K_\lambda=
 \begin{pmatrix}
  0&-i\Omega_\lambda\\
  iH_\lambda&-i\Omega_\lambda
 \end{pmatrix},
 \qquad
 Y_\lambda'=\mathcal K_\lambda Y_\lambda,
 \quad Y_\lambda(R)=\binom{\Id}{0}.
 \label{eq:centrifugal-doubled-system}
\end{equation}
Extend this system through the free tail $r>R$ and introduce the reciprocal
amplitude balance
\begin{equation}
 \widehat{\mathsf A}_\lambda
   =\Omega_\lambda^{-1/2}\mathsf A_\lambda,
 \qquad
 \widehat{\mathsf C}_\lambda
   =\Omega_\lambda^{1/2}\mathsf C_\lambda^{\rm ph},
 \qquad
 B_\lambda
   =\Omega_\lambda^{-1/2}(\Omega_\lambda^{1/2})'
   =\frac12\Omega_\lambda^{-1}\Omega_\lambda'.
 \label{eq:reciprocal-amplitude-balance}
\end{equation}
Then, with $\widehat Y_\lambda
=(\widehat{\mathsf A}_\lambda,\widehat{\mathsf C}_\lambda)^\top$,
\begin{equation}
 \widehat Y_\lambda'
 =\widehat{\mathcal K}_\lambda\widehat Y_\lambda,
 \qquad
 \widehat{\mathcal K}_\lambda=
 \begin{pmatrix}
  -B_\lambda&-i\Id\\
  iG_\lambda&B_\lambda-i\Omega_\lambda
 \end{pmatrix},
 \qquad
 \widehat Y_\lambda(\infty)
 =\binom{(2k)^{-1/2}\Id}{0}.
 \label{eq:fully-balanced-doubled-system}
\end{equation}
The terminal vector is independent of $\lambda$, and
\begin{equation}
 \boxed{
 \partial_\lambda\widehat{\mathcal K}_\lambda
 =\begin{pmatrix}
  -\partial_\lambda B_\lambda&0\\
  -[G_\lambda,Z_\lambda]&
  \partial_\lambda B_\lambda-i\partial_\lambda\Omega_\lambda
 \end{pmatrix}.}
 \label{eq:fully-balanced-generator-variation}
\end{equation}
Thus no derivative of the amplitude-weighted interaction $H_\lambda$ remains.
More precisely, the two diagonal reference terms obey
\begin{equation}
 B_\lambda=-\Re(F_{\lambda,L}^{-1}F_{\lambda,L}'),
 \qquad
 \partial_\lambda B_\lambda=-(\Re\zeta_\lambda)',
 \qquad
 \partial_\lambda\Omega_\lambda=-2\Omega_\lambda\Re\zeta_\lambda.
 \label{eq:diagonal-reference-drift-identities}
\end{equation}
In particular, both $\partial_\lambda B_\lambda$ and
$\partial_\lambda\Omega_\lambda$ annihilate $e_0$.

Let $\widehat{\mathcal P}_\lambda(r,s)$ be the propagator of
\eqref{eq:fully-balanced-doubled-system} and let
$\Pi_+=(\Id,0)$.  Differentiation in $\lambda$ gives the exact Duhamel formula
\begin{equation}
 \partial_\lambda\widehat{\mathsf A}_\lambda(1,k)
 =-\Pi_+\int_1^R
   \widehat{\mathcal P}_\lambda(1,r)
   (\partial_\lambda\widehat{\mathcal K}_\lambda(r))
   \widehat Y_\lambda(r)\dd r
 -\Pi_+\int_R^\infty
   \widehat{\mathcal P}_\lambda(1,r)
   (\partial_\lambda\widehat{\mathcal K}_\lambda(r))
   \widehat Y_\lambda(r)\dd r.
 \label{eq:balanced-centrifugal-A-duhamel}
\end{equation}
Equivalently the two integrals form one integral over $[1,\infty)$; the split
only displays that the diagonal inverse-square reference tail remains present
after the potential has vanished.

There is a further exact cancellation.  Put
\begin{equation}
 P_\lambda^{\rm ph}(r,k)=e^{-ikr}U_\lambda(r,k),
 \qquad
 \widetilde Y_\lambda
 =\begin{pmatrix}P_\lambda^{\rm ph}&0\\0&P_\lambda^{\rm ph}\end{pmatrix}
  \widehat Y_\lambda.
 \label{eq:outgoing-phase-gauge}
\end{equation}
Writing
\begin{equation}
 \Gamma_\lambda
 =(P_\lambda^{\rm ph})'(P_\lambda^{\rm ph})^*
 =i\left(\frac12\Omega_\lambda-k\Id\right),
 \label{eq:phase-gauge-connection}
\end{equation}
the phase-gauged system is
\begin{equation}
 \widetilde Y_\lambda'
 =\widetilde{\mathcal K}_\lambda\widetilde Y_\lambda,
 \qquad
 \widetilde{\mathcal K}_\lambda=
 \begin{pmatrix}
  \Gamma_\lambda-B_\lambda&-i\Id\\
  iW_L^\circ&\Gamma_\lambda+B_\lambda-i\Omega_\lambda
 \end{pmatrix},
 \qquad
 \widetilde Y_\lambda(\infty)
 =\binom{(2k)^{-1/2}\Id}{0}.
 \label{eq:phase-gauged-doubled-system}
\end{equation}
In particular the interaction $W_L^\circ$ is now independent of $\lambda$.
The Wronskian identity gives
\begin{equation}
 (\Im\zeta_\lambda)'
 =\frac12\partial_\lambda\Omega_\lambda
 =-\Omega_\lambda\Re\zeta_\lambda.
 \label{eq:phase-amplitude-wronskian-identity}
\end{equation}
Together with~\eqref{eq:diagonal-reference-drift-identities}, this proves the
pure radial-gauge formula
\begin{equation}
 \boxed{
 \partial_\lambda\widetilde{\mathcal K}_\lambda
 =\partial_r\mathcal Q_\lambda,
 \qquad
 \mathcal Q_\lambda
 =\begin{pmatrix}\zeta_\lambda&0\\0&-\zeta_\lambda\end{pmatrix}.}
 \label{eq:pure-radial-gauge-variation}
\end{equation}
Let $\widetilde{\mathcal P}_\lambda(r,s)$ be the propagator of
\eqref{eq:phase-gauged-doubled-system}.  Since
$\mathcal Q_\lambda(\infty)=0$, integration by parts in the exact Duhamel
formula yields
\begin{align}
 \partial_\lambda\widetilde Y_\lambda(1)
 &=\mathcal Q_\lambda(1)\widetilde Y_\lambda(1)
   +\int_1^\infty
    \widetilde{\mathcal P}_\lambda(1,r)
    [\mathcal Q_\lambda(r),\widetilde{\mathcal K}_\lambda(r)]
    \widetilde Y_\lambda(r)\dd r,
 \label{eq:phase-gauge-integrated-duhamel}\\
 [\mathcal Q_\lambda,\widetilde{\mathcal K}_\lambda]
 &=\begin{pmatrix}
    0&-2i\zeta_\lambda\\
    -i(\zeta_\lambda W_L^\circ+W_L^\circ\zeta_\lambda)&0
   \end{pmatrix}.
 \label{eq:charged-refocusing-curvature}
\end{align}
Thus all bare centrifugal variation is a removable endpoint gauge.  The only
interior term is the charged off-diagonal curvature in
\eqref{eq:charged-refocusing-curvature}; it records conversion and reciprocal
return in the same integral.
If
\begin{equation}
 x_\lambda(k)=\mathsf A_\lambda(1,k)^{-1}e_0,
 \qquad
 \ell_\lambda(k)=-\log\|x_\lambda(k)\|^2,
 \label{eq:centrifugal-selected-loss}
\end{equation}
and $\widehat x_\lambda
=\widehat{\mathsf A}_\lambda(1,k)^{-1}e_0$, then
\begin{equation}
 x_\lambda=\omega_0(1,k)^{-1/2}\widehat x_\lambda,
 \qquad
 \ell_\lambda(k)=\log\omega_0(1,k)
 -\log\|\widehat x_\lambda(k)\|^2,
 \label{eq:selected-loss-balanced-equivalence}
\end{equation}
where $\Omega_\lambda e_0=\omega_0e_0$ and $\omega_0$ is independent of
$\lambda$.
If also
\begin{equation}
 \widetilde x_\lambda
 =\widetilde{\mathsf A}_\lambda(1,k)^{-1}e_0,
 \qquad
 \mathcal R_\lambda(k)
 =\Pi_+\int_1^\infty
    \widetilde{\mathcal P}_\lambda(1,r)
    [\mathcal Q_\lambda(r),\widetilde{\mathcal K}_\lambda(r)]
    \widetilde Y_\lambda(r)\dd r,
 \label{eq:charged-refocusing-remainder}
\end{equation}
then $\|\widetilde x_\lambda\|=\|\widehat x_\lambda\|$.  Moreover the
endpoint term in~\eqref{eq:phase-gauge-integrated-duhamel} contributes zero to
the selected logarithmic derivative, because
$\zeta_\lambda e_0=0$ and
$\widetilde{\mathsf A}_\lambda\widetilde x_\lambda=e_0$.  Hence, if
\begin{equation}
 \mathfrak S_{\kappa,L}^{(\lambda)}
 :=\int_0^\infty\varpi_\kappa(E)
       \ell_\lambda(\sqrt E)\dd E.
 \label{eq:lambda-signed-radiation-functional}
\end{equation}
then
\begin{equation}
 \partial_\lambda\ell_\lambda(k)
 =2\Re\frac{\left\langle \widetilde x_\lambda,
  \widetilde{\mathsf A}_\lambda(1,k)^{-1}
  \mathcal R_\lambda(k)\widetilde x_\lambda
  \right\rangle}{\|\widetilde x_\lambda\|^2}.
 \label{eq:centrifugal-selected-loss-variation}
\end{equation}
Equivalently, define the right dressed trajectory and the normalized left
refocusing trajectory by
\begin{align}
 \binom{a_\lambda(r,k)}{c_\lambda(r,k)}
 &:=\widetilde Y_\lambda(r,k)\widetilde x_\lambda(k),
 \label{eq:right-curvature-trajectory}\\
 \binom{p_\lambda(r,k)}{q_\lambda(r,k)}
 &:=\frac{1}{\|\widetilde x_\lambda(k)\|^2}
   \widetilde{\mathcal P}_\lambda(1,r)^*\Pi_+^*
   \widetilde{\mathsf A}_\lambda(1,k)^{-*}
   \widetilde x_\lambda(k).
 \label{eq:left-curvature-trajectory}
\end{align}
They solve the reciprocal equations
$v_\lambda'=\widetilde{\mathcal K}_\lambda v_\lambda$ and
$w_\lambda'=-\widetilde{\mathcal K}_\lambda^*w_\lambda$.  Formula
\eqref{eq:centrifugal-selected-loss-variation} is therefore exactly
\begin{equation}
 \boxed{
 \partial_\lambda\ell_\lambda(k)
 =2\Re\int_1^\infty\left[
  -2i\langle p_\lambda,\zeta_\lambda c_\lambda\rangle
  -i\left\langle q_\lambda,
    (\zeta_\lambda W_L^\circ+W_L^\circ\zeta_\lambda)
    a_\lambda\right\rangle
 \right]\dd r.}
 \label{eq:explicit-signed-curvature-work}
\end{equation}
No undifferentiated centrifugal term or derivative of $W_L^\circ$ occurs.
The first integrand is charged through the return state $c_\lambda$, which
vanishes when $W_L^\circ=0$; the second displays the charged scalar multiplier
explicitly.  Both must remain paired with the reciprocal left trajectory.
In particular, for the rational weight in
\eqref{eq:rational-second-trace-weight},
\begin{equation}
 \mathfrak S_{\kappa,L}^{(1)}
 =\mathfrak S_{\kappa,L}^{(0)}
 +\int_0^1\int_0^\infty
   \varpi_\kappa(E)\,
   \partial_\lambda\ell_\lambda(\sqrt E)
   \dd E\dd\lambda.
 \label{eq:exact-signed-centrifugal-interpolation}
\end{equation}
The transport-free estimate gives
\begin{equation}
 \mathfrak S_{\kappa,L}^{(0)}
 \leq C_{\kappa,M}
 \left(1+\frac1{4\pi}
  \int_1^R\|V(r,\cdot)\|_2^2\dd r\right),
 \label{eq:transport-free-rational-signed-bound}
\end{equation}
uniformly in $L$, $R$, and the angular bandwidth.  Consequently the sole
remaining $N=1$ estimate is the signed centrifugal interpolation condition
$\mathrm{SCI}_1$:
\begin{equation}
 \liminf_{j\to\infty}
 \int_0^1\int_0^\infty
  \varpi_\kappa(E)\,
  \partial_\lambda\ell_{\lambda,L_j}(\sqrt E)
  \dd E\dd\lambda
 \leq C_{\kappa,M}\bigl(1+\|V\|_{X_3}^2\bigr).
 \label{eq:SCI1}
\end{equation}
Along a double-buffered fixed-potential diagonal,
\eqref{eq:SCI1} implies~\eqref{eq:signed-radiation-column-endpoint}, and hence
the fixed-coupling $N=1$ spectral-density lower bound.
\end{proposition}

\begin{proof}
All diagonal entries of $F_{\lambda,L}$ are nonzero on the positive real
axis by their Wronskians, so the displayed derivatives are legitimate.
Differentiating the first two definitions in
\eqref{eq:lambda-balanced-coefficients} gives
\eqref{eq:centrifugal-omega-variation} and
\eqref{eq:centrifugal-H-variation}.  Since
$U_\lambda^*\partial_\lambda U_\lambda=iZ_\lambda$, differentiation of
$G_\lambda=U_\lambda^*W_L^\circ U_\lambda$ gives
\eqref{eq:centrifugal-balanced-commutator}.  The degree-zero reference equation
does not contain $\lambda$, hence $\zeta_\lambda e_0=Z_\lambda e_0=0$.

Equation~\eqref{eq:centrifugal-doubled-system} is
\eqref{eq:bessel-balanced-refocusing-system} with its $\lambda$ dependence
shown explicitly.  Substitution of
\eqref{eq:reciprocal-amplitude-balance} gives
\eqref{eq:fully-balanced-doubled-system}.  Because
$\Omega_\lambda(r,k)\to2k\Id$ as $r\to\infty$ for fixed $L$, its terminal
vector is common to all $\lambda$.  Differentiating its generator and using
\eqref{eq:centrifugal-balanced-commutator} gives
\eqref{eq:fully-balanced-generator-variation}.  The first identity in
\eqref{eq:diagonal-reference-drift-identities} follows from
$\Omega_\lambda=2k|F_{\lambda,L}|^{-2}$; the other two follow by
differentiation.  Variation of constants from infinity gives
\eqref{eq:balanced-centrifugal-A-duhamel}; convergence on the free tail follows
from the inverse-square reference coefficients.  Since
$\mathsf A_\lambda=\Omega_\lambda^{1/2}
\widehat{\mathsf A}_\lambda$ and the degree-zero entry of $\Omega_\lambda$ is
independent of $\lambda$, we obtain
\eqref{eq:selected-loss-balanced-equivalence}.

For~\eqref{eq:phase-gauge-connection}, write each diagonal reference solution
as $f_{\lambda,n}=|f_{\lambda,n}|e^{i\theta_{\lambda,n}}$.  Its outgoing
Wronskian gives
$\theta_{\lambda,n}'=k|f_{\lambda,n}|^{-2}
=\frac12(\Omega_\lambda)_{nn}$, proving both
\eqref{eq:phase-gauge-connection} and
\eqref{eq:phase-amplitude-wronskian-identity}.  Direct substitution then gives
\eqref{eq:phase-gauged-doubled-system} and
\eqref{eq:pure-radial-gauge-variation}.  Differentiating the terminal-value
problem, integrating by parts, and using
\begin{equation*}
 \partial_r\!\left(
  \widetilde{\mathcal P}_\lambda(1,r)\mathcal Q_\lambda(r)
  \widetilde Y_\lambda(r)\right)
 =\widetilde{\mathcal P}_\lambda(1,r)
  \bigl(\mathcal Q_\lambda'
       +[\mathcal Q_\lambda,\widetilde{\mathcal K}_\lambda]\bigr)
  \widetilde Y_\lambda(r)
\end{equation*}
proves~\eqref{eq:phase-gauge-integrated-duhamel}; block multiplication proves
\eqref{eq:charged-refocusing-curvature}.  Finally,
$\partial_\lambda\widetilde x_\lambda
=-\widetilde{\mathsf A}_\lambda^{-1}
(\partial_\lambda\widetilde{\mathsf A}_\lambda)\widetilde x_\lambda$.
The top endpoint block is
$\zeta_\lambda(1)\widetilde{\mathsf A}_\lambda(1)$, and its contribution is
zero because it sends $\widetilde x_\lambda$ to
$\zeta_\lambda(1)e_0=0$.  Differentiating
$-\log\|\widetilde x_\lambda\|^2$ proves
\eqref{eq:centrifugal-selected-loss-variation}.  Move the propagator in
\eqref{eq:charged-refocusing-remainder} to the first slot of the inner product
and insert~\eqref{eq:charged-refocusing-curvature}; this gives
\eqref{eq:explicit-signed-curvature-work}.  Integration in $\lambda$ proves
\eqref{eq:exact-signed-centrifugal-interpolation}.

At $\lambda=0$ the reference is scalar in angle, so the operator is the direct
integral used in Proposition~\ref{prop:decomposable-scalar-entropy}.  Jensen's
inequality followed by the global Deift--Killip transmission bound, now
multiplied by
$2k^2/(k^2+\kappa^2)^3\leq2\kappa^{-6}k^2$, proves
\eqref{eq:transport-free-rational-signed-bound}.  Finally combine
\eqref{eq:exact-signed-centrifugal-interpolation},
\eqref{eq:transport-free-rational-signed-bound}, and Proposition
\ref{prop:signed-radiation-column-reduction}.
\end{proof}

\begin{remark}[Absolute centrifugal variation is the wrong target]
\label{rem:absolute-centrifugal-variation-false}
The commutator in~\eqref{eq:centrifugal-balanced-commutator} must remain inside
the signed state-correlated integral.  Indeed, already on $\mathbb C^2$, take
\begin{equation}
 U_N(\lambda)=\operatorname{diag}(1,e^{iN\lambda}),
 \qquad
 W_N=\delta_N
 \begin{pmatrix}0&1\\1&0\end{pmatrix},
 \qquad G_N=U_N^*W_NU_N.
\end{equation}
Then
\begin{equation}
 \|G_Ne_0\|^2=\delta_N^2,
 \qquad
 \int_0^1\|\partial_\lambda G_N(\lambda)e_0\|\dd\lambda
 =N\delta_N.
\end{equation}
With $\delta_N=N^{-1}$ the charged quadratic column tends to zero while the
absolute variation stays equal to one.  This is an algebraic diagnostic, not a
scalar Schr\"odinger counterexample.

The stronger absolute-curvature failure does have a scalar realization.  Take
zero radial mean and let
\begin{equation}
 p_n(\theta,\phi)=\sin^n\theta\cos(n\phi),
 \qquad
 W_n^\circ(r,\omega)=\chi(r)p_n(\omega),
 \label{eq:highest-weight-curvature-test}
\end{equation}
where $\chi=1$ on a fixed nontrivial shell.  Then $p_n$ is a pure real harmonic
of degree $n$, $\|p_n\|_\infty=1$, and
$\|p_n\|_2^2\asymp n^{-1/2}$.  At any fixed $k>0$ and on that shell, the
degree-$n$ diagonal entry of $\zeta_\lambda$ satisfies
\begin{equation}
 \int_0^1|\zeta_{\lambda,n}(r,k)|\dd\lambda
 \geq\left|\log|f_{1,n}(r,k)|-\log|f_{0,n}(r,k)|\right|
 \gtrsim_{r,k}n\log n.
 \label{eq:centrifugal-log-amplitude-growth}
\end{equation}
Indeed $f_{0,n}=e^{ikr}$, whereas the fixed-argument spherical Hankel
asymptotic gives
$|f_{1,n}(r,k)|\asymp(2n-1)!!(kr)^{-n}$.  Since
$\zeta_\lambda e_0=0$ and $W_n^\circ e_0=(4\pi)^{-1/2}\chi p_n$, the lower-left
block of~\eqref{eq:charged-refocusing-curvature} gives
\begin{equation}
 \int_0^1\int_{\operatorname{supp}\chi}
 \left\|[\mathcal Q_\lambda,
  \widetilde{\mathcal K}_\lambda]\binom{e_0}{0}\right\|
 \dd r\dd\lambda
 \gtrsim n^{3/4}\log n,
 \qquad
 \int\|W_n^\circ e_0\|^2\dd r\asymp n^{-1/2}.
 \label{eq:absolute-charged-curvature-failure}
\end{equation}
Thus even for uniformly bounded compactly supported real scalar multipliers,
no absolute curvature estimate at the charged scale is possible.  Taking an
absolute value before the signed, dressed $\lambda$ and energy integrations
destroys exactly the high-phase and reciprocal-refocusing cancellation which
\eqref{eq:SCI1} must retain.
\end{remark}

\begin{proposition}[Charged negative-energy centrifugal action]
\label{prop:negative-energy-centrifugal-action}
In the finite $N=1$ model of Proposition
\ref{prop:exact-centrifugal-interpolation}, decompose
$P_0\oplus Q_0$ and write
\begin{equation}
 H_{\lambda,L}=\begin{pmatrix}
  h_q&B_L^*\\ B_L&H_{Q,\lambda,L}
 \end{pmatrix},
 \qquad
 H_{Q,\lambda,L}=H_{Q,0,L}+\lambda\frac{D_{Q,L}}{r^2}.
 \label{eq:lambda-star-Feshbach-blocks}
\end{equation}
Assume $\|V\|_\infty\leq M$, choose $\kappa^2>M+1$, and put
\begin{equation}
 R_{Q,\lambda}=(H_{Q,\lambda,L}+\kappa^2)^{-1},
 \qquad
 R_{00}^{(\lambda)}
 =P_0(H_{\lambda,L}+\kappa^2)^{-1}P_0,
 \qquad
 R_q=(h_q+\kappa^2)^{-1}.
 \label{eq:lambda-compressed-negative-resolvents}
\end{equation}
Then
\begin{equation}
 R_{00}^{(\lambda)}
 =\left(h_q+\kappa^2-B_L^*R_{Q,\lambda}B_L\right)^{-1}
 \label{eq:lambda-Feshbach-negative-resolvent}
\end{equation}
is decreasing in $\lambda$ and satisfies the exact derivative identity
\begin{equation}
 -\partial_\lambda R_{00}^{(\lambda)}
 =T_\lambda^*\frac{D_{Q,L}}{r^2}T_\lambda\succeq0,
 \qquad
 T_\lambda=R_{Q,\lambda}B_LR_{00}^{(\lambda)}.
 \label{eq:negative-energy-centrifugal-action-density}
\end{equation}
Moreover,
\begin{align}
 0&\preceq R_{00}^{(0)}-R_{00}^{(1)}
 \preceq R_{00}^{(0)}-R_q,
 \label{eq:negative-energy-centrifugal-order}\\
 \int_0^1\left\|
   \left(\frac{D_{Q,L}}{r^2}\right)^{1/2}T_\lambda
  \right\|_{\mathfrak S_2}^2\dd\lambda
 &=\operatorname{Tr}\bigl(R_{00}^{(0)}-R_{00}^{(1)}\bigr)
 \leq C_{\kappa,M}\int_1^\infty\|b_L(r)\|^2\dd r
 \leq\frac{C_{\kappa,M}}{4\pi}\|V^\circ\|_{X_3}^2.
 \label{eq:charged-negative-energy-action-bound}
\end{align}
The constants are independent of $L$, the angular bandwidth, and the radial
support.
\end{proposition}

\begin{proof}
For finite $L$, $D_{Q,L}/r^2$ is bounded on the exterior half-line, so the
resolvent is norm differentiable in $\lambda$ and
\begin{equation}
 \partial_\lambda R_{Q,\lambda}
 =-R_{Q,\lambda}\frac{D_{Q,L}}{r^2}R_{Q,\lambda}.
\end{equation}
Differentiate the Schur complement in
\eqref{eq:lambda-Feshbach-negative-resolvent}.  This proves
\eqref{eq:negative-energy-centrifugal-action-density} and hence monotonicity.
Because $B_L^*R_{Q,\lambda}B_L\succeq0$, inverse monotonicity also gives
$R_{00}^{(\lambda)}\succeq R_q$.  Consequently
\eqref{eq:negative-energy-centrifugal-order} follows.

The right side of~\eqref{eq:negative-energy-centrifugal-order} is positive
trace class, and Proposition~\ref{prop:radial-mean-star-Feshbach}, applied at
$\lambda=0$, bounds its trace by the charged column.  Integrating
\eqref{eq:negative-energy-centrifugal-action-density}, taking traces, and using
monotone convergence proves~\eqref{eq:charged-negative-energy-action-bound}.
Approximation removes the temporary smoothness needed to justify pointwise
kernel manipulations.
\end{proof}

\begin{remark}[Resolvent order alone cannot yield the entropy]
\label{rem:resolvent-order-not-entropy}
Proposition~\ref{prop:negative-energy-centrifugal-action} is a genuine charged
estimate, but it is not by itself a proof of $\mathrm{SCI}_1$.  The failure is
already visible for abstract scalar measures.  On $[0,1]$, let the reference
density be $h_0=1$ and set, with $\epsilon_n=e^{-n}$,
\begin{equation}
 h_n(E)=\begin{cases}
  \epsilon_n,&0<E<\frac12,\\
  2-\epsilon_n,&\frac12<E<1.
 \end{cases}
 \label{eq:monotone-Stieltjes-entropy-test}
\end{equation}
The masses agree.  For every $\kappa>0$,
\begin{equation}
 \int_0^1\frac{h_0(E)-h_n(E)}{E+\kappa^2}\dd E
 =(1-\epsilon_n)\left(
  \int_0^{1/2}\frac{\dd E}{E+\kappa^2}
  -\int_{1/2}^1\frac{\dd E}{E+\kappa^2}\right)>0,
 \label{eq:all-negative-energy-Stieltjes-order}
\end{equation}
and the left side is uniformly bounded.  The affine path
$h_{n,\lambda}=(1-\lambda)h_0+\lambda h_n$ therefore has decreasing Stieltjes
transform at every negative energy.  Nevertheless,
\begin{equation}
 \int_0^1\log^-h_n(E)\dd E=\frac n2,
 \qquad
 \int_0^1\varpi_\kappa(E)\log^-h_n(E)\dd E
 \geq n\int_0^{1/2}\varpi_\kappa(E)\dd E\longrightarrow\infty.
 \label{eq:Stieltjes-order-entropy-failure}
\end{equation}
Thus even monotonicity at every negative energy plus a uniformly bounded total
resolvent action does not control boundary entropy for general Herglotz data.
Any implication from~\eqref{eq:charged-negative-energy-action-bound} to
$\mathrm{SCI}_1$ must use the local radial Schr\"odinger evolution, not only
operator order or abstract spectral-measure positivity.
\end{remark}

\begin{proposition}[Exact angular heat-flow interpolation]
\label{prop:angular-heat-flow-interpolation}
Let $A_3=-\Delta_\omega$, write $V=q+V^\circ$ with zero angular mean
$V^\circ$, and define
\begin{equation}
 V_t(r,\omega)=q(r)+e^{-tA_3/r^2}V^\circ(r,\cdot)(\omega),
 \qquad t\geq0.
 \label{eq:angular-heat-flow-potential}
\end{equation}
For smooth compact finite-band data this path is real, has the same radial
support, obeys $\|V_t\|_\infty\leq\|V\|_\infty$, and satisfies the exact heat
budget
\begin{equation}
 \int_0^\infty\int_1^R\frac1{r^2}
  \|\nabla_\omega V_t(r,\cdot)\|_2^2\dd r\dd t
 =\frac12\int_1^R\|V^\circ(r,\cdot)\|_2^2\dd r.
 \label{eq:exact-angular-heat-budget}
\end{equation}

Let $F_{t,L}(r,k)$ be the outgoing Jost matrix for
\begin{equation}
 H_{t,L}=-\partial_r^2+\frac{D_L}{r^2}+P_LM_{V_t}P_L,
 \qquad J_{t,L}(k)=F_{t,L}(1,k),
 \label{eq:heat-flow-finite-operator}
\end{equation}
and let $j_q(k)$ be the degree-zero radial-mean Jost factor.  Put
\begin{equation}
 x_t(k)=J_{t,L}(k)^{-1}e_0,
 \qquad
 \ell_t^{\rm h}(k)
 =-\log\bigl(|j_q(k)|^2\|x_t(k)\|^2\bigr).
 \label{eq:heat-flow-selected-loss}
\end{equation}
Then $\ell_\infty^{\rm h}=0$.  To state its derivative, write the phase-space
system for $(F_{t,L},F_{t,L}')^\top$ as
$\mathbb Y_t'=\mathbb K_t\mathbb Y_t$, let
$\mathbb P_t(r,s)$ be its propagator, and set $\Pi_1=(\Id,0)$.  Define right
and left reciprocal trajectories by
\begin{align}
 \binom{\alpha_t(r,k)}{\delta_t(r,k)}
 &=\mathbb Y_t(r,k)x_t(k),
 \label{eq:heat-right-trajectory}\\
 \binom{\gamma_t(r,k)}{\beta_t(r,k)}
 &=\frac1{\|x_t(k)\|^2}\,
   \mathbb P_t(1,r)^*\Pi_1^*J_{t,L}(k)^{-*}x_t(k).
 \label{eq:heat-left-trajectory}
\end{align}
One has the exact identities
\begin{align}
 \partial_t\ell_t^{\rm h}(k)
 &=2\Re\int_1^R
   \langle\beta_t(r,k),M_{r^{-2}A_3V_t^\circ(r)}
          \alpha_t(r,k)\rangle\dd r
 \label{eq:heat-loss-variation}\\
 &=2\Re\int_1^R\frac1{r^2}\int_{\mathbb S^2}
   \nabla_\omega V_t^\circ\cdot
   \nabla_\omega(\overline{\beta_t}\alpha_t)
   \dd\omega\dd r,
 \label{eq:heat-loss-product-gradient}
\end{align}
and consequently
\begin{equation}
 \ell_0^{\rm h}(k)
 =-2\Re\int_0^\infty\int_1^R\frac1{r^2}
   \int_{\mathbb S^2}
   \nabla_\omega V_t^\circ\cdot
   \nabla_\omega(\overline{\beta_t}\alpha_t)
   \dd\omega\dd r\dd t.
 \label{eq:exact-heat-flow-work}
\end{equation}

For the rational weight~\eqref{eq:rational-second-trace-weight}, define the
heat-flow reciprocal Fisher action
\begin{equation}
 \mathfrak D_{\kappa,L}^{\rm h}
 :=\int_0^\infty\int_0^\infty\int_1^R
   \frac{\varpi_\kappa(E)}{r^2}
   \|\nabla_\omega(
     \overline{\beta_t(r,\sqrt E)}\alpha_t(r,\sqrt E))\|_2^2
   \dd r\dd E\dd t.
 \label{eq:heat-reciprocal-Fisher-action}
\end{equation}
Pointwise on the sphere,
\begin{equation}
 \|\nabla_\omega(\overline\beta\alpha)\|_2^2
 \leq2\int_{\mathbb S^2}
  \left(|\beta|^2|\nabla_\omega\alpha|^2
       +|\alpha|^2|\nabla_\omega\beta|^2\right)\dd\omega.
 \label{eq:reciprocal-product-density-action}
\end{equation}
Thus $\mathfrak D_{\kappa,L}^{\rm h}$ is controlled by the sum of the two
reciprocal density actions, with exactly the physical centrifugal weight
$r^{-2}$ and no support-radius loss.
The trajectories have additional structure not visible in
\eqref{eq:reciprocal-product-density-action}.  Set
\begin{align}
 \mathfrak j_t(r,k,\omega)
 &=\overline{\beta_t}\,\partial_r\alpha_t
   -\overline{\partial_r\beta_t}\,\alpha_t
   =\overline{\beta_t}\delta_t+\overline{\gamma_t}\alpha_t,
 \label{eq:reciprocal-Wronskian-density}\\
 \mathfrak J_t(r,k,\omega)
 &=\overline{\beta_t}\nabla_\omega\alpha_t
   -\alpha_t\nabla_\omega\overline{\beta_t}.
 \label{eq:reciprocal-angular-current}
\end{align}
Hamiltonian duality and scalarity of $V_t$ give the exact local conservation
law
\begin{equation}
 \partial_r\mathfrak j_t
 +\frac1{r^2}\operatorname{div}_\omega\mathfrak J_t=0,
 \qquad
 \int_{\mathbb S^2}\mathfrak j_t(r,k,\omega)\dd\omega=1.
 \label{eq:reciprocal-local-Wronskian-law}
\end{equation}
Thus the product gradient in~\eqref{eq:heat-reciprocal-Fisher-action} is the
symmetric partner of a divergence-controlled reciprocal current; this is the
compensation absent from the arbitrary $H^1$ diagnostic below.
Since $\int_0^\infty\varpi_\kappa(E)\dd E=\pi/(8\kappa^3)$,
Cauchy--Schwarz and~\eqref{eq:exact-angular-heat-budget} give
\begin{equation}
 \left|\int_0^\infty\varpi_\kappa(E)
       \ell_0^{\rm h}(\sqrt E)\dd E\right|
 \leq C_\kappa\|V^\circ\|_{X_3}
      (\mathfrak D_{\kappa,L}^{\rm h})^{1/2}.
 \label{eq:heat-flow-entropy-Fisher-bound}
\end{equation}
Therefore the cutoff-compatible estimate
\begin{equation}
 \liminf_{j\to\infty}\mathfrak D_{\kappa,L_j}^{\rm h}
 \leq C_{\kappa,M}\bigl(1+\|V^\circ\|_{X_3}^2\bigr)
 \label{eq:HFI1}
\end{equation}
along the double-buffered fixed-potential diagonal implies the $N=1$ relative
entropy bound and hence the fixed-coupling spectral-density lower bound.
\end{proposition}

\begin{proof}
The spherical heat semigroup is Markov and self-adjoint, proving the reality
and $L^\infty$ assertions.  Expanding $V^\circ$ in nonconstant spherical
harmonics gives
\begin{equation}
 \int_0^\infty\frac1{r^2}
  \|A_3^{1/2}e^{-tA_3/r^2}V^\circ(r)\|_2^2\dd t
 =\frac12\|V^\circ(r)\|_2^2;
\end{equation}
integration in $r$ proves~\eqref{eq:exact-angular-heat-budget}.

The terminal Jost data beyond $R$ are independent of $t$.  Since
$\partial_tV_t^\circ=-r^{-2}A_3V_t^\circ$, variation of constants in the phase-space
system gives
\begin{equation}
 \partial_t\mathbb Y_t(1)
 =-\int_1^R\mathbb P_t(1,r)
   \begin{pmatrix}0&0\\-M_{r^{-2}A_3V_t^\circ(r)}&0\end{pmatrix}
   \mathbb Y_t(r)\dd r.
\end{equation}
Differentiate~\eqref{eq:heat-flow-selected-loss}, move the propagator to the
first slot, and use~\eqref{eq:heat-right-trajectory}--
\eqref{eq:heat-left-trajectory}; this proves~\eqref{eq:heat-loss-variation}.
Angular integration by parts proves~\eqref{eq:heat-loss-product-gradient}.
As $t\to\infty$, $V_t\to q$ and
$J_{t,L}^{-1}e_0=j_q^{-1}e_0$, so
$\ell_\infty^{\rm h}=0$ and integration in $t$ proves
\eqref{eq:exact-heat-flow-work}.  Finally apply Cauchy--Schwarz in
$(t,E,r,\omega)$ and use the heat budget.  Estimate~\eqref{eq:HFI1} then gives
a uniform signed relative entropy bound; Proposition
\ref{prop:radial-mean-relative-reduction} and the diagonal completion theorem
\ref{thm:buffered-diagonal-completion} finish the asserted implication.
\end{proof}

\begin{remark}[Why the heat endpoint is genuinely cutoff-compatible]
\label{rem:heat-endpoint-cutoff-compatible}
For each buffered approximant $V_j$, the whole heat path has angular bandwidth
at most that of $V_j$, while~\eqref{eq:exact-angular-heat-budget} is uniform in
$j$.  Thus the degree-$2L_j$ cutoff-edge lift cannot reappear during the
interpolation.  A bound such as~\eqref{eq:HFI1} is false with uniform
quantifiers over arbitrary unbuffered multiplier families, because together
with~\eqref{eq:heat-flow-entropy-Fisher-bound} it would imply the already false
$\mathrm{PSR}_1$.  Its live content is precisely the fixed-potential,
bandwidth-before-cutoff diagonal order.
\end{remark}

\begin{remark}[Separate angular-energy bounds do not control the heat action]
\label{rem:H1-product-criticality}
The reciprocal relation in~\eqref{eq:heat-left-trajectory} is essential.
Indeed $H^1(\mathbb S^2)$ is not an algebra.  In a fixed polar coordinate disk,
put $L_\epsilon=\log(1/\epsilon)$ and, up to harmless smooth cutoffs, let
\begin{equation}
 f_\epsilon(r)=\begin{cases}
  \sqrt{L_\epsilon},&0<r\leq\epsilon,\\[1mm]
  \displaystyle\frac{\log(1/r)}{\sqrt{L_\epsilon}},
     &\epsilon<r<1,\\[2mm]
 0,&r\geq1.
 \end{cases}
 \label{eq:Moser-product-test}
\end{equation}
Then $\|f_\epsilon\|_2=O(L_\epsilon^{-1/2})$ and
\begin{equation}
 \int|\nabla f_\epsilon|^2=2\pi+o(1),
 \qquad
 \int|\nabla(f_\epsilon^2)|^2
 =\frac{8\pi}{3}L_\epsilon+O(1).
 \label{eq:Moser-product-divergence}
\end{equation}
After adding a constant and normalizing, one obtains smooth
$\alpha_\epsilon=\beta_\epsilon$ with bounded $L^2$ norms and bounded separate
angular energies, but
$\|\nabla(\overline{\beta_\epsilon}\alpha_\epsilon)\|_2\to\infty$.
Consequently~\eqref{eq:HFI1} cannot follow from independent centrifugal-energy
bounds for $\alpha_t$ and $\beta_t$.  A proof must use their exact right/left
Schr\"odinger reciprocity, or establish additional integrability generated by
the scalar multiplication dynamics.  This functional diagnostic is not a
Schr\"odinger counterexample to~\eqref{eq:HFI1}.
\end{remark}

\begin{remark}[What the literature trace formulas actually supply]
\label{rem:matrix-szego-literature-audit}
The established finite-dimensional matrix Szeg\H{o} sum rules are formulated
through $\log\det$ of the matrix spectral weight, equivalently the trace of
its operator logarithm; see, for example, the matrix orthogonal-polynomial
survey of Damanik--Pushnitski--Simon~\cite{DamanikPushnitskiSimon2008}.
Laptev--Naboko--Safronov instead select the scalar radial Weyl measure, but
their second trace develops the complex high-minor pole contribution
described above~\cite{LNS2005}.  Neither result contains the fixed-vector
inequality~\eqref{eq:star-fixed-diagonal-log-endpoint}.  Proposition
\ref{prop:negative-energy-row-no-go} also shows why that missing inequality
cannot be recovered by evaluating a forward row at one point of the upper
half-plane: the uncontrolled exponential type is already present in a
commutative fixed-band scalar model.
\end{remark}

\begin{remark}[What the signed reduction does and does not remove]
\label{rem:signed-vector-versus-determinant}
Estimate~\eqref{eq:signed-radiation-column-endpoint} is strictly weaker than
the original positive-work endpoint, but it is not a consequence of the
star-relative determinant trace.  At the algebraic two-port level, take
\begin{equation}
 \mathsf T=\cosh\rho\,\Id,
 \qquad
 \mathsf R=\sinh\rho\,\operatorname{diag}(1,-1),
 \qquad \Theta=\Id.
\end{equation}
Then $\mathsf T^*\mathsf T-\mathsf R^*\mathsf R=\Id$ and
$\mathsf A=\mathsf T+\Theta\mathsf R=\operatorname{diag}(e^\rho,e^{-\rho})$.
Thus $\det\mathsf A=1$, while for the first coordinate
$x=\mathsf A^{-1}e_0=e^{-\rho}e_0$ and
$-\log\|x\|=\rho$.  A determinant can therefore hide the selected loss by an
opposite high-channel gain even under the exact flux law.

Likewise, positivity and trace control at one negative resolvent point are not
enough at the level of abstract measures.  On $[0,1]$, compare the reference
density $1$ with
\begin{equation}
 h_n(E)=
 \begin{cases}
  2-e^{-n},&0<E<1/2,\\
  e^{-n},&1/2<E<1.
 \end{cases}
\end{equation}
The masses are equal and, for every decreasing Stieltjes weight
$(E+\kappa^2)^{-1}$, the weighted difference is positive and uniformly
bounded, but $\int_0^1\log^-h_n\,\dd E=n/2$.  This is only an abstract
diagnostic, not a Schr\"odinger counterexample; it proves that the local
pole-free propagation structure must enter the passage from
\eqref{eq:charged-compressed-resolvent-trace-bound} to
\eqref{eq:signed-radiation-column-endpoint}.
\end{remark}

\begin{corollary}[A cutoff-compatible refocusing endpoint is sufficient]
\label{cor:reciprocal-refocusing-endpoint}
For the double-buffered approximants, suppose that for every
$I\Subset(0,\infty)$,
\begin{equation}
 \liminf_{j\to\infty}\int_I
  \log\!\left(1+\mathfrak Q_{L_j}(\sqrt E)\right)\dd E
 \leq C_I\left(1+\|V^\circ\|_{X_3}^2\right).
 \label{eq:reciprocal-refocusing-endpoint}
\end{equation}
Then~\eqref{eq:relative-radiation-column-entropy-target} holds for $N=1$, and
the limiting exterior density has rank at least one almost everywhere on
$(0,\infty)$.  Equivalently, the desired relative entropy itself is exactly
twice the positive energy integral of the work in
~\eqref{eq:exact-refocusing-work}.  Proposition
~\ref{prop:signed-radiation-column-reduction} shows that this sufficient
condition is stronger than necessary: after rational energy localization, a
charged bound on the \emph{signed} work already suffices.
\end{corollary}

\begin{proof}
Combine~\eqref{eq:reciprocal-reflection-entropy-bound},
~\eqref{eq:relative-radiation-column-density}, and Proposition
~\ref{prop:radial-mean-relative-reduction}.
\end{proof}

\begin{corollary}[A positive directional-reflection endpoint is sufficient]
\label{cor:directional-reflection-endpoint}
For the double-buffered approximants, suppose that for every
$I\Subset(0,\infty)$,
\begin{equation}
 \liminf_{j\to\infty}\int_I
 \operatorname{arsinh}\!\left(
  \|\mathsf R_{L_j}(\sqrt E)\xi_{L_j}(\sqrt E)\|
 \right)\dd E
 \leq C_I\left(1+\|V^\circ\|_{X_3}^2\right).
 \label{eq:directional-reflection-endpoint}
\end{equation}
Then~\eqref{eq:relative-radiation-column-entropy-target} holds for $N=1$, and
the limiting exterior density has rank at least one almost everywhere on
$(0,\infty)$.
\end{corollary}

\begin{proof}
Combine~\eqref{eq:directional-reflection-entropy},
~\eqref{eq:relative-radiation-column-density}, and Proposition
~\ref{prop:radial-mean-relative-reduction}.
\end{proof}

\begin{proposition}[Exact hyperbolic-angle identity and bilinear density target]
\label{prop:directional-reflection-hyperbolic-angle}
Fix real $k>0$ and keep the unit vector $\xi_L(k)$ from
~\eqref{eq:directional-reflection-state} fixed as $r$ varies.  Put
\begin{equation}
 t(r,k)=\widetilde{\mathsf T}_L(r,k)\xi_L(k),
 \qquad
 s(r,k)=\widetilde{\mathsf R}_L(r,k)\xi_L(k),
 \qquad
 \rho(r,k)=\operatorname{arsinh}\|s(r,k)\|.
 \label{eq:directional-hyperbolic-state}
\end{equation}
Then $\|t\|^2-\|s\|^2=1$.  At points where $s\ne0$, define
\begin{equation}
 \eta_+(r,k)=\frac{t(r,k)}{\|t(r,k)\|},
 \qquad
 \eta_-(r,k)=\frac{s(r,k)}{\|s(r,k)\|}.
\end{equation}
The reflection-only system~\eqref{eq:reflection-only-two-port-evolution}
gives the exact identity
\begin{equation}
 \partial_r\rho(r,k)
 =\Re\langle\eta_+(r,k),\mathsf C_L(r,k)\eta_-(r,k)\rangle.
 \label{eq:directional-hyperbolic-derivative}
\end{equation}
Consequently, with $R$ beyond the support, the terminal condition gives the
exact state-correlated return work
\begin{equation}
 \rho(1,k)
 =-\int_1^R
 \Re\langle\eta_+(r,k),\mathsf C_L(r,k)\eta_-(r,k)\rangle\dd r
 \geq0.
 \label{eq:directional-state-correlated-return-work}
\end{equation}
In particular,
\begin{equation}
 \rho(1,k)
 \leq\frac1{2k}\int_1^R
 \|V^\circ(r,\cdot)\|_2
 \|\Psi_+(r,k)\Psi_-(r,k)\|_2\dd r,
 \label{eq:directional-bilinear-density-bound}
\end{equation}
where
\begin{equation}
 \Psi_+(r,k)=F_{q,L}(r,k)U_+(r,k)\eta_+(r,k),
 \qquad
 \Psi_-(r,k)=\overline{F_{q,L}(r,k)}U_-(r,k)\eta_-(r,k).
 \label{eq:directional-dressed-waves}
\end{equation}
The product in~\eqref{eq:directional-bilinear-density-bound} is pointwise on
$\mathbb S^2$.  Values of $\eta_-$ on the set $s=0$ may be chosen arbitrarily;
they do not contribute to the integrated identity.

In particular, define along the buffered diagonal
\begin{equation}
 \mathfrak B_{I,j}
 :=\int_I\int_1^{R_j}\int_{\mathbb S^2}
 |\Psi_{+,j}(r,\sqrt E,\omega)|^2
 |\Psi_{-,j}(r,\sqrt E,\omega)|^2
 \dd\omega\dd r\dd E.
 \label{eq:directional-bilinear-density-action}
\end{equation}
If
\begin{equation}
 \limsup_{j\to\infty}\mathfrak B_{I,j}
 \leq C_I\left(1+\|V^\circ\|_{X_3}^2\right),
 \label{eq:directional-bilinear-density-target}
\end{equation}
then~\eqref{eq:directional-reflection-endpoint} holds and the limiting
exterior density has rank at least one almost everywhere.
\end{proposition}

\begin{proof}
The flux identity gives $\|t\|^2-\|s\|^2=1$.  Using
$s'=\mathsf C_L^*t$ and $\|t\|=\sqrt{1+\|s\|^2}$ gives
\begin{equation}
 \partial_r\operatorname{arsinh}\|s\|
 =\Re\langle\eta_-,\mathsf C_L^*\eta_+\rangle,
\end{equation}
which is~\eqref{eq:directional-hyperbolic-derivative}.  Since $s(R,k)=0$,
integrating inward proves
~\eqref{eq:directional-state-correlated-return-work}; taking the absolute
value inside the integral gives an upper bound.  From
~\eqref{eq:reflection-only-two-port-evolution}, scalar multiplication gives
\begin{align}
 |\langle\eta_+,\mathsf C_L\eta_-\rangle|
 &=\frac1{2k}\left|
   \int_{\mathbb S^2}V^\circ(r,\omega)
   \overline{\Psi_+(r,k,\omega)}\Psi_-(r,k,\omega)\dd\omega
  \right|\notag\\
 &\leq\frac1{2k}\|V^\circ(r,\cdot)\|_2
       \|\Psi_+(r,k)\Psi_-(r,k)\|_2.
\end{align}
This proves~\eqref{eq:directional-bilinear-density-bound}.  Finally integrate
over $E\in I$, use $k\geq\sqrt{\inf I}$, and apply Cauchy--Schwarz on
$I\times(1,R_j)\times\mathbb S^2$:
\begin{equation}
 \int_I\rho(1,\sqrt E)\dd E
 \leq C_I\sqrt{|I|}\,\|V^\circ\|_{X_3}
       \mathfrak B_{I,j}^{1/2}.
\end{equation}
Under~\eqref{eq:directional-bilinear-density-target}, the right-hand side is
bounded by $C_I(1+\|V^\circ\|_{X_3}^2)$.  Apply Corollary
~\ref{cor:directional-reflection-endpoint}.
\end{proof}

\begin{proposition}[An ideal scalar shell exposes the evanescent-coordinate artifact]
\label{prop:ideal-shell-reflection-artifact}
Fix an observation radius $a>1$ and a real $k>0$.  In a finite diagonal
reference, insert the ideal shell
$W(r)=g\delta_a(r)W_0$, where $W_0=W_0^*$, and normalize the coefficient
matrices to $(\mathsf T,\mathsf R)=(\Id,0)$ just outside the shell.  Then the
coefficients just inside it are exactly
\begin{align}
 \mathsf T(a^-)&=\Id-\frac{g}{2ik}
   \overline{F(a)}W_0F(a),
 &
 \mathsf R(a^-)&=\frac{g}{2ik}F(a)W_0F(a),
 \label{eq:ideal-shell-coefficients}
\end{align}
and therefore
\begin{equation}
 \mathsf A(a^-)=\mathsf T(a^-)+\Theta(a)\mathsf R(a^-)=\Id.
 \label{eq:ideal-shell-exact-cancellation}
\end{equation}
Below the shell the coefficient matrices are constant and
\begin{equation}
 \mathsf A(r)=\Id+\left(\Theta(r)-\Theta(a)\right)\mathsf R(a^-),
 \qquad 1\leq r<a.
 \label{eq:ideal-shell-refocusing-drift}
\end{equation}

This cancellation occurs for scalar angular multiplication.  Indeed, take the
free reference, let $e_0=Y_{00}$, choose a real normalized spherical harmonic
$Y_{\ell0}$, and set
$w_\ell=\sqrt{4\pi}Y_{\ell0}$ and
$W_0=P_LM_{w_\ell}P_L$ with $L\geq\ell$.  Then
$W_0e_0=Y_{\ell0}$ and
\begin{equation}
 \|\mathsf R(a^-)e_0\|
 =\frac{|g|}{2k}|f_0(a,k)f_\ell(a,k)|
 \longrightarrow\infty
 \qquad(\ell\to\infty),
 \label{eq:ideal-shell-raw-reflection-diverges}
\end{equation}
whereas at the observation plane $x=e_0$ and the relative density ratio is
exactly one.  Thus raw directional reflection, and any estimate that replaces
the exact boundary interference by its hyperbolic angle, can arbitrarily
overcount a harmless evanescent coefficient.
\end{proposition}

\begin{proof}
Continuity across the shell and the derivative jump
$F_{V,L}'(a^+)-F_{V,L}'(a^-)=gW_0F_{V,L}(a)$ give
~\eqref{eq:ideal-shell-coefficients}; equivalently, integrate the coefficient
system across the atom.  Its atomic generator squares to zero because
$(F,\overline F)(\overline F,-F)^{\mathsf T}=0$.  Since
$\Theta(a)F(a)=\overline{F(a)}$, the two shell terms cancel exactly in
~\eqref{eq:ideal-shell-exact-cancellation}.  Equation
~\eqref{eq:ideal-shell-refocusing-drift} follows because
$\mathsf T$ and $\mathsf R$ are constant where $W=0$.

For the scalar example, $w_\ell e_0=Y_{\ell0}$.  Hence the only component of
$\mathsf R(a^-)e_0$ is the displayed degree-$\ell$ component.  With
$f_n(r,k)=krh_n^{(1)}(kr)$, the positive-coefficient identity
~\eqref{eq:spherical-Hankel-amplitude} gives
\begin{equation}
 |f_\ell(a,k)|^2
 \geq \frac{((2\ell-1)!!)^2}{(ka)^{2\ell}},
\end{equation}
which diverges for fixed $ka$.  On the other hand,
~\eqref{eq:ideal-shell-exact-cancellation} gives $x=e_0$ at the observation
plane.  This proves the claim.
\end{proof}

\begin{remark}[Scope of the shell diagnostic]
The atom $g\delta_a w_\ell$ is not an admissible $X_3$ potential, so
Proposition~\ref{prop:ideal-shell-reflection-artifact} is neither a
counterexample to~\eqref{eq:reciprocal-refocusing-endpoint} nor a disproof of
Simon's conjecture.  It is an exact conditioning test.  It proves that the
unweighted endpoint~\eqref{eq:directional-reflection-endpoint} is not the
intrinsic quantity to attack.  A proof may first eliminate elliptic channels
by the outgoing Feshbach estimate, or work directly with the exact reciprocal
factor $\Theta'=-2ikF^{-2}$ in
~\eqref{eq:exact-refocusing-work}.  Either route retains the cancellation in
~\eqref{eq:ideal-shell-exact-cancellation}.
\end{remark}

\begin{remark}[What remains after the exact gauge]
The directional endpoint~\eqref{eq:directional-reflection-endpoint} remains a
valid sufficient condition, and its hyperbolic derivative contains neither a
signed logarithm nor a forward WKB phase.  It is nevertheless too strong to be
the canonical target.  Proposition~\ref{prop:ideal-shell-reflection-artifact}
shows that it can count the conditioning of an elliptic Bessel basis rather
than physical loss.  The same warning applies to the bilinear density action
~\eqref{eq:directional-bilinear-density-action}.

The exact cutoff-compatible target is instead
~\eqref{eq:exact-refocusing-entropy}: density loss is the positive part of one
state-correlated reciprocal work.  Its factor
$|\Theta_n'|=2k/|f_{q,n}|^2$ automatically suppresses channels below their
turning radius and preserves boundary refocusing.  The buffered high-channel
coercivity still removes the Galerkin edge.  Proposition
~\ref{prop:signed-radiation-column-reduction} makes the remaining theorem
strictly sharper: prove a scalar charged bound on the rationally weighted
\emph{signed} refocusing work inside the growing active block.  Density
amplification pays for its negative part automatically.  Bounding the two
factors in the work separately would again destroy the correlation which the
exact identity retains.
\end{remark}

Proposition~\ref{prop:zero-background-NMSR} does \emph{not} prove
~\eqref{eq:NMSR} for the physical finite models: their reference operator is
$-\partial_r^2+D_L/r^2$, and Proposition
~\ref{prop:naive-Bessel-cross-term} shows that ordinary background subtraction
still destroys dimension uniformity.  Proposition
~\ref{prop:radial-mean-Bessel-cancellation} removes that algebraic divergence.
For $N=1$, Proposition~\ref{prop:radial-mean-star-Feshbach} sharpens the next
positive target to the relative-radiation-column estimate
~\eqref{eq:relative-radiation-column-entropy-target}.  For general $N$ the
target remains the radial-mean-relative estimate
~\eqref{eq:radial-mean-relative-entropy-target}, or its dimension-normalized
version followed by fixed-subspace tightness~\eqref{eq:TIGHT}.  None of these
conclusions follows from strong resolvent convergence, one off-axis
trace-class bound, or the signed trace coefficient alone.

\section{The first full-space converter obstruction}
\label{app:scalar-converter-obstruction}

The preceding corollary leaves open whether angular dynamics repairs the
non-invariance of the compressed lift.  It is useful first to remove that
dynamics completely.  The resulting model has an exact obstruction which is
independent of the radial profile and of its cost.

Before doing so, one can quantify exactly what the minimum lift loses into the
complement.  This already rejects the most literal adiabatic use of the
compressed rotator.

\begin{proposition}[Complementary leakage of the minimum lift]
\label{prop:minimum-lift-complement-leakage}
Let $S=4\pi$, let $h_\ell$ be~\eqref{eq:real-highest-weight}, and let
$\Pi_\ell$ be the projection onto $\operatorname{span}\{e_0,h_\ell\}$.  For
$v_{a,b,d}$ in~\eqref{eq:minimum-scalar-lift},
\begin{align}
 \|(I-\Pi_\ell)M_{v_{a,b,d}}e_0\|_2^2
 &=\frac{(d-a)^2}{S G_\ell},
 \label{eq:minimum-lift-low-leakage}\\
 \|(I-\Pi_\ell)M_{v_{a,b,d}}h_\ell\|_2^2
 &=S b^2G_\ell
   +\frac{(d-a)^2}{G_\ell^2}
      \bigl(H_{6,\ell}-H_{4,\ell}^2\bigr).
 \label{eq:minimum-lift-high-leakage}
\end{align}
Moreover,
\begin{equation}
 \frac{H_{6,\ell}-H_{4,\ell}^2}{G_\ell^2}
 \longrightarrow \frac{20}{9\sqrt3}-1.
 \label{eq:minimum-lift-diagonal-leakage-limit}
\end{equation}
For the lifted rotator in Corollary
~\ref{cor:scalar-rotator-barrier-lift}, on every fixed interior portion of the
mixing region on which $|\sin\theta\cos\theta|\geq c>0$,
\begin{equation}
 \|(I-\Pi_\ell)M_{v_{a,b,d}}h_\ell\|_2
 \geq c'\delta_\Lambda\ell^{1/4}.
 \label{eq:minimum-lift-gap-failure}
\end{equation}
Thus the complementary coupling divided by the intended two-channel gap
$\delta_\Lambda$ diverges like $\ell^{1/4}$.
\end{proposition}

\begin{proof}
Write $g_\ell=h_\ell^2-S^{-1}$ and
\begin{equation}
 w_\ell=g_\ell h_\ell-G_\ell h_\ell.
\end{equation}
Direct substitution in~\eqref{eq:minimum-scalar-lift} gives
\begin{align}
 (I-\Pi_\ell)M_{v_{a,b,d}}e_0
 &=\frac{d-a}{\sqrt S\,G_\ell}g_\ell,
 \label{eq:minimum-lift-low-vector}\\
 (I-\Pi_\ell)M_{v_{a,b,d}}h_\ell
 &=\sqrt S\,b\,g_\ell+\frac{d-a}{G_\ell}w_\ell.
 \label{eq:minimum-lift-high-vector}
\end{align}
The azimuthal frequencies of $g_\ell$ are $0$ and $2\ell$, while those of
$w_\ell$ are $\ell$ and $3\ell$, so the two terms in
~\eqref{eq:minimum-lift-high-vector} are orthogonal.  Also
\begin{equation}
 \|w_\ell\|_2^2
 =\|g_\ell h_\ell\|_2^2-G_\ell^2
 =H_{6,\ell}-H_{4,\ell}^2.
\end{equation}
This proves~\eqref{eq:minimum-lift-low-leakage} and
~\eqref{eq:minimum-lift-high-leakage}.  The asymptotics in
~\eqref{eq:highest-H4}--\eqref{eq:highest-H6} give
\eqref{eq:minimum-lift-diagonal-leakage-limit}.  Finally,
$G_\ell\asymp\sqrt\ell$ and $|b|\geq c\delta_\Lambda$ on the stated region,
so the first term of~\eqref{eq:minimum-lift-high-leakage} proves
~\eqref{eq:minimum-lift-gap-failure}.
\end{proof}

The preceding growth is specific to the minimum lift, but non-negligible
leakage is forced for every bounded scalar lift.  The reason is pointwise:
near the large nodal region of a concentrated highest-weight state,
multiplication cannot turn $h_\ell$ into a nonzero constant component.

\begin{proposition}[Nodal-set obstruction to a two-state scalar lift]
\label{prop:nodal-set-lift-obstruction}
Let $(X,\mu)$ have measure $S$, let $e_0=S^{-1/2}$, let $h$ be real with
$\|h\|_2=1$ and $h\perp e_0$, and let $\Pi$ project onto
$\operatorname{span}\{e_0,h\}$.  For every real $v\in L^\infty(X)$, put
\begin{equation}
 b=\langle e_0,M_vh\rangle,
 \qquad d=\langle h,M_vh\rangle,
 \qquad B=\|v\|_\infty.
\end{equation}
Then, for every $t>0$,
\begin{equation}
 \|(I-\Pi)M_vh\|_2^2
 \geq \mu\{|h|\leq t\}
 \left(\frac{|b|}{\sqrt S}-(B+|d|)t\right)_+^2.
 \label{eq:nodal-set-general-bound}
\end{equation}
If $b\neq0$ and
\begin{equation}
 t_b=\frac{|b|}{2\sqrt S(B+|d|)},
 \label{eq:nodal-threshold}
\end{equation}
then
\begin{equation}
 \|(I-\Pi)M_vh\|_2^2
 \geq\frac{b^2}{4S}\mu\{|h|\leq t_b\}.
 \label{eq:nodal-set-half-bound}
\end{equation}
For $h=h_\ell$, if
\begin{equation}
 \log\frac{B_\ell+|d_\ell|}{|b_\ell|}=o(\ell),
 \label{eq:subexponential-amplitude-ratio}
\end{equation}
then $\mu\{|h_\ell|\leq t_{b_\ell}\}\to4\pi$ and hence
\begin{equation}
 \|(I-\Pi_\ell)M_{v_\ell}h_\ell\|_2
 \geq\left(\frac12-o(1)\right)|b_\ell|.
 \label{eq:nodal-highest-weight-bound}
\end{equation}
\end{proposition}

\begin{proof}
Since
\begin{equation}
 (I-\Pi)M_vh=vh-be_0-dh,
\end{equation}
the reverse triangle inequality gives, on $\{|h|\leq t\}$,
\begin{equation}
 |vh-be_0-dh|
 \geq\frac{|b|}{\sqrt S}-(B+|d|)t.
\end{equation}
Integration proves~\eqref{eq:nodal-set-general-bound}, and the choice
~\eqref{eq:nodal-threshold} proves~\eqref{eq:nodal-set-half-bound}.

For the highest-weight state, fix $\eta>0$.  On
$\{|\theta-\pi/2|\geq\eta\}$,
\begin{equation}
 |h_\ell(\theta,\phi)|
 \leq C\ell^{1/4}(\cos\eta)^\ell.
\end{equation}
Condition~\eqref{eq:subexponential-amplitude-ratio} says
$t_{b_\ell}=e^{-o(\ell)}$ up to fixed factors.  Thus the displayed region is
contained in $\{|h_\ell|\leq t_{b_\ell}\}$ for all large $\ell$.  Its area
tends to $4\pi$ as $\eta\downarrow0$, proving the measure assertion and
~\eqref{eq:nodal-highest-weight-bound}.
\end{proof}

For the lifted rotating block, $|b|\asymp\delta_\Lambda$ during a fixed
fraction of the quarter-turn.  The minimum lift has
$B=O(\delta_\Lambda\ell^{1/4})$, so
~\eqref{eq:subexponential-amplitude-ratio} holds.  It also holds for any
uniformly bounded alternative lift when
$\ell=(\log\Lambda)^p$ with $p>1$, because
$\log(B/|b|)=O(\log\Lambda)=o(\ell)$.  Consequently no such scalar realization
makes $\operatorname{span}\{e_0,h_\ell\}$ asymptotically invariant during the
mixing stage.  A successful converter, if it exists, must deliberately enter
the complementary angular sector and later refocus it.

\begin{proposition}[Decomposable scalar converters cannot concentrate]
\label{prop:decomposable-converter-no-go}
Let $(X,\mu)$ have finite measure $S$, let $e_0=S^{-1/2}$, and let
$h\in L^2(X)$ have norm one.  Suppose that a lossless two-port scatterer is
decomposable over $X$:
\begin{equation}
 (\mathcal Sf)(x)=\mathsf S(x)f(x),
 \qquad \mathsf S(x)\in U(2)\quad\text{for a.e. }x.
 \label{eq:decomposable-scatterer}
\end{equation}
If $\mathbf e_+$ denotes incidence in the first port and the same symbol
denotes the desired transmitted port, then, for every constant phase $\gamma$,
\begin{equation}
 \left\|\mathcal S(e_0\mathbf e_+)
       -e^{i\gamma}h\mathbf e_+\right\|_{L^2(X;\C^2)}^2
 \geq 2-\frac{2}{\sqrt S}\|h\|_{L^1(X)}.
 \label{eq:decomposable-conversion-lower-bound}
\end{equation}
In particular, deleting $A_3/r^2$ from a real scalar angular Schr\"odinger
converter makes~\eqref{eq:decomposable-scatterer} hold with
$X=\mathbb S^2$: every angular point is then an independent one-dimensional
lossless scattering problem.
\end{proposition}

\begin{proof}
Both vectors in~\eqref{eq:decomposable-conversion-lower-bound} have norm one.
If $s_{++}(x)$ is the desired output component of
$\mathsf S(x)\mathbf e_+$, unitarity gives $|s_{++}(x)|\leq1$.  Hence
\begin{align}
 \left|\left\langle e^{i\gamma}h\mathbf e_+,
                   \mathcal S(e_0\mathbf e_+)\right\rangle\right|
 &=\left|S^{-1/2}\int_X e^{-i\gamma}\overline{h(x)}
                   s_{++}(x)\dd\mu(x)\right|\notag\\
 &\leq S^{-1/2}\|h\|_1.
\end{align}
Expanding the squared distance proves the claim.  If $A_3/r^2$ is absent,
the real scalar radial equation contains no angular derivatives.  Its
direction-resolved scattering matrix is therefore the direct integral of the
unitary scalar matrices $\mathsf S(x)$.
\end{proof}

For the real highest-weight harmonic~\eqref{eq:real-highest-weight},
\begin{equation}
 \|h_\ell\|_1
 =4C_\ell\sqrt\pi\,
   \frac{\Gamma((\ell+2)/2)}{\Gamma((\ell+3)/2)}
 =4\sqrt2\,\pi^{-1/4}\ell^{-1/4}(1+o(1)).
 \label{eq:highest-weight-L1}
\end{equation}
Indeed, the azimuthal integral of $|\cos(\ell\phi)|$ is $4$, and the polar
integral is
$\int_0^\pi\sin^{\ell+1}\theta\dd\theta
=\sqrt\pi\,\Gamma((\ell+2)/2)/\Gamma((\ell+3)/2)$.
Consequently Proposition~\ref{prop:decomposable-converter-no-go} gives
\begin{equation}
 \inf_{\mathcal S,\gamma}
 \left\|\mathcal S(e_0\mathbf e_+)
       -e^{i\gamma}h_\ell\mathbf e_+\right\|^2
 \geq
 2-4\sqrt2\,\pi^{-3/4}\ell^{-1/4}(1+o(1)).
 \label{eq:highest-decomposable-error}
\end{equation}
Thus the best decomposable modal error tends to $\sqrt2$, not to zero.  If
$\mathcal S_\ell^{\rm full}$ is a proposed full converter and
$\mathcal S_\ell^{\rm dec}$ is the scatterer obtained by deleting angular
transport, the triangle inequality yields the useful perturbative test
\begin{equation}
 \left\|(\mathcal S_\ell^{\rm full}
              -\mathcal S_\ell^{\rm dec})e_0\mathbf e_+\right\|
 \geq \sqrt{2-\frac{\|h_\ell\|_1}{\sqrt\pi}}
       -\operatorname{err}_{\rm conv}(\ell).
 \label{eq:angular-transport-order-one}
\end{equation}
Every successful converter must therefore use the centrifugal term at order
one; a comparison which makes it $o(1)$ disproves the candidate construction.

The same necessity has a local conservation-law formulation.  It also
identifies the geometric transport which a positive estimate must control.

\begin{proposition}[Exact angular-current requirement]
\label{prop:angular-current-requirement}
Let $v$ be real and let $u$ solve on $[R_1,R_2]\times\mathbb S^2$
\begin{equation}
 -\partial_r^2u-\frac1{r^2}\Delta_{\mathbb S^2}u+vu=k^2u.
 \label{eq:full-scalar-converter-equation}
\end{equation}
Define the radial and angular currents
\begin{equation}
 j_r=\Im(\overline u\,\partial_ru),
 \qquad
 J_\omega=\frac1{r^2}\Im(\overline u\,\nabla_\omega u).
 \label{eq:radial-angular-currents}
\end{equation}
Then
\begin{equation}
 \partial_rj_r+\operatorname{div}_\omega J_\omega=0.
 \label{eq:angular-current-continuity}
\end{equation}
If the boundary data are exact forward modes $e_0^+$ at $R_1$ and
$h_\ell^+$ at $R_2$, then
\begin{equation}
 \int_{R_1}^{R_2}\frac1{r^2}
   \int_{\mathbb S^2}|u|\,|\nabla_\omega u|\dd\omega\dd r
 \geq k\left(\frac\pi2-1-C\ell^{-1/2}\right).
 \label{eq:angular-current-lower-bound}
\end{equation}
The same conclusion holds with an $O_K(\varepsilon)$ loss when the two
directional boundary modes have phase-space error at most $\varepsilon$.
\end{proposition}

\begin{proof}
Taking the imaginary part of
$\overline u$ times~\eqref{eq:full-scalar-converter-equation} proves
~\eqref{eq:angular-current-continuity}; the real scalar potential disappears
identically.  Integrate in $r$ and test against the one-Lipschitz function
\begin{equation}
 \varphi(\theta,\phi)=|\theta-\pi/2|.
\end{equation}
Since a forward mode $e^{ikr}f(\omega)$ has radial current $k|f|^2$, integration
by parts on the sphere gives
\begin{align}
 k\left|\int_{\mathbb S^2}\varphi
       \left(h_\ell^2-\frac1{4\pi}\right)\dd\omega\right|
 &\leq \int_{R_1}^{R_2}\int_{\mathbb S^2}|J_\omega|
       \dd\omega\dd r\notag\\
 &\leq \int_{R_1}^{R_2}\frac1{r^2}
       \int_{\mathbb S^2}|u|\,|\nabla_\omega u|
       \dd\omega\dd r.
 \label{eq:tested-angular-current}
\end{align}
The uniform spherical measure satisfies
\begin{equation}
 \frac1{4\pi}\int_{\mathbb S^2}|\theta-\pi/2|\dd\omega
 =\int_0^{\pi/2}(\pi/2-\theta)\sin\theta\dd\theta
 =\frac\pi2-1.
\end{equation}
The polar marginal of $h_\ell^2\dd\omega$ is
$J_\ell^{-1}\sin^{2\ell+1}\theta\dd\theta$, whose first absolute moment
about the equator is $O(\ell^{-1/2})$.  Substitution in
~\eqref{eq:tested-angular-current} proves~\eqref{eq:angular-current-lower-bound}.
The phase-space stability statement follows by expanding the boundary currents
and applying Cauchy--Schwarz.
\end{proof}

For example, if along the converter
$\|u(r)\|_2\leq C_0$ and
$\|\nabla_\omega u(r)\|_2\leq C_1L_{\rm eff}$, then
\begin{equation}
 \int_{R_1}^{R_2}\frac1{r^2}
   \int|u|\,|\nabla_\omega u|
 \leq C_0C_1L_{\rm eff}
       \left(\frac1{R_1}-\frac1{R_2}\right).
 \label{eq:effective-bandwidth-current}
\end{equation}
Therefore a converter supported far beyond its effective angular bandwidth is
impossible.  A construction must instead operate near the centrifugal scale
$R_1=O(L_{\rm eff})$ or create a genuine cascade to angular frequencies much
larger than the target degree.  This is the sharpened scalar fork: the
compressed rotator cannot simply be moved to a region where $A_3/r^2$ is
negligible and then justified perturbatively.

The scalar cost alone does not bound $L_{\rm eff}$.  A long weak slab can
create an arbitrarily oscillatory angular phase at vanishing cost.  This is a
second necessary falsification test for any proposed positive converter bound.

\begin{proposition}[Cheap scalar phase masks]
\label{prop:cheap-scalar-phase-mask}
Let $K\Subset(0,\infty)$, let $\tau>0$, and let
$w\in L^\infty(\mathbb S^2;\mathbb R)$ with $\|w\|_\infty\leq1$.  Delete
$A_3/r^2$ and put the scalar potential
\begin{equation}
 v_\epsilon(r,\omega)=\epsilon w(\omega)
 \mathbf1_{[R,R+\tau/\epsilon]}(r).
 \label{eq:weak-phase-mask-slab}
\end{equation}
Its full constant-column cost is
\begin{equation}
 \int\|M_{v_\epsilon(r)}e_0\|_2^2\dd r
 =\frac{\tau\epsilon}{4\pi}\|w\|_2^2\longrightarrow0.
 \label{eq:weak-phase-mask-cost}
\end{equation}
Nevertheless, after removal of the free propagation phase, its two-port
scattering matrix has reflection $O_K(\epsilon)$ and transmission multiplier
\begin{equation}
 t_\epsilon(k,\omega)
 =\exp\!\left(-\frac{i\tau}{2k}w(\omega)\right)+O_K(\epsilon),
 \label{eq:weak-phase-mask-limit}
\end{equation}
uniformly for $k\in K$ and almost every $\omega$.  The same statement holds in
the reverse direction.

In particular, for the smooth bounded masks
\begin{equation}
 w_L(\theta,\phi)=\cos(L\theta)=T_L(\cos\theta),
 \label{eq:chebyshev-phase-mask}
\end{equation}
one has $\|w_L\|_2\asymp1$ but
$\|\nabla_\omega w_L\|_2\asymp L$.  The limiting transmitted wave from $e_0$
therefore has angular gradient comparable to $L$, while the cost in
~\eqref{eq:weak-phase-mask-cost} can tend to zero independently of $L$.
\end{proposition}

\begin{proof}
For each fixed $\omega$, the slab has internal wave number
\begin{equation}
 p_\epsilon(k,\omega)=\sqrt{k^2-\epsilon w(\omega)}
 =k-\frac{\epsilon}{2k}w(\omega)+O_K(\epsilon^2).
\end{equation}
The reflection coefficient at either interface is $O_K(\epsilon)$, uniformly
in $\omega$.  The exact rectangular-slab formula therefore gives total
reflection $O_K(\epsilon)$ and, relative to free propagation over the same
length,
\begin{equation}
 t_\epsilon(k,\omega)
 =\exp\!\left(i(p_\epsilon-k)\frac{\tau}{\epsilon}\right)
  \bigl(1+O_K(\epsilon)\bigr),
\end{equation}
which proves~\eqref{eq:weak-phase-mask-limit}.  Reciprocity gives the reverse
statement.  The cost is immediate.  Finally,
\begin{align}
 \|w_L\|_2^2
 &=2\pi\left(1-\frac1{4L^2-1}\right),\\
 \|\nabla_\omega w_L\|_2^2
 &=2\pi L^2\left(1+\frac1{4L^2-1}\right),
\end{align}
so differentiating the phase in~\eqref{eq:weak-phase-mask-limit} proves the
last assertion.
\end{proof}

Proposition~\ref{prop:cheap-scalar-phase-mask} shows why
~\eqref{eq:angular-current-lower-bound} is a necessary condition but not yet a
converter lower bound.  Vanishing column cost can seed arbitrarily large
angular frequency, after which the centrifugal propagator may transport
angular density.  The remaining issue is not generation of a cascade but its
broadband refocusing into one prescribed highest-weight mode.

There is, however, a radial incompatibility hidden by the idealized product
notation.  A weak slab is a clean phase mask only where the centrifugal
evolution during that slab is negligible.  That is incompatible with making
the same slab cheap at the angular-transport scale.

\begin{proposition}[Phase-mask/transport incompatibility]
\label{prop:phase-mask-transport-incompatibility}
Let $w_L$ be real with
$0<c_0\leq\|w_L\|_2\leq C_0$ and $\|w_L\|_\infty\leq C_0$.  Put
\begin{equation}
 v_L(r,\omega)=\epsilon_Lw_L(\omega)
 \mathbf1_{[R_L,R_L+T_L]}(r),
 \qquad \tau_L=\epsilon_LT_L.
\end{equation}
Assume $0<c_1\leq\tau_L\leq C_1$.  Its constant-column cost and the maximal
centrifugal phase scale on degrees at most $L$ are
\begin{align}
 \mathcal C_L
 &:=\int\|M_{v_L(r)}e_0\|_2^2\dd r
 =\frac{\tau_L^2}{4\pi T_L}\|w_L\|_2^2,
 \label{eq:mask-cost-duration}\\
 \mathcal D_L
 &:=L(L+1)\int_{R_L}^{R_L+T_L}\frac{\dd r}{r^2}
 =\frac{L(L+1)T_L}{R_L(R_L+T_L)}.
 \label{eq:mask-centrifugal-action}
\end{align}
If $R_L\leq C L$ and $\mathcal C_L\to0$, then
$T_L\to\infty$ and $\mathcal D_L\to\infty$.  Conversely, under the same
radius condition, $\mathcal D_L\to0$ forces $T_L\to0$ and
$\mathcal C_L\to\infty$.  Hence a nontrivial phase mask cannot be both cheap
and centrifugally clean at the scale where degree-$L$ angular transport is
available.
\end{proposition}

\begin{proof}
Equations~\eqref{eq:mask-cost-duration} and
~\eqref{eq:mask-centrifugal-action} are direct calculations.  Vanishing cost
with $\tau_L$ and $\|w_L\|_2$ bounded below implies $T_L\to\infty$.  If
$T_L\leq R_L$, then
\begin{equation}
 \mathcal D_L\geq
 \frac{L(L+1)}{2R_L^2}T_L\geq cT_L.
\end{equation}
If $T_L\geq R_L$, then
\begin{equation}
 \mathcal D_L\geq\frac{L(L+1)}{2R_L}\geq cL.
\end{equation}
This proves the first assertion.  If $\mathcal D_L\to0$, the second case is
impossible and the first estimate forces $T_L\to0$; then
~\eqref{eq:mask-cost-duration} diverges.
\end{proof}

Thus the ideal masks in~\eqref{eq:weak-phase-mask-limit} are available cheaply
only in a far-field regime relative to the frequencies which remain active
during the slab.  But Proposition~\ref{prop:angular-current-requirement} and
~\eqref{eq:effective-bandwidth-current} show that such a regime cannot change
the angular density unless the slab simultaneously creates frequencies much
larger than its radius.  The putative converter must therefore use a coupled
scalar--centrifugal cascade during the weak slabs themselves; it is not a
literal succession of independent optical elements.

The linearized Fourier budget has an exact nonlinear counterpart for the
standard pure-detuning two-level ensemble.  Recording it separates a genuine
nonlinear obstruction from the additional energy scaling in the radial
Schr\"odinger problem.

\begin{proposition}[Nonlinear Plancherel budget for a pure-detuning ensemble]
\label{prop:SU2-nonlinear-Plancherel}
Let $q\in C_c^\infty(\mathbb R;\mathbb C)$ and, for $\xi\in\mathbb R$, let
$G(\xi,x)$ solve
\begin{equation}
 \partial_xG(\xi,x)=
 \begin{pmatrix}
 0&q(x)e^{2i\xi x}\\
 -\overline{q(x)}e^{-2i\xi x}&0
 \end{pmatrix}G(\xi,x),
 \qquad G(\xi,-\infty)=\Id.
 \label{eq:SU2-nonlinear-Fourier-system}
\end{equation}
Write
\begin{equation}
 G(\xi,+\infty)=
 \begin{pmatrix}a(\xi)&b(\xi)\\
 -\overline{b(\xi)}&\overline{a(\xi)}
 \end{pmatrix}.
\end{equation}
Then $|a|^2+|b|^2=1$ and
\begin{equation}
 -\int_{\mathbb R}\log\bigl(1-|b(\xi)|^2\bigr)\dd\xi
 \leq \pi\int_{\mathbb R}|q(x)|^2\dd x.
 \label{eq:SU2-nonlinear-Plancherel}
\end{equation}
More precisely, if $z_j$ are the zeros, counted with multiplicity, of the
analytic continuation of $a$ to $\mathbb C_+$, then
\begin{equation}
 -\int_{\mathbb R}\log\bigl(1-|b(\xi)|^2\bigr)\dd\xi
 =\pi\|q\|_2^2-4\pi\sum_j\Im z_j.
 \label{eq:SU2-trace-with-zeros}
\end{equation}
Consequently, if $|b(\xi)|\geq\eta$ on an interval $J$, then
\begin{equation}
 \|q\|_2^2\geq
 \frac{|J|}{\pi}\log\frac1{1-\eta^2}.
 \label{eq:SU2-fixed-band-cost}
\end{equation}
\end{proposition}

\begin{proof}
For real $\xi$, the generator in
~\eqref{eq:SU2-nonlinear-Fourier-system} is skew-Hermitian and traceless, so
$G(\xi,+\infty)\in\mathrm{SU}(2)$.  The Volterra series gives an analytic
continuation of $a$ to $\mathbb C_+$ and the high-imaginary-energy expansion
\begin{equation}
 \log a(z)=-\frac{i}{2z}\|q\|_2^2+O(z^{-2}).
 \label{eq:SU2-high-energy-expansion}
\end{equation}
Factor $a=B a_{\rm out}$ into its Blaschke and outer parts.  With
\begin{equation}
 B(z)=\prod_j\frac{z-z_j}{z-\overline{z_j}},
 \qquad
 \log a_{\rm out}(z)=\frac1{\pi i}
 \int_{\mathbb R}\frac{\log|a(\xi)|}{\xi-z}\dd\xi,
\end{equation}
their expansions at infinity are
\begin{equation}
 \log B(z)=-\frac{2i}{z}\sum_j\Im z_j+O(z^{-2}),
 \qquad
 \log a_{\rm out}(z)=\frac{i}{\pi z}
 \int_{\mathbb R}\log|a(\xi)|\dd\xi+O(z^{-2}).
\end{equation}
Comparison with~\eqref{eq:SU2-high-energy-expansion}, followed by
$2\log|a|=\log(1-|b|^2)$, proves
~\eqref{eq:SU2-trace-with-zeros}.  Dropping the nonnegative zero sum proves
~\eqref{eq:SU2-nonlinear-Plancherel}; restricting its left side to $J$ proves
~\eqref{eq:SU2-fixed-band-cost}.  This is the continuous nonlinear-Fourier
analogue of the nonlinear Plancherel identities discussed in
~\cite{TaoThiele2012}.
\end{proof}

This proposition does not yet settle the physical frozen converter.  After
subtracting a scalar phase, that two-state forward model has the form
\begin{equation}
 i\partial_x\psi_s
 =s\begin{pmatrix}0&\beta_\ell u(x)\\
                  \beta_\ell u(x)&c_\ell
   \end{pmatrix}\psi_s,
 \qquad s=(2k)^{-1},\qquad c_\ell=\frac{\ell(\ell+1)}{R_\ell^2}.
 \label{eq:physical-scaled-two-level-model}
\end{equation}
In the interaction picture the spectral parameter is
$\xi=sc_\ell/2$, but the nonlinear-Fourier potential is
$q_s=-is\beta_\ell u$ and therefore varies with the same parameter.  The data
occupy a diagonal curve in $(\xi,q)$, whereas
Proposition~\ref{prop:SU2-nonlinear-Plancherel} integrates $\xi$ with one fixed
$q$.  Thus~\eqref{eq:SU2-fixed-band-cost} is the correct exact model theorem,
but applying it to~\eqref{eq:physical-scaled-two-level-model} requires an
additional scaled-potential or second-order Schr\"odinger argument.

That missing step is not formal: the corresponding scaled-ensemble inequality
is false for arbitrary matrix Hamiltonians.

\begin{proposition}[Scaled nonlinear Plancherel fails for matrix controls]
\label{prop:scaled-matrix-ensemble-no-go}
Fix $0<s_0<s_1$.  There are real symmetric $2\times2$ Hamiltonians
$H_\Lambda(x)$ supported on $[0,\Lambda]$ such that
\begin{equation}
 \int_0^\Lambda\|H_\Lambda(x)e_1\|^2\dd x
 =\frac12\Lambda^{-1/3}\longrightarrow0,
 \label{eq:scaled-matrix-control-cost}
\end{equation}
while the solutions of
\begin{equation}
 i\partial_x\psi_s=sH_\Lambda(x)\psi_s,
 \qquad \psi_s(0)=e_1,\qquad s\in[s_0,s_1],
 \label{eq:scaled-matrix-ensemble}
\end{equation}
satisfy
\begin{equation}
 \sup_{s\in[s_0,s_1]}
 \inf_{\gamma\in\mathbb R}
 \|\psi_s(\Lambda)-e^{i\gamma}e_2\|
 =O(\Lambda^{-1/3}).
 \label{eq:scaled-matrix-transfer}
\end{equation}
Consequently the fixed-band transition entropy diverges although the charged
cost vanishes:
\begin{equation}
 \int_{s_0}^{s_1}
 -\log\bigl(1-|\langle e_2,\psi_s(\Lambda)\rangle|^2\bigr)\dd s
 \geq c\log\Lambda-C.
 \label{eq:scaled-matrix-entropy-failure}
\end{equation}
\end{proposition}

\begin{proof}
Use the first-order version of the rotator in Appendix~\ref{app:rotator}:
\begin{equation}
 H_\Lambda(x)=\delta_\Lambda
 R(\omega_\Lambda x)\diag(0,1)R(\omega_\Lambda x)^*,
 \quad
 \delta_\Lambda=\Lambda^{-2/3},
 \quad
 \omega_\Lambda=\frac\pi{2\Lambda}.
\end{equation}
The cost is
$\delta_\Lambda^2\int_0^\Lambda\sin^2(\omega_\Lambda x)\dd x$, which is
~\eqref{eq:scaled-matrix-control-cost}.  Writing
$\psi_s=R(\omega_\Lambda x)y_s$ gives the constant equation
\begin{equation}
 i\partial_xy_s=
 \bigl[s\delta_\Lambda\diag(0,1)-i\omega_\Lambda\mathsf K\bigr]y_s.
\end{equation}
Its lower spectral projection differs from $e_1e_1^*$ by
$O(\omega_\Lambda/\delta_\Lambda)=O(\Lambda^{-1/3})$, uniformly in $s$.
Since $R(\pi/2)e_1=e_2$, this proves~\eqref{eq:scaled-matrix-transfer}.
The complementary transition probability is $O(\Lambda^{-2/3})$, and taking
its negative logarithm proves~\eqref{eq:scaled-matrix-entropy-failure}.
\end{proof}

The principal nonlinear escape from a Fourier budget would be robust adiabatic
passage.  Such passage is effective in finite quantum ladders and qubit
ensembles~\cite{LeghtasSarletteRouchon2011,RobinEtAl2021}.  Scalar angular
Schr\"odinger operators, however, possess an order structure absent from an
arbitrary matrix Hamiltonian.

\begin{proposition}[Ground-state positivity excludes a scalar adiabatic swap]
\label{prop:scalar-ground-adiabatic-no-go}
Let
\begin{equation}
 \mathcal L(r)=r^{-2}A_3+M_{v(r)},
 \qquad v(r,\cdot)\in L^\infty(\mathbb S^2;\mathbb R),
 \label{eq:instantaneous-angular-Hamiltonian}
\end{equation}
with $v$ continuous in $r$ in the $L^\infty$ norm.  Its lowest eigenvalue is
simple and has a strictly positive normalized eigenfunction $g_r$.  Hence any
continuous rank-one isolated eigenprojection of $\mathcal L(r)$ which equals
$e_0e_0^*$ at a free left endpoint remains the ground-state projection.  If the
right endpoint is also free, it equals $e_0e_0^*$ there and cannot converge to
$h_\ell h_\ell^*$.

Quantitatively, if at some radius $R$ one has
$\|v(R,\cdot)\|_\infty\leq\epsilon$ and $g_R$ is the normalized ground state,
then
\begin{equation}
 |\langle h_\ell,g_R\rangle|^2
 \leq\frac{2\epsilon R^2}{\ell(\ell+1)}.
 \label{eq:ground-highest-weight-overlap}
\end{equation}
In particular, for $R\leq C\ell$ and $\epsilon=o(1)$, an adiabatically followed
ground mode has $o(1)$ overlap with $h_\ell$.
\end{proposition}

\begin{proof}
The heat semigroup of $r^{-2}A_3+M_v$ is positivity improving on the connected
sphere.  Equivalently, the strong maximum principle and the variational
characterization show that the ground eigenvalue is simple and its eigenfunction
may be chosen strictly positive.  A continuous eigenvalue branch which starts
at the bottom cannot leave the ground branch without meeting a second
eigenvalue; simplicity of the ground state forbids that meeting.  At a free
endpoint the ground eigenspace is $\mathbb C e_0$.

For the quantitative assertion, let $E_R$ be the ground energy.  Testing on
$e_0$ gives $E_R\leq\epsilon$, while
\begin{equation}
 E_R\geq R^{-2}\langle g_R,A_3g_R\rangle-\epsilon.
\end{equation}
Thus $\langle g_R,A_3g_R\rangle\leq2\epsilon R^2$.  Since
$A_3h_\ell=\ell(\ell+1)h_\ell$ and $A_3\geq0$, the spectral theorem gives
~\eqref{eq:ground-highest-weight-overlap}.
\end{proof}

Proposition~\ref{prop:scalar-ground-adiabatic-no-go} rules out the direct scalar
analogue of the slowly rotating lowest-eigenvector construction in
Appendix~\ref{app:rotator}, including a vanishing endpoint-matching potential
at the centrifugal scale.  It does not exclude a chirped rotating-frame branch,
an excited-state path with deliberate transitions, or a genuinely nonadiabatic
multimode cascade.  Those mechanisms must leave the positive ground bundle;
the problem is now to charge that departure on a fixed energy band.

The most direct rotating-frame escape is a monotone chirped highest-weight
pulse.  The continuum of energies makes its usual adiabatic scaling expensive.

\begin{proposition}[A monotone chirped scalar passage is not cheap on a band]
\label{prop:chirped-ensemble-cost}
Let $\varphi_\ell$ be the normalized complex highest-weight harmonic,
$L_z\varphi_\ell=\ell\varphi_\ell$, and put
\begin{equation}
 M_\ell^{\rm c}=\|\varphi_\ell\|_\infty\asymp\ell^{1/4},
 \qquad
 w_{\ell,\Phi}(\omega)=
 \frac{\Re(e^{-i\Phi}\varphi_\ell(\omega))}{M_\ell^{\rm c}}.
\end{equation}
Consider the forward angular equation
$i\partial_rU_{\ell,s}=s(R_\ell^{-2}A_3+M_{v_\ell(r)})U_{\ell,s}$ at
$R_\ell=C\ell$ with
\begin{equation}
 v_\ell(r,\omega)=\epsilon_\ell
 w_{\ell,\Phi_\ell(r)}(\omega),
 \qquad
 g_\ell=\frac{\epsilon_\ell}
 {2\sqrt{4\pi}M_\ell^{\rm c}}.
 \label{eq:chirped-highest-weight-control}
\end{equation}
Let $\nu_\ell=\Phi_\ell'$ be monotone and sweep the complete resonance interval
\begin{equation}
 I_\ell=c_\ell[s_0,s_1],
 \qquad c_\ell=\frac{\ell(\ell+1)}{R_\ell^2}\asymp1.
\end{equation}
If the standard uniform adiabatic parameter
\begin{equation}
 \eta_\ell:=
 \sup_{\{r:\nu_\ell(r)\in I_\ell\}}
 \frac{|\nu_\ell'(r)|}{g_\ell^2}
 \longrightarrow0,
 \label{eq:chirped-adiabatic-parameter}
\end{equation}
then the charged cost on the sweep obeys
\begin{equation}
 \int\|M_{v_\ell(r)}e_0\|_2^2\dd r
 \geq\frac{2|I_\ell|}{\eta_\ell}\longrightarrow\infty.
 \label{eq:chirped-ensemble-cost}
\end{equation}
In particular, a uniformly adiabatic monotone chirp cannot produce a
vanishing-cost converter on a fixed nondegenerate energy band.
\end{proposition}

\begin{proof}
Rotation by $e^{-i\Phi_\ell L_z/\ell}$ fixes the angular mask and adds
$-\nu_\ell L_z/\ell$ to the transformed generator.  On
$\operatorname{span}\{e_0,\varphi_\ell\}$ the detuning is therefore
$sc_\ell-\nu_\ell(r)$ and the off-diagonal coupling has size $s g_\ell$.
Every $s\in[s_0,s_1]$ crosses resonance, so uniform adiabaticity requires
~\eqref{eq:chirped-adiabatic-parameter} throughout the complete interval
$I_\ell$.  Monotonicity gives a sweep duration
\begin{equation}
 T_\ell\geq\frac{|I_\ell|}{\eta_\ell g_\ell^2}.
\end{equation}
Since
$\|w_{\ell,\Phi}\|_2^2=(2(M_\ell^{\rm c})^2)^{-1}$, the cost density is
\begin{equation}
 \|M_{v_\ell}e_0\|_2^2
 =\frac{\epsilon_\ell^2}{8\pi(M_\ell^{\rm c})^2}
 =2g_\ell^2.
\end{equation}
Multiplication by $T_\ell$ proves~\eqref{eq:chirped-ensemble-cost}.
\end{proof}

This proposition is deliberately restricted to a single monotone chirp and
the usual uniform adiabatic regime.  It does not exclude nonmonotone composite
pulses, rotating frames involving several angular momenta, or nonadiabatic
multimode transfer.  It proves that the most immediate ensemble-control
construction cannot be the missing cheap converter.

A seemingly cheaper alternative is to use a rotating low-degree scalar field
and climb the highest-weight ladder one unit at a time.  The centrifugal
spacing between adjacent rungs is only $O(\ell^{-2})$ at $R\asymp\ell$, and a
complete sweep through $\ell$ rungs has width $O(\ell^{-1})$.  The fixed energy
band nevertheless prevents an adiabatic ground ladder from ending on one
common rung.

\begin{proposition}[Rotating-frame ground ladders fan out on an energy band]
\label{prop:rotating-ground-ladder-fanout}
Let $L_z=-i\partial_\phi$, let
$\varphi_n=Y_n^n$, and consider any rotating-frame scalar angular generator
whose control vanishes at its right endpoint and whose endpoint operator is
\begin{equation}
 \mathcal H_{s,\Omega}^{\rm end}
 =sR^{-2}A_3-\Omega L_z,
 \qquad s\in[s_0,s_1].
 \label{eq:rotating-frame-free-endpoint}
\end{equation}
For $\Omega\geq0$, the vector $\varphi_\ell$ can be a ground state of
~\eqref{eq:rotating-frame-free-endpoint} only if
\begin{equation}
 \frac{2s\ell}{R^2}\leq\Omega
 \leq\frac{2s(\ell+1)}{R^2}.
 \label{eq:rotating-ladder-target-window}
\end{equation}
Consequently no single $\Omega=\Omega_\ell$ makes $\varphi_\ell$ a ground
state for every $s\in[s_0,s_1]$ once
\begin{equation}
 \ell>\frac{s_0}{s_1-s_0}.
 \label{eq:rotating-ladder-band-threshold}
\end{equation}
Thus a uniformly adiabatic construction which starts from the ground state
$e_0$, follows an isolated rotating-frame ground branch, turns off the scalar
control, and aims at one common highest-weight mode cannot work on a fixed
nondegenerate energy band.
\end{proposition}

\begin{proof}
At the free endpoint,
\begin{equation}
 \mathcal H_{s,\Omega}^{\rm end}Y_n^m
 =\left(\frac{s n(n+1)}{R^2}-\Omega m\right)Y_n^m.
\end{equation}
For $\Omega\geq0$ and fixed $n$, the lowest value occurs at $m=n$.  Comparing
the energies of $\varphi_\ell$ with $\varphi_{\ell-1}$ and
$\varphi_{\ell+1}$ gives exactly
~\eqref{eq:rotating-ladder-target-window}.  A common endpoint value would
therefore have to satisfy
\begin{equation}
 \frac{2s_1\ell}{R^2}\leq\Omega_\ell
 \leq\frac{2s_0(\ell+1)}{R^2}.
\end{equation}
The interval is empty precisely when
$s_1\ell>s_0(\ell+1)$, which is
~\eqref{eq:rotating-ladder-band-threshold}.  An isolated ground branch remains
the ground branch under uniform adiabatic following, so it cannot terminate at
$\varphi_\ell$ for the complete band.
\end{proof}

For $\Omega<0$ the same argument selects $Y_n^{-n}$.  A real highest-weight
packet, which combines the two signs of $m$, is even less compatible with a
nonzero endpoint rotation rate.  Proposition
~\ref{prop:rotating-ground-ladder-fanout} does not exclude an excited-branch
pump or a deliberately nonadiabatic ladder, but it closes the most economical
adiabatic supercascade suggested by the shrinking adjacent-level spacings.

The natural next attempt is to keep the $e_0$--$h_\ell$ centrifugal gap
cancelled and replace the chirp by a finite composite sequence.  The scalar
cost and the exact angular current are incompatible with every such
bounded-depth construction.  More precisely, gap cancellation can be cheap
only for a very short time, while changing the angular density in that time
requires frequencies far above the target frequency.

\begin{proposition}[A gap-cancelled converter must supercascade]
\label{prop:gap-cancelled-supercascade}
Let $S=4\pi$, $R_\ell\asymp\ell$, and
$c_\ell=\ell(\ell+1)/R_\ell^2$.  Fix
$s\in[s_0,s_1]\Subset(0,\infty)$, and let the normalized solution of
\begin{equation}
 i\partial_x\psi_{\ell,s}
 =s\bigl(R_\ell^{-2}A_3+M_{v_\ell(x)}\bigr)\psi_{\ell,s},
 \qquad \psi_{\ell,s}(0)=e_0,
 \label{eq:gap-cancelled-evolution}
\end{equation}
be defined on $[0,T_\ell]$.  Suppose that, for almost every $x$,
\begin{equation}
 \left|c_\ell
 +\langle h_\ell,M_{v_\ell(x)}h_\ell\rangle
 -\langle e_0,M_{v_\ell(x)}e_0\rangle\right|
 \leq\frac{c_\ell}{2}.
 \label{eq:gap-cancellation-hypothesis}
\end{equation}
There are constants $\varepsilon_*,C,c>0$, independent of $\ell$, such that
if
\begin{equation}
 \inf_{\gamma\in\mathbb R}
 \|\psi_{\ell,s}(T_\ell)-e^{i\gamma}h_\ell\|_2
 \leq\varepsilon_*,
 \label{eq:gap-cancelled-conversion}
\end{equation}
and
\begin{equation}
 \mathcal E_\ell
 :=\int_0^{T_\ell}\|M_{v_\ell(x)}e_0\|_2^2\dd x,
\end{equation}
then
\begin{align}
 \mathcal E_\ell&\geq c\,T_\ell\ell^{-1/2},
 \label{eq:gap-cancelled-cost-time}\\
 \int_0^{T_\ell}\|A_3^{1/2}\psi_{\ell,s}(x)\|_2\dd x
 &\geq cR_\ell^2,
 \label{eq:gap-cancelled-current-time}\\
 \sup_{0\leq x\leq T_\ell}
 \|A_3^{1/2}\psi_{\ell,s}(x)\|_2
 &\geq \frac{c\ell^{3/2}}{\mathcal E_\ell}.
 \label{eq:gap-cancelled-supercascade}
\end{align}
In particular, vanishing charged cost forces angular frequency
$\gg\ell^{3/2}$ somewhere during the converter.
\end{proposition}

\begin{proof}
Put $a_\ell(x)=\langle e_0,M_{v_\ell(x)}e_0\rangle$ and
$d_\ell(x)=\langle h_\ell,M_{v_\ell(x)}h_\ell\rangle$.  Since
\begin{equation}
 d_\ell-a_\ell
 =\int_{\mathbb S^2}v_\ell
 \left(h_\ell^2-\frac1S\right)\dd\omega,
\end{equation}
Cauchy--Schwarz,~\eqref{eq:gap-cancellation-hypothesis}, and
$G_\ell=\|h_\ell^2-S^{-1}\|_2^2\asymp\sqrt\ell$ give
\begin{equation}
 \|M_{v_\ell(x)}e_0\|_2^2
 =\frac1S\|v_\ell(x)\|_2^2
 \geq\frac{|d_\ell(x)-a_\ell(x)|^2}{SG_\ell}
 \geq c\ell^{-1/2}.
 \label{eq:gap-cancelled-pointwise-cost}
\end{equation}
Integration proves~\eqref{eq:gap-cancelled-cost-time}.

The density $\rho=|\psi_{\ell,s}|^2$ satisfies the exact continuity equation
\begin{equation}
 \partial_x\rho
 +\frac{2s}{R_\ell^2}
 \operatorname{div}_\omega
 \Im(\overline{\psi_{\ell,s}}\nabla_\omega\psi_{\ell,s})=0;
 \label{eq:frozen-angular-continuity}
\end{equation}
the real scalar potential again disappears.  Test against
$\varphi(\theta,\phi)=|\theta-\pi/2|$.  The uniform density has
$\varphi$-moment $\pi/2-1$, whereas the $h_\ell^2$ moment is
$O(\ell^{-1/2})$.  Also
\begin{equation}
 \bigl\||\psi|^2-h_\ell^2\bigr\|_1
 \leq2\inf_\gamma\|\psi-e^{i\gamma}h_\ell\|_2.
\end{equation}
Choosing $\varepsilon_*$ small and integrating
\eqref{eq:frozen-angular-continuity} therefore yields
\begin{equation}
 c\leq\frac{2s}{R_\ell^2}
 \int_0^{T_\ell}\int_{\mathbb S^2}
 |\psi_{\ell,s}|\,|\nabla_\omega\psi_{\ell,s}|\dd\omega\dd x
 \leq\frac{C}{R_\ell^2}
 \int_0^{T_\ell}\|A_3^{1/2}\psi_{\ell,s}\|_2\dd x,
\end{equation}
which proves~\eqref{eq:gap-cancelled-current-time}.  Finally,
\eqref{eq:gap-cancelled-cost-time} gives
$T_\ell\leq C\mathcal E_\ell\sqrt\ell$.  Combining this with
\eqref{eq:gap-cancelled-current-time} and $R_\ell\asymp\ell$ proves
\eqref{eq:gap-cancelled-supercascade}.
\end{proof}

This has a direct finite-depth consequence once the layers have the usual
bounded phase-gradient budget.  A quadratic-form calculation in
\eqref{eq:gap-cancelled-evolution} gives
\begin{equation}
 \|A_3^{1/2}\psi_{\ell,s}(x)\|_2
 \leq s\int_0^x\|\nabla_\omega v_\ell(t)\|_\infty\dd t.
 \label{eq:continuous-gradient-growth}
\end{equation}
The discrete counterpart is that free angular propagation preserves the
$H^1$ seminorm, while a phase mask satisfies
\begin{equation}
 \|A_3^{1/2}(e^{-isw}f)\|_2
 \leq\|A_3^{1/2}f\|_2+s\|\nabla_\omega w\|_\infty\|f\|_2.
 \label{eq:phase-mask-gradient-growth}
\end{equation}
Hence a gap-cancelled implementation with $m_\ell$ layers satisfying
\begin{equation}
 \int_{I_{j,\ell}}\|\nabla_\omega v_\ell(t)\|_\infty\dd t
 \leq C\ell
 \label{eq:layer-gradient-budget}
\end{equation}
for each layer---or its ideal mask/free analogue with
$\|\nabla w_{j,\ell}\|_\infty\leq C\ell$---has intermediate frequency at
most $C(1+m_\ell)\ell$.  Under the hypotheses of Proposition
~\ref{prop:gap-cancelled-supercascade}, it must satisfy
\begin{equation}
 (1+m_\ell)\mathcal E_\ell\geq c\sqrt\ell.
 \label{eq:composite-depth-cost}
\end{equation}
Thus fixed-depth designs fail, and if $\mathcal E_\ell\to0$ then necessarily
$m_\ell/\sqrt\ell\to\infty$.  The surviving counterexample mechanism must be
a genuinely deep super-target cascade, not an ordinary composite pulse.

The same conclusion explains why a norm-only argument cannot close the
problem.  The exact block dynamics has a useful normal form which isolates the
remaining high-sector rotation.

\begin{proposition}[Exact scalar star-frame reduction]
\label{prop:scalar-star-frame}
Let $Q=\Id-e_0e_0^*$ and write
\begin{equation}
 \begin{aligned}
 m_0(x)&=\langle e_0,M_{v(x)}e_0\rangle,
 &q(x)&=QM_{v(x)}e_0,\\
 K(x)&=Q(D+M_{v(x)})Q.
 \end{aligned}
\end{equation}
Let $W_s$ solve $i\partial_xW_s=sK(x)W_s$, $W_s(0)=\Id_{Q\mathcal K}$,
and decompose a solution as $\psi_s=z_se_0+\chi_s$.  Set
\begin{equation}
 \phi_0(x)=\int_0^xm_0(t)\dd t,
 \qquad \alpha_s=e^{is\phi_0}z_s,
 \qquad \eta_s=W_s^*\chi_s,
 \qquad r_s=e^{-is\phi_0}W_s^*q.
 \label{eq:scalar-star-frame-variables}
\end{equation}
Then
\begin{equation}
 i\partial_x
 \begin{pmatrix}\alpha_s\\ \eta_s\end{pmatrix}
 =s
 \begin{pmatrix}0&r_s^*\\ r_s&0\end{pmatrix}
 \begin{pmatrix}\alpha_s\\ \eta_s\end{pmatrix},
 \qquad
 \|r_s(x)\|_2=\|q(x)\|_2,
 \label{eq:scalar-star-frame-system}
\end{equation}
and the charged density is exactly
\begin{equation}
 \|M_{v(x)}e_0\|_2^2=|m_0(x)|^2+\|r_s(x)\|_2^2.
 \label{eq:scalar-star-frame-cost}
\end{equation}
\end{proposition}

\begin{proof}
Projection of the equation onto $\mathbb Ce_0\oplus Q\mathcal K$ gives
\begin{align}
 iz_s'&=s(m_0z_s+q^*\chi_s),\\
 i\chi_s'&=s(qz_s+K\chi_s).
\end{align}
Substitution of~\eqref{eq:scalar-star-frame-variables} proves
\eqref{eq:scalar-star-frame-system}.  Unitarity of $W_s$ gives the norm
identity, and orthogonality of $m_0e_0$ and $q$ gives
\eqref{eq:scalar-star-frame-cost}.
\end{proof}

The star frame also upgrades the scale obstruction.  Reaching an orthogonal
state requires an order-one pulse area in the charged vector, irrespective of
how violently its direction rotates in the complementary sector.

\begin{corollary}[Pulse-area and scalar-roughness toll]
\label{cor:scalar-pulse-area-roughness}
Under the hypotheses of Proposition~\ref{prop:gap-cancelled-supercascade}, put
\begin{equation}
 q_\ell(x)=QM_{v_\ell(x)}e_0,
 \qquad
 \mathcal G_\ell=
 \int_0^{T_\ell}\|\nabla_\omega v_\ell(x)\|_\infty\dd x
 \label{eq:scalar-gradient-action}
\end{equation}
whenever the second quantity is finite.  Then
\begin{align}
 \int_0^{T_\ell}\|q_\ell(x)\|_2\dd x&\geq c,
 \label{eq:charged-pulse-area}\\
 T_\ell\mathcal E_\ell&\geq c,
 \label{eq:time-cost-uncertainty}\\
 \mathcal E_\ell&\geq c\ell^{-1/4}.
 \label{eq:gap-cancelled-quarter-cost}
\end{align}
If in addition $v_\ell(x,\cdot)\in W^{1,\infty}(\mathbb S^2)$ for almost
every $x$, then
\begin{align}
 \mathcal G_\ell
 &\geq \frac{c\ell^2}{T_\ell}
 \geq \frac{c\ell^{3/2}}{\mathcal E_\ell},
 \label{eq:gradient-action-toll}\\
 \frac{\mathcal G_\ell}{T_\ell}
 &\geq \frac{c\ell}{\mathcal E_\ell^2}.
 \label{eq:average-gradient-toll}
\end{align}
Consequently, if
$\operatorname*{ess\,sup}_x\|\nabla_\omega v_\ell(x)\|_\infty\leq B_\ell$,
then
\begin{equation}
 \mathcal E_\ell\geq c\sqrt{\frac{\ell}{B_\ell}}.
 \label{eq:Lipschitz-cost-floor}
\end{equation}
In particular, uniformly bounded scalar controls whose angular Lipschitz scale
is $B_\ell=O(\ell)$ have a strictly positive cost floor.  If each
$v_\ell(x,\cdot)$ is a spherical polynomial of degree at most $L_\ell$ and
$\|v_\ell\|_\infty\leq B$, the spherical Bernstein inequality gives
\begin{equation}
 \mathcal E_\ell\geq
 c\sqrt{\frac{\ell}{B L_\ell}}.
 \label{eq:bandwidth-cost-floor}
\end{equation}
Thus a bounded band-limited converter with $\mathcal E_\ell\to0$ must satisfy
$L_\ell\gtrsim \ell/\mathcal E_\ell^2\gg\ell$.
\end{corollary}

\begin{proof}
Use the variables of Proposition~\ref{prop:scalar-star-frame}.  Since
$|\alpha_s|=|\langle e_0,\psi_{\ell,s}\rangle|$, set
$\vartheta_s=\arccos|\alpha_s|$.  Equation
~\eqref{eq:scalar-star-frame-system} and
$\|\eta_s\|_2^2=1-|\alpha_s|^2$ imply, in the a.e. metric-derivative sense,
\begin{equation}
 |\vartheta_s'(x)|\leq s\|r_s(x)\|_2
 =s\|q_\ell(x)\|_2.
 \label{eq:star-frame-quantum-speed}
\end{equation}
Initially $\vartheta_s(0)=0$, while
~\eqref{eq:gap-cancelled-conversion} and $h_\ell\perp e_0$ give
$|\alpha_s(T_\ell)|\leq\varepsilon_*$.  After decreasing the fixed
$\varepsilon_*$ if necessary, integration proves
~\eqref{eq:charged-pulse-area}.  Cauchy--Schwarz and
~\eqref{eq:scalar-star-frame-cost} give
\begin{equation}
 c\leq
 \left(\int_0^{T_\ell}\|q_\ell\|_2\dd x\right)^2
 \leq T_\ell\mathcal E_\ell,
\end{equation}
which is~\eqref{eq:time-cost-uncertainty}.  On the other hand,
~\eqref{eq:gap-cancelled-cost-time} says
$T_\ell\leq C\mathcal E_\ell\sqrt\ell$.  Combining the last two inequalities
proves~\eqref{eq:gap-cancelled-quarter-cost}.

Write $g_\ell(x)=\|\nabla_\omega v_\ell(x)\|_\infty$.  Integrating
~\eqref{eq:continuous-gradient-growth} in $x$ and reversing the order of
integration gives
\begin{equation}
 \int_0^{T_\ell}\|A_3^{1/2}\psi_{\ell,s}(x)\|_2\dd x
 \leq s\int_0^{T_\ell}(T_\ell-t)g_\ell(t)\dd t
 \leq sT_\ell\mathcal G_\ell.
\end{equation}
Together with~\eqref{eq:gap-cancelled-current-time}, this yields the first
bound in~\eqref{eq:gradient-action-toll}; the cost--time estimate yields the
second and then~\eqref{eq:average-gradient-toll}.  The bound
~\eqref{eq:Lipschitz-cost-floor} follows by comparing the last display with
$\mathcal G_\ell/T_\ell\leq B_\ell$.  Finally, a spherical polynomial of
degree $L_\ell$ restricts to a trigonometric polynomial of that degree on each
great circle, so Bernstein's inequality gives
$\|\nabla_\omega v_\ell\|_\infty\leq C L_\ell\|v_\ell\|_\infty$.
This proves~\eqref{eq:bandwidth-cost-floor}.
\end{proof}

The conclusion is stronger than a state-frequency warning.  A cheap
gap-cancelled converter cannot be assembled from bounded scalar profiles at
the target angular scale: the scalar control itself must contain a
super-target rough component.  At the extremal pulse-area scale
$\mathcal E_\ell\asymp\ell^{-1/4}$, its degree must already be at least of
order $\ell^{3/2}$.  The remaining question is whether such a rough component
can be populated and refocused uniformly in energy.

Thus the magnitude of the off-diagonal star coupling is fully charged.  The
only uncharged freedom is its direction
$W_s(x)^*q(x)/\|q(x)\|$, which is rotated by the complementary scalar dynamics
and depends on the same ensemble parameter $s$.  Generic bilinear quantum
systems can transfer eigenstates with arbitrarily small $L^p$ control for
$p>1$, and related ensemble systems exploit averaging and noncommuting
brackets~\cite{BoussaidCaponigroChambrion2012,Chambrion2013}.  Those theorems
do not apply directly to the common scaling and scalar-multiplication geometry
here, but they rule out any argument based only on the smallness of
\eqref{eq:scalar-star-frame-cost}.  A proof of $\mathrm{SSNP}_\ell$ must charge
the energy-dependent directional rotation and its reciprocal refocusing.

\begin{remark}[Resonant-drive scaling and the first broadband barrier]
There is a natural coupled mechanism which survives the preceding
incompatibility at one energy.  Put
\begin{equation}
 M_\ell=\|h_\ell\|_\infty\asymp\ell^{1/4},
 \qquad w_\ell=M_\ell^{-1}h_\ell,
 \qquad \beta_\ell=(\sqrt{4\pi}M_\ell)^{-1}\asymp\ell^{-1/4}.
 \label{eq:resonant-drive-beta}
\end{equation}
Then $\|w_\ell\|_\infty=1$ and, for every real radial envelope $u$,
\begin{equation}
 \langle h_\ell,M_{u w_\ell}e_0\rangle=\beta_\ell u,
 \qquad
 \|M_{u w_\ell}e_0\|_2^2=\beta_\ell^2|u|^2.
 \label{eq:resonant-drive-coupling-cost}
\end{equation}
Freeze the forward angular evolution on a slab beginning at
$R_\ell=C\ell$.  Its $e_0$--$h_\ell$ gap is
\begin{equation}
 \Delta_\ell(k)=\frac{\ell(\ell+1)}{2kR_\ell^2}
 \longrightarrow\frac1{2C^2k}.
 \label{eq:frozen-resonant-gap}
\end{equation}
At a single $k_0$, the resonant choice
$u(x)=\epsilon\cos(\Delta_\ell(k_0)x)\mathbf1_{[0,T]}(x)$ with
$T\asymp(\beta_\ell\epsilon)^{-1}$ has order-one first-order transition
amplitude but charged cost
\begin{equation}
 \beta_\ell^2\int_0^T|u(x)|^2\dd x
 \asymp\beta_\ell\epsilon\longrightarrow0.
 \label{eq:single-energy-resonant-cost}
\end{equation}
For fixed $\epsilon>0$, one also has $T=O(\ell^{1/4})=o(R_\ell)$, so the
centrifugal gap is nearly frozen.  This scaling is therefore a serious
candidate mechanism, not a clean far-field phase mask.

The same calculation exposes its first obstruction.  In the linearized
frozen model the transition amplitude is
\begin{equation}
 c_\ell(\Delta)=-i\beta_\ell\widehat u(\Delta),
 \qquad
 \widehat u(\Delta)=\int_\mathbb R u(x)e^{i\Delta x}\dd x.
 \label{eq:born-resonant-amplitude}
\end{equation}
If $|c_\ell(\Delta)|\geq\eta$ throughout an interval $J$, Plancherel gives
the exact Born-level lower bound
\begin{equation}
 \int_\mathbb R\|M_{u(x)w_\ell}e_0\|_2^2\dd x
 =\beta_\ell^2\|u\|_2^2
 \geq\frac{\eta^2|J|}{2\pi}.
 \label{eq:born-broadband-lower-bound}
\end{equation}
For any fixed nondegenerate $K$, the interval
$J_\ell=\Delta_\ell(K)$ has length bounded below as $\ell\to\infty$.
Consequently the cheap resonant pulse is intrinsically single-energy at first
order.  A full converter must either prove a nonlinear composite-pulse
evasion of~\eqref{eq:born-broadband-lower-bound} through complementary modes,
or a positive proof must upgrade this Fourier budget to the exact reciprocal
radial evolution.  The statement here is only a rigorous diagnostic for the
linearized frozen model, not such an upgrade.
\end{remark}

The Fourier obstruction becomes stronger, rather than weaker, if the rough
control tries to launch the supercascade in one step.  Let $P_{\geq L}$ denote
the projection onto spherical harmonics of degrees at least $L$.

\begin{proposition}[No broadband one-jump supercascade at Born order]
\label{prop:born-high-band-recoil}
In the frozen laboratory~\eqref{eq:scaled-scalar-angular-laboratory}, let
$\psi_s^{(1)}(T)$ be the first Born correction to the free evolution from
$e_0$.  For every integer $L\geq1$,
\begin{equation}
 \int_{s_0}^{s_1}
 \|P_{\geq L}\psi_s^{(1)}(T)\|_2^2\dd s
 \leq C_S\frac{R^2}{L(L+1)}
 \int_0^T\|M_{v(x)}e_0\|_2^2\dd x.
 \label{eq:born-high-band-recoil}
\end{equation}
Consequently, if the first Born term has norm at least $\eta$ in degrees
$\geq L$ for every $s\in S$, then
\begin{equation}
 \int_0^T\|M_{v(x)}e_0\|_2^2\dd x
 \geq c_S\eta^2\frac{L^2}{R^2}.
 \label{eq:born-one-jump-cost}
\end{equation}
For $R\asymp\ell$, a direct broadband launch to the roughness scale
$L\gtrsim\ell/\mathcal E_\ell^2$ forced by
Corollary~\ref{cor:scalar-pulse-area-roughness} is therefore incompatible with
$\mathcal E_\ell\to0$ at Born order.
\end{proposition}

\begin{proof}
Let $\{Y_{n,m}\}$ be an orthonormal spherical-harmonic basis and set
\begin{equation}
 q_{n,m}(x)=\langle Y_{n,m},M_{v(x)}e_0\rangle,
 \qquad \lambda_n=\frac{n(n+1)}{R^2}.
\end{equation}
Up to an irrelevant terminal free phase, the first Born coefficient is
\begin{equation}
 b_{n,m}^{(1)}(s)
 =-is\int_0^Tq_{n,m}(x)e^{is\lambda_nx}\dd x.
\end{equation}
With the Fourier convention
$\widehat q(\xi)=\int q(x)e^{i\xi x}\dd x$, the change of variables
$\xi=s\lambda_n$ and Plancherel give
\begin{equation}
 \int_{s_0}^{s_1}|b_{n,m}^{(1)}(s)|^2\dd s
 \leq \frac{2\pi s_1^2}{\lambda_n}
 \int_0^T|q_{n,m}(x)|^2\dd x.
\end{equation}
Sum over $n\geq L$ and $|m|\leq n$, use
$\lambda_n\geq L(L+1)/R^2$, and apply Parseval in the angular variable.  This
proves~\eqref{eq:born-high-band-recoil}; the uniform lower bound gives
~\eqref{eq:born-one-jump-cost}.
\end{proof}

This leaves only a nonlinear bootstrap: a putative counterexample must first
concentrate substantial mass using lower modes, then use a small rough scalar
component where that mass is already concentrated, and finally undo the
entire excursion on a fixed energy band.  Neither a direct rough kick nor a
target-scale composite can do the job.

The phrase ``first concentrate'' can itself be charged exactly.  Write
$P_{\leq N}$ for the sum of the spherical-harmonic sectors of degrees at most
$N$, $Q_{>N}=\Id-P_{\leq N}$, and
\begin{equation}
 a(v):=\|M_ve_0\|_2=|\mathbb S^2|^{-1/2}\|v\|_2.
 \label{eq:scalar-charged-amplitude}
\end{equation}

\begin{proposition}[The charged gate precedes nonlinear high-sector amplification]
\label{prop:scalar-preconcentration-gate}
For every bounded real $v$ and every $N\geq0$,
\begin{equation}
 \|M_vP_{\leq N}\|_{\mathrm{op}}
 \leq \|M_vP_{\leq N}\|_{\mathrm{HS}}
 =(N+1)a(v).
 \label{eq:scalar-low-band-gate}
\end{equation}
Moreover, every normalized $f\in H^1(\mathbb S^2)$ satisfies
\begin{equation}
 \|M_vf\|_2^2
 \leq C\|v\|_\infty a(v)
 \bigl(1+\|A_3^{1/2}f\|_2\bigr).
 \label{eq:scalar-preconcentration-inequality}
\end{equation}
In particular, if $\|v\|_\infty\leq B$ and $\|M_vf\|_2\geq\eta$, then
\begin{equation}
 1+\|A_3^{1/2}f\|_2\geq
 \frac{c\eta^2}{B\,a(v)}.
 \label{eq:scalar-preconcentration-toll}
\end{equation}
If $\psi_s$ solves the frozen laboratory
\eqref{eq:scaled-scalar-angular-laboratory} with
$\|\psi_s(0)\|_2=1$ and $\psi_s(0)\in\Ran P_{\leq N}$, then, without any restriction on the evolution
inside $\Ran Q_{>N}$,
\begin{equation}
 \arcsin\|Q_{>N}\psi_s(T)\|_2
 \leq s(N+1)\int_0^Ta(v(x))\dd x.
 \label{eq:scalar-band-exit-angle}
\end{equation}
Thus $\|Q_{>N}\psi_s(T)\|_2\geq\eta$ implies
\begin{equation}
 \int_0^Ta(v(x))\dd x\geq
 \frac{\arcsin\eta}{s_1(N+1)},
 \qquad
 T\mathcal E_\ell\geq
 \frac{(\arcsin\eta)^2}{s_1^2(N+1)^2}.
 \label{eq:scalar-band-exit-cost}
\end{equation}
\end{proposition}

\begin{proof}
The addition theorem gives, pointwise,
\begin{equation}
 \sum_{n=0}^N\sum_{m=-n}^n|Y_{n,m}(\omega)|^2
 =\frac{(N+1)^2}{|\mathbb S^2|}.
\end{equation}
Consequently
\begin{equation}
 \|M_vP_{\leq N}\|_{\mathrm{HS}}^2
 =\frac{(N+1)^2}{|\mathbb S^2|}\|v\|_2^2,
\end{equation}
which proves~\eqref{eq:scalar-low-band-gate}.  Next,
$\|v\|_4^2\leq\|v\|_\infty\|v\|_2$ and the two-dimensional
Gagliardo--Nirenberg inequality on the sphere gives
\begin{equation}
 \|f\|_4^2\leq C\|f\|_2
 \bigl(\|A_3^{1/2}f\|_2+\|f\|_2\bigr).
\end{equation}
Hölder's inequality and $\|f\|_2=1$ prove
\eqref{eq:scalar-preconcentration-inequality}.

For the dynamical claim, split the Hamiltonian into its block-diagonal part
relative to $P_{\leq N}\oplus Q_{>N}$ and its off-diagonal part.  Passing to
the unitary interaction frame generated by the block-diagonal part leaves an
off-diagonal block of norm
$s\|Q_{>N}M_{v(x)}P_{\leq N}\|_{\mathrm{op}}$.  The standard two-subspace
angle estimate therefore gives
\begin{equation}
 \left|\frac{\dd}{\dd x}
 \arcsin\|Q_{>N}\psi_s(x)\|_2\right|
 \leq s\|Q_{>N}M_{v(x)}P_{\leq N}\|_{\mathrm{op}}.
\end{equation}
Now use~\eqref{eq:scalar-low-band-gate} and integrate.  The last assertion
follows from Cauchy--Schwarz and
$\mathcal E_\ell=\int_0^Ta(v(x))^2\dd x$.
\end{proof}

Estimate~\eqref{eq:scalar-band-exit-angle} is nonlinear: the entire
high--high block, including arbitrarily rough scalar amplification and
refocusing, has been retained in the block-diagonal propagator.  It says that
this amplification cannot occur before charged mass has crossed the chosen
spectral cut.  It does not sum over $N$, because one charged kick may cross
many cuts at once.  In fact, the transport-free scalar model disproves every
attempt to recover the quadratic cost merely by adding these gates.

\begin{proposition}[Wrapped phases disprove charged-gate additivity]
\label{prop:wrapped-phase-gate-additivity-no-go}
Fix $s_*>0$ and $0<\delta<1$, and put
$S_\delta=[s_*(1-\delta),s_*(1+\delta)]$.  In azimuthal coordinates define the
bounded real mask
\begin{equation}
 \Theta_L(\theta,\phi)
 :=\operatorname{Arg}(e^{iL\phi})\in(-\pi,\pi],
 \qquad L\geq1.
 \label{eq:wrapped-azimuthal-mask}
\end{equation}
Delete the centrifugal term and solve
\begin{equation}
 i\partial_x\psi_{\epsilon,L,s}
 =sM_{v_{\epsilon,L}(x)}\psi_{\epsilon,L,s},
 \qquad
 v_{\epsilon,L}(x,\omega)
 =\epsilon\Theta_L(\omega)
 \mathbf1_{[0,(s_*\epsilon)^{-1}]}(x),
 \qquad
 \psi_{\epsilon,L,s}(0)=e_0.
 \label{eq:wrapped-mask-evolution}
\end{equation}
Then
\begin{equation}
 \mathcal E_{\epsilon,L}
 =\int\|M_{v_{\epsilon,L}(x)}e_0\|_2^2\dd x
 =\frac{\pi^2}{3s_*}\epsilon\longrightarrow0,
 \label{eq:wrapped-mask-vanishing-cost}
\end{equation}
uniformly in $L$, but at the terminal time
$T_\epsilon=(s_*\epsilon)^{-1}$,
\begin{equation}
 \inf_{s\in S_\delta}
 \|P_{\geq L}\psi_{\epsilon,L,s}(T_\epsilon)\|_2
 \geq
 \frac{\sin(\pi\delta)}{\pi\delta}>0.
 \label{eq:wrapped-mask-high-band}
\end{equation}
Thus no exact nonlinear high-band estimate, nor any summation of the
first-crossing gates, can be controlled solely by the charged quadratic cost
after centrifugal transport has been deleted.
\end{proposition}

\begin{proof}
The evolution is explicit:
\begin{equation}
 \psi_{\epsilon,L,s}(T_\epsilon)
 =e^{-i(s/s_*)\Theta_L}e_0.
\end{equation}
Also
$|\mathbb S^2|^{-1}\|\Theta_L\|_2^2=\pi^2/3$, which proves
\eqref{eq:wrapped-mask-vanishing-cost}.  Set $a=s/s_*$.  The azimuthal Fourier
coefficient of $e^{-ia\Theta_L}$ at frequency $-L$ is
\begin{equation}
 c_{-L}(a)
 =\frac1{2\pi}\int_{-\pi}^{\pi}e^{-i(a-1)y}\dd y
 =\frac{\sin(\pi(a-1))}{\pi(a-1)}.
 \label{eq:wrapped-mask-sinc}
\end{equation}
Every spherical harmonic with azimuthal frequency $-L$ has degree at least
$L$.  Orthogonality in $\phi$ therefore gives
$\|P_{\geq L}\psi_{\epsilon,L,s}(T_\epsilon)\|_2\geq|c_{-L}(a)|$.
For $|a-1|\leq\delta<1$ this proves
\eqref{eq:wrapped-mask-high-band}.
\end{proof}

This is an exact fixed-band, all-orders failure, not a Born calculation.  It
does not contradict Proposition~\ref{prop:scalar-preconcentration-gate}:
\eqref{eq:scalar-band-exit-angle} charges the nonvanishing pulse area
$\int a(v_{\epsilon,L})\dd x$, whereas
\eqref{eq:wrapped-mask-vanishing-cost} makes only its quadratic cost vanish.
It also does not produce the converter, because the density remains uniform.
The centrifugal/density coupling is therefore indispensable.

There is a second exact toll.  Scalar multiplication changes phase but not
density, so all density concentration is generated by the centrifugal flow.

\begin{proposition}[Exact density-concentration action]
\label{prop:density-concentration-action}
Let $\psi_s$ be a normalized $H^1$ solution of the frozen laboratory
\eqref{eq:scaled-scalar-angular-laboratory}, starting from $e_0$, with the
action below finite, and assume
that for some fixed $\varepsilon<1$,
\begin{equation}
 \inf_{\gamma\in\mathbb R}
 \|\psi_s(T)-e^{i\gamma}h_\ell\|_2\leq\varepsilon.
 \label{eq:density-action-target}
\end{equation}
Then, uniformly for $s\in S$,
\begin{equation}
 \int_0^T\int_{\mathbb S^2}
 |\psi_s(x,\omega)|^2|\nabla_\omega\psi_s(x,\omega)|^2
 \dd\omega\dd x
 \geq c_{S,\varepsilon}R^2\ell^{1/2}.
 \label{eq:density-concentration-action}
\end{equation}
For the variable-coefficient first-order model
\begin{equation}
 i\partial_r\psi_s
 =s\bigl[r^{-2}A_3+M_{v(r)}\bigr]\psi_s,
 \label{eq:variable-angular-laboratory}
\end{equation}
the corresponding necessary condition is
\begin{equation}
 \int_{R_1}^{R_2}\frac1{r^2}
 \int_{\mathbb S^2}|\psi_s|^2|\nabla_\omega\psi_s|^2
 \dd\omega\dd r
 \geq c_{S,\varepsilon}\ell^{1/2}.
 \label{eq:variable-density-concentration-action}
\end{equation}
In particular, ideal bounded phase masks contribute exactly zero to the
increase of $\|\psi_s\|_4^4$; every phase-mask construction must create the
required concentration during its free-propagation pieces.
\end{proposition}

\begin{proof}
Put $\rho=|\psi_s|^2$ and
$J=\Im(\overline{\psi_s}\nabla_\omega\psi_s)$.  Since the scalar potential is
real, it cancels from the continuity equation:
\begin{equation}
 \partial_x\rho+\frac{2s}{R^2}\operatorname{div}_\omega J=0.
\end{equation}
For $Y(x)=\int_{\mathbb S^2}\rho(x)^2$, integration by parts and the pointwise
bounds $|\nabla\rho|\leq2|\psi_s||\nabla\psi_s|$ and
$|J|\leq|\psi_s||\nabla\psi_s|$ give
\begin{equation}
 |Y'(x)|\leq\frac{8s}{R^2}
 \int_{\mathbb S^2}|\psi_s|^2|\nabla_\omega\psi_s|^2.
 \label{eq:L4-density-derivative}
\end{equation}
For the real highest-weight harmonic, direct integration of
$|\sin^\ell\theta\cos(\ell\phi)|^p$ shows
\begin{equation}
 \|h_\ell\|_p\asymp_p
 \ell^{1/4-1/(2p)},
 \qquad 1\leq p<\infty.
 \label{eq:highest-weight-Lp}
\end{equation}
Choosing the phase in~\eqref{eq:density-action-target} and applying Hölder,
\begin{equation}
 1-\frac{\varepsilon^2}{2}
 \leq|\langle h_\ell,\psi_s(T)\rangle|
 \leq\|h_\ell\|_{4/3}\|\psi_s(T)\|_4,
\end{equation}
so $Y(T)=\|\psi_s(T)\|_4^4\geq c_\varepsilon\ell^{1/2}$, whereas
$Y(0)=|\mathbb S^2|^{-1}$.  Integrating
\eqref{eq:L4-density-derivative} proves
\eqref{eq:density-concentration-action}.  Repeating the computation with
$R^{-2}$ replaced by $r^{-2}$ proves
\eqref{eq:variable-density-concentration-action}.  Finally, multiplication by
a unimodular scalar phase leaves $\rho$ and hence $Y$ unchanged.
\end{proof}

For a smooth ideal product~\eqref{eq:phase-mask-product}, let
$\psi_{j,k}(\tau)$ denote the state during the $j$th free leg,
$0\leq\tau\leq\sigma_{j,\ell}$.  The same calculation, now with angular
coefficient $(2k)^{-1}$, shows that
\begin{equation}
 \sum_{j=1}^{m}\int_0^{\sigma_{j,\ell}}
 \int_{\mathbb S^2}|\psi_{j,k}|^2|\nabla_\omega\psi_{j,k}|^2
 \dd\omega\dd\tau
 \geq c_{K,\varepsilon}\ell^{1/2}
 \label{eq:phase-mask-density-action}
\end{equation}
whenever the product reaches the target with error at most $\varepsilon$,
uniformly for $k$ in a compact interval $K\Subset(0,\infty)$.  Thus a wrapped
phase may cheaply create enormous angular derivatives, but focusing still has
to convert those derivatives into the density action in
\eqref{eq:phase-mask-density-action}.  This condition does not exclude the
construction: the same cheap masks can make the integrand enormous during a
short free leg.  The statement extends to rough masks whenever the left-hand
side is finite by regular approximation; an infinite action already satisfies
the necessary condition.

The centrifugal term also gives a monotone radial energy.  This rules out the
stationary long weak slab even when its large centrifugal action is retained
rather than treated as an error.

\begin{proposition}[Radial energy/BV obstruction]
\label{prop:radial-energy-BV-obstruction}
\leavevmode\par\noindent
Let $v$ belong to
$W^{1,1}([R_1,R_2];L^\infty(\mathbb S^2;\mathbb R))$, and let
$\psi_s$ be a normalized $H^1$ solution of
\eqref{eq:variable-angular-laboratory} with $\psi_s(R_1)=e_0$.  Put
\begin{equation}
 \Gamma_v
 :=\|v(R_1)\|_\infty
   +\sup_{R_1\leq r\leq R_2}\|v(r)\|_\infty
   +\int_{R_1}^{R_2}\|\partial_rv(r)\|_\infty\dd r.
 \label{eq:radial-potential-variation}
\end{equation}
Then, uniformly in $s>0$,
\begin{align}
 \frac1{r^2}\|A_3^{1/2}\psi_s(r)\|_2^2
 &\leq \Gamma_v,
 \label{eq:radial-energy-bound}\\
 \int_{R_1}^{R_2}\frac1{r^2}
 \int_{\mathbb S^2}|\psi_s|\,|\nabla_\omega\psi_s|
 \dd\omega\dd r
 &\leq \sqrt{\Gamma_v}\log\frac{R_2}{R_1}.
 \label{eq:radial-BV-current-bound}
\end{align}
Consequently, fix
$\Sigma=[s_0,s_1]\Subset(0,\infty)$.  If $s\in\Sigma$ and the endpoint is
within a sufficiently small fixed $L^2$ error of $h_\ell$, then, for all
sufficiently large $\ell$,
\begin{equation}
 \Gamma_v\log^2\frac{R_2}{R_1}\geq c_\Sigma>0.
 \label{eq:radial-BV-conversion-toll}
\end{equation}

In particular, consider stationary weak slabs
\begin{equation}
 v_n(r,\omega)=\epsilon_nw_n(\omega),
 \qquad R_n\leq r\leq R_n+T_n,
 \label{eq:stationary-weak-angular-slab}
\end{equation}
where $R_n\geq R_*>0$,
$\|w_n\|_\infty\leq B$, and $a(w_n)\geq b>0$.  If
\begin{equation}
 \mathcal E_n=\epsilon_n^2T_na(w_n)^2\longrightarrow0,
 \label{eq:stationary-weak-slab-cost}
\end{equation}
then the left-hand side of
\eqref{eq:radial-BV-current-bound} tends to zero.  Hence no stationary long
weak slab, nor a version with switch-on and switch-off ramps of total
$L^\infty$ variation $O(\epsilon_n)$, can perform the scalar conversion in the
variable-coefficient first-order model.
\end{proposition}

\begin{proof}
Define the instantaneous angular energy
\begin{equation}
 \mathcal H_s(r)
 :=\left\langle\psi_s(r),
 \bigl[r^{-2}A_3+M_{v(r)}\bigr]\psi_s(r)\right\rangle.
\end{equation}
Differentiating along the Schr\"odinger flow cancels the commutator with the
Hamiltonian and gives the exact identity
\begin{equation}
 \mathcal H_s'(r)
 =-\frac2{r^3}\|A_3^{1/2}\psi_s(r)\|_2^2
 +\langle\psi_s(r),M_{\partial_rv(r)}\psi_s(r)\rangle.
 \label{eq:radial-energy-identity}
\end{equation}
The factor $s$ cancels.  Since
$\mathcal H_s(R_1)=\langle e_0,M_{v(R_1)}e_0\rangle$, integration of
\eqref{eq:radial-energy-identity}, followed by
$|\langle\psi,M_v\psi\rangle|\leq\|v\|_\infty$, proves
\eqref{eq:radial-energy-bound}.  Cauchy--Schwarz on the sphere then gives
\begin{equation}
 \frac1{r^2}\int|\psi_s|\,|\nabla_\omega\psi_s|
 \leq \frac1{r^2}\|A_3^{1/2}\psi_s\|_2
 \leq\frac{\sqrt{\Gamma_v}}r,
\end{equation}
and integration proves~\eqref{eq:radial-BV-current-bound}.

For the conversion toll, apply the continuity equation associated with
\eqref{eq:variable-angular-laboratory} to the one-Lipschitz test function
$\varphi(\theta,\phi)=|\theta-\pi/2|$.  The uniform density and the
highest-weight density have moments differing by
$\pi/2-1-O(\ell^{-1/2})$.  Endpoint $L^2$ error controls the density error in
$L^1$, while the continuity equation gives
\begin{equation}
 \left|\int_{\mathbb S^2}\varphi
 \bigl(|\psi_s(R_2)|^2-|e_0|^2\bigr)\dd\omega\right|
 \leq 2s\int_{R_1}^{R_2}\frac1{r^2}
 \int_{\mathbb S^2}|\psi_s|\,|\nabla_\omega\psi_s|
 \dd\omega\dd r.
\end{equation}
Since $s\leq s_1$, a sufficiently accurate conversion forces the current in
\eqref{eq:radial-BV-current-bound} to be bounded below uniformly for
$s\in\Sigma$.  This proves~\eqref{eq:radial-BV-conversion-toll}.

For~\eqref{eq:stationary-weak-angular-slab},
$\Gamma_{v_n}\leq2B\epsilon_n$.  Put
$\eta_n:=\mathcal E_n/b^2$.  Since $a(w_n)\geq b$,
$T_n\leq\eta_n/\epsilon_n^2$.  Therefore, with
$t_n:=\eta_n/(R_*\epsilon_n^2)$,
\begin{equation}
 \sqrt{\Gamma_{v_n}}
 \log\left(1+\frac{T_n}{R_n}\right)
 \leq C_B\left(\frac{\eta_n}{R_*}\right)^{1/4}
 t_n^{-1/4}\log(1+t_n)
 \leq C_{B,R_*}\eta_n^{1/4}\longrightarrow0,
\end{equation}
where $\sup_{t>0}t^{-1/4}\log(1+t)<\infty$.
Equation~\eqref{eq:radial-BV-current-bound} completes the proof.  Adding ramps
of total variation $O(\epsilon_n)$ leaves $\Gamma_{v_n}=O(\epsilon_n)$.
\end{proof}

\begin{corollary}[Radial superoscillation toll]
\label{cor:radial-superoscillation-toll}
Assume the gap-cancelled converter hypotheses of Proposition
~\ref{prop:gap-cancelled-supercascade}, now in the variable-coefficient
first-order model on $[R_\ell,R_\ell+T_\ell]$, with
$R_\ell\asymp\ell$, and assume the same duration bound
\begin{equation}
 T_\ell\leq C\mathcal E_\ell\sqrt\ell.
 \label{eq:radial-gap-duration}
\end{equation}
If the conversion succeeds with fixed sufficiently small error, then
\begin{equation}
 \Gamma_{v_\ell}
 \geq c\frac{R_\ell^2}{T_\ell^2}
 \geq c\frac{\ell}{\mathcal E_\ell^2}.
 \label{eq:radial-superoscillation-toll}
\end{equation}
If in addition $\|v_\ell\|_\infty\leq B$ and its radial dependence has
$m_\ell$ bounded-amplitude monotone layers, then
\begin{equation}
 m_\ell\geq c_B\frac{\ell}{\mathcal E_\ell^2}.
 \label{eq:radial-layer-count-toll}
\end{equation}
Thus every vanishing-cost gap-cancelled converter is forced to be super-rough
both angularly and radially.
\end{corollary}

\begin{proof}
Equation~\eqref{eq:radial-BV-conversion-toll} and
$\log(1+T_\ell/R_\ell)\leq T_\ell/R_\ell$ give the first inequality in
\eqref{eq:radial-superoscillation-toll}; use
\eqref{eq:radial-gap-duration} and $R_\ell\asymp\ell$ for the second.
A bounded profile with $m_\ell$ monotone layers has
$\Gamma_{v_\ell}\leq C B(1+m_\ell)$, which proves
\eqref{eq:radial-layer-count-toll}.
\end{proof}

The preceding lower bound concerns raw variation.  Raw variation by itself is
not a dynamical resource and cannot be placed on the cost side of the desired
estimate.

\begin{proposition}[Raw radial variation is not a cost proxy]
\label{prop:raw-radial-variation-no-go}
Fix $I=[R_1,R_2]$, a compact
$\Sigma=[s_0,s_1]\Subset(0,\infty)$, and a real
$v\in W^{1,1}(I;L^\infty(\mathbb S^2))$.  Let $U_s^v$ denote the propagator
of~\eqref{eq:variable-angular-laboratory}.  There are real smooth
perturbations $z_n$, supported away from the endpoints of $I$, such that
\begin{align}
 \int_I a(z_n(r))^2\dd r&\longrightarrow0,
 &
 \int_I\|z_n(r)\|_\infty\dd r&\longrightarrow0,
 \label{eq:inert-radial-oscillation-small}\\
 \int_I\|\partial_rz_n(r)\|_\infty\dd r&\longrightarrow\infty,
 &
 \int_I\|\nabla_\omega z_n(r)\|_2^2\dd r&\longrightarrow0,
 \label{eq:inert-radial-oscillation-rough}
\end{align}
while
\begin{equation}
 \sup_{s\in\Sigma}
 \|U_s^{v+z_n}(R_2,R_1)-U_s^v(R_2,R_1)\|_{\mathrm{op}}
 \longrightarrow0.
 \label{eq:inert-radial-oscillation-propagator}
\end{equation}
Moreover,
\begin{equation}
 \left|
 \int_Ia(v(r)+z_n(r))^2\dd r
 -\int_Ia(v(r))^2\dd r
 \right|\longrightarrow0.
\end{equation}
Consequently,
$\Gamma_{v+z_n}$ can be inflated without appreciable charged cost or dynamical
effect.  No coercive bridge from raw radial variation to the constant-column
cost can close the converter problem.
\end{proposition}

\begin{proof}
Choose nonzero real $w\in C^\infty(\mathbb S^2)$ and
$0\leq\chi\in C_c^\infty((R_1,R_2))$, and put
\begin{equation}
 z_n(r,\omega)
 =\delta_n\chi(r)\sin(\nu_nr)w(\omega),
 \qquad
 \delta_n\downarrow0,
 \qquad
 \delta_n\nu_n\longrightarrow\infty.
\end{equation}
The first, second, and fourth assertions follow directly from the factor
$\delta_n$.  Periodic averaging gives
$\int_I\chi(r)|\cos(\nu_nr)|\dd r\to
(2/\pi)\int_I\chi(r)\dd r$; hence the term
$\delta_n\nu_n\chi(r)\cos(\nu_nr)w$ in $\partial_rz_n$ proves the third
assertion.  Since $v$ has finite radial variation, the reverse triangle
inequality also gives $\Gamma_{v+z_n}\to\infty$.
Bounded-perturbation Duhamel gives
\begin{equation}
 \sup_{s\in\Sigma}
 \|U_s^{v+z_n}(R_2,R_1)-U_s^v(R_2,R_1)\|_{\mathrm{op}}
 \leq s_1\int_I\|z_n(r)\|_\infty\dd r,
\end{equation}
which proves~\eqref{eq:inert-radial-oscillation-propagator}.  The cost
assertion follows from Cauchy--Schwarz and
$\int_Ia(z_n)^2\to0$.
\end{proof}

The dynamically relevant replacement is the radial work tested against the
evolving density.  It admits an exact scalar continuity representation and
therefore yields a first, genuinely coupled cost--action bridge.

\begin{proposition}[Effective radial-work cost--action bridge]
\label{prop:effective-radial-work-bridge}
Let
\begin{equation}
 v\in W^{1,1}([R_1,R_2];L^\infty(\mathbb S^2))
 \cap L^2([R_1,R_2];H^1(\mathbb S^2))
\end{equation}
be real, with $v(R_1)=v(R_2)=0$, and let $\psi_s$ solve
~\eqref{eq:variable-angular-laboratory} from $e_0$.  Define
\begin{align}
 \mathfrak W_s
 &:=\int_{R_1}^{R_2}
 \langle\psi_s,M_{\partial_rv}\psi_s\rangle\dd r,
 \label{eq:effective-radial-work}\\
 \mathfrak R_v
 &:=\int_{R_1}^{R_2}\|\nabla_\omega v(r)\|_2^2\dd r,
 &
 \mathfrak D_s
 &:=\int_{R_1}^{R_2}\frac1{r^2}
 \int_{\mathbb S^2}|\psi_s|^2|\nabla_\omega\psi_s|^2
 \dd\omega\dd r.
 \label{eq:radial-work-factors}
\end{align}
Assume $\mathfrak D_s<\infty$; the lower bounds below are automatic if it is
infinite.
Then the exact identities
\begin{align}
 \mathfrak W_s
 &=\frac1{R_2^2}\|A_3^{1/2}\psi_s(R_2)\|_2^2
   +2\int_{R_1}^{R_2}\frac1{r^3}
       \|A_3^{1/2}\psi_s(r)\|_2^2\dd r,
 \label{eq:radial-work-energy-identity}\\
 \mathfrak W_s
 &=-2s\int_{R_1}^{R_2}\frac1{r^2}
 \int_{\mathbb S^2}
 \nabla_\omega v\cdot
 \Im(\overline{\psi_s}\nabla_\omega\psi_s)
 \dd\omega\dd r
 \label{eq:radial-work-current-identity}
\end{align}
hold.  Fix $\Sigma=[s_0,s_1]\Subset(0,\infty)$ and constants
$0<c_R<C_R<\infty$.  If
\begin{equation}
 c_R\ell\leq R_1\leq R_2\leq C_R\ell,
 \qquad s\in\Sigma,
 \qquad
 \inf_{\gamma\in\mathbb R}
 \|\psi_s(R_2)-e^{i\gamma}h_\ell\|_2\leq\varepsilon_0
 \label{eq:radial-work-conversion-hypothesis}
\end{equation}
for a sufficiently small fixed $\varepsilon_0$, then
\begin{equation}
 \mathfrak W_s\geq c>0,
 \qquad
 \mathfrak R_v\mathfrak D_s\geq cR_1^2.
 \label{eq:radial-work-product-toll}
\end{equation}
If, in addition, $v(r)$ has angular harmonic degree at most $L\geq1$ for
almost every $r$, then its charged cost
\begin{equation}
 \mathcal E(v):=\int_{R_1}^{R_2}a(v(r))^2\dd r
\end{equation}
satisfies
\begin{equation}
 \mathcal E(v)\mathfrak D_s
 \geq c\frac{R_1^2}{L(L+1)}.
 \label{eq:bandlimited-cost-action-bridge}
\end{equation}
Thus, at $R_1\asymp\ell$ and target-scale bandwidth $L\lesssim\ell$, a
vanishing-cost converter would have to create density action at least
$\mathfrak D_s\gtrsim1/\mathcal E(v)$.
\end{proposition}

\begin{proof}
Integrating~\eqref{eq:radial-energy-identity}, using the endpoint conditions
on $v$, $A_3e_0=0$, and $\psi_s(R_1)=e_0$, proves
~\eqref{eq:radial-work-energy-identity}.  Next put
$\rho_s=|\psi_s|^2$.  Its exact continuity equation is
\begin{equation}
 \partial_r\rho_s
 +\frac{2s}{r^2}\operatorname{div}_\omega
 \Im(\overline{\psi_s}\nabla_\omega\psi_s)=0.
\end{equation}
Integrating $\int\rho_s\partial_rv$ by parts first in $r$ and then on the
sphere proves~\eqref{eq:radial-work-current-identity}.

The endpoint hypothesis gives
$|\langle h_\ell,\psi_s(R_2)\rangle|\geq
1-\varepsilon_0^2/2$.  Therefore
\begin{equation}
 \|A_3^{1/2}\psi_s(R_2)\|_2^2
 \geq\ell(\ell+1)
 |\langle h_\ell,\psi_s(R_2)\rangle|^2,
\end{equation}
and~\eqref{eq:radial-work-energy-identity} gives
$\mathfrak W_s\geq c$.  By~\eqref{eq:radial-work-current-identity} and
Cauchy--Schwarz,
\begin{equation}
 |\mathfrak W_s|
 \leq2s_1\sqrt{\mathfrak R_v}
 \left(
 \int_{R_1}^{R_2}\frac1{r^4}
 \int_{\mathbb S^2}|\psi_s|^2|\nabla_\omega\psi_s|^2
 \dd\omega\dd r
 \right)^{1/2}
 \leq\frac{2s_1}{R_1}
 \sqrt{\mathfrak R_v\mathfrak D_s}.
\end{equation}
This proves~\eqref{eq:radial-work-product-toll}.  Finally, angular
bandlimiting gives
\begin{equation}
 \mathfrak R_v
 \leq L(L+1)\int_{R_1}^{R_2}\|v(r)\|_2^2\dd r
 =4\pi L(L+1)\mathcal E(v),
\end{equation}
which proves~\eqref{eq:bandlimited-cost-action-bridge}.
\end{proof}

These results expose the nonlinear bootstrap without closing it.  A
small charged component can act strongly only after the state has acquired
large angular regularity by~\eqref{eq:scalar-preconcentration-inequality}; the
first escape through every fixed spectral cut is nevertheless charged by
\eqref{eq:scalar-band-exit-angle}; and matching the highest-weight state forces
the independent density action~\eqref{eq:density-concentration-action}.
Proposition~\ref{prop:wrapped-phase-gate-additivity-no-go} eliminates
multiscale additivity based only on those gates, while
Proposition~\ref{prop:radial-energy-BV-obstruction} excludes stationary and
radially low-variation weak slabs once centrifugal transport is restored.
Corollary~\ref{cor:radial-superoscillation-toll} moreover forces radial
variation and bounded-amplitude layer count
$\gtrsim\ell/\mathcal E_\ell^2$ under gap cancellation, but Proposition
~\ref{prop:raw-radial-variation-no-go} proves that this raw BV quantity can be
inflated at negligible cost and with negligible dynamical effect.
Proposition~\ref{prop:effective-radial-work-bridge} replaces it by the exact
state-correlated work and proves the product bridge
~\eqref{eq:bandlimited-cost-action-bridge}.  What is still missing is a
nonlinear estimate controlling this effective work when both the scalar
bandwidth and the evolving density action are super-rough.  Without such an
estimate, this phase is a sharper reduction, not a proof of
$\mathrm{SSNP}_\ell$.

\begin{remark}[Controllability warning]
The refusal to promote the preceding tolls to a no-go theorem is essential.
In other bilinear Schr\"odinger systems, directions lost by the linearized or
quadratic dynamics can be recovered by a cubic mechanism for rougher controls
\cite{Bournissou2022}, and phase control can combine with Lie brackets to
generate small-time approximate density transport by diffeomorphisms
\cite{BeauchardPozzoli2025}.  Those results do not provide the charged
$L^2$-cost decay, fixed-band uniformity, or reciprocal conversion required
here.  They do show that Born recoil plus a first-crossing estimate cannot be
declared nonlinear-complete.  Conversely, known bilinear Schr\"odinger
obstructions use spacetime smoothing effects \cite{ChambrionThomann2020},
which supports the present conclusion that the missing bridge should be a
uniform smoothing or nonlinear-Plancherel estimate that controls the
state-correlated work~\eqref{eq:effective-radial-work}, not raw radial
variation.
\end{remark}

\begin{problem}[Physically scaled scalar nonlinear-Plancherel and refocusing test]
\label{prob:achromatic-phase-mask-refocusing}
First consider the frozen angular laboratory
\begin{equation}
 \begin{aligned}
 i\partial_xU_{\ell,s}(x)
 &=s\bigl[D_{\ell,R}+M_{v(x)}\bigr]U_{\ell,s}(x),
 &U_{\ell,s}(0)&=\Id,\\
 D_{\ell,R}&=R^{-2}A_3.
 \end{aligned}
 \label{eq:scaled-scalar-angular-laboratory}
\end{equation}
where $s\in S=[s_0,s_1]\Subset(0,\infty)$, $v$ is real and bounded, and its
endpoint values vanish.  Put
\begin{equation}
 B_{\ell,v}(s)=\langle h_\ell,U_{\ell,s}(T)e_0\rangle.
\end{equation}
Decide whether the scalar-multiplication estimate
\begin{equation}
 \int_S-\log\bigl(1-|B_{\ell,v}(s)|^2\bigr)\dd s
 \leq C_S\int_0^T\|M_{v(x)}e_0\|_2^2\dd x
 \label{eq:scaled-scalar-nonlinear-Plancherel}
\end{equation}
denoted $\mathrm{SSNP}_\ell$, holds uniformly in $\ell$, $R\asymp\ell$, and
$T$.  Proposition
~\ref{prop:scaled-matrix-ensemble-no-go} proves that scalar multiplication is
essential, while Proposition~\ref{prop:scalar-ground-adiabatic-no-go} excludes
the direct ground-state adiabatic counterexample.  A proof of
~\eqref{eq:scaled-scalar-nonlinear-Plancherel}, followed by a controlled
freezing and directional-mode argument, excludes the converter.  A failure
must exhibit a nonadiabatic or rotating-frame scalar multimode mechanism.

As a constructive stress test, let
\begin{equation}
 \mathcal P_{w,k}f=e^{-iw/(2k)}f,
 \qquad
 \mathcal F_{\sigma,k}f=e^{-i\sigma A_3/(2k)}f.
\end{equation}
Decide whether bounded real masks $w_{j,\ell}$ and nonnegative propagation
times $\sigma_{j,\ell}$ can produce products
\begin{equation}
 \mathcal U_{\ell,k}
 =\mathcal P_{w_{m,\ell},k}\mathcal F_{\sigma_{m,\ell},k}
  \cdots
  \mathcal P_{w_{1,\ell},k}\mathcal F_{\sigma_{1,\ell},k}
 \label{eq:phase-mask-product}
\end{equation}
such that, on a fixed positive-measure energy interval,
\begin{equation}
 \sup_{k\in K}\inf_{\gamma(k)}
 \|\mathcal U_{\ell,k}e_0-e^{i\gamma(k)}h_\ell\|_2=o(1),
 \label{eq:phase-mask-conversion-target}
\end{equation}
with the reciprocal directional conversion needed by the barrier argument.
If so, realize the masks by separated weak slabs as in Proposition
~\ref{prop:cheap-scalar-phase-mask} and justify the full second-order radial
error including $A_3/r^2$ and the incompatibility in Proposition
~\ref{prop:phase-mask-transport-incompatibility}.  Equivalently, replace the
ideal product by a coupled weak-slab evolution and prove refocusing directly.
 If this cannot be done, prove an obstruction which uses uniformity in $k$ or
 bidirectionality together with the radial action budget.  Angular Sobolev
growth alone cannot work because of Proposition~\ref{prop:cheap-scalar-phase-mask},
but Proposition~\ref{prop:gap-cancelled-supercascade} shows that angular growth
combined with the scalar cost already excludes every shallow gap-cancelled
composite.
 Corollary~\ref{cor:scalar-pulse-area-roughness} now excludes every
 target-scale Lipschitz realization and forces any cheap gap-cancelled control
 into angular degrees $\gtrsim\ell/\mathcal E_\ell^2$.  Proposition
 ~\ref{prop:born-high-band-recoil} excludes direct fixed-band injection into
 that scale at first order, while Proposition
 ~\ref{prop:rotating-ground-ladder-fanout} excludes the natural adiabatic
 adjacent-rung bootstrap.  Propositions
 ~\ref{prop:scalar-preconcentration-gate} and
 ~\ref{prop:density-concentration-action} now charge the first spectral escape,
 strong action by a small charged component, and the density concentration
 needed to reach $h_\ell$.  Proposition
 ~\ref{prop:wrapped-phase-gate-additivity-no-go} shows that the charged gates
 cannot be made additive without centrifugal transport, while Proposition
 ~\ref{prop:radial-energy-BV-obstruction} excludes stationary and radially
 low-variation weak slabs when that transport is restored.  Corollary
 ~\ref{cor:radial-superoscillation-toll} shows that any surviving
 gap-cancelled profile is doubly super-rough: its radial variation or bounded
 layer count, as well as its angular bandwidth, is at least
 $\gtrsim\ell/\mathcal E_\ell^2$.  Proposition
 ~\ref{prop:raw-radial-variation-no-go} shows that this raw variation can be
 dynamically inert, whereas Proposition
 ~\ref{prop:effective-radial-work-bridge} proves the exact bandlimited bridge
 ~\eqref{eq:bandlimited-cost-action-bridge}.  The remaining stronger
 full-space target is to control the state-correlated work
 ~\eqref{eq:effective-radial-work} by the charged cost for a super-bandwidth
 radially oscillatory or moving-support profile, using fixed-band uniformity
 and reciprocal refocusing to prevent simultaneous blow-up of
 $\mathfrak R_v$ and $\mathfrak D_s$.  The estimate must remain uniform through
 the scalar--centrifugal excursion and its reverse passage.  By
 Propositions~\ref{prop:scaled-matrix-ensemble-no-go} and
 ~\ref{prop:scalar-ground-adiabatic-no-go}, its proof must use scalar
 multiplication, the cost of leaving the ground bundle, and the directional
 rotation isolated by Proposition~\ref{prop:scalar-star-frame}; no abstract
 ensemble inequality can suffice.
\end{problem}

\begin{remark}[Finite-grid diagnostic only]
The idealized product~\eqref{eq:phase-mask-product} was stress-tested on a
periodic one-dimensional angular grid with generator $-\partial_x^2$.  The
target was a normalized localized oscillatory packet of nominal degree $5$,
and the masks and free-propagation times were optimized simultaneously at five
 values of the common dispersion parameter $s=1/(2k)$.  On a dense $41$-point
 check over $s\in[0.95,1.05]$, a six-mask run reached minimum fidelity $0.987$
 and mean fidelity $0.995$; over the wider interval $[0.75,1.25]$, a four-mask
 run reached minimum fidelity $0.779$ and mean fidelity $0.910$.  The optimizer
 hit its iteration limit, so these are neither certified optima nor asymptotic
evidence.  They serve only as a falsification warning: modest-band dependence
on $1/k$ does not by itself give an obvious algebraic no-go.  This model has no
radial coordinate, is not $\mathbb S^2$, and does not satisfy the action budget
of Proposition~\ref{prop:phase-mask-transport-incompatibility}; it therefore
provides no evidence for a full scalar converter.  The reproducible diagnostic
is recorded in
\nolinkurl{tmp/numerical_verification/phase_mask_circle_ensemble.py}.
\end{remark}

\begin{remark}[Full-modal barrier/composite diagnostic only]
The scale obstruction in Proposition~\ref{prop:gap-cancelled-supercascade} was
also checked against an explicit compressed composite pulse.  Alternating
$X$ and $Z$ rotations with seven recorded pulse areas give minimum fidelity
$0.999996$ on a dense check of $s\in[0.9,1.1]$ in the ideal two-state model.
The same pulse areas were then realized by a normalized equatorial scalar
barrier and the real highest-weight mixer $h_\ell/\|h_\ell\|_\infty$ in the
spherical-harmonic sector consisting of degrees $n\leq4\ell$ and azimuthal
orders in $\ell\mathbb Z$.  At $\ell=8$ this is a $137$-dimensional Galerkin
system.  Its minimum, mean, and maximum fidelities were respectively
$0.000203$, $0.0375$, and $0.184$.  This dramatic collapse is exactly the
complementary-sector effect erased by compression.

The calculation is not a lower bound: the barrier is not asymptotic at
$\ell=8$, the recorded pulse was optimized only in the compressed model, and
a different deep cascade could behave differently.  It is therefore only a
falsification check on the ordinary finite-depth composite idea, consistent
with the rigorous frequency requirement
\eqref{eq:gap-cancelled-supercascade}.  The reproducible calculation is
recorded in
\nolinkurl{tmp/numerical_verification/barrier_composite_sphere.py}.
\end{remark}

\section{A fixed-potential Riesz--Talbot counterexample to the buffered
$N=1$ entropy route}
\label{app:riesz-talbot-fixed-counterexample}

The preceding converter restrictions leave one genuinely nonlinear escape:
concentrate the selected angular state on a set of exponentially small area,
apply a translated reflector there, and use reciprocity to refocus the
transmitted state.  This section proves that the escape is real.  It gives a
single bounded zonal scalar potential in $X_3$ for which the selected
$N=1$ entropy diverges along every cofinal positively smoothed, radially
truncated, double-buffered diagonal.  The result disproves the proposed
fixed-potential selected-channel theorem.  It does not determine the spectral
multiplicity of the limiting operator and therefore does not disprove Simon's
conjecture.

\subsection{Translated barriers and the exact Hardy isometry}

\begin{lemma}[Translated-barrier Hardy isometry]
\label{lem:translated-barrier-hardy}
Let $\mathcal H$ be finite dimensional.  Suppose that $A,B\in
\mathcal B(\mathcal H)$ satisfy
\begin{equation}
 A^*A-B^*B=\Id,\qquad A\text{ is invertible},\qquad
 \|BA^{-1}\|<1,
 \label{eq:barrier-flux-data}
\end{equation}
and let $S$ be unitary.  For $|z|<1$ put
\begin{equation}
 F(z)=(A+zSB)^{-1}.
 \label{eq:barrier-hardy-function}
\end{equation}
Then
\begin{equation}
 \frac1{2\pi}\int_0^{2\pi}
 F(e^{i\phi})^*F(e^{i\phi})\dd\phi=\Id,
 \label{eq:barrier-hardy-isometry}
\end{equation}
and consequently, for every $y\ne0$,
\begin{equation}
 \frac1{2\pi}\int_0^{2\pi}
 \log\frac{\|y\|}{\|F(e^{i\phi})y\|}\dd\phi\geq0.
 \label{eq:barrier-hardy-nonnegative-loss}
\end{equation}

Let $P$ be an orthogonal projection fixed by conjugation and let
\begin{equation}
 A=\Id+(\cosh\rho-1)P,\qquad B=(\sinh\rho)P,\qquad \rho>0.
 \label{eq:localized-hyperbolic-barrier}
\end{equation}
If $S\bar y=y$ and, for some $|\lambda|=1$,
\begin{equation}
 \|(\Id-P)y\|+\|\bar y-\lambda y\|
 \leq\epsilon\|y\|,
 \label{eq:localized-phase-real-state}
\end{equation}
then
\begin{equation}
 \frac1{2\pi}\int_0^{2\pi}
 \log\frac{\|y\|}{\|F(e^{i\phi})y\|}\dd\phi
 \geq\log\cosh\rho-C_\rho\epsilon.
 \label{eq:localized-barrier-positive-loss}
\end{equation}
\end{lemma}

\begin{proof}
Set $C=BA^{-1}$ and $T=SC$.  Factoring on the right gives
\begin{equation}
 F(z)=A^{-1}(\Id+zT)^{-1}
 =A^{-1}\sum_{n\geq0}(-zT)^n.
\end{equation}
The flux identity gives
\begin{equation}
 A^{-*}A^{-1}=\Id-C^*C=\Id-T^*T.
\end{equation}
Parseval and telescoping therefore yield
\begin{equation}
 \frac1{2\pi}\int F^*F
 =\sum_{n\geq0}T^{*n}(\Id-T^*T)T^n=\Id.
\end{equation}
Jensen applied to $\log\|F(e^{i\phi})y\|^2$ proves
\eqref{eq:barrier-hardy-nonnegative-loss}.

If $Py=y$ and $\bar y=\lambda y$, then $Sy=\bar\lambda y$ and
\begin{equation}
 F(e^{i\phi})y
 =\frac{y}{\cosh\rho+
   e^{i\phi}\bar\lambda\sinh\rho}.
\end{equation}
The phase mean of the logarithm of the denominator is
$\log\cosh\rho$.  In the approximate case the inverse in
\eqref{eq:barrier-hardy-function} has norm at most $e^\rho$.
The resolvent identity and~\eqref{eq:localized-phase-real-state} compare
$F(e^{i\phi})y$ uniformly with the preceding scalar expression with error
$C_\rho\epsilon\|y\|$.  Both norms lie between fixed multiples of $\|y\|$,
so the logarithm is uniformly Lipschitz and
\eqref{eq:localized-barrier-positive-loss} follows.
\end{proof}

\subsection{Radial realization at summable charged cost}

The angular operators constructed later in this section are realized by one
scalar radial cell as follows.
The order in which the scales are chosen is important.  First fix the finite
incoming angular space, the number of Riesz factors, the desired error, and
the finite integer ratios between their frequencies.  Next choose the weak
slab lengths.  Only then choose a common odd carrier so large that all of the
frequencies are simultaneously far above the inverse radial scale while the
entire cell remains far below the square of the smallest frequency.  This
leaves room for both the forward and reverse Talbot flights.

\begin{proposition}[A scalar reciprocal Riesz cell]
\label{prop:scalar-reciprocal-riesz-cell}
Let $K\Subset(0,\infty)$, let $\mathcal K$ be a finite-dimensional,
conjugation-invariant zonal space containing $e_0$, and fix $m\geq1$,
$\rho>0$, and $\epsilon>0$.  Outside any prescribed radius $R_-$ there is a
bounded, compactly supported, real zonal potential $v_{m,\epsilon}$ with
vanishing spherical mean at every radius such that the following hold.

\begin{enumerate}
 \item At a chosen tuning point $k_\star\in\operatorname{int}K$, its forward
 Riesz train, localized hyperbolic reflector, and reverse train satisfy the
 hypotheses of Proposition~\ref{prop:robust-reciprocal-cascade} on
 $\mathcal K$, with cell error at most $\epsilon$.
 \item If $s=k_\star/k$, the error is at most
 \begin{equation}
   C_{K,\mathcal K}m|s-1|+\epsilon.
   \label{eq:physical-riesz-cell-chromatic-error}
 \end{equation}
 \item Inserting a sufficiently distant free collar before the cell changes
 its transfer, up to error $\epsilon$, only through the scalar circle phase
 $z=e^{i\vartheta(k)}$; the derivative of $\vartheta$ may be made arbitrarily
 large, uniformly on $K$.
 \item Its charged cost obeys
 \begin{equation}
  \int\!\int_{\mathbb S^2}|v_{m,\epsilon}(r,\omega)|^2
       \dd\omega\dd r
  \leq C_\rho e^{-c m}+\epsilon.
  \label{eq:physical-riesz-cell-cost}
 \end{equation}
\end{enumerate}
The supports of any prescribed countable collection of such cells can be
chosen disjoint.
\end{proposition}

\begin{proof}
We give the scale ledger, including the two approximations used repeatedly
below.  On a slab $[R,R+T]$ put $T^{-1}w(\omega)$, where $w$ is a fixed
finite-band real multiplier of degree at most $D$.  Variation of constants,
uniformly for $k\in K$, gives
\begin{align}
 \mathsf T_{T,w}(k)
 &=\exp\!\left(-\frac{i}{2k}w\right)
   +O_K\!\left(T^{-1}+\frac{D^2T}{R^2}+\frac{T}{R}\right),
 \label{eq:long-weak-mask-transmission}\\
 \mathsf R_{T,w}(k)
 &=O_K\!\left(T^{-1}+\frac{D^2T}{R^2}+\frac{T}{R}\right),
 \label{eq:long-weak-mask-reflection}\\
 \int_R^{R+T}\|T^{-1}w\|_2^2\dd r
 &=T^{-1}\|w\|_2^2.
 \label{eq:long-weak-mask-cost}
\end{align}
The first error is the one-dimensional long-weak-slab error, the second comes
from the angular generator on the slab, and the third from replacing $r$ by
$R$.  The estimates also hold after one $k$ derivative.  On a free collar
$[r_0,r_1]$, the outgoing evolution on every fixed angular block is
\begin{equation}
 \exp\!\left[-\frac{iA_3}{2k}
       \left(\frac1{r_0}-\frac1{r_1}\right)\right]
 +O_K(r_0^{-1}),
 \label{eq:free-collar-angular-flight}
\end{equation}
with the error uniform after the block has been fixed.  These are the standard
Volterra and outgoing Hankel asymptotics; their uniformity is exactly the
finite-block uniformity imposed in the statement.

Apply Proposition~\ref{prop:talbot-riesz-concentration} and first record its
finite odd frequency ratios $1=q_1<\cdots<q_m=Q$.  Choose the mask-slab
lengths so large that the sum of~\eqref{eq:long-weak-mask-cost}, except for
the reflector, is below $\epsilon/10$.  Now choose a common odd carrier
$L_\ast$ so large that, with $L_t=L_\ast q_t$ and with $T_{\max}$ the largest
slab length, there is a radial interval satisfying
\begin{equation}
 R_-+\operatorname{length}(\text{cell})\ll R,\qquad
 QL_\ast\sqrt{T_{\max}}\ll R\ll L_\ast^2.
 \label{eq:riesz-cell-scale-separation}
\end{equation}
This interval is nonempty once $L_\ast\gg Q\sqrt{T_{\max}}$.  Thus the
centrifugal error in every weak mask slab tends to zero, while
\eqref{eq:free-collar-angular-flight} can realize both
$\pi/(2L_t^2)$ and $3\pi/(2L_t^2)$ in the forward and reverse trains.  A
successive enlargement of $L_\ast$, $R$, and the slab lengths makes the sum
of all transmission, reflection, spherical-Talbot, and outgoing errors below
$\epsilon$.

At the focus insert a fixed smooth radial bump $Q_0$ multiplying
\begin{equation}
 \chi_m-\langle\chi_m\rangle_{\mathbb S^2}.
 \label{eq:mean-zero-localized-reflector}
\end{equation}
Choose the bump strength so that the induced two-port coefficient on the
concentrated state is~\eqref{eq:localized-hyperbolic-barrier}.  Subtracting
the spherical mean changes only the scalar channel; a compensating
reflectionless scalar phase slab absorbs that change.  Equations
\eqref{eq:riesz-concentration}--\eqref{eq:riesz-small-support-cost} and the
stability of finite-dimensional transfer matrices show that the resulting
cell has error $O(\epsilon+e^{-cm})$ and reflector cost $C_\rho e^{-cm}$.
All masks have zero mean because the frequencies are odd, and
\eqref{eq:mean-zero-localized-reflector} has zero mean by construction.

The reverse train gives~\eqref{eq:riesz-refocusing}; differentiating the
finite product in $s=k_\star/k$ gives
\eqref{eq:physical-riesz-cell-chromatic-error}.  A free collar of length
comparable to its starting radius conjugates the reflection coefficient by
$e^{2ikd}$, while its angular action is
$O((1/r_0-1/r_1)D^2)=o(1)$ on the already fixed block.  Rechoose the common
carrier beyond that collar and absorb this error into $\epsilon$.  This proves
the phase assertion.  Finally choose the next cell beyond the end of the
preceding one.
This proves disjointness and~\eqref{eq:physical-riesz-cell-cost}.
\end{proof}

\subsection{Uniform lacunary placement and the fixed potential}

The radial translations must work for a continuum of energies and for all
later positive smoothings and Galerkin cutoffs.  The required uniformity is a
compact-family form of the Riemann--Lebesgue lemma.

\begin{lemma}[Uniform torus placement]
\label{lem:uniform-torus-placement}
Let $I$ be a compact interval and let
$\mathcal F\Subset L^1(I;C(\mathbb T^q))$.  Given $\delta>0$ and $M>0$,
there are real frequencies $M<n_1<\cdots<n_q$ such that, uniformly for
$F\in\mathcal F$,
\begin{equation}
 \left|\int_I F(k,e^{in_1k},\ldots,e^{in_qk})\dd k
 -\int_I\int_{\mathbb T^q}F(k,z)\dd z\dd k\right|<\delta.
 \label{eq:uniform-torus-placement}
\end{equation}
The frequencies may be chosen successively: once $n_1,\ldots,n_{q-1}$ have
been fixed to meet earlier finite lists of inequalities, $n_q$ can be taken
arbitrarily large without disturbing those inequalities.  The same conclusion
holds simultaneously for finitely many compact families.
\end{lemma}

\begin{proof}
Approximate the compact family by finitely many functions which are
trigonometric polynomials in $z$ with $L^1(I)$ coefficients.  Choose $n_1$
so that every nonconstant $z_1$ Fourier mode is small.  Having chosen
$n_1,\ldots,n_{t-1}$, every Fourier mode involving a nonzero power of $z_t$
has frequency $\ell_t n_t+O(1)$ and tends to zero as $n_t\to\infty$ by the
Riemann--Lebesgue lemma.  Its zero power is the $z_t$ average.  Induct on $t$
and absorb the finite-net errors.  Earlier inequalities do not contain the
new frequency and are therefore unchanged.  Taking the maximum of finitely
many thresholds proves the final assertion.
\end{proof}

\begin{theorem}[Failure of fixed-potential buffered selected entropy]
\label{thm:fixed-potential-riesz-counterexample}
There are a bounded real zonal potential
$V\in X_3$, with
\begin{equation}
 \int_{\mathbb S^2}V(r,\omega)\dd\omega=0
 \quad\text{for a.e. }r,
 \label{eq:fixed-counterexample-zero-mean}
\end{equation}
and a compact interval $K\Subset(0,\infty)$ with the following property.
For every choice
\begin{equation}
 B_n,R_n,L_n\longrightarrow\infty,qquad
 \frac{L_n}{B_n+R_n+1}\longrightarrow\infty,
 \label{eq:arbitrary-buffered-diagonal}
\end{equation}
let
\begin{equation}
 V_n=\mathbf1_{[1,R_n]}\mathsf S_{B_n}V,
 \label{eq:fixed-counterexample-approximants}
\end{equation}
and let $J_n(k)$ be the Jost matrix of
$H_{R_n,L_n}(V_n)$.  Then
\begin{equation}
 \int_K\log^-\|J_n(k)^{-1}e_0\|^2\dd k
 \longrightarrow+\infty.
 \label{eq:fixed-potential-selected-entropy-divergence}
\end{equation}
In particular, the $N=1$ hypothesis
\eqref{eq:buffered-diagonal-PSR} is false even when one is free to choose the
cofinal positive smoothing, radial truncation, and double-buffered Galerkin
diagonal.
\end{theorem}

\begin{proof}
Fix $k_\star>0$ and a small compact interval $K$ about it.  Take
$m_j=\lceil C_0\log(j+1)\rceil$, with $C_0$ large enough that
\begin{equation}
 \sum_j e^{-cm_j}<\infty.
 \label{eq:riesz-cell-cost-summable}
\end{equation}
Use Proposition~\ref{prop:scalar-reciprocal-riesz-cell} with errors
$\epsilon_j$ satisfying $\sum_j\epsilon_j^\beta<\infty$, and put the resulting
cells on successive disjoint shells.  Their amplitudes are uniformly bounded;
by~\eqref{eq:physical-riesz-cell-cost} and
\eqref{eq:riesz-cell-cost-summable},
\begin{equation}
 \|V\|_{X_3}^2
 =\sum_j\|v_{m_j,\epsilon_j}\|_{L^2_{r,\omega}}^2<\infty.
 \label{eq:fixed-riesz-potential-cost}
\end{equation}
The mean-zero and reality assertions hold cell by cell.

Choose the free radial translations successively by
Lemma~\ref{lem:uniform-torus-placement}.  At the $j$th choice the earlier
cells, a finite range of smoothing degrees, all radial cuts through the
current cell, and all Galerkin cutoffs which resolve its finite angular band
form finitely many compact transfer-matrix families.  A diagonal choice with
error at most $2^{-j}$ therefore makes the physical energy average agree with
the independent circle-phase average, uniformly for every later member of
those families.  Increasing the finite range at each step covers every
positive smoothing degree, every radial cut, and every finite cutoff.  The
high-channel Schur estimate
\eqref{eq:buffered-edge-self-energy-vanishes} supplies the remaining uniform
passage from a resolved finite active block to the full double-buffered
Galerkin model.

Group the cells into dyadic blocks
\begin{equation}
 \mathcal B_q=\{j:2^{q-1}<j\leq2^q\}.
\end{equation}
For $N=2^q$, Proposition~\ref{prop:robust-reciprocal-cascade} applies on
\begin{equation}
 |s-1|\leq w_N
 =\frac{c}{N\log(N+1)\log\log(N+2)}.
\end{equation}
Its conditional version, obtained by starting the proof at the left endpoint
of $\mathcal B_q$, gives circle-phase mean loss at least $c2^q-C$.  Because
$s=k_\star/k$ is bi-Lipschitz near $k_\star$, the corresponding energy window
has length comparable to $w_{2^q}$.  On the complement, the conditional phase
mean is nonnegative by Lemma~\ref{lem:translated-barrier-hardy}.  Hence every
resolved dyadic block contributes at least
\begin{equation}
 c2^qw_{2^q}-C2^{-q}
 \geq \frac{c'}{q\log(q+2)}-C2^{-q}
 \label{eq:dyadic-block-entropy-gain}
\end{equation}
to the energy-averaged selected loss.  The uniform placement error is included
in the summable final term.

A smoothing or radial cut may leave one cell only partially resolved.  Its
circle-phase mean cannot decrease the accumulated loss, again by
\eqref{eq:barrier-hardy-nonnegative-loss}; uniform placement bounds its actual
negative contribution by the assigned summable error.  Cofinality of
\eqref{eq:arbitrary-buffered-diagonal} implies that every fixed dyadic block
is eventually fully resolved and remains far below the Galerkin edge.
Therefore~\eqref{eq:dyadic-block-entropy-gain} and
\begin{equation}
 \sum_{q\geq2}\frac1{q\log q}=+\infty
\end{equation}
give divergence of the signed selected loss.

Finally, pointwise for $t>0$ one has $-\log t\leq\log^-t$.  Apply this with
$t=\|J_n(k)^{-1}e_0\|^2$.  The divergent lower bound for the signed selected
loss therefore implies~\eqref{eq:fixed-potential-selected-entropy-divergence}
directly.  This proves the theorem.
\end{proof}

\begin{corollary}[What the counterexample decides]
\label{cor:fixed-potential-route-closed}
The universal endpoint estimates $\mathrm{SCI}_1$ and $\mathrm{HFI}_1$ are
false in every form which implies the fixed selected-channel entropy bound
\eqref{eq:buffered-diagonal-PSR}.  The exact signed centrifugal and angular
heat interpolation identities remain valid, but their proposed uniform
endpoint bounds do not.  Theorem
\ref{thm:fixed-potential-riesz-counterexample} does not disprove Simon's
conjecture: it controls one fixed probe only, while absolutely continuous
spectral density may survive in moving angular directions.  Any continuation
of the conjecture must therefore be basis-free (for example, a moving-probe
or full-operator entropy), or use a mechanism other than selected-channel
entropy.
\end{corollary}

\subsection{Reciprocal phase synchronization}

The earlier obstruction to stacking localized barriers was a possible
exponential amplification of a small non-real tail.  Reciprocal barriers do the
opposite in phase average.  Let $a=\cosh\rho$, $b=\sinh\rho$, and let $H$ be
the symmetric unitary reciprocal phase on an active block.  An ideal refocused
barrier acts by
\begin{equation}
 \Phi_z(H)=(a\Id+zbH)^{-1}(aH+\bar z b\Id),\qquad |z|=1.
 \label{eq:reciprocal-mobius-map}
\end{equation}
For $h,\lambda\in\mathbb T$, direct cross multiplication gives
\begin{equation}
 |\Phi_z(h)-\Phi_z(\lambda)|
 =\frac{|h-\lambda|}
 {|a+zbh|\,|a+zb\lambda|}.
 \label{eq:mobius-chord-identity}
\end{equation}
At a synchronized phase the projective derivative is
\begin{equation}
 L_z(\lambda)=|a+zb\lambda|^{-2},\qquad
 \frac1{2\pi}\int_0^{2\pi}L_{e^{i\phi}}(\lambda)\dd\phi=1.
 \label{eq:mobius-critical-derivative}
\end{equation}
For every $0<\beta<1$, strict Jensen gives
\begin{equation}
 q_\beta:=\frac1{2\pi}\int_0^{2\pi}
 L_{e^{i\phi}}(\lambda)^\beta\dd\phi<1.
 \label{eq:mobius-fractional-contraction}
\end{equation}
The value is independent of $\lambda$.

\begin{proposition}[Robust reciprocal-cascade entropy]
\label{prop:robust-reciprocal-cascade}
Consider $N$ successive real reciprocal cells on nested finite active blocks.
Suppose that, after the forward and reverse converter gauges, the $j$th cell
has transfer coefficients
\begin{equation}
 A_j(s)=a\Id+E_{j,+}(s),\qquad
 B_j(s)=b\Id+E_{j,-}(s)
 \label{eq:near-scalar-cell-coefficients}
\end{equation}
on the active block, with the exact flux relations and
\begin{equation}
 \|E_{j,+}(s)\|+\|E_{j,-}(s)\|
 \leq\eta_j(s),\qquad
 \eta_j(s)\leq C m_j|s-1|+\epsilon_j,
 \label{eq:near-scalar-cell-error}
\end{equation}
where $m_j\leq C_0\log(j+1)$ and
$\sum_j\epsilon_j^\beta$ is sufficiently small.  Starting from the free
constant channel, let $J_N(s,z)$ be the resulting Jost block.  There are
$c_0,C_1>0$ such that, for
\begin{equation}
 |s-1|\leq w_N,\qquad
 w_N=\frac{c_0}{N\log(N+1)\log\log(N+2)},
 \label{eq:dyadic-coherence-window}
\end{equation}
one has
\begin{equation}
 \int_{\mathbb T^N}
 -\log\|J_N(s,z)^{-1}e_0\|\dd z
 \geq c_0N-C_1.
 \label{eq:robust-cascade-linear-loss}
\end{equation}
Every circle carries normalized Haar measure.
\end{proposition}

\begin{proof}
With zero error and $s=1$, the refocused active Jost block has the scalar
recurrence
\begin{equation}
 d_j=a d_{j-1}+z_jb\overline{d_{j-1}},\qquad d_0=1.
 \label{eq:scalar-hyperbolic-recurrence}
\end{equation}
Thus
\begin{equation}
 \int_{\mathbb T}\log|a d+zb\bar d|\dd z
 =\log|d|+\log a,
 \label{eq:scalar-barrier-mean-gain}
\end{equation}
and the exact mean loss is $N\log a$.

For the perturbed cascade put $H_j=J_j^{-1}\overline{J_j}$ in the refocused
active frame.  Its ideal update is~\eqref{eq:reciprocal-mobius-map}; the
operator-norm perturbation is $O(\eta_j)$.  If the spectrum of $H_j$ lies in
an arc $I$, then the ideal image lies in $\Phi_z(I)$ and
\begin{equation}
 |\Phi_z(I)|=\int_I|a+zbe^{i\vartheta}|^{-2}\dd\vartheta.
 \label{eq:mobius-arc-length}
\end{equation}
Its phase mean is $|I|$.  Strict concavity and
\eqref{eq:mobius-fractional-contraction} give, uniformly for arcs of length at
most a fixed $\ell_0<2\pi$,
\begin{equation}
 \mathbb E_z|\Phi_z(I)|^\beta\leq q|I|^\beta,\qquad q<1.
 \label{eq:arc-fractional-contraction}
\end{equation}

There is only one way a short arc can make a macroscopic excursion: the random
repelling point of $\Phi_z$ must fall within $O(|I|+\eta_j)$ of it.
Equivalently, apply the maximal inequality to the nonnegative arc-length
martingale generated by~\eqref{eq:mobius-arc-length}.  An error arc injected
with length $C\eta_j$ has probability at most
$C\eta_j/\ell_0$ of ever reaching $\ell_0$.  Hence the excursion probability
before time $N$ is at most
\begin{equation}
 C\sum_{j\leq N}\eta_j(s)
 \leq \frac{C}{\log\log(N+2)}+C\sum_j\epsilon_j,
 \label{eq:mobius-excursion-probability}
\end{equation}
where~\eqref{eq:dyadic-coherence-window} and
$\sum_{j\leq N}m_j=O(N\log N)$ were used.  On the complement,
\eqref{eq:arc-fractional-contraction} and subadditivity of $t^\beta$ keep the
reciprocal phase in an arbitrarily short arc, except for a set of phase
histories of arbitrarily small measure.

If $H_j$ lies in an arc of length $o(1)$ around $\lambda_j$, then
$H_j\bar y_j=y_j$ implies
\begin{equation}
 \|\bar y_j-\bar\lambda_j y_j\|=o(1)\|y_j\|.
\end{equation}
Lemma~\ref{lem:translated-barrier-hardy} and
\eqref{eq:near-scalar-cell-error} give a conditional phase-mean loss
$\log a-o(1)$ at such a stage.  At the exceptional stages the conditional
mean remains nonnegative by~\eqref{eq:barrier-hardy-nonnegative-loss}.
Choose the summable error budget small, discard finitely many initial stages,
and sum the conditional means.  This proves
\eqref{eq:robust-cascade-linear-loss}.
\end{proof}

For a physical inner Jost matrix $J$, a reflectionless real converter with
transmission $U$, and $x=J^{-1}e_0$, the state at the reflector is $y=Ux$ and
the reciprocal phase is
\begin{equation}
 S=UJ^{-1}\overline J\,U^T.
 \label{eq:riesz-reciprocal-phase}
\end{equation}
The half-line scattering identity makes $S$ unitary and gives the exact
outgoing-square relation
\begin{equation}
 S\bar y=y.
 \label{eq:riesz-outgoing-square}
\end{equation}

\subsection{The zonal quarter-Talbot Riesz factor}

Put $D=-\partial_x^2$ on the circle and fix $0<a<\pi/4$.  Define
\begin{equation}
 g_s(x)=e^{-i\pi sD/2}e^{-ias\cos x}.
 \label{eq:flat-talbot-factor}
\end{equation}
Even and odd Fourier modes have multipliers $1$ and $-i$ at $s=1$, so
\begin{align}
 g_1(x)&=\cos(a\cos x)-\sin(a\cos x)>0,
 \label{eq:positive-talbot-factor}\\
 p(x)&=g_1(x)^2=1-\sin(2a\cos x),\qquad
 \frac1{2\pi}\int_0^{2\pi}p(x)\dd x=1.
 \label{eq:talbot-riesz-density}
\end{align}
Strict Jensen gives the positive tilted entropy
\begin{equation}
 D_p:=\frac1{2\pi}\int_0^{2\pi}p\log p\dd x>0.
 \label{eq:talbot-positive-entropy}
\end{equation}

\begin{proposition}[Zonal spherical Talbot factor and reciprocal inverse]
\label{prop:zonal-talbot-factor}
Let $\mathcal K\subset C^\infty(\mathbb S^2)$ be finite dimensional,
zonal, and invariant under conjugation.  Uniformly for $s$ in a fixed compact
neighbourhood of $1$ and odd integers $L\to\infty$,
\begin{equation}
 e^{-i\pi sA_3/(2L^2)}
   \bigl(e^{-ias\cos(L\theta)}f\bigr)
 =g_s(L\theta)f+\mathcal R_{L,s}f,\qquad
 \|\mathcal R_{L,s}\|_{\mathcal K\to L^2}
 \leq C_{\mathcal K,a}\frac{\sqrt{\log L}}L.
 \label{eq:spherical-talbot-asymptotic}
\end{equation}
A sign-reversed mask following a free flight of scaled duration
$3\pi/(2L^2)$ gives an operator $W_{L,s}$ such that, for the forward factor
$U_{L,s}$ above,
\begin{align}
 \|W_{L,1}U_{L,1}-\Id\|_{\mathcal K\to L^2}
 &\leq C_{\mathcal K,a}\frac{\sqrt{\log L}}L,
 \label{eq:talbot-reciprocal-inverse}\\
 \|W_{L,s}U_{L,s}-W_{L,1}U_{L,1}\|_{\mathcal K\to L^2}
 &\leq C_{\mathcal K,a}|s-1|.
 \label{eq:talbot-chromatic-bound}
\end{align}
\end{proposition}

\begin{proof}
The identity $\cos(L\theta)=T_L(\cos\theta)$ makes the mask globally smooth.
On zonal functions
\begin{equation}
 A_3=-\partial_\theta^2-\cot\theta\,\partial_\theta.
\end{equation}
Compare its scaled evolution with that of $-\partial_\theta^2$ by Duhamel.
For every smooth periodic $H$,
\begin{equation}
 \|\cot\theta\,\partial_\theta H(L\theta)\|_{L^2(\mathbb S^2)}
 \leq C_H L\sqrt{\log L}.
\end{equation}
The evolution time is $O(L^{-2})$, giving
\eqref{eq:spherical-talbot-asymptotic}; derivatives falling on the fixed
envelope cost only $O(L^{-1})$.  In the flat fast variable
$e^{-i3\pi D/2}=e^{i\pi D/2}$.  Thus the sign-reversed mask after the
$3\pi/2$ flight is the inverse factor at $s=1$.  The same Duhamel comparison
proves~\eqref{eq:talbot-reciprocal-inverse}; differentiation in $s$ on the
scaled finite fast-frequency block proves~\eqref{eq:talbot-chromatic-bound}.
\end{proof}

\begin{proposition}[Uniform Riesz concentration on a finite incoming block]
\label{prop:talbot-riesz-concentration}
Fix $0<\alpha<D_p$, a finite-dimensional conjugation-invariant zonal space
$\mathcal K$, an integer $m$, and $\epsilon>0$.  There are odd,
superlacunary frequencies $L_1<\cdots<L_m$, a real finite-band function
$0\leq\chi_m\leq1$, and forward and reverse Talbot products $U_{m,s}$ and
$W_{m,s}$ such that, with
\begin{equation}
 G_m(\theta)=\prod_{t=1}^m g_1(L_t\theta),\qquad
 E_m=\left\{\sum_{t=1}^m\log p(L_t\theta)\geq\alpha m\right\},
 \label{eq:riesz-product-typical-set}
\end{equation}
one has
\begin{align}
 \|U_{m,1}f-G_mf\|_2
 &\leq\epsilon\|f\|_2,
 \label{eq:riesz-forward-approximation}\\
 \|(1-\chi_m)U_{m,1}f\|_2
 &\leq(\epsilon+Ce^{-cm})\|f\|_2,
 \label{eq:riesz-concentration}\\
 \|\chi_m\|_2^2&\leq Ce^{-cm},
 \label{eq:riesz-small-support-cost}\\
 \|W_{m,1}U_{m,1}-\Id\|_{\mathcal K\to L^2}
 &\leq\epsilon,
 \label{eq:riesz-refocusing}\\
 \|W_{m,s}U_{m,s}-W_{m,1}U_{m,1}\|
 +\|U_{m,s}-U_{m,1}\|
 &\leq C_{\mathcal K,a}m|s-1|.
 \label{eq:riesz-chromatic-error}
\end{align}
The last two norms are from $\mathcal K$ to $L^2$.
\end{proposition}

\begin{proof}
Choose the frequencies successively.  Riemann--Lebesgue, uniformly on the unit
sphere of $\mathcal K$, makes every exponential moment used below as close as
desired to the product moment of independent copies of $\log p$.  On $E_m$,
$\prod p(L_t\theta)\geq e^{\alpha m}$, hence $|E_m|\leq Ce^{-\alpha m}$.
Under the tilted density $p(x)\dd x/(2\pi)$, the mean of $\log p$ is
$D_p>\alpha$.  Chernoff therefore gives, uniformly for $f\in\mathcal K$,
\begin{equation}
 \int_{E_m^c}|G_m|^2|f|^2\dd\omega
 \leq Ce^{-cm}\|f\|_2^2.
 \label{eq:riesz-tilted-large-deviation}
\end{equation}
Apply the positive finite-band spherical smoother
\eqref{eq:positive-angular-smoother} to a slightly enlarged typical set.
For sufficiently large smoothing degree it gives $0\leq\chi_m\leq1$,
\eqref{eq:riesz-small-support-cost}, and the same weighted-tail estimate.
Proposition~\ref{prop:zonal-talbot-factor}, with still more separated
frequencies, proves~\eqref{eq:riesz-forward-approximation} and
\eqref{eq:riesz-concentration}.  Apply the reciprocal factors in reverse order
to obtain~\eqref{eq:riesz-refocusing}.  Telescoping the two products and using
\eqref{eq:talbot-chromatic-bound} proves~\eqref{eq:riesz-chromatic-error}.
\end{proof}


\section{Notation and normalization checks}

\begin{enumerate}[label=(\roman*)]
 \item $P_N\leq P_L$ means inclusion of ranges; it is not equivalent to the numerical statement $L\geq N$ unless the harmonic enumeration has been fixed accordingly.
 \item The physical probe contains the factor $r^{-(d-1)/2}$ in~\eqref{eq:physical-probe}; the flattened probe does not.
 \item The Jost normalization is fixed by~\eqref{eq:jost-normalization}.  Calling $e^{ikr}$ ``unit flux'' without a factor $(2k)^{-1/2}$ would be inconsistent.
 \item The exterior matrix measure is not identified with the full-space measure until the multiplicity preserving bridge in Theorem~\ref{thm:completion}(d) is supplied.
 \item The proposed FCEN/PSR statement is a sufficient strengthening of Simon's conjecture, not a known equivalent reformulation.
\end{enumerate}

\section*{Acknowledgements and AI-use disclosure}

The author is not a trained mathematician.  GPT-5.6 Sol developed the
mathematical constructions in this manuscript from the ground up under the
author's direction.

\begin{thebibliography}{99}

\bibitem{Simon2019}
B.~Simon,
\emph{Tosio Kato's work on non-relativistic quantum mechanics, Part 2},
Bull. Math. Sci. \textbf{9} (2019), 1950005; Conjecture~20.2.
\url{https://math.caltech.edu/SimonPapers/R62.pdf}

\bibitem{DeiftKillip1999}
P.~Deift and R.~Killip,
\emph{On the absolutely continuous spectrum of one-dimensional Schr\"odinger operators with square summable potentials},
Commun. Math. Phys. \textbf{203} (1999), 341--347.
\url{https://doi.org/10.1007/s002200050615}

\bibitem{Hunt1968}
R.~A.~Hunt,
\emph{On the convergence of Fourier series},
in \emph{Orthogonal Expansions and Their Continuous Analogues},
Southern Illinois University Press, Carbondale, 1968, 235--255.

\bibitem{ChristKiselev1998}
M.~Christ and A.~Kiselev,
\emph{Absolutely continuous spectrum for one-dimensional Schr\"odinger operators with slowly decaying potentials: some optimal results},
J. Amer. Math. Soc. \textbf{11} (1998), 771--797.
\url{https://doi.org/10.1090/S0894-0347-98-00276-8}

\bibitem{DenisovSurvey2012}
S.~A.~Denisov,
\emph{Multidimensional $L^2$ conjecture: a survey},
arXiv:1203.4255 (2012).
\url{https://arxiv.org/abs/1203.4255}

\bibitem{DenisovWKB2009}
S.~A.~Denisov,
\emph{An evolution equation as the WKB correction in long-time asymptotics of Schr\"odinger dynamics},
Comm. Partial Differential Equations \textbf{33} (2008), 307--319.
\url{https://doi.org/10.1080/03605300701249655}

\bibitem{MolchanovVainberg2004}
S.~A.~Molchanov and B.~R.~Vainberg,
\emph{Schr\"odinger operators with matrix potentials: transition from the absolutely continuous to the singular spectrum},
J. Funct. Anal. \textbf{215} (2004), 111--129.
\url{https://arxiv.org/abs/math-ph/0308012}

\bibitem{LNS2005}
A.~Laptev, S.~Naboko, and O.~Safronov,
\emph{A Szeg\H{o} condition for a multidimensional Schr\"odinger operator},
J. Funct. Anal. \textbf{219} (2005), 285--305.
\url{https://www.ma.ic.ac.uk/~alaptev/Papers/schr.pdf}

\bibitem{DamanikPushnitskiSimon2008}
D.~Damanik, A.~Pushnitski, and B.~Simon,
\emph{The analytic theory of matrix orthogonal polynomials},
Surv. Approx. Theory \textbf{4} (2008), 1--85.
\url{https://arxiv.org/abs/0711.2703}

\bibitem{Perelman2005}
G.~Perelman,
\emph{Stability of the absolutely continuous spectrum for multidimensional Schr\"odinger operators},
Int. Math. Res. Not. \textbf{2005} (2005), 2289--2313.
\url{https://doi.org/10.1155/IMRN.2005.2289}

\bibitem{Lotoreichik2013}
V.~Lotoreichik,
\emph{Singular Values and Trace Formulae for Resolvent Power Differences of Self-Adjoint Elliptic Operators},
Monographic Series TU Graz: Computation in Engineering and Science, vol.~22,
Verlag der Technischen Universit\"at Graz, Graz, 2013; Theorem~4.17 and Corollary~4.19.
\url{https://doi.org/10.3217/978-3-85125-304-7}

\bibitem{GPS2007}
F.~Gesztesy, A.~Pushnitski, and B.~Simon,
\emph{On the Koplienko spectral shift function, I. Basics},
J. Math. Phys. Anal. Geom. \textbf{4} (2008), 63--107.
\url{https://arxiv.org/abs/0705.3629}

\bibitem{Denisov2022}
S.~A.~Denisov,
\emph{Spatial asymptotics of Green's function and applications},
operator-valued Schr\"odinger analysis (2022).
\url{https://people.math.wisc.edu/~denissov/final.pdf}

\bibitem{LaptevSafronov2024}
A.~Laptev and O.~Safronov,
\emph{Absolutely continuous spectrum of a typical Schr\"odinger operator with an operator-valued potential},
St. Petersburg Math. J. \textbf{35} (2024), 171--183.
\url{https://doi.org/10.1090/spmj/1799}

\bibitem{Denisov2025}
S.~A.~Denisov,
\emph{Wave packet decomposition for Schr\"odinger evolution with rough potential and generic value of parameter},
arXiv:2408.03470v4 (2025).
\url{https://arxiv.org/abs/2408.03470}

\bibitem{TaoThiele2012}
T.~Tao and C.~Thiele,
\emph{Nonlinear Fourier analysis},
arXiv:1201.5129 (2012).
\url{https://arxiv.org/abs/1201.5129}

\bibitem{LeghtasSarletteRouchon2011}
Z.~Leghtas, A.~Sarlette, and P.~Rouchon,
\emph{Adiabatic passage and ensemble control of quantum systems},
J. Phys. B: At. Mol. Opt. Phys. \textbf{44} (2011), 154017.
\url{https://doi.org/10.1088/0953-4075/44/15/154017}

\bibitem{RobinEtAl2021}
R.~Robin, N.~Augier, U.~Boscain, and M.~Sigalotti,
\emph{Ensemble qubit controllability with a single control via adiabatic and rotating wave approximations},
arXiv:2003.05831v2 (2021).
\url{https://arxiv.org/abs/2003.05831}

\bibitem{BoussaidCaponigroChambrion2012}
N.~Boussa\"id, M.~Caponigro, and T.~Chambrion,
\emph{Which notion of energy for bilinear quantum systems?},
Proc. 4th IFAC Workshop on Lagrangian and Hamiltonian Methods for Non Linear
Control (2012), 226--230.
\url{https://arxiv.org/abs/1302.0815}

\bibitem{Chambrion2013}
T.~Chambrion,
\emph{A sufficient condition for partial ensemble controllability of bilinear
Schr\"odinger equations with bounded coupling terms},
52nd IEEE Conference on Decision and Control (2013), 3708--3713.
\url{https://arxiv.org/abs/1303.0298}

\bibitem{Bournissou2022}
M.~Bournissou,
\emph{Small-time local controllability of the bilinear Schr\"odinger equation,
despite a quadratic obstruction, thanks to a cubic term},
arXiv:2203.03955v2 (2022).
\url{https://arxiv.org/abs/2203.03955}

\bibitem{BeauchardPozzoli2025}
K.~Beauchard and E.~Pozzoli,
\emph{Small-time approximate controllability of bilinear Schr\"odinger
equations and diffeomorphisms},
arXiv:2410.02383v2 (2025).
\url{https://arxiv.org/abs/2410.02383}

\bibitem{ChambrionThomann2020}
T.~Chambrion and L.~Thomann,
\emph{On the bilinear control of the Gross--Pitaevskii equation},
arXiv:1810.09792v2 (2020).
\url{https://arxiv.org/abs/1810.09792}

\end{thebibliography}

\end{document}
